\documentclass[11pt,a4paper]{article}
\usepackage[T1]{fontenc}
\usepackage{lmodern,microtype}
\usepackage[margin=27mm]{geometry}
\usepackage{amsmath,amssymb,amsthm,mathtools,booktabs,array}
\PassOptionsToPackage{hyphens}{url}
\usepackage[hidelinks]{hyperref}
\emergencystretch=2em
\newtheorem{theorem}{Theorem}[section]
\newtheorem{proposition}[theorem]{Proposition}
\newtheorem{lemma}[theorem]{Lemma}
\newtheorem{corollary}[theorem]{Corollary}
\theoremstyle{definition}
\newtheorem{definition}[theorem]{Definition}
\theoremstyle{remark}
\newtheorem{remark}[theorem]{Remark}
\newcommand{\GT}{\widehat{\mathrm{GT}}}
\newcommand{\GTz}{\widehat{\mathrm{GT}}_0}
\newcommand{\Zh}{\widehat{\mathbf Z}}
\newcommand{\Z}{\mathbf Z}
\newcommand{\Q}{\mathbf Q}
\newcommand{\F}{\mathbf F}
\newcommand{\GQ}{G_{\Q}}
\newcommand{\ab}{\mathrm{ab}}
\newcommand{\cts}{\mathrm{cts}}
\newcommand{\cl}{\overline}
\DeclareMathOperator{\Hom}{Hom}
\DeclareMathOperator{\Gal}{Gal}
\DeclareMathOperator{\Aut}{Aut}
\DeclareMathOperator{\Out}{Out}
\DeclareMathOperator{\Br}{Br}
\DeclareMathOperator{\SL}{SL}
\DeclareMathOperator{\PSL}{PSL}
\DeclareMathOperator{\Ad}{Ad}
\DeclareMathOperator{\res}{res}
\DeclareMathOperator{\infla}{inf}
\newcommand{\workbenchurl}{https://github.com/KokunoYumeto/k-theory-and-galois-groups-workbench}
\newcommand{\recordurl}{https://github.com/KokunoYumeto/mathematics-workbenches}
\title{The arithmetic quotient of the profinite Grothendieck--Teichm\"uller group}
\author{}
\date{25 November 2026}
\begin{document}
\maketitle
\begin{abstract}
Ihara extended the Kummer cocycles to the profinite Grothendieck--Teichm\"uller group, Enriquez proved their distribution relations, and Nakamura and Schneps related the cocycle of two to the Teichm\"uller tower. Starting from the established decomposition into normalized and Galois elements, we identify the arithmetic quotient and the obstruction to extending invariant characters. Dividing the integral Kummer cocycles by two gives independent invariant quadratic characters with quartic Kummer restrictions; the invariant space and its restriction image both have countably infinite dimension. We then remove all rational Kummer coordinates and compute successive exact two-adic monodromy images through weight ten, already attained by the Galois subgroup. Their common kernel has no Magnus term below degree eleven. An integral beta measure gives exact arithmetic images in every odd weight and removes the entire dyadic metabelian multiplier. Combining these reductions and removing a primitive mixed degree-eleven coordinate gives a smaller group carrying the same relative abelian obstruction. We determine which of the resulting quadratic characters extend through the earlier normalizations. The full restriction image and relative kernel remain open.

The two corrected adelic modular frames have exactly their joint Galois image, giving the quotient associated with the maximal abelian extensions of the Gaussian and Eisenstein fields. We determine its overlap with rational radicals and settle the complete prime-two and exceptional prime-three compatibility. Furusho's nonabelian dyadic identities prove the normalized harmonic comparison. Finite-subgroup centralizers, genus-zero rigidity and the exact cubic distribution relation prove the normalized cubic comparison. The remaining conditions are explicit tame tests. A finite-cover determinant gives another explicit relative character whose nonvanishing is undecided. We also compute a local Brauer invariant and the Hecke conductors at the quartic Kummer layer. The mathematics and text were produced by GPT-6 Astra Pro (ChatGPT Pro mode), with independent same-model re-proofs and targeted checks; no human has checked them.
\end{abstract}

\section{The question and the results}
The comparison between the absolute Galois group $\GQ$ and Drinfeld's profinite Grothendieck--Teichm\"uller group $\GT$ is a central question in Grothendieck's programme. Drinfeld introduced the group by the two-cycle, three-cycle and pentagon equations; Ihara developed the arithmetic characters of its Galois action. We use the standard tangential-base-point inclusion $\GQ\subset\GT$. All groups, homomorphisms and cohomology in this paper are profinite or continuous unless another meaning is stated.

Let $\lambda:\GT\to\Zh^\times$ be the multiplier, and let
$\rho_2:\GT\to\Zh(1)$ be the cocycle extending the action on the compatible positive roots of $2$. Thus
\[
 \rho_2(gh)=\rho_2(g)+\lambda(g)\rho_2(h).
\]
Put
\[
 K_2=\Q(\mu_\infty,2^{1/\infty}),\qquad
 N=\ker\lambda\cap\ker\rho_2,\qquad H=\GQ,\qquad K=G_{K_2}.
\]
The notation $N$ denotes the normalized subgroup also written $\Lambda_0$. It lies in the group $\Lambda$ of Hatcher, Lochak and Schneps, whose formulas act on the pure Teichm\"uller tower. The input used here is the decomposition
\begin{equation}\label{eq:decomposition}
 \GT=NH.
\end{equation}
It is Theorem 2 of \cite{H17}; Lemma 7 there gives the cocycle law and evenness of $\rho_2$ when $\lambda=1$. The same note proves relation (IV), relation $(\mathrm{III}')$ and the resulting identification of the Nakamura--Schneps subgroup with $\GT$ in Theorems 1 and 20 and Corollary 21. These statements are cited background, with their stated AI checks and no human verification. The gentle group $\GTz$ is defined by the two- and three-cycle equations. Invertibility and a unit multiplier follow from these equations by \cite[Theorems A and E]{Monoid}; no extra invertibility test is imposed below.

For a prime $p$, write
\[
 X_p=H^1(N,\F_p)^{\GT},\qquad
 Y_p=H^1(K,\F_p)^H.
\]
Coefficients have trivial action. The restriction in the question lands in $Y_p$, although its codomain can of course be written $H^1(K,\F_p)$.

\begin{theorem}\label{thm:main-structural}
Assume \eqref{eq:decomposition} and the preceding normalization of $\rho_2$. Then $N$ is closed and normal, $N\cap H=K$, and
\[
 \GT/N\simeq A:=\Gal(K_2/\Q).
\]
In the affine coordinates $(u,t)$, with product
$(u,t)(v,s)=(uv,t+us)$, this quotient is
\[
 A=\{(u,t)\in\Zh^\times\ltimes\Zh:
 t\equiv\varepsilon_8(u)\pmod2\},
\]
where $\varepsilon_8(u)=0$ for $u\equiv1,7\pmod8$ and $1$ for $u\equiv3,5\pmod8$.
There is a canonical topological isomorphism
\[
 \GT^{\ab}\simeq\Zh^\times\times B,\qquad
 B=N/\overline{[\GT,N]K}.
\]
For every prime $p$, restriction identifies
\[
 \Hom(B,\F_p)\simeq\ker(X_p\longrightarrow Y_p).
\]
Consequently the relative kernel vanishes for every prime if and only if the multiplier identifies $\GT^{\ab}$ with $\Zh^\times$.
\end{theorem}

In words, the multiplier and the cocycle of two describe precisely this Galois quotient. All extra abelian characters belong to $B$. This is the precise form of Lochak and Schneps's Question 1.2 \cite[\S1, Question 1.2]{LS2006}. It is weaker than deciding whether $\GT=\GQ$: a group can have trivial relative abelianization and still have a nontrivial relative nonabelian quotient.

\begin{theorem}\label{thm:main-dimension}
With the same notation,
\[
 \dim_{\F_2} X_2=\dim_{\F_2}\operatorname{im}(X_2\longrightarrow Y_2)=\aleph_0.
\]
For every positive integer $d$, Ihara's integral cocycle $\rho_d^I$ is even on the multiplier-one subgroup of $\GTz$. For every odd prime $q$,
\[
 \psi_q=\frac{\rho_q^I|_N}{2}\pmod2
\]
is an invariant continuous character of $N$. Its restriction is the character of $K_2(q^{1/4})/K_2$, and these restrictions are linearly independent as $q$ varies. For every prime $p$, including odd $p$, $\dim_{\F_p}Y_p=\aleph_0$.
\end{theorem}

Thus the answer to the dimension question is infinite even if $\GT=\GQ$. It is the restriction, rather than the existence of invariant characters, that tests equality. In particular this theorem supplies no nonzero character of $\GT$ trivial on $\GQ$.

The normalization can be made substantially stronger. Let
\[
 L_{\mathrm{rad}}=\Q(\mu_\infty,q^{1/n}:q\text{ prime},\ n\geq1),
 \qquad
 N_{\mathrm{rad}}=N\cap\bigcap_{q\text{ odd prime}}\ker\rho_q^I.
\]
Here all radicals are the coherent positive roots.

\begin{theorem}\label{thm:main-rational-normalization}
With the preceding assumptions and notation, the joint map
\[
 N\longrightarrow\prod_{q\text{ odd prime}}2\Zh,
 \qquad n\longmapsto(\rho_q^I(n))_q
\]
is surjective already on $K$. Its kernel $N_{\mathrm{rad}}$ is
a proper closed normal subgroup, and
\[
 \GT=N_{\mathrm{rad}}\GQ,\qquad
 N_{\mathrm{rad}}\cap\GQ=G_{L_{\mathrm{rad}}},\qquad
 \GT/N_{\mathrm{rad}}\simeq\Gal(L_{\mathrm{rad}}/\Q).
\]
The latter quotient consists of the affine coordinates
$(u,(t_q)_{q\text{ prime}})$ satisfying
$t_q\equiv\epsilon_q(u)\pmod2$, with multiplication
$(u,t)(v,s)=(uv,t+us)$; here $\epsilon_q$ is the character of
$\Q(\sqrt q)$. The relative abelian quotient is unchanged:
\[
 B\simeq N_{\mathrm{rad}}/
       \overline{[\GT,N_{\mathrm{rad}}]G_{L_{\mathrm{rad}}}}.
\]
The invariant quadratic characters of $N$ that factor through all
rational Kummer coordinates are exactly the finite linear combinations
of the $\psi_q$ for odd primes $q$, and restriction is injective on
this entire subspace. No nonzero invariant character with values in
$\F_p$, for odd $p$, factors through these coordinates.
\end{theorem}

This removes all rational-radical data from the open problem at once.
An extra character would have to be nonzero on the smaller subgroup
$N_{\mathrm{rad}}$. The theorem does not show that its relative
abelianization is zero. Enriquez's distribution theorem is essential
for describing the whole rational Kummer subspace: it proves
$\rho_{ab}^I=\rho_a^I+\rho_b^I$ integrally on $\GT$
\cite[Theorem~8.1 and Remark~8.6]{Enriquez2007}.

The normalization also controls every modular congruence representation
and identifies the two individual elliptic-frame signs.

\begin{theorem}\label{thm:main-congruence-frames}
Let $g=(1,f)\in\GT$ satisfy $\rho_2(g)=0$, and put
\[
 U=\begin{pmatrix}1&1\\0&1\end{pmatrix},\qquad
 V=\begin{pmatrix}1&0\\-1&1\end{pmatrix}.
\]
In the signed adelic braid representation,
\[
 f(U^2,V^2)=I\quad\text{in }\operatorname{SL}_2(\Zh).
\]
For the signed cubic and harmonic frame bits $a_3(g),a_2(g)$
of \cite[Lemmas 12--15]{H17}, with the positive signed lifts
specified in Section~\ref{sec:global-frame-character}, one has
\[
 a_2(g)=a_3(g)=\rho_3^I(g)/2\pmod2.
\]
In particular both signs vanish on $N_{\mathrm{rad}}$.
\end{theorem}

The first assertion says that every normalized element induces the
identity on every signed modular congruence quotient. The second says
that these two particular frame coordinates are already rational
Kummer data. Neither assertion identifies the action on the full
profinite modular group with its congruence action. The proofs are
Theorem~\ref{thm:individual-frame-kummer} and
Proposition~\ref{prop:normalized-congruence-triviality}.

The rational Kummer characters do not exhaust even the explicitly
accessible arithmetic part of $X_2$.  The first two odd-weight coefficients of pro-$2$
monodromy give further exact arithmetic normalizations.

\begin{theorem}\label{thm:main-cubic-normalization}
For $g=(1,f)$ put $m_3(g)=[S^2T]f$ in the Magnus expansion
$x=1+S$, $y=1+T$.  Then
\[
 m_3(N)\subseteq8\Z_2,\qquad
 m_3(N_{\mathrm{rad}})=m_3(G_{L_{\mathrm{rad}}})=8\Z_2.
\]
Thus $\eta_3=m_3/8\pmod2$ belongs to $X_2$, and its restriction
to $G_{L_{\mathrm{rad}}}$ is nonzero.  It is linearly independent
of the entire rational Kummer subspace.  The joint map
\[
 N\longrightarrow
 \left(\prod_{q\text{ odd prime}}2\Zh\right)\times8\Z_2,
 \qquad n\longmapsto\bigl((\rho_q^I(n))_q,m_3(n)\bigr)
\]
is surjective already on $K$.  Its kernel
$N_{\mathrm{rad},3}=N_{\mathrm{rad}}\cap\ker m_3$ is a proper
closed normal subgroup of $N_{\mathrm{rad}}$ and of $\GT$, and
\[
 \GT=N_{\mathrm{rad},3}\GQ,\qquad
 B\simeq N_{\mathrm{rad},3}/
 \overline{[\GT,N_{\mathrm{rad},3}]
                 (N_{\mathrm{rad},3}\cap\GQ)}.
\]
The kernel of restriction from $X_2$ to
$H^1(N_{\mathrm{rad},3},\F_2)$ is precisely
\[
 \left(\bigoplus_{q\text{ odd prime}}\F_2\psi_q\right)
       \oplus\F_2\eta_3.
\]
On $N_{\mathrm{rad},3}$ put $m_5(g)=[S^4T]f$ and
$K_{\mathrm{rad},3}=N_{\mathrm{rad},3}\cap\GQ$.  Then
\[
 m_5(N_{\mathrm{rad},3})=m_5(K_{\mathrm{rad},3})=32\Z_2.
\]
Thus $N_{\mathrm{rad},3,5}=N_{\mathrm{rad},3}\cap\ker m_5$
is again proper and closed normal, with
\[
 \GT=N_{\mathrm{rad},3,5}\GQ,\qquad
 B\simeq N_{\mathrm{rad},3,5}/
 \overline{[\GT,N_{\mathrm{rad},3,5}]
                  (N_{\mathrm{rad},3,5}\cap\GQ)}.
\]
The divided character $\eta_5=m_5/32\pmod2$ on
$N_{\mathrm{rad},3}$ does not extend to an invariant quadratic
character of $N_{\mathrm{rad}}$ or of $N$.  The kernel of restriction
from $X_2$ to $H^1(N_{\mathrm{rad},3,5},\F_2)$ is still precisely
the displayed rational Kummer and cubic span.
\end{theorem}

This sharpens the subgroup on which the open relative character must
be detected.  It also distinguishes two arithmetic sources of invariant
characters: roots of rational numbers and cubic monodromy.
Soul\'e's characters and their extensions to GT are classical;
Marin obtains the degree-three and degree-five extensions from the
Burau representation \cite[\S6.3, p.~780]{Marin2008}.
The argument here determines the exact integral cubic and weight-five
images after the respective normalizations and gives concrete Galois witnesses.
The proofs are in Sections~\ref{sec:cubic-magnus-normalization}
and~\ref{sec:weight-five-normalization}.  An arithmetic commutator
obstructs extending the divided weight-five character to the original
invariant character problem.

The next odd weight gives another exact arithmetic normalization.

\begin{theorem}\label{thm:main-weight-seven-alternative}
Put $D_3=N_{\mathrm{rad},3}$, $D_5=N_{\mathrm{rad},3,5}$,
and $J_5=D_5\cap H$.
Every $g\in D_5$ starts in Magnus degree seven, and
$m_7(g)=[S^6T]f_g$ is an additive weight-seven coordinate with
\[
 m_7(D_5)=m_7(J_5)=128\Z_2.
\]
The proper closed normal subgroup $D_7=D_5\cap\ker m_7$ satisfies
\[
 \GT=D_7H,\qquad
 B\simeq D_7/\overline{[\GT,D_7](D_7\cap H)}.
\]
The invariant character $\eta_7=m_7/128\pmod2$ on $D_5$ has
nonzero arithmetic restriction. It has no invariant extension to
$N_{\mathrm{rad},3}$, $N_{\mathrm{rad}}$ or $N$.
The kernel of restriction from $X_2$ to $H^1(D_7,\F_2)$ is still
exactly the rational Kummer span together with $\F_2\eta_3$.
\end{theorem}

Section~\ref{sec:weight-seven-partial} gives the full integral
coefficient certificates and a Galois element with valuation seven.
The upper bound uses two coefficients of the full pentagon through
weight nine, in addition to the leading relation. This third
normalization gives a smaller subgroup on which the unresolved
relative character must be detected; it supplies no character
separating $\GT$ from its Galois subgroup.

The following layer is an even-weight bracket, so one cannot continue
by considering only the depth-one coefficients in odd weights.

\begin{theorem}\label{thm:main-weight-eight-normalization}
Every $g\in D_7$ has Magnus expansion starting in degree at least eight.
Put
\[
 m_8(g)=[S^4TST^2]f_g,\qquad J_7=D_7\cap H.
\]
This is an additive coordinate of conjugation weight $\lambda_2^8$,
and
\[
 m_8(D_7)=m_8(J_7)=\Z_2.
\]
Its kernel $D_8$ is proper, closed and normal in $\GT$, satisfies
$\GT=D_8H$, and gives
\[
 B\simeq D_8/\overline{[\GT,D_8](D_8\cap H)}.
\]
Every element of $D_8$ starts in Magnus degree at least nine.
The invariant quadratic character $\eta_8=m_8\pmod2$ on $D_7$
extends invariantly to $D_5$ by the same raw coefficient modulo two,
but has no invariant extension to $D_3$, $N_{\mathrm{rad}}$ or $N$.
More precisely, no nonzero combination of its extension and $\eta_7$
on $D_5$ extends invariantly to $D_3$, and
\[
 \ker\!\left(H^1(D_3,\F_2)^{\GT}\longrightarrow H^1(D_8,\F_2)\right)
 =\F_2\eta_5.
\]
The kernel of restriction from $X_2$ to $H^1(D_8,\F_2)$ remains
exactly the rational Kummer span together with $\F_2\eta_3$.
\end{theorem}

The primitive leading line is twice the Ihara bracket of the rationally
normalized weight-three and weight-five directions. An arithmetic
commutator has odd coefficient on that line. This proves the full
integral image without lifting a finite jet. Section~\ref{sec:weight-eight-normalization}
gives the coefficient certificate and the arithmetic construction.

The next odd coordinate gives a further exact normalization.

\begin{theorem}\label{prop:main-weight-nine-bounds}
Put $J_j=D_j\cap H$ for $j=7,8$, and
$m_9(g)=[S^8T]f_g$ for $g\in D_7$.
This is a continuous homomorphism of conjugation weight $\lambda_2^9$,
and
\[
 m_9(D_7)=m_9(J_7)=m_9(D_8)=m_9(J_8)=512\Z_2.
\]
Its integral kernel $D_9=D_8\cap\ker m_9$ is proper, closed and
normal in $\GT$, and
\[
 \GT=D_9H,\qquad
 B\simeq D_9/\overline{[\GT,D_9](D_9\cap H)}.
\]
Every $D_9$ word starts in degree at least ten.
The joint map $(m_8,m_9):D_7\to\Z_2\times512\Z_2$
is surjective already on $J_7$ and has kernel $D_9$.
The invariant character $\eta_9=m_9/512\pmod2$ on $D_7$
has nonzero arithmetic restriction; its restriction to $D_8$
has no invariant extension to $D_5$.
The kernel of restriction from $X_2$ to $H^1(D_9,\F_2)$
remains exactly the rational Kummer span together with $\F_2\eta_3$.
\end{theorem}

Section~\ref{sec:weight-nine-bounds} proves the divisibility by an
integral measure and gives a Galois word attaining it. The smaller subgroup $D_9$ is therefore an unconditional
arithmetic normalization. These successive normalizations do not
settle the remaining relative character problem.

The tenth layer is again a mixed bracket coordinate.

\begin{theorem}\label{thm:main-weight-ten-normalization}
On $D_9$, put $m_{10}(g)=[S^5TST^3]f_g$ and $J_9=D_9\cap H$.
This is an additive coordinate of conjugation weight $\lambda_2^{10}$,
with
\[
 m_{10}(D_9)=m_{10}(J_9)=\Z_2.
\]
Its kernel $D_{10}$ is proper, closed and normal, and
\[
 \GT=D_{10}H,\qquad
 B\simeq D_{10}/\overline{[\GT,D_{10}](D_{10}\cap H)}.
\]
Every element of $D_{10}$ starts in Magnus degree at least eleven.
The invariant character $\eta_{10}=m_{10}\pmod2$ extends invariantly
to $D_8$, but not to $D_7$ or any of $D_5,D_3,N_{\mathrm{rad}},N$. The kernel of
restriction from $X_2$ to $H^1(D_{10},\F_2)$ is still exactly the
rational Kummer span together with $\F_2\eta_3$.
\end{theorem}

Section~\ref{sec:weight-ten-normalization} identifies the leading
line as the Ihara bracket of the primitive cubic and seventh
lines. A commutator of arithmetic elements has a unit coordinate
on it. The remaining relative obstruction can therefore be studied
on $D_{10}$.

The odd metabelian coordinates admit a uniform continuation. For
$g\in\ker\lambda$, put $a_j(g)=[S^jT]f_g$ and
\[
 c_{2r+1}(g)=[u^{2r}]\sum_{j\ge1}a_j(g)(e^u-1)^j
 \qquad(r\ge1).
\]
These rational two-adic coordinates are the logarithmic boundary
coefficients of the dyadic beta multiplier.

\begin{theorem}\label{thm:main-uniform-dyadic-normalization}
Put $C_0=D=N_{\mathrm{rad}}$ and
\[
 C_r=D\cap\bigcap_{j=1}^r\ker c_{2j+1},\qquad
 C_\infty=\bigcap_{r\ge0}C_r.
\]
For every $r\ge1$, the coefficient
$m_{2r+1}=[S^{2r}T]f$ is additive on $C_{r-1}$, has conjugation
weight $\lambda_2^{2r+1}$, and satisfies
\[
 m_{2r+1}(C_{r-1})=m_{2r+1}(C_{r-1}\cap H)
                     =2^{2r+1}\Z_2,\qquad
 \ker(m_{2r+1}|_{C_{r-1}})=C_r.
\]
All these subgroups are closed and normal, and
\[
 G=C_rH=C_\infty H,\qquad
 B\simeq C_\infty/
       \overline{[G,C_\infty](C_\infty\cap H)}.
\]
The subgroup $C_\infty$ is exactly the kernel in $D$ of the dyadic
beta multiplier. The kernel of restriction from $X_2$ to
$H^1(C_\infty,\F_2)$ is the rational Kummer span together with
$\F_2\eta_3$.
Finally, $E=C_\infty\cap D_{10}$ has infinite index in $D_{10}$,
and
\[
 G=EH,\qquad B\simeq E/\overline{[G,E](E\cap H)}.
\]
Its elements have no pro-$2$ Magnus term below degree eleven and
have trivial dyadic beta multiplier. The kernel of restriction
from $X_2$ to $H^1(E,\F_2)$ is again exactly the rational Kummer
span together with $\F_2\eta_3$.
The mixed coordinate $b_{11}=[S^6TST^3]f$ on $E$ is additive,
has conjugation weight $\lambda_2^{11}$, and has full and arithmetic
image $\Z_2$. Its kernel $E'$ is a proper closed normal subgroup,
$G=E'H$, and
\[
 B\simeq E'/\overline{[G,E'](E'\cap H)}.
\]
The kernel of restriction from $X_2$ to $H^1(E',\F_2)$ is the
same rational Kummer span together with $\F_2\eta_3$.
\end{theorem}

Section~\ref{sec:uniform-dyadic-normalization} proves this by combining
integer-valued polynomials with arithmetic cyclotomic differences.
Compactness passes from the finite stages to their intersection.
This identifies a smaller group on which the original relative
character question is decided; the character question itself remains
open. Ichimura and Kaneko characterize the arithmetic beta image at
odd primes in terms of Vandiver's conjecture
\cite[Section~1.2, pp.~314--315, Theorem~1]{IchimuraKaneko1989}.
Their theorem assumes an odd prime. The dyadic reduction here uses
the normalized duplication identity of Furusho and gives its exact
successive coefficient lattices directly.

There is a second arithmetic normalization using the two elliptic
points of the modular curve.  Let $G_1=\ker\lambda\subset\GT$.
Use the signed braid representation $q$, the unique pure frames
$A_g,\beta_g$ at the cubic and harmonic elliptic points, and the
normalized upper-triangular intertwiner $D_g$ defined in
Section~\ref{sec:adelic-cm-frames}.  Put
\[
 T_g=q(A_g)^{-1}D_g,\qquad W_g=q(\beta_g)^{-1}D_g
                       \quad(g\in G_1).
\]

\begin{theorem}\label{thm:main-cm-normalization}
The corrected pair $(T,W)$ has the same full adelic image on $G_1$
and on $G_{\Q^{\ab}}$.  Its kernel $C_{\mathrm{CM}}$ in $G_1$
is closed and normal in $\GT$.  If
$K_3=\Q(\sqrt{-3})$ and $K_4=\Q(i)$, then
\[
 \GT=C_{\mathrm{CM}}\GQ,\qquad
 C_{\mathrm{CM}}\cap\GQ=G_{K_3^{\ab}K_4^{\ab}},\qquad
 \GT/C_{\mathrm{CM}}\simeq
       \Gal(K_3^{\ab}K_4^{\ab}/\Q).
\]
Every continuous abelian character of $\GT$ which is trivial on
$\GQ$ and whose restriction to $G_1$ factors through $(T,W)$ is zero.
\end{theorem}

The complete adelic frame pair therefore contributes no further
relative abelian character.  The proof identifies its Galois image
by complex multiplication and the exact signed Tate-module comparison
on the Legendre family.  The two CM division fields intersect in
$\Q^{\ab}(3^{1/4})$.  The common field of their compositum and the rational-radical field is
\[
 L_{\mathrm{rad}}\cap K_3^{\ab}K_4^{\ab}
 =\Q^{\ab}\bigl(q^{1/12}:q\text{ prime},\,2^{1/8},3^{1/8}\bigr).
\]
Theorem~\ref{thm:radical-cm-intersection} proves this formula by an
explicit derived-subgroup calculation and division points.  Combining
the two normalizations is therefore equivalent to a specified crossed
cocycle with values in
$(\Z/12)^2\times\prod_{q\ne2,3}\Z/6$ being zero.
The equivalence is Proposition~\ref{prop:radical-cm-joint-gate}.
The entire coordinate belonging to $2$ vanishes: on every multiplier-one
gentle element the cubic and eighth-root actions of the frames both agree
with $\rho_2/2$, combined modulo twelve
(Corollary~\ref{cor:complete-two-compatibility}).
The remaining compatibility conditions have explicit finite forms.
On $D=N_{\mathrm{rad}}$, put
\[
 W_{g,3}\bmod3=(-I)^{w_3(g)},\qquad
 T_{g,2}\bmod8=(1+4\tau_3(g))I.
\]
Proposition~\ref{prop:lr-exceptional-three-character} proves that
$\chi_3=w_3+\tau_3$ is $\GT$-invariant and trivial on $D\cap H$;
it therefore defines an element of $\Hom(B,\F_2)$. Its vanishing
is exactly the exceptional eighth-root-of-three compatibility.
For each prime $\ell\ge5$, the other tests are
\[
 B_{g,\ell}(\omega,\omega)^{(\ell-\chi_{-3}(\ell))/3}=1,
 \qquad
 W_{g,\ell}^{(\ell-\chi_{-4}(\ell))/2}=1
\]
in the corresponding residue norm-one tori, for every $g\in D$.
Here $\omega$ is a primitive cube root, and
$\chi_{-3},\chi_{-4}$ are the two quadratic Dirichlet characters.
The wild cubic condition at three already holds.
Section~\ref{sec:lr-relative-gate} proves these equivalences, including
the integral cubic-frame comparison. The character $\chi_3$ is zero
by the two comparison theorems below. An explicit homomorphism from a
192-element group modulo eight to $\mathrm{SL}_2(\mathbf F_3)$
compares this character with two finite harmonic and cubic frame
defects (Proposition~\ref{prop:pair-of-gates}). Each defect is $\GT$-invariant and extends uniquely to a
quadratic character of $\GT$ trivial on $H$, and the two
extensions sum to $\chi_3$
(Proposition~\ref{prop:two-relative-extensions}).
The harmonic defect is zero on full GT by the following stronger
nonabelian identity.

\begin{theorem}[Harmonic comparison after Furusho]
\label{thm:main-harmonic-dyadic}
For every $g=(1,f_g)\in\GT$ with $\rho_2(g)=0$, the pure
harmonic frame satisfies
\[
 \beta_g(u^2,v^2)^{-1}=f_g(u^2,uvu)
 \quad\text{in }\widehat B_3.
\]
Consequently $w_3=\eta$ on $D$, the relative character
$\epsilon_h$ is zero, and $\chi_3=\epsilon_c=\tau_3+\eta$.
\end{theorem}

Section~\ref{sec:harmonic-dyadic-comparison} proves the identity
from Furusho's full dyadic word identities, the established
relation $(\mathrm{IV})$, and Enriquez's theorem on $(K_2)$. 

\begin{theorem}[Cubic comparison from genus-zero rigidity]
\label{thm:main-cubic-normalized}
For every $g=(1,f_g)\in\GT$ with $\rho_3^I(g)=0$, the pure
cubic frame satisfies
\[
 A_g(u^2,v^2)^{-1}=f_g(u^2,uv)
 \quad\text{in }\widehat B_3.
\]
Consequently $\epsilon_c=0$ and $\chi_3=0$ on full GT.
The exceptional eighth-root-of-three compatibility is therefore
satisfied. Joint radical--CM compatibility is equivalent to the
remaining tame cubic and harmonic tests displayed above.
\end{theorem}

Section~\ref{sec:normalized-cubic-identity} proves the comparison.
Boggi--Zalesskii's finite-subgroup and centralizer theorems,
Asada's genus-zero faithfulness and
Hoshi--Minamide--Mochizuki's description of the automorphisms
of the five-point group make the cyclic-cover quotient available
with its full profinite topology. The exact $(K_3)$ relation
identifies its GT parameter. Forgetting the two branch points
then gives the cubic comparison with all inner corrections
accounted for. The tame tests are not decided here.

A different finite homology representation gives a quadratic character
in the remaining relative kernel: it uses the ten character lines
$(\xi,\omega)$ with $\xi$ a nontrivial eleventh root of unity and
$\omega$ a primitive cube root, in characteristic seventeen, in the homology of the genus-ten curve
$u^{11}+v^3=1$.
Corollary~\ref{cor:concrete-radical-anderson-character} gives its
determinant formula. Its vanishing on the Galois subgroup is proved;
whether it is nonzero on $\GT$ remains open.

\subsection*{Credit and scope}
The Kummer cocycles used here are Ihara's, and their full distribution relations are Enriquez's Theorem 8.1. The Galois description of rational radicals is classical Kummer theory in the explicit form of Moree and Stevenhagen's Theorem 2.5. The group-theoretic obstruction is the continuous Hochschild--Serre transgression. The structural use of normalized elements starts from \cite[Theorem 2]{H17}. The paper writes out the quotient, restriction, cyclic-cover parity and lattice arguments in the form needed for the character question. The kernel of restriction remains unresolved, as does the invariant character space for odd primes.

Lochak's Theorem 5.10 identifies the inertia-preserving automorphisms of generic procongruence mapping class groups with his group of compatible pairs \cite[Definition 5.4, Theorem 5.10]{Lochak2025}. Relation (IV) identifies every such pair with a diagonal pair, and the resulting group is $\GT$ \cite[Corollary 22]{H17}. This answers his closing questions about these groups. It does not by itself identify separately defined tower actions, or prove that preservation of nonseparating twist inertia implies preservation of every twist inertia.

\section{Proof overview}
The arithmetic quotient is obtained by combining the cocycle with the cyclotomic character. A positive radical of $2$ lying in an abelian extension can have denominator at most two, so the Kummer translation subgroup over the cyclotomic field is exactly $2\Zh$. This determines the quotient, including its parity condition.

The extension argument uses the factorization $G=NH$: an invariant character of $N$ extends by zero on $H$ exactly when it vanishes on $N\cap H$. To identify the full abelianization one also uses Kronecker--Weber, which makes the image of $H^{\ab}$ a canonical section of the multiplier.

The dimension argument uses a different feature. A positive peripheral automorphism acts on each cyclic cover of the three-punctured sphere. In its class-two quotient, preservation of the boundary word makes the coefficients of the conjugator symmetric under reversal of the cyclic indices. Passing from a cover of degree $2d$ to one of degree $d$ proves the parity of $\rho_d^I$. For each odd rational prime, adjoining its fourth root to $K_2$ gives a nontrivial quadratic extension, and these extensions are independent. The divided cocycles give the desired lifts. Conjugation acts by the odd scalar $\lambda$, so their reductions modulo two are invariant. Compact metrizability gives the matching countable upper bound. This identifies an infinite subspace on which restriction is injective; it does not exhaust the invariant character space.

The simultaneous normalization uses the same arithmetic idea for all rational primes. A positive radical in an abelian extension of $\Q$ has at most two conjugates, since all its conjugates must be real. Finite-character separation then makes every joint Kummer image exactly the even lattice. Compactness prescribes all coordinates at once, so their common kernel still multiplies with $G_\Q$ to give $\GT$. Enriquez's distribution theorem identifies the cocycles at composite indices with the sums of those at prime indices. This shows that the entire rational Kummer character space restricts injectively; the relative obstruction is unchanged on their common kernel.

The cubic refinement uses an integral feature invisible after tensoring
with $\Q_2$.  The two- and three-cycle equations determine enough
coefficients through degree five to make a signed modular matrix entry
congruent to $80m_3$ modulo $128$.  Congruence triviality forces divisibility
by eight.  Choose a cyclotomic Galois element whose Kummer coordinate of two
is two and whose other rational Kummer coordinates vanish.  Its cubic
coefficient is odd.  Conjugation by a Galois element with $2$-adic
multiplier three, followed by the inverse multiplier power, kills every
rational Kummer coordinate and multiplies the cubic coefficient by $24$.  This proves both independence and the exact integral image.

The weight-five refinement uses the same matrix representation, now
modulo $2048$.  Explicit integral coefficient identities in degrees
five, seven and nine force the first weight-five coefficient to be a
multiple of $32$.  The Alexander cofactor supplies an exact weight-five
homomorphism, and two multiplier-difference operations on a Galois
element kill the rational and cubic coordinates while producing a
generator of that lattice.  The proof prints the finite recursions and
coefficient certificates needed for its divisibility argument.

At weight seven, the leading pentagon first restricts the integral
Hall vector to a primitive line. Two coefficients of the full Artin
pentagon, together with the normalized signed matrix identity, then
force divisibility by $128$. A third multiplier-difference operation
on the same arithmetic starting element produces a value of valuation
seven. Arithmetic commutators also show why the divided weight-five
and weight-seven characters cannot extend to invariant characters on
the earlier normalized groups.

At weight eight the normalized leading space is a primitive integral
line generated by an Ihara bracket. A commutator of two multiplier-one
Galois elements kills every preceding additive coordinate. Comparing
one element with a cyclotomic conjugate gives an odd value on the new
line, so its image is the full two-adic lattice. The same commutator
also obstructs extending its quadratic reduction to the original group.
The exact weight-seven divisibility allows that reduction to extend to
$D_5$. A cyclotomic difference and a complex-conjugation norm of arithmetic
elements give independent evaluations of this extension and $\eta_7$,
ruling out every nonzero combination as the restriction of an invariant
character on $D_3$.

At weight nine, an integral multiple of the logarithmic beta coefficient
gives a homomorphism before normalization. Four cyclotomic differences
and correction by the weight-eight commutator produce an arithmetic
value of valuation nine. An integral measure obtained from the beta
boundary series has zero lower even moments and obeys the dyadic
distribution identity. Integrating an integer-valued polynomial gives
the matching upper bound, uniformly for every odd coefficient whose
predecessors vanish. The full and arithmetic
images are therefore equal, giving the normalization $D_9$. The
weight-eight extension controls the correction exponent and prevents
the divided ninth character from extending invariantly to $D_5$.
At weight ten a primitive Ihara bracket has an odd mixed coefficient.
A commutator with a weight-seven arithmetic element realizes that
unit, giving the further kernel $D_{10}$.

For the finite harmonic comparison, Furusho's four-variable
dyadic identity specializes to a factorization of $f$ as a word
times the inverse of its interchange. The swap involution has
trivial fixed subgroup in the relevant free profinite group, so
this factor is exactly the pure harmonic frame. Specialization
to the signed braid group gives the identity directly, since the
central argument has total exponent zero. Its finite reduction removes the harmonic comparison
character. For the cubic comparison, the order-three centralizer
on an eleven-point sphere has the completed five-point quotient.
Genus-zero rigidity gives a GT parameter for that quotient.
Collision multitwists determine its multiplier and labelling, and
$(K_3)$ identifies it with the original parameter.
This extends the scope of the normalized cubic relation proved by
Schneps for her special-homomorphism subgroup
\cite[Theorem~3.4 and equation~(3.10), p.~667]{Schneps2006Special}. The other
filling then identifies the pure cubic conjugator. Exponent sums
remove the possible central discrepancy in the signed braid group.
Together the two comparisons remove the exceptional character at three.

For the CM normalization, correcting each pure frame by the same
upper-triangular intertwiner produces homomorphisms into the norm-one
tori of $\Q(\sqrt{-3})$ and $\Q(i)$.  Ordinary paths in the Legendre
family identify these signed integral matrices with the Tate actions
at its two CM fibres.  Class field theory and explicit division points
compute the two images and their common quadratic subfield.  The frame
sign equality on GT then forces its joint image to equal the Galois image.

\section{The arithmetic quotient}
\label{sec:arithmetic-quotient}

We first identify the quotient through which the multiplier and the Kummer
cocycle of~$2$ factor. Put
\[
 G=\widehat{\mathrm{GT}},\quad H=G_{\mathbb Q},\quad
 N=\Lambda _0,\quad
 K_2=\mathbb Q(\mu _\infty,2^{1/\infty}),\quad K=G_{K_2},
 \qquad U=\widehat{\mathbb Z}^{\times}.
\]
The embedding of $H$ in $G$ is understood throughout. We use the decomposition
$G=NH$ of \cite[Theorem~2]{H17}, the cocycle identity
\[
 \rho _2(gg')=\rho _2(g)+\lambda(g)\rho _2(g'),
\]
and the fact that $\rho _2$ restricts to the Kummer cocycle of the positive
coherent roots of~$2$; see \cite[Definition~5.1]{NS2000} and
\cite[Lemma~7]{H17}. Define the additive quadratic character
$\epsilon _8:U\to\mathbb F _2$ by
\[
 \epsilon _8(a)=
 \begin{cases}
 0,&a\equiv 1,7\pmod 8,\\
 1,&a\equiv 3,5\pmod 8.
 \end{cases}
\]

\begin{proposition}\label{prop:arithmetic-affine}
The subgroup $N$ is closed and normal in $G$, and $N\cap H=K$.
The map $(\lambda,\rho _2)$ identifies $G/N$ with
\[
 A=\{(a,r)\in U\ltimes\widehat{\mathbb Z}:
                 r\bmod 2=\epsilon _8(a)\},
 \qquad
 (a,r)(b,s)=(ab,r+as).
\]
This identifies $G/N$ canonically with $\operatorname{Gal}(K_2/\mathbb Q)$.
In particular $\rho _2\bmod 2=\epsilon _8\circ\lambda$ on $G$.
\end{proposition}

\begin{proof}
The cocycle identity makes $(\lambda,\rho _2)$ a continuous homomorphism to
the affine semidirect product. Its kernel is $N$, so $N$ is closed and normal.
On $H$ the simultaneous kernel fixes every root of unity and every coherent
root of~$2$, and is consequently $K$.

We need the precise Kummer image over $\mathbb Q^{\mathrm{ab}}
=\mathbb Q(\mu _\infty)$. It is a closed subgroup $J$ of
$\widehat{\mathbb Z}$. For a positive integer $d$, the inclusion
$J\subset d\widehat{\mathbb Z}$ is equivalent to
$2^{1/d}\in\mathbb Q^{\mathrm{ab}}$. This happens exactly for $d=1,2$.
Indeed $X^d-2$ is Eisenstein; if $d>2$, the field generated by its positive
real root is real and misses a nonreal conjugate, so is not normal over
$\mathbb Q$. Every finite subfield of an abelian Galois extension is normal.
Closed subgroups of $\widehat{\mathbb Z}$ have the form
$\prod_p p^{e_p}\mathbb Z_p$ with $e_p\in\mathbb N\cup\{\infty\}$,
and the inclusions in $d\widehat{\mathbb Z}$ determine all the $e_p$.
Thus $J=2\widehat{\mathbb Z}$.

Given $g\in G$, choose $h\in H$ with $\lambda(h)=\lambda(g)$.
The element $gh^{-1}$ has multiplier one, and
\cite[Lemma~7(c)]{H17} gives $\rho _2(gh^{-1})\in2\widehat{\mathbb Z}$.
The cocycle identity therefore gives
$\rho _2(g)\equiv\rho _2(h)\pmod2$. On $H$ this parity is the action
on $\sqrt2=\zeta _8+\zeta _8^{-1}$, namely $\epsilon _8(\lambda(h))$.
For each fixed multiplier, the Kummer values on $H$ fill the corresponding
coset of $2\widehat{\mathbb Z}$, by the preceding calculation. Thus the
image of $H$ is exactly $A$. The decomposition $G=NH$ shows that $G$ has
the same image. Finally the continuous bijection $H/K\to G/N$ is a
homeomorphism, since its source is compact and its target Hausdorff.
\end{proof}

\begin{proposition}\label{prop:arithmetic-commutators}
Let $T=\{(1,2t):t\in\widehat{\mathbb Z}\}\subset A$.
Then
\[
 \overline{[A,A]}=T,\qquad
 \overline{[A,T]}=\{(1,4t):t\in\widehat{\mathbb Z}\}.
\]
Every finite $p$-group quotient of $A$ is abelian when $p$ is odd.
Every finite quotient of $A$ which is a $2$-group of nilpotency class at
most two has commutator subgroup of order at most two.
\end{proposition}

\begin{proof}
The quotient $A/T=U$ is abelian. Choose $a\in U$ whose $2$-adic
component is~$3$ and whose odd-primary components are~$1$.
Then $(a,1)\in A$, and its commutator with $c=(-1,0)$ is $(1,2)$.
The closed subgroup generated by this translation is $T$, proving the first
equality. Conjugation by $(a,r)$ multiplies a translation coordinate by~$a$.
Thus $[A,T]$ consists of coordinates generated by $(a-1)2t$.
They are all divisible by~$4$ at the prime~$2$, and $a=-1$ produces
$(1,-4)$, whose closed cyclic subgroup is $4\widehat{\mathbb Z}$.
This proves the second equality.

The element $c$ has order two and acts on $T$ by inversion. It maps to the
identity in an odd-primary finite group, so the image of $T$ in such a
group has exponent dividing two and is trivial. The quotient is therefore
abelian. A quotient of nilpotency class at most two kills $[A,T]$, and its
commutator subgroup is consequently a quotient of
$T/\overline{[A,T]}\simeq C_2$.
\end{proof}

\section{Characters and the cyclotomic abelianization}
\label{sec:characters-abelianization}

The decomposition $G=NH$ separates the characters obtained from the multiplier
from the characters that could distinguish $G$ from $H$. All characters
in this section are continuous, and all abelianizations are profinite.

\begin{lemma}\label{lem:character-gluing}
Let $G=NH$ be a profinite group with $N$ closed and normal and $H$ closed,
and put $K=N\cap H$. Let $M$ be a finite abelian group with trivial action.
An invariant character $\psi:N\to M$ and a character $\eta:H\to M$
extend to a character of $G$ if and only if $\psi|_K=\eta|_K$.
The extension is unique and is given by
\[
 \chi(nh)=\psi(n)+\eta(h).
\]
\end{lemma}

\begin{proof}
Necessity follows by restriction. If $nh=n'h'$, then
$n'^{-1}n=h'h^{-1}\in K$, so the proposed two values coincide.
The identity
\[
 (n_1h_1)(n_2h_2)=n_1(h_1n_2h_1^{-1})h_1h_2
\]
and invariance of $\psi$ prove additivity. The multiplication map
$N\times H\to G$ is a continuous surjection from a compact space to a
Hausdorff space and hence is a quotient map. The displayed function on
$N\times H$ is continuous and constant on its fibres, proving continuity
on $G$. Uniqueness follows from $G=NH$.
\end{proof}

In our arithmetic situation define
\[
 C_N=\overline{\langle[g,n],k:g\in G,\ n\in N,\ k\in K\rangle},
 \qquad B=N/C_N.
\]
The group $B$ is abelian, since $C_N$ contains $[N,N]$; it is also the
largest quotient of $N$ on which conjugation by $G$ and the image of $K$
are both trivial.

\begin{theorem}\label{thm:abelian-splitting}
There is a canonical isomorphism of profinite abelian groups
\[
 G^{\mathrm{ab}}\simeq U\times B
\]
under which the multiplier is the first projection. Moreover
\[
 C_N=N\cap\overline{[G,G]},\qquad
 B\simeq\bigl(G/\overline{\langle\!\langle H\rangle\!\rangle}\bigr)^{\mathrm{ab}}.
\]
For every prime $p$, put
\[
 X_p=H^1(N,\mathbb F_p)^G,
 \qquad r_p:X_p\longrightarrow H^1(K,\mathbb F_p).
\]
Restriction gives canonical identifications
\begin{align*}
 \ker r_p&\simeq\operatorname{Hom}_{\mathrm{cts}}(B,\mathbb F_p),\\
 H^1(G,\mathbb F_p)&\simeq
 \lambda^*H^1(U,\mathbb F_p)\oplus\ker r_p.
\end{align*}
Consequently $\ker r_p=0$ if and only if every continuous
$\mathbb F_p$-character of $G$ factors through the multiplier.
The multiplier identifies $G^{\mathrm{ab}}$ with $U$ if and only if
$N\subset\overline{[G,G]}$, and this holds if and only if $\ker r_p=0$
for every prime~$p$.
\end{theorem}

\begin{proof}
Kronecker--Weber identifies $H^{\mathrm{ab}}$ with $U$ through the
cyclotomic character. Since $K$ fixes $\mathbb Q^{\mathrm{ab}}$, it lies
in $\overline{[H,H]}$ and hence in $\overline{[G,G]}$.
The map $N\to G^{\mathrm{ab}}$ therefore factors through $B$.
Conversely the rule
\[
 nh\longmapsto(\lambda(h),n\bmod C_N)
\]
is well-defined: changing the decomposition changes $n$ by an element of
$K$, which is zero in $B$. It is continuous by the compact quotient
argument of Lemma~\ref{lem:character-gluing}. It is a homomorphism because
conjugation by $H$ is trivial on $B$. The maps
$H^{\mathrm{ab}}=U\to G^{\mathrm{ab}}$ and $B\to G^{\mathrm{ab}}$
give the inverse of the resulting homomorphism
$G^{\mathrm{ab}}\to U\times B$: both composites are the identity on the
images of $N$ and $H$, which generate the relevant groups.
This proves the splitting and the formula for $C_N$.
Killing $H$ and then abelianizing imposes exactly the relations defining
$C_N$, which proves the second description of~$B$.

Apply Lemma~\ref{lem:character-gluing} with $\eta=0$ to identify the
restriction kernel. Every character of $H$ is the pullback of a unique
character of $U$ and vanishes on $K$. Subtracting that pullback from a
character of $G$ gives the asserted direct sum, or it follows directly
from the abelian splitting. Finally a nonzero profinite abelian group has
a nonzero finite abelian quotient, and that quotient has a quotient
$C_p$ for some prime $p$. Thus $B=0$ if and only if it has no character
to $\mathbb F_p$ for any prime~$p$.
\end{proof}

The expression ``the abelianization is $U$'' in this theorem specifies
the multiplier map. The kernel assertion does not follow from an
unmarked abstract isomorphism of profinite abelian groups.

\subsection{The transgression describes the image}

The kernel of restriction and the image of restriction involve different
obstructions. The next proposition states both in one exact sequence.
Write
\[
 Y_p=H^1(K,\mathbb F_p)^H.
\]
Restriction $r_p$ lands in $Y_p$, and $A=G/N=H/K$ acts on $K$ through
conjugation by $H$.

\begin{proposition}\label{prop:transgression-comparison}
The transgression
$d_H:Y_p\to H^2(A,\mathbb F_p)$ is injective and identifies $Y_p$ with
\[
 \ker\bigl(H^2(A,\mathbb F_p)\longrightarrow H^2(H,\mathbb F_p)\bigr).
\]
There is an exact sequence
\[
 0\longrightarrow\operatorname{Hom}_{\mathrm{cts}}(B,\mathbb F_p)
 \longrightarrow X_p
 \xrightarrow{\ d_Hr_p\ }
 \ker\bigl(H^2(A,\mathbb F_p)\longrightarrow H^2(G,\mathbb F_p)\bigr)
 \longrightarrow0.
\]
Thus, under the injection $d_H$, the image of restriction is precisely
the inflation kernel for~$G$.
\end{proposition}

\begin{proof}
For completeness, we describe the central-extension form of the
five-term argument. Given a conjugation-invariant character
$\psi:N\to\mathbb F_p$, its pushout extension is
\[
 E_\psi=(\mathbb F_p\times G)/
       \{(-\psi(n),n):n\in N\}.
\]
The subgroup being divided out is closed and normal, by the invariance
of $\psi$. The quotient is a central extension of $A$ by $\mathbb F_p$;
its class is $d_G\psi$, with a consistent choice of the transgression
sign. A splitting of this extension is equivalent to a character of
$G$ extending $\psi$: in one direction subtract the splitting from the
natural lift $g\mapsto[0,g]$, and in the other use the character to
define a splitting independent of the lift from~$A$.

More generally, a central extension of $A$ becomes trivial after
inflation to $G$ if and only if it admits a continuous homomorphism
$G\to E$ lifting $G\to A$. Restricting such a lift to $N$ gives an
invariant character $\psi$, and the map
$[z,g]\mapsto z\,\widetilde g$ identifies its pushout with $E$.
This proves exactness at $X_p$ and identifies the image of $d_G$ with
the indicated inflation kernel. The construction respects sums of
characters and Baer sums of central extensions, so these are linear
maps.

Apply the same argument to $H$ and $K$. A character of $K$ in the
kernel of $d_H$ extends to $H$. Every character of $H$ factors through
$U$ and therefore vanishes on $K$, so $d_H$ is injective. The preceding
lifting argument gives its image. Finally pulling the pushout extension
back from $G/N$ along $H/K$ is the same construction with $\psi|_K$.
Thus $d_G=d_Hr_p$. Theorem~\ref{thm:abelian-splitting} identifies the
kernel of $d_G$, proving the displayed exact sequence.
\end{proof}

\section{The arithmetic invariant characters}
\label{sec:arithmetic-target}

The target allowed by Galois invariance is already infinite dimensional.
Its size does not determine which characters extend across $K\subset N$.

\begin{theorem}\label{thm:arithmetic-target-dimension}
For every prime $p$,
\[
 \dim_{\mathbb F_p}H^1(G_{K_2},\mathbb F_p)^{G_{\mathbb Q}}=\aleph_0.
\]
For $p=2$, an explicit independent family is given, for the odd rational
primes $q$, by the quadratic extensions $K_2(q^{1/4})/K_2$.
Equivalently these are the Kummer classes of $\sqrt q$ in
$K_2^\times/(K_2^\times)^2$.
\end{theorem}

\begin{proof}
For an odd prime $q$, the square of the positive fourth root $q^{1/4}$
belongs to $\mathbb Q^{\mathrm{ab}}\subset K_2$. The extension in the
statement is therefore at most quadratic. It is normal over $\mathbb Q$,
because its other conjugates are $i^j q^{1/4}$ and $i\in K_2$.
Its character, if nonzero, is invariant under $H$, because conjugation
acts trivially on a group of order two.

Suppose that a nonempty finite product of these characters is zero.
Let $a$ be the product of the corresponding distinct odd primes. Then
$a^{1/4}\in K_2$. The extension
$\mathbb Q^{\mathrm{ab}}(a^{1/4})/\mathbb Q^{\mathrm{ab}}$ is a
nontrivial quadratic extension. To see nontriviality, $X^4-a$ is
Eisenstein at any prime dividing $a$, and its positive real root
generates a nonnormal real field, which cannot lie in
$\mathbb Q^{\mathrm{ab}}$. By Proposition~\ref{prop:arithmetic-affine},
$\operatorname{Gal}(K_2/\mathbb Q^{\mathrm{ab}})=2\widehat{\mathbb Z}$
has exactly one subgroup of index two. Its quadratic subfield is
$\mathbb Q^{\mathrm{ab}}(2^{1/4})$. The two quadratic fields are
therefore equal, so the product $a^{1/4}2^{1/4}=(2a)^{1/4}$ lies in
$\mathbb Q^{\mathrm{ab}}$: their common nontrivial automorphism negates
both factors. But $X^4-2a$ is Eisenstein at~$2$, and its positive root
again generates a real nonnormal field. This contradiction proves
nontriviality and independence of the whole quadratic family.

Now let $p$ be odd and let $m\geq1$. Put
$P=\operatorname{UT}_3(\mathbb F_p)^m$, a direct product of upper
unitriangular matrix groups. Its centre and commutator subgroup agree
and are isomorphic to $(C_p)^m$, and its abelianization is
$(C_p)^{2m}$. By \v{S}afarevi\v{c}'s theorem, in the precise form
\cite[Theorem~1, p.~1]{SchmidtWingberg}, there is a finite Galois
extension $E/\mathbb Q$ with group~$P$.
Let $L=E\cap K_2$. Its Galois group over $\mathbb Q$ is a finite
$p$-group quotient of $A$, and is abelian by
Proposition~\ref{prop:arithmetic-commutators}. Hence
$L\subset E^{[P,P]}$. The reverse inclusion follows from
$E^{[P,P]}\subset\mathbb Q^{\mathrm{ab}}\subset K_2$. Thus
\[
 \operatorname{Gal}(EK_2/K_2)=\operatorname{Gal}(E/L)=[P,P]
                  \simeq(C_p)^m.
\]
This subgroup is central in $\operatorname{Gal}(EK_2/\mathbb Q)$,
because $[P,P]$ is central in~$P$. Its characters give $m$ independent
members of $Y_p$.

Finally the algebraic field $\overline{\mathbb Q}$ is countable.
Every cyclic degree-$p$ extension of $K_2$ inside it is generated by
one element, so there are at most countably many such extensions.
Each has finitely many characters. This gives the countable upper
bound, while the preceding arguments give arbitrarily large finite
independent families for every~$p$.
\end{proof}

\section{Infinitely many invariant quadratic characters of the normalized group}
\label{sec:quadratic-dimension}

Ihara's Kummer cocycles give explicit invariant quadratic characters and
determine the dimension of their restriction image. They do not settle
the kernel of restriction. We write
$\rho_d^I$ for these cocycles to distinguish their construction from the
frame cocycle used to define~$N$.

Let $d$ be a positive integer and put
\[
 F=\widehat F_2=\langle x,y\rangle,\qquad
 H_d=\ker(F\longrightarrow\mathbb Z/d\mathbb Z),
 \quad x\longmapsto1,\quad y\longmapsto0.
\]
The Schreier basis of $H_d$ consists of $x^d$ and
$y_j=x^jyx^{-j}$ for $0\leq j<d$. Let
$\ell_d:H_d\to\widehat{\mathbb Z}$ send $y_0$ to~$1$ and the other
basis elements to~$0$. For $g=(\lambda,f)\in\GTz$ define
\[
 \rho_d^I(g)=-\ell_d(f).
\]
This is Ihara's cocycle $-\Psi_d^{(0)}$; the construction and its
identification with the Kummer cocycle of $d$ on $G_{\mathbb Q}$ are
recalled in \cite[\S2.2]{Furusho2003}, following Ihara's theory
\cite{Ihara1999}. We use that arithmetic identification as a published
input. The following direct verification fixes the cocycle convention.

\begin{lemma}\label{lem:ihara-cocycle}
For $g,g'\in\GTz$,
\[
 \rho_d^I(gg')=\rho_d^I(g)+\lambda(g)\rho_d^I(g').
\]
If $n\in N$, then
\[
 \rho_d^I(gng^{-1})=\lambda(g)\rho_d^I(n).
\]
\end{lemma}

\begin{proof}
The automorphism $\Phi_g:x\mapsto x^\lambda$,
$y\mapsto f^{-1}y^\lambda f$ preserves $H_d$. On its abelianization,
\[
 [x^d]\longmapsto\lambda[x^d],\qquad
 [y_j]\longmapsto\lambda[y_{\lambda j\bmod d}].
\]
The two conjugating contributions of $f$ cancel. Replacing an exponent
$\lambda j$ by its residue modulo $d$ only conjugates by a power of
$x^d$, which is also invisible in the abelianization. Since $\lambda$
is a unit modulo $d$, index zero has no preimage other than zero.
Consequently $\ell_d\Phi_g=\lambda\ell_d$.
With the convention $\Phi_{gg'}=\Phi_g\Phi_{g'}$, the second coordinate
is $f_{gg'}=f_g\Phi_g(f_{g'})$. Applying $-\ell_d$ proves the first
formula. Applying it twice to $gng^{-1}$, and using $\lambda(n)=1$,
cancels the terms involving $\rho_d^I(g)$ and proves the second.
\end{proof}

\begin{lemma}\label{lem:cyclic-cover-symmetry}
Let $F=\widehat F_2=\langle x,y\rangle$. Suppose a continuous
automorphism $\phi$ satisfies
\[
 \phi(x)=x,\qquad \phi(y)=f^{-1}yf,\qquad
 f\in\overline{[F,F]},
\]
and $\phi(xy)$ is conjugate to $xy$. For a positive integer $n$, use
the free basis
\[
 X=x^n,\qquad Y_j=x^jyx^{-j}\quad(0\leq j<n)
\]
of $H_n$. If $[f]=a[X]+\sum_{j=0}^{n-1}c_j[Y_j]$ in
$H_n^{\mathrm{ab}}$, then
\[
 a=0,\qquad \sum_{j=0}^{n-1}c_j=0,\qquad
 c_j=c_{-j\bmod n}\quad(0\leq j<n).
\]
\end{lemma}

\begin{proof}
The map $H_n^{\mathrm{ab}}\to F^{\mathrm{ab}}$ takes $[X]$ to
$n[x]$ and every $[Y_j]$ to $[y]$. Hence $na=0$ and $\sum c_j=0$.
Multiplication by $n$ is injective on $\widehat{\mathbb Z}$, so $a=0$.
The automorphism preserves $H_n$, fixes $X$, and sends
\[
 Y_i\longmapsto f_i^{-1}Y_if_i,\qquad f_i=x^ifx^{-i}.
\]
Conjugation by $x$ cyclically permutes the $[Y_j]$ in the abelianization
and fixes $[X]$, so $[f_i]=\sum_jc_{j-i}[Y_j]$, with indices modulo $n$.

The boundary word
\[
 P=Y_0Y_1\cdots Y_{n-1}X=(yx)^n
\]
is a power of a conjugate of $xy$, and therefore
$\phi(P)=tPt^{-1}$ for some $t\in F$. Multiplying $t$ on the right
by a suitable integer power of $yx$ makes $t$ belong to $H_n$ without
changing this equality: $yx$ centralizes $P$ and generates $F/H_n$.

Work in the quotient of $H_n$ of nilpotency class two, and use the
commutator convention $[u,v]=uvu^{-1}v^{-1}$. The commutator pairing
identifies its degree-two group with the exterior square of
$H_n^{\mathrm{ab}}$ over $\widehat{\mathbb Z}$. In particular the
commutators of distinct basis elements are independent. This independence
can also be checked in finite unitriangular $3$-by-$3$ matrix groups:
map one chosen pair of generators to the two elementary generators and
all other generators to the identity, for every finite modulus.
Using the same symbols for degree-one vectors, the boundary identity gives
\[
 \sum_{i=0}^{n-1}[Y_i,[f_i]]
       =\left[[t],X+\sum_{i=0}^{n-1}Y_i\right].
\]
Here each outer bracket is the bilinear commutator pairing; $[f_i]$
and $[t]$ denote degree-one vectors. Write
$[t]=t_XX+\sum_i t_iY_i$. The coefficient of $[X,Y_i]$ on the left
is zero because every $[f_i]$ has zero $X$-coordinate. On the right it
is $t_X-t_i$. Thus every $t_i=t_X$, and the right side is zero.
For $i<j$, the left side has coefficient $c_{j-i}-c_{i-j}$ at
$[Y_i,Y_j]$. Taking $i=0$ gives the asserted symmetry.
\end{proof}

\begin{theorem}\label{thm:all-kummer-parity}
For every positive integer $d$, Ihara's integral Kummer cocycle is even
on the multiplier-one subgroup of the gentle group:
\[
 \rho_d^I\bigl(\ker(\lambda:\widehat{\mathrm{GT}}_0\to U)\bigr)
       \subset2\widehat{\mathbb Z}.
\]
On the whole gentle group it satisfies
\[
 \rho_d^I\bmod2=\epsilon_d\circ\lambda,
\]
where $\epsilon_d:U\to\mathbb F_2$ is the additive quadratic
Dirichlet character of $\mathbb Q(\sqrt d)$.
\end{theorem}

\begin{proof}
The two cycle equations at multiplier one give the positive peripheral
preservation in Lemma~\ref{lem:cyclic-cover-symmetry}. Apply that lemma
with $n=2d$. Pairing indices $j$ and $-j$ gives
\[
 0=\sum_{j=0}^{2d-1}c_j
   =c_0+c_d+2\sum_{j=1}^{d-1}c_j.
\]
Under $H_{2d}\subset H_d$, only $Y_0$ and $Y_d$ among these basis
elements have nonzero $y$-coordinate in $H_d^{\mathrm{ab}}$, and that
coordinate is one for each. Indeed $Y_{j+d}=x^dY_jx^{-d}$.
The element $x^{2d}$ has zero $y$-coordinate as well. Consequently
\[
 \rho_d^I=-\ell_d(f)=-(c_0+c_d)\in2\widehat{\mathbb Z}.
\]
For a general gentle element $g$, choose $h\in G_{\mathbb Q}$ with
$\lambda(h)=\lambda(g)$. The cocycle identity and the first assertion
give $\rho_d^I(g)\equiv\rho_d^I(h)\pmod2$. The latter is the action
of $h$ on $\sqrt d$, and hence is $\epsilon_d(\lambda(g))$.
\end{proof}

\begin{theorem}\label{thm:quadratic-dimension}
The vector space of invariant continuous quadratic characters of the
normalized subgroup has dimension
\[
 \dim_{\mathbb F_2}H^1(\Lambda_0,\mathbb F_2)^{\widehat{\mathrm{GT}}}
 =\aleph_0.
\]
For every positive integer $d$ the formula
\[
 \psi_d^N=\frac{\rho_d^I|_N}{2}\bmod2
\]
defines a continuous invariant quadratic character. For the odd rational
primes $q$, the characters $\psi_q^N$ are linearly independent and
restrict to the characters of $K_2(q^{1/4})/K_2$. Thus
\[
 \dim_{\mathbb F_2}\operatorname{im}r_2=\aleph_0.
\]
More precisely, for distinct odd primes $q_1,\ldots,q_m$, let
$R:N\to\mathbb Z_2^m$ be the vector of the $2$-adic components of
$\rho_{q_1}^I,\ldots,\rho_{q_m}^I$. Then
\[
 R(K)=R(N)=2\mathbb Z_2^m,
\]
and restriction is injective on the $m$-dimensional character space
obtained by pullback from this lattice.
\end{theorem}

\begin{proof}
Theorem~\ref{thm:all-kummer-parity} permits division by two on $N$.
Since $\lambda=1$ there, the result is a continuous homomorphism.
The conjugation formula of Lemma~\ref{lem:ihara-cocycle} and the fact
that every multiplier is odd make its reduction modulo two invariant.
On $K$, the character $\psi_q^N$ is the ordinary fourth-root Kummer
character because $K$ fixes $\sqrt q$. Independence follows from
Theorem~\ref{thm:arithmetic-target-dimension}.

The vector $R$ is a continuous homomorphism. On $K$ it is the vector of ordinary Kummer
characters of the $q_j$. All their values are even, since $K$ fixes
$\sqrt{q_j}$. Modulo four, half their values are precisely the
independent quadratic characters of
Theorem~\ref{thm:arithmetic-target-dimension}. Thus the image of $R(K)$
modulo $4\mathbb Z_2^m$ is all of
$2\mathbb Z_2^m/4\mathbb Z_2^m$.

A closed additive subgroup of $\mathbb Z_2^m$ is a
$\mathbb Z_2$-submodule, because integer multiplication extends by
continuity to $2$-adic scalar multiplication. A closed submodule
$M\subset2\mathbb Z_2^m$ whose reduction modulo four is onto equals
$2\mathbb Z_2^m$. Explicitly, approximate any vector modulo four by
an element of $M$, then correct the error modulo eight by twice an
element of $M$, and continue. The convergent sum belongs to $M$ by
closedness. Apply this to the compact image $R(K)$.

The parity theorem gives $R(N)\subset2\mathbb Z_2^m$, and the reverse
inclusion follows from $K\subset N$. Hence $R(N)=R(K)$.
Conjugation acts on this lattice by the odd scalar $\lambda(g)$, so
all its quadratic characters are invariant. Its character space has
dimension $m$, and restriction to $K$ is injective because $R(K)=R(N)$.

There is also a countable upper bound. The finitely generated free
profinite group $\widehat F_2$ has a countable base, as there are only
countably many finite quotients with specified generator images.
Thus $U\times\widehat F_2$, its closed subspace $G$, and the closed
subgroup $N$ are compact metrizable spaces. Such a space has countably
many clopen subsets: each is a finite union of members of a countable
base. A continuous function to $\mathbb F_2$ is determined by its
zero fibre. Consequently $H^1(N,\mathbb F_2)$, and hence its invariant
subspace, is countable. The lower bounds for arbitrary $m$ prove the
dimension formula.

The same upper bound applies to the restriction image, while the
independent quartic restrictions give its countably infinite dimension.
\end{proof}

The dimension computation identifies the explicit quartic subspace. To determine its complete extent, and to normalize all rational Kummer coordinates simultaneously, we next use the distribution theorem.

\section{The classical distribution relations and their arithmetic consequences}
\label{sec:kummer-distribution}

The multiplicative relations among Ihara's Kummer cocycles already hold
integrally on the full profinite Grothendieck--Teichm\"uller group.
They follow from Enriquez's distribution theorem, which is used here as
published background.

For the cyclic-cover subgroup $H_n\subset\widehat F_2$, define
\[
 d_n:H_n\longrightarrow\widehat F_2,\qquad
 x^n\longmapsto x,\quad y_0\longmapsto y,\quad
 y_j\longmapsto1\quad(0<j<n).
\]
Let $c_n=\operatorname{exp}_y\circ d_n$; thus $c_n=\ell_n$ in
our Schreier coordinates. Enriquez proves that, for every
$g=(\lambda,f)\in\widehat{\mathrm{GT}}$ and every $n\geq1$,
\begin{equation}
 d_n(f)=y^{c_n(f)}f.
 \label{eq:full-kummer-distribution}
\end{equation}
This is equation (50) and Theorem 8.1 of \cite{Enriquez2007}; in
particular his subgroup $\mathrm{GTK}$ is the full profinite group.
The map is the filling map on the cyclic cover, so this is an identity
in the full free profinite group, beyond its abelian or congruence images.

\begin{proposition}\label{prop:kummer-multiplicativity}
For positive integers $a,b$ and every $g\in\widehat{\mathrm{GT}}$,
\[
 \rho_{ab}^I(g)=\rho_a^I(g)+\rho_b^I(g).
\]
Consequently
\[
 \rho_n^I=\sum_{q\mid n}v_q(n)\rho_q^I,
 \qquad
 \frac{\rho_{ab}^I-\rho_a^I-\rho_b^I}{2}\bmod2=0.
\]
The last expression therefore supplies no nonzero character in the
restriction kernel.
\end{proposition}
\begin{proof}
On $H_{ab}$ one has $d_b(H_{ab})\subset H_a$ and
$d_a d_b=d_{ab}$. Indeed $x^{ab}$ maps successively to $x^a,x$;
for $y_j=x^jyx^{-j}$, the first map kills it unless $b\mid j$,
and otherwise sends it to $y_{j/b}$. The second map retains only
$j=0$, for $0\leq j<ab$.
Applying \eqref{eq:full-kummer-distribution} twice now gives
\[
 y^{c_{ab}(f)}f=d_{ab}(f)
   =d_a\bigl(y^{c_b(f)}f\bigr)
   =y^{c_b(f)+c_a(f)}f.
\]
The procyclic subgroup generated by $y$ is isomorphic to
$\widehat{\mathbb Z}$, so $c_{ab}(f)=c_a(f)+c_b(f)$.
This is also the observation in \cite[Remark 8.6]{Enriquez2007}.
Since $\rho_n^I=-c_n(f)$, the asserted formulas follow, with the
prime factorization formula obtained by iteration.
\end{proof}

The cocycle normalization agrees with the one used for the modular
frame: \cite[\S2.2, footnote on p.~100]{Furusho2003} identifies the
Nakamura--Schneps cocycle $\rho_2$ with $-\Psi_2^{(0)}$, and
\cite[Lemma 7(a) and Proposition 11]{H17} identifies that cocycle
with the pure-frame definition. Hence $\rho_2^I=\rho_2$ on
$\widehat{\mathrm{GT}}$. The negative-index Schreier generators in
Furusho's convention give the same coefficient $c_n$: for $0<j<n$,
$x^{-j}yx^j$ is conjugate within $H_n$ to $y_{n-j}$, so both
coefficient maps vanish on it.

\begin{corollary}\label{cor:dyadic-anderson}
Every element of $\widehat{\mathrm{GT}}$ satisfies all dyadic
Anderson multiplication relations $(A_{2^j})$, $j\geq1$, in the
notation of \cite[Definition 2.3]{Furusho2003}.
\end{corollary}
\begin{proof}
Relation (IV) holds on the group by \cite[Theorem 1]{H17}, and
$(K_2)$ holds by \eqref{eq:full-kummer-distribution}.
\cite[Theorem 3.1]{Furusho2003} therefore gives the duplication
relation; its dyadic iterates follow from that paper's Lemma 3.2.
Equivalently, combine
$\mathrm{GTK}=\widehat{\mathrm{GT}}$,
$\mathrm{I}\Gamma=\widehat{\mathrm{GT}}$
\cite[Corollary 21]{H17}, and
\cite[Theorem 3.4]{Furusho2003}.
\end{proof}

These are consequences of the cited distribution and multiplication
theorems. They do not determine the full restriction kernel: they show
that the displayed multiplicative defects vanish and impose the
specified dyadic identities.

\section{Simultaneous normalization of all rational Kummer coordinates}
\label{sec:rational-radical-normalization}

The arithmetic part of every rational Kummer coordinate can be removed
at once. This replaces the normalized subgroup by a proper closed
subgroup without changing the relative abelian obstruction.
The Galois calculation used here is the classical description of radical
extensions; see \cite[\S2, especially Theorem~2.5, author pp.~4--6]{MoreeStevenhagen}.
We give the finite-coordinate argument to make its simultaneous use and
the resulting GT quotient explicit.

\begin{lemma}\label{lem:simultaneous-kummer-image}
For distinct rational primes $q_1,\ldots,q_m$, the joint Kummer map
\[
 G_{\mathbb Q^{\mathrm{ab}}}\longrightarrow\widehat{\mathbb Z}^{\,m},
 \qquad \sigma\longmapsto
       (\kappa_{q_1}(\sigma),\ldots,\kappa_{q_m}(\sigma))
\]
has image exactly $2\widehat{\mathbb Z}^{\,m}$.
Consequently, for distinct odd primes $q_1,\ldots,q_m$, its restriction
to $K=G_{K_2}$ also has image $2\widehat{\mathbb Z}^{\,m}$.
\end{lemma}

\begin{proof}
Denote the first image by $J$. It is a compact, hence closed, subgroup
of $\widehat{\mathbb Z}^{\,m}$. Every continuous character of that
product with values in $\mathbb Q/\mathbb Z$ has the form
\[
 (t_1,\ldots,t_m)\longmapsto\frac1n\sum_{i=1}^m a_it_i\pmod{\mathbb Z}
\]
for integers $n>0$ and $a_i$. Such a character vanishes on $J$ exactly
when the positive real radical
\[
 \alpha=\prod_{i=1}^m q_i^{a_i/n}
\]
belongs to $\mathbb Q^{\mathrm{ab}}$. Indeed the Kummer action on
$\alpha$ multiplies it by the corresponding $n$th root of unity.

If $\alpha\in\mathbb Q^{\mathrm{ab}}$, its field over $\mathbb Q$
is normal and is contained in $\mathbb R$. Every conjugate of
$\alpha$ is a real root of $X^n-\alpha^n$, and is therefore one
of $\alpha,-\alpha$. It follows that $\alpha^2\in\mathbb Q$.
Writing this last rational number in its unique prime factorization
gives $2a_i/n\in\mathbb Z$ for each $i$. Conversely these
integrality conditions make $\alpha$ a square root of a positive
rational number, and hence an element of $\mathbb Q^{\mathrm{ab}}$.
Thus the annihilator of $J$ is precisely the group of characters
of exponent at most two. It is also the annihilator of
$2\widehat{\mathbb Z}^{\,m}$, so the two closed subgroups agree.
For this last implication, a proper closed subgroup of a profinite
abelian group is separated from any point outside it by a character
of a finite quotient.

For the second assertion, apply the first to the list
$2,q_1,\ldots,q_m$. Its image is the full product
$2\widehat{\mathbb Z}^{\,m+1}$. Every point of that product whose
first coordinate is exactly zero has a preimage fixing all roots
of two as well as all roots of unity. Such a preimage belongs to
$K$, proving the assertion with no compatibility condition left over.
\end{proof}

Let $\mathcal P$ be the set of rational primes and put
\[
 L_{\mathrm{rad}}=
 \mathbb Q\bigl(\mu_\infty,q^{1/n}:q\in\mathcal P,
                                      \ n\geq1\bigr).
\]
The roots in this expression are positive real roots; adjoining
$\mu_\infty$ also adjoins all their conjugates.
Define
\[
 N_{\mathrm{rad}}=N\cap
       \bigcap_{\substack{q\in\mathcal P\\q\ne2}}\ker\rho_q^I,
 \qquad K_{\mathrm{rad}}=G_{L_{\mathrm{rad}}}.
\]
The frame cocycle $\rho_2$ and $\rho_2^I$ agree on $G$;
the normalization is recorded in \cite[\S2.2, p.~100, footnote]{Furusho2003}
and \cite[Proposition~11]{H17}. Thus $N_{\mathrm{rad}}$ is the
simultaneous kernel of the multiplier and all rational-prime Kummer
cocycles.

\begin{theorem}\label{thm:rational-radical-normalization}
The joint map
\[
 \mathcal R:N\longrightarrow
       \prod_{\substack{q\in\mathcal P\\q\ne2}}2\widehat{\mathbb Z},
 \qquad n\longmapsto(\rho_q^I(n))_q
\]
is surjective already on $K$. Its kernel is $N_{\mathrm{rad}}$,
and
\[
 N=N_{\mathrm{rad}}K,\qquad
 G=N_{\mathrm{rad}}H,\qquad
 N_{\mathrm{rad}}\cap H=K_{\mathrm{rad}}.
\]
In particular $N_{\mathrm{rad}}$ is a proper closed normal subgroup
of $N$, and
\[
 N/N_{\mathrm{rad}}\simeq
       \prod_{\substack{q\in\mathcal P\\q\ne2}}2\widehat{\mathbb Z}.
\]
There are canonical topological isomorphisms
\begin{align*}
 G/N_{\mathrm{rad}}&\simeq
            \operatorname{Gal}(L_{\mathrm{rad}}/\mathbb Q),\\
 B&\simeq N_{\mathrm{rad}}/
       \overline{[G,N_{\mathrm{rad}}]K_{\mathrm{rad}}}.
\end{align*}
The first quotient is explicitly
\[
 A_{\mathrm{rad}}=
 \left\{(u,(t_q)_{q\in\mathcal P}):
 \begin{array}{l}
 u\in\widehat{\mathbb Z}^{\times},\quad
 t_q\in\widehat{\mathbb Z},\\
 t_q\equiv\epsilon_q(u)\pmod2\quad(q\in\mathcal P)
 \end{array}\right\},
\]
with multiplication $(u,t)(v,s)=(uv,t+us)$, where
$\epsilon_q$ is the quadratic character of $\mathbb Q(\sqrt q)$.
\end{theorem}

\begin{proof}
Each coordinate of $\mathcal R$ is a homomorphism on $N$ and is
even by Theorem~\ref{thm:all-kummer-parity}. The joint map is
continuous for the product topology. By
Lemma~\ref{lem:simultaneous-kummer-image}, every finite-coordinate
projection of $\mathcal R(K)$ is the whole corresponding product.
Hence $\mathcal R(K)$ is dense in the full product. It is compact
and therefore closed, proving
\[
 \mathcal R(K)=\mathcal R(N)=\prod_{q\ne2}2\widehat{\mathbb Z}.
\]
Equivalently, one can prescribe successively larger finite lists
of coordinates and use the finite intersection property in $K$.
This compactness step prescribes the whole profinite value of every
coordinate, not only its residue at one finite level.

Given $n\in N$, choose $k\in K$ with $\mathcal R(k)=\mathcal R(n)$.
Then $nk^{-1}\in N_{\mathrm{rad}}$, proving
$N=N_{\mathrm{rad}}K$ and then $G=N_{\mathrm{rad}}H$.
The conjugation identity for the Kummer cocycles shows that
$N_{\mathrm{rad}}$ is normal in $G$; it is closed by definition.
Its Galois intersection fixes precisely the roots of unity and
the roots of all rational primes, so it is $K_{\mathrm{rad}}$.
The quotient isomorphism $N/N_{\mathrm{rad}}\simeq\prod_{q\ne2}
2\widehat{\mathbb Z}$ follows from compactness. This nontrivial
quotient also proves that $N_{\mathrm{rad}}$ is proper.

The factorization $G=N_{\mathrm{rad}}H$ identifies
$G/N_{\mathrm{rad}}$ with $H/K_{\mathrm{rad}}$.
On the Galois group, each translation coordinate has parity
$\epsilon_q(u)$ because it records the action on $\sqrt q$.
For a fixed multiplier choose one Galois lift. The first assertion
of Lemma~\ref{lem:simultaneous-kummer-image}, followed by compactness
over all finite prime lists, allows an arbitrary even translation
vector to be added to this lift. Thus the image is exactly
$A_{\mathrm{rad}}$, with the displayed cocycle multiplication.

Finally apply the proof of Theorem~\ref{thm:abelian-splitting}
to $G=N_{\mathrm{rad}}H$. The Galois intersection
$K_{\mathrm{rad}}$ lies in $\overline{[H,H]}$, since it fixes
$\mathbb Q^{\mathrm{ab}}$. That proof identifies the new relative
abelian quotient with
$(G/\overline{\langle\!\langle H\rangle\!\rangle})^{\mathrm{ab}}$,
which is also $B$. This gives the asserted canonical isomorphism.
\end{proof}

\begin{corollary}\label{cor:rational-kummer-characters}
The invariant quadratic characters of $N$ that factor through all
rational Kummer coordinates form exactly
\[
 X_2^{\mathrm{Kum}}=
 \bigoplus_{\substack{q\in\mathcal P\\q\ne2}}
       \mathbb F_2\,\psi_q^N.
\]
Restriction to $K$ is injective on this subspace. For an odd prime
$p$, no nonzero invariant $\mathbb F_p$-character of $N$ factors
through the rational Kummer coordinates.
For every prime $p$, the unresolved relative kernel is equivalently
\[
 \operatorname{Hom}(B,\mathbb F_p)=
 \left\{\chi\in H^1(N_{\mathrm{rad}},\mathbb F_p)^G:
                       \chi|_{K_{\mathrm{rad}}}=0\right\}.
\]
\end{corollary}

\begin{proof}
Enriquez's distribution theorem gives
$\rho_{ab}^I=\rho_a^I+\rho_b^I$ on $G$ for positive integers
$a,b$; see \cite[Theorem~8.1 and Remark~8.6, author pp.~30--32]{Enriquez2007}.
Since $\rho_2^I$ is zero on $N$, all rational Kummer coordinates
on $N$ are generated by the odd-prime coordinates.
A continuous character of the product in
Theorem~\ref{thm:rational-radical-normalization} to a finite group
depends on finitely many coordinates: its open kernel contains
a basic product neighborhood of zero. Each factor
$2\widehat{\mathbb Z}$ has a one-dimensional quadratic character
space, generated by division by two followed by reduction modulo
two. This proves the direct-sum formula. Every multiplier is odd,
so every one of these characters is $G$-invariant. Surjectivity
already on $K$ proves injectivity of restriction.

Complex conjugation acts on the product by multiplication by $-1$.
If $p$ is odd and a character $\chi$ is invariant, then
$\chi(t)=\chi(-t)=-\chi(t)$, so it is zero. The final identification
is the character-gluing lemma applied to
$G=N_{\mathrm{rad}}H$.
\end{proof}

\begin{proposition}\label{prop:radical-character-pullback}
For every prime $p$, put
\[
 X_p^{\mathrm{rad}}=H^1(N_{\mathrm{rad}},\mathbb F_p)^G,
 \qquad
 Y_p^{\mathrm{rad}}=H^1(K_{\mathrm{rad}},\mathbb F_p)^H.
\]
The restriction maps identify $X_p$ with the fiber product
\[
 X_p\simeq X_p^{\mathrm{rad}}
                  \mathop{\times}_{Y_p^{\mathrm{rad}}}Y_p.
\]
The kernels of $X_p\to X_p^{\mathrm{rad}}$ and
$Y_p\to Y_p^{\mathrm{rad}}$ are identified by restriction.
They are zero for odd $p$, and for $p=2$ they are respectively
$X_2^{\mathrm{Kum}}$ and its restriction to $K$.
\end{proposition}
\begin{proof}
A character of $N$ restricts to the indicated compatible pair,
and the pair determines it because $N=N_{\mathrm{rad}}K$.
Conversely, let $\alpha\in X_p^{\mathrm{rad}}$ and
$\beta\in Y_p$ have the same restriction to $K_{\mathrm{rad}}$.
Define
\[
 \chi(nk)=\alpha(n)+\beta(k)
       \qquad(n\in N_{\mathrm{rad}},\ k\in K).
\]
If $nk=n'k'$, then $n'^{-1}n=k'k^{-1}$ lies in
$N_{\mathrm{rad}}\cap K=K_{\mathrm{rad}}$; the equality of
the two restrictions proves well-definedness. The invariance of
$\alpha$ under $K$ makes this formula a homomorphism.
It is continuous: the multiplication map from the compact space
$N_{\mathrm{rad}}\times K$ onto the Hausdorff group $N$ is
a quotient map, and its pullback is the continuous sum
$\alpha+\beta$. Conjugation by $H$ preserves the formula,
since $\alpha$ is $G$-invariant and $\beta$ is $H$-invariant.
Conjugation by $N$ preserves any abelian character of $N$.
Since $G=NH$, the resulting character belongs to $X_p$.

The kernel of restriction to $N_{\mathrm{rad}}$ is the invariant
character space of $N/N_{\mathrm{rad}}$; the corresponding kernel
on the Galois side is that of $K/K_{\mathrm{rad}}$.
The joint Kummer map identifies both quotients with
$\prod_{q\ne2}2\widehat{\mathbb Z}$, with the same scalar action
of the multiplier. The preceding corollary computes their
invariant character spaces, proving the final assertions.
\end{proof}

The subgroup $N_{\mathrm{rad}}$ gives a strictly smaller location
for the open obstruction. A character of $G$ which vanishes on
$H$ is zero if and only if its restriction to this subgroup is zero.
The theorem does not determine its relative abelianization.

\section{A concrete quadratic character from the eleventh multiplication formula}
\label{sec:odd-anderson-characters}

This section constructs the explicit continuous quadratic character
$\Theta_{17,11}$, which vanishes on the arithmetic subgroup.  Its
nonvanishing on the full
profinite Grothendieck--Teichm\"uller group is not established.  A separate
formal example shows why beta factorization, the unitary relation, and
the dyadic multiplication formulas do not by themselves eliminate these
characters.  The example is not asserted to lift to a pentagon element.

Let $G_0$ be the profinite gentle group and let
$G=\widehat{\mathrm{GT}}\subset G_0$.  Write
$F=\widehat F(x,y)$ and $z=(xy)^{-1}$.  We use the convention
\[
 \phi_g(x)=x^{\lambda_g},\qquad
 \phi_g(y)=f_g^{-1}y^{\lambda_g}f_g,
 \qquad f_g\in F'.
\]
Let $B_g(X,Y)$ be the normalized Alexander coefficient, so that the
semilinear coefficient on $F'/F''$ is
\[
 A_g(X,Y)=S_{\lambda_g}(X)S_{\lambda_g}(Y)B_g(X,Y),
 \qquad S_a(T)=\frac{T^a-1}{T-1}.
\]
Composition gives
\begin{equation}
 B_{gh}(X,Y)=B_g(X,Y)
                 B_h(X^{\lambda_g},Y^{\lambda_g}).
 \label{eq:odd-beta-cocycle}
\end{equation}
These are the conventions of \cite[\S2.3, p.~101]{Furusho2003}.

\subsection{The boundary characters in the unitary relation}

\begin{lemma}\label{lem:all-finite-unitarity}
Suppose that $\lambda_g=1$.  Let $p$ be an odd prime and let
$a,b\in\overline{\mathbb F}_p^\times$ have finite orders prime to $p$.
Then
\[
 B_g(a,b)B_g(a^{-1},b^{-1})=1.
\]
In fact $B_g(a,b)=1$ if $a=1$, $b=1$, or $ab=1$.
\end{lemma}

\begin{proof}
Choose $m$ prime to $p$ divisible by the orders of $a$ and $b$, and
put $U_m=\ker(F\to(\mathbb Z/m)^2)$.  Since $f_g\in F'$, the
automorphism $\phi_g$ induces the identity on $F/\gamma_3F$.  Choose
$c\in F$ such that $\phi_g(z)=c^{-1}zc$, and write
$[c]=(u,v)\in F^{\mathrm{ab}}=\widehat{\mathbb Z}^2$.  In the
class-two quotient the commutator subgroup is freely generated over
$\widehat{\mathbb Z}$ by $[x,y]$.  The equality $c^{-1}zc=z$ in this
quotient implies $u=v$.

The cover with fundamental group $U_m$ has deck group $(\mathbb Z/m)^2$.
Its boundary subrepresentation in $H_1(U_m,\mathbb F_p)$ is the sum
of the three permutation representations induced from the cyclic
subgroups generated by the images of $x$, $y$, and $z$, with one
trivial relation among the boundary classes.  On the boundaries above
$x$ the conjugator is $1$; above $y$ it is $f_g$, whose deck image is
$1$; above $z$ it is $c$, whose deck image is diagonal.  Thus $\phi_g$
fixes the boundary representation pointwise.  More explicitly, a
character $(a,b)$ occurring above $z$ satisfies $ab=1$, and its
conjugator scalar is $a^u b^v=(ab)^u=1$.  The corresponding assertions
above $x$ and $y$ are immediate.  Every nontrivial deck character
occurs in $H_1(U_m,\overline{\mathbb F}_p)$ with multiplicity one.
The coefficient $B_g(a,b)$ is its scalar action.  This proves the
last assertion for all nontrivial boundary characters.  At the trivial
character $B_g(1,1)=1$, since the class-two commutator is fixed.

For completeness, the complement of the boundary subrepresentation
carries the usual nondegenerate alternating form.  One can see its
preservation without a topological realization of $\phi_g$.  A surface
presentation of $U_m$ has the form
\[
 \prod_{j=1}^{(m-1)(m-2)/2}[a_j,b_j]
       \prod_{i=1}^{3m}c_i=1.
\]
In the class-two exponent-$p$ quotient, project the degree-two
commutators to the exterior square of homology modulo the boundary
classes.  The relation yields the bivector
$\sum_j[a_j]\wedge[b_j]$.  Conjugating a boundary word changes its
degree-two coordinate by a wedge containing its boundary class, so
this change vanishes after projection.  Positive boundary preservation
therefore fixes the bivector and the associated alternating form.

The deck action preserves the form, so the character spaces at
$(a,b)$ and $(a^{-1},b^{-1})$ pair perfectly whenever
$a,b,ab\ne1$.  Preservation of the form makes the product of their
scalar actions equal to one.  The boundary calculation proves the
same identity in the remaining cases.
\end{proof}

\subsection{An explicit family of arithmetic-trivial characters}

\begin{theorem}\label{thm:odd-anderson-character}
Let $p\equiv5\pmod{12}$ be prime, and let $n\geq1$ satisfy
$\gcd(n,3p)=1$.  Choose a primitive cube root
$\omega\in\mathbb F_{p^2}$.  The expression
\begin{equation}
 E_{p,n}(g)=
 \frac{\displaystyle\prod_{k=0}^{n-1}B_g(\omega,\omega^k)}
      {\displaystyle\prod_{\xi^n=1}B_g(\xi,\omega)}
 \label{eq:odd-anderson-defect}
\end{equation}
belongs to the norm-one subgroup of $\mathbb F_{p^2}^\times$.
Here all evaluations are reduced modulo $p$, and the denominator
ranges over all $n$th roots in $\overline{\mathbb F}_p$.
The rule
\begin{equation}
 (-1)^{\Theta_{p,n}(g)}=E_{p,n}(g)^{(p+1)/2}
 \label{eq:odd-anderson-sign}
\end{equation}
defines a continuous homomorphism $G_0\to\mathbb F_2$, independent
of the choice between $\omega$ and $\omega^2$.  It vanishes on
$G_{\mathbb Q}$.

In particular its restriction to $G$ is an explicit candidate for a
nonzero element of
$\operatorname{Hom}_{\mathrm{cts}}(G/\langle\!\langle
G_{\mathbb Q}\rangle\!\rangle,\mathbb F_2)$.
The theorem does not assert that this restriction is nonzero.
\end{theorem}

\begin{proof}
The evaluations are nonzero because they are coefficients of
invertible semilinear actions on the appropriate character spaces.
The $p^2$-power Frobenius fixes $\omega$ and permutes the $n$th
roots.  Hence the expression lies in $\mathbb F_{p^2}^\times$.
Equation~\eqref{eq:odd-beta-cocycle} and reindexing the denominator
give
\begin{equation}
 E_{p,n}(gh)=E_{p,n}(g)\,
             \sigma_{\lambda_g}\bigl(E_{p,n}(h)\bigr),
 \label{eq:odd-anderson-crossed}
\end{equation}
where $\sigma_\lambda$ is the identity for $\lambda\equiv1\pmod3$
and the $p$-power Frobenius for $\lambda\equiv-1\pmod3$.
Indeed $\lambda$ permutes all $n$th roots, and its only remaining
effect in the displayed expression is its action on $\omega$.

If $\lambda_g=1$, Lemma~\ref{lem:all-finite-unitarity} gives
$E_{p,n}(g)^p=E_{p,n}(g)^{-1}$.  In the denominator one reindexes
the $p$th powers of the $n$th roots as their inverses; in the numerator
$\omega^p=\omega^{-1}$.  For arbitrary $g$, choose
$h\in G_{\mathbb Q}$ with $\lambda_h=\lambda_g$ and write $g=uh$,
where $\lambda_u=1$.  The arithmetic Anderson multiplication formula
\cite[Corollary~8.6.3]{Anderson1989}, in the beta form stated explicitly
in \cite[Remark~2.4, p.~101]{Furusho2003}, gives
\begin{equation}
 E_{p,n}(h)=\omega^{n\Psi_n(h)}\in\mu_3.
 \label{eq:odd-anderson-arithmetic}
\end{equation}
Its exact sign convention is immaterial here; in the convention above
$\Psi_n=-\rho_n$.  Since $p\equiv-1\pmod3$, the group $\mu_3$ also
has norm one.  Equation~\eqref{eq:odd-anderson-crossed} now proves
that $E_{p,n}(g)$ has norm one.

The norm-one subgroup is cyclic of order $p+1$, so
\eqref{eq:odd-anderson-sign} is a sign.  Frobenius acts on that
subgroup by inversion and therefore fixes its quadratic sign.
Consequently \eqref{eq:odd-anderson-crossed} makes $\Theta_{p,n}$
an ordinary homomorphism.  It is continuous because it uses finitely
many finite character evaluations.  Since $3$ divides $(p+1)/2$,
\eqref{eq:odd-anderson-arithmetic} makes this homomorphism zero on
$G_{\mathbb Q}$.  Replacing $\omega$ by $\omega^2$ applies Frobenius
and leaves the sign unchanged.
\end{proof}

For a pentagon element, Ihara's beta factorization gives
$B_g(s,t)=\Gamma_g(s)\Gamma_g(t)/\Gamma_g(s+t)$ in additive
notation.  Consequently
\begin{equation}
 E_{p,n}(g)=
 \frac{\displaystyle\prod_{c\in\frac1n\mathbb Z/\mathbb Z}
                    \Gamma_g(\tfrac13+c)}
      {\displaystyle\Gamma_g(\tfrac n3)
          \prod_{c\in\frac1n\mathbb Z/\mathbb Z}\Gamma_g(c)}.
 \label{eq:odd-anderson-gamma}
\end{equation}
This expression is independent of the character ambiguity in
$\Gamma_g$.  If the Anderson relation $(A_n)$ holds for $g$, then
$\Theta_{p,n}(g)=0$.  Thus the known dyadic multiplication formulas
eliminate the cases $n=2^a$.  Enriquez's theorem that
$\mathrm{GTK}=\widehat{\mathrm{GT}}$ also eliminates every proposed
character obtained from a multiplicative defect of the integral
Kummer coordinates: his Remark~8.6 gives
$\rho_{ab}=\rho_a+\rho_b$ before any reduction or division.
The exact reference is \cite[Theorem~8.1 and Remark~8.6]{Enriquez2007}.
It does not assert the equality
$\mathrm{GTA}=\widehat{\mathrm{GT}}$; the relation $(K_n)$ is recorded
as implying the logarithmic derivative of $(A_n)$ in
\cite[p.~102]{Furusho2003}.

\begin{corollary}\label{cor:concrete-radical-anderson-character}
Let $H=G_{\mathbb Q}$ and let $N_{\mathrm{rad}}$ be the simultaneous
kernel in $G$ of the multiplier and the rational-prime Kummer cocycles.
Put
\[
 L_{\mathrm{rad}}=\mathbb Q(\mu_\infty,q^{1/m}:
                          q\text{ prime},\ m\geq1),
 \qquad K_{\mathrm{rad}}=G_{L_{\mathrm{rad}}}.
\]
Assume the simultaneous normalization
$G=N_{\mathrm{rad}}H$, with
$N_{\mathrm{rad}}\cap H=K_{\mathrm{rad}}$, and identify the relative
abelian quotient as
\[
 B=N_{\mathrm{rad}}/
       \overline{[G,N_{\mathrm{rad}}]K_{\mathrm{rad}}}.
\]
For a primitive cube root $\omega\in\mathbb F_{289}$ the formula
\begin{equation}
 (-1)^{\psi_{17,11}(u)}=
    \left(\prod_{\substack{\xi^{11}=1\\\xi\ne1}}
                        B_u(\xi,\omega)\right)^9,
 \qquad u\in N_{\mathrm{rad}},
 \label{eq:concrete-radical-anderson-sign}
\end{equation}
defines a continuous $G$-invariant quadratic character vanishing on
$K_{\mathrm{rad}}$.  It therefore gives a precisely specified element
of $\operatorname{Hom}_{\mathrm{cts}}(B,\mathbb F_2)$.
Its unique extension to $G$ vanishing on $H$ is $\Theta_{17,11}$.
In particular $\psi_{17,11}$ is nonzero if and only if
$\Theta_{17,11}|_G$ is nonzero; neither nonvanishing assertion is
proved here.
\end{corollary}

\begin{proof}
Take $p=17$, $n=11$ in Theorem~\ref{thm:odd-anderson-character}.
For $u\in N_{\mathrm{rad}}$ the multiplier is one.
In the numerator of \eqref{eq:odd-anderson-defect}, four values of
$k$ are $1$ modulo $3$, four are $0$ modulo $3$, and three are $2$
modulo $3$.  Lemma~\ref{lem:all-finite-unitarity} gives
\[
 \prod_{k=0}^{10}B_u(\omega,\omega^k)=B_u(\omega,\omega)^4,
 \qquad B_u(\omega,\omega)^{18}=1.
\]
Its ninth-power contribution is therefore one.  The denominator's
$\xi=1$ term is also one.  Since inversion fixes a sign, the
restriction of \eqref{eq:odd-anderson-sign} is exactly
\eqref{eq:concrete-radical-anderson-sign}.

Being a restriction of an ordinary character of $G$, this map is
$G$-invariant; being zero on $H$, it is zero on $K_{\mathrm{rad}}$.
In fact the arithmetic formula gives $E_{17,11}(u)=1$ on
$K_{\mathrm{rad}}$, since $\Psi_{11}(u)=0$ there.
It consequently factors through the displayed quotient $B$.
The factorization $G=N_{\mathrm{rad}}H$ proves both uniqueness
of the extension and the equivalence of the two nonvanishing claims.
\end{proof}

There is also a finite-dimensional interpretation of the product in
\eqref{eq:concrete-radical-anderson-sign}.  Take the degree-$33$ cover
with deck group $C_{11}\times C_3$, given algebraically by
$u^{11}+v^3=1$ and the map $t=u^{11}$.  There are three boundary
components above $t=0$, eleven above $t=1$, and one above infinity.
Euler characteristic therefore gives genus ten for its compactification.
Its homology has the ten character lines $(\xi,\omega)$ with
$\xi$ a nontrivial $11$th root, together with their inverse lines.
The individual evaluations may lie in $\mathbb F_{17^{10}}$;
$\mu_{11}$ is not contained in $\mathbb F_{289}$.
The sum descends to a ten-dimensional space over $\mathbb F_{289}$,
since the $17^2$-power Frobenius permutes these ten characters.  It is
the whole $\omega$-eigenspace of the cubic deck action on the compact
homology.  The product is the determinant of an element of
$N_{\mathrm{rad}}$ on this space, and
\eqref{eq:concrete-radical-anderson-sign} is its quadratic norm-one
sign.  Thus the proposed test involves one finite homology
representation of a genus-ten curve.

\subsection{A formal sign ambiguity which survives duplication}

\begin{proposition}\label{prop:formal-odd-sign}
There exists a Frobenius-equivariant sign-valued formal gamma function
at the prime $17$ whose beta coboundary satisfies beta associativity,
symmetry, the unitary relation, and the duplication formula, but for
which the expression in \eqref{eq:odd-anderson-gamma} with $n=11$
equals $-1$.  It may be taken constant in the pro-$17$ directions,
so all its logarithmic derivatives vanish.

This is a non-implication statement for the listed formal identities.
It does not provide a solution of the full nonabelian relations
$(K_n)$ or of the pentagon equation.
\end{proposition}

\begin{proof}
Let $T=(\mathbb Q/\mathbb Z)_{17'}$ be the subgroup of elements of
order prime to $17$.  In $(\mathbb Z/33)^\times$ put
\[
 S=\langle2\rangle
   =\{1,2,4,8,16,32,31,29,25,17\}.
\]
Define $\gamma:T\to\mathbb F_2$ by
\[
 \gamma(s)=
 \begin{cases}
  1,&s=a/33\text{ for }a\in S,\\
  0,&\text{otherwise},
 \end{cases}
 \qquad \Gamma(s)=(-1)^{\gamma(s)}.
\]
For an arbitrary rational torsion argument $s$, the definition means
$\gamma(s)=\gamma(s_{17'})$, where $s_{17'}$ is its prime-to-$17$
component.  Thus the value is constant in the $17$-primary variable;
it is not declared to be $1$ merely because the denominator is
divisible by $17$.  In particular $\gamma$ is zero whenever the
prime-to-$17$ component has even order.
The set $S$ is stable under $2$, $-1=2^5$, and $17=2^9$.
Thus $\Gamma(-s)=\Gamma(s)$ and $\Gamma(17s)=\Gamma(s)$.

We check the duplication relation
\begin{equation}
 \gamma(s)+\gamma(s+\tfrac12)=\gamma(2s)+\gamma(\tfrac12).
 \label{eq:formal-sign-duplication}
\end{equation}
The last term is zero.  If $s$ has odd order, then $s+1/2$ has
even order and the assertion is the invariance under doubling on
odd torsion.  If the $2$-primary component of $s$ has order exactly
two, put $v=s+1/2$.  Then $v$ has odd order, $2s=2v$, and the same
invariance proves the identity.  If the $2$-primary component has
order at least four, all three possibly nonzero arguments on the
left and right have even order, so all terms vanish.  This proves
\eqref{eq:formal-sign-duplication}, or equivalently
\[
 \Gamma(s)\Gamma(s+\tfrac12)=\Gamma(2s)\Gamma(\tfrac12).
\]
Iteration gives all dyadic multiplication identities with zero Kummer
correction.

Put $B(s,t)=\Gamma(s)\Gamma(t)/\Gamma(s+t)$.  It is symmetric,
satisfies beta associativity, is one on the axes and on $s+t=0$,
and satisfies $B(s,t)B(-s,-t)=1$.  These assertions follow directly
from the coboundary definition and the fact that all values are signs.
For the prime-to-$17$ spectral components, the Frobenius invariance
allows the assignment to descend to $\mathbb Z_{17}$ coefficients.
Make it constant on each pro-$17$ formal neighborhood.  Equivalently,
use the product decomposition of
$\mathbb Z_{17}[[\widehat{\mathbb Z}]]$ into the unramified formal
power-series rings indexed by Frobenius orbits of prime-to-$17$
characters.  The specified constants define a unit in this completed
group algebra.  The beta coboundary is a corresponding unit in the
two-variable algebra, and its logarithmic derivatives are zero.

Finally $\Gamma(11/3)=\Gamma(2/3)=1$, and $\Gamma(c)=1$ for
every $c\in\frac1{11}\mathbb Z/\mathbb Z$.  The numerators of the
points $1/3+k/11$, expressed with denominator $33$, are $11+3k$.
Their intersection with $S$ is
\[
 \{2,8,17,29,32\},
\]
which has five elements.  Therefore the numerator of
\eqref{eq:odd-anderson-gamma} is $(-1)^5=-1$ and its denominator
is $1$.  At $p=17$ the quadratic exponent is $(17+1)/2=9$, so
the associated formal quadratic sign is also $-1$.
\end{proof}

The remaining question is a full-group realization or elimination
problem.  Neither the finite sign computation nor the fact that
logarithmic derivatives forget this sign permits one to infer an
element of $\widehat{\mathrm{GT}}$ with nonzero
$\Theta_{17,11}$.  Conversely, eliminating these candidates requires
an additional consequence of the full relations, rather than only
the four formal identities in Proposition~\ref{prop:formal-odd-sign}.

\subsection{Two limitations of a direct eleventh-Kummer comparison}

Put $D=\ker\lambda\cap\bigcap_q\ker\rho_q^I$ in the full
profinite Grothendieck--Teichm\"uller group.  By \cite[Theorem~8.1 and Remark~8.6]{Enriquez2007}, the exact
Kummer relation, with zero Kummer coordinate on $D$, is
\begin{equation}\label{eq:eleventh-local-system-K}
d_n(f_g)=f_g\qquad(g\in D).
\end{equation}
Here $H_n\leq\widehat F(x,y)$ is the kernel of $x\mapsto1,y\mapsto0$
in $\mathbb Z/n$, freely generated by
\[
 X=x^n,\qquad Y_j=x^{-j}yx^j\quad(0\leq j<n),
\]
and $d_n:H_n\to\widehat F(x,y)$ sends $X\mapsto x$, $Y_0\mapsto y$
and $Y_j\mapsto1$ for $j\ne0$.  We apply the full nonabelian relation \eqref{eq:eleventh-local-system-K}.

\begin{lemma}\label{lem:kn-different-local-systems}
For $n=11$ and a primitive cube root $\omega$ in
$\mathbb F_{17^2}$, consider characters over $\overline{\mathbb F}_{17}$.
The character of $H_{11}$ obtained by pulling
back $x\mapsto1,y\mapsto\omega$ along $d_{11}$ is different from
the restriction of every character
$x\mapsto\xi,y\mapsto\omega$ with $\xi^{11}=1$.
Nor is either kind a deck translate of the other.
\end{lemma}
\begin{proof}
The pullback along $d_{11}$ has values
\[
 X\mapsto1,\qquad Y_0\mapsto\omega,\qquad
 Y_j\mapsto1\ (j\ne0).
\]
Each character of the second kind has values
\[
 X\mapsto\xi^{11}=1,\qquad Y_j\mapsto\omega
 \quad\hbox{for every }j.
\]
The deck generator cyclically permutes the conjugacy classes of
the $Y_j$.  It therefore permutes the eleven characters having a
single nontrivial $Y_j$-value and fixes the character whose eleven
$Y_j$-values are all $\omega$.  These two orbits are disjoint.
\end{proof}

Consequently the direct operation ``apply \eqref{eq:eleventh-local-system-K}, then pass to the
pulled-back rank-one local system'' does not identify the product
\begin{equation}\label{eq:eleventh-local-system-P}
P_g=\prod_{\substack{\xi^{11}=1\\\xi\ne1}}B_g(\xi,\omega)
\end{equation}
with the determinant on that pulled-back local system.  This
observation is a limitation of this proposed argument, not a claim
that the full relation \eqref{eq:eleventh-local-system-K} cannot constrain \eqref{eq:eleventh-local-system-P} by another route.

Here is a separate, explicit limitation of a pairing-only argument.
Let $C$ be the smooth compact curve with affine equation
$u^{11}+v^3=1$, and let
$A=C_{11}\times C_3$ be its deck group over the $t$-line,
$t=u^{11}$.  We use the geometric finite homology description:
over an algebraic closure of $\mathbb F_{17}$, its $20$-dimensional
homology has exactly the character lines
\begin{equation}\label{eq:eleventh-local-system-H}
(\xi^a,\omega^b),\qquad
 a\in(\mathbb Z/11)^\times,\quad b\in\{1,2\}.
\end{equation}
Every line is one-dimensional.  The alternating intersection form
pairs a line only with its inverse-character line.
For completeness, the character description follows from the
two-loop cellular complex of the three-punctured sphere.  On a
nontrivial character the boundary from the two one-cells to the
one zero-cell has rank one, so the open cover has a single
homology line.  Compactification kills precisely the peripheral
lines whose character is trivial on one of the three inertia
groups.  The inertia values here are $\xi^a$, $\omega^b$ and
$(\xi^a\omega^b)^{-1}$.  Since eleven and three are coprime,
all three are nontrivial exactly for the pairs in \eqref{eq:eleventh-local-system-H}.
The trivial character contributes no compact homology.

\begin{proposition}\label{prop:kn-linear-sign-model}
There is an $\mathbb F_{17}$-linear automorphism of the space \eqref{eq:eleventh-local-system-H}
that commutes with the whole deck group $A$, preserves its
alternating intersection form, and has determinant $-1$ on the
ten-dimensional $\omega$-eigenspace over $\mathbb F_{17^2}$.
Its quadratic norm-one determinant sign is therefore $-1$.
This is a linear model, not a GT element or a solution of the
pentagon or of \eqref{eq:eleventh-local-system-K}.
\end{proposition}
\begin{proof}
The square of Frobenius acts on the nonzero exponents modulo
eleven by multiplication by $17^2\equiv3$.  Its two orbits are
\[
 A_0=\{1,3,9,5,4\},\qquad
 A_1=\{2,6,7,10,8\}=-A_0.
\]
Ordinary Frobenius acts on the pairs $(a,b)$ in \eqref{eq:eleventh-local-system-H} as
\[
 (a,b)\longmapsto(6a,-b).
\]
It thus preserves the set
\[
 \mathcal O=(A_0\times\{1\})\cup(A_1\times\{2\}).
\]
Indeed $6A_0=A_1$ and $6A_1=A_0$.
The set $\mathcal O$ is also stable under inversion
$(a,b)\mapsto(-a,-b)$, as is its complement.

Define the automorphism to be $-1$ on the ten character lines in
$\mathcal O$ and $+1$ on the ten remaining lines.  Frobenius
stability makes it descend to $\mathbb F_{17}$, and its scalar
action on each character line makes it commute with $A$.
The intersection form pairs only mutually inverse characters:
if $v,w$ have characters $\chi,\psi$, invariance gives
$\langle v,w\rangle=\chi(a)\psi(a)\langle v,w\rangle$ for every
deck element $a$, which forces $\psi=\chi^{-1}$ whenever the
pairing is nonzero.  Our automorphism has equal signs on these
paired lines, so it preserves that form.

There are exactly five negative lines in the $\omega$-eigenspace,
namely those indexed by $A_0$.  Its determinant there is
$(-1)^5=-1$.  This value belongs to the norm-one group of
$\mathbb F_{17^2}^\times$, and raising it to
$(17+1)/2=9$ leaves $-1$, as asserted.
\end{proof}

The exact remaining question is whether the full nonabelian
relations, including \eqref{eq:eleventh-local-system-K} and the pentagon, exclude this sign in
the GT action.  Neither Lemma~\ref{lem:kn-different-local-systems}
nor Proposition~\ref{prop:kn-linear-sign-model} answers it.
In particular, no assertion about nonvanishing of
$\Theta_{17,11}|_{\widehat{\mathrm{GT}}}$ follows from these two statements.


\section{The comparison with Galois and the gentle group}
The factorization also makes the nonabelian equality question and the
role of gentle characters precise.

\begin{proposition}\label{prop:normalized-comparison}
There is a natural homeomorphism of right-coset spaces
\[
 N/K\longrightarrow G/H,\qquad nK\longmapsto nH.
\]
In particular $G=H$ if and only if $N=K$.
Put $N_0=\ker(\lambda:\GTz\to U)\cap\ker\rho_2$.
Then $\GTz=N_0H$. If a continuous homomorphism from $\GTz$
to any topological group is trivial on $H$, its restriction to $G$
is trivial if and only if its restriction to $N$ is trivial.
\end{proposition}
\begin{proof}
Surjectivity of the coset map is $G=NH$. Two elements of $N$ have
the same right $H$-coset precisely when their quotient lies in
$N\cap H=K$. The map is continuous; compactness of the source and
the Hausdorff property of the target make it a homeomorphism.
There is a single coset exactly when $N=K$, proving the equivalence.

For $g\in\GTz$, choose $h\in H$ with $\lambda(h)=\lambda(g)$.
By \cite[Lemma~7(c)]{H17}, the value $\rho_2(gh^{-1})$ is even;
that lemma uses only the three-cycle equation. The Galois Kummer
image over $\Q^{\mathrm{ab}}$ is $2\Zh$, as proved in
Proposition~\ref{prop:arithmetic-affine}. Choose
$k\in G_{\Q^{\mathrm{ab}}}$ with
$\rho_2(k)=\rho_2(gh^{-1})$. The cocycle law gives
$gh^{-1}k^{-1}\in N_0$, proving the gentle factorization.
Finally $G=NH$ shows directly that a homomorphism killing $H$
kills $G$ exactly when it kills $N$.
\end{proof}

The equivalence is about equality of the two full profinite groups.
The quotient $B$ records its abelian consequence. In particular a
gentle character trivial on Galois proves strictness of $G\supset H$
only if it is nonzero on $N$; a character vanishing on the pentagon
group does not provide such a witness.

\section{Comparison of the two higher-genus actions}
The explicit action on the groups with zero or one marked point agrees,
as an outer action, with the action constructed from quilts. The following
calculation gives the change of representative.

\begin{proposition}\label{prop:ns-lns-comparison}
Use the compatible standard markings of Nakamura--Schneps
\cite[Theorem 11.2]{NS2000} and Lochak--Nakamura--Schneps
\cite[Theorem 1.1]{LNS}. For $F=(\lambda,f)\in\widehat{\operatorname{GT}}$,
write $\rho=\rho_2(F)$ and
\[
 w_1=1,\qquad w_j=(a_1\cdots a_{j-1})^j\quad(2\leq j\leq2g),
 \qquad
 K_F=\prod_{j=2}^{2g}w_j^{\,(-1)^{j+1}4\rho}.
\]
Let $A_F^{\mathrm{NS2000}}$ denote the standard quilt representative and
$A_F^{\mathrm{LNS}}$ the representative given by the explicit chain
formulas. On each of $\widehat\Gamma_{g,0}$ and
$\widehat\Gamma_{g,1}$ one has
\[
 A_F^{\mathrm{NS2000}}=\operatorname{Ad}(K_F)\circ A_F^{\mathrm{LNS}}.
\]
In particular the two outer actions are equal. If $\rho_2(F)=0$, these
standard representatives themselves are equal.
\end{proposition}
\begin{proof}
The equality of the Nakamura--Schneps subgroup with
$\widehat{\operatorname{GT}}$ is used as established background
\cite[Corollary 21]{H17}; it ensures that their formulas apply to every
$F$ in the statement.

The elements $a_1,\ldots,a_{2g}$ satisfy the braid relations. The full
braid twist $w_j$ is central in the subgroup generated by
$a_1,\ldots,a_{j-1}$. Hence the $w_j$ commute with one another, and
$w_j$ commutes with $a_i$ whenever $j\ne i$. The usual full-twist
factorization gives
\[
 w_i=y_iw_{i-1},\qquad
 y_i=a_{i-1}\cdots a_1^2\cdots a_{i-1}.
\]
The element $w_{i-1}$ commutes both with $a_i$ and with $y_i$.
If $z$ centralizes $a,b$ and $f$ belongs to the closed commutator
subgroup of a free profinite group on two generators, then
$f(a,bz)=f(a,b)$. Indeed the homomorphism from
$\langle a,b\rangle\times\langle z\rangle$ given by multiplication
sends the second coordinate of a commutator word to the identity;
the identity extends from commutator words by continuity.
Consequently
\[
 f(a_i^2,w_i)=f(a_i^2,y_i),\qquad
 f(w_i,a_i^2)=f(y_i,a_i^2).
\]

Set $c_i=f(a_i^2,y_i)$. All $w_j$ with $j\ne i$ commute with
$c_i$ and with $a_i$. The explicit chain action is
$A_F^{\mathrm{LNS}}(a_i)=c_i a_i^\lambda c_i^{-1}$.
Conjugating by $K_F$ therefore gives
\[
 w_i^{\,(-1)^{i+1}4\rho}
 f(a_i^2,w_i)a_i^\lambda f(w_i,a_i^2)
 w_i^{-(-1)^{i+1}4\rho},
\]
which is the standard quilt formula. For $a_1$ both maps give
$a_1^\lambda$, and $K_F$ commutes with $a_1$.

For $g\geq2$, the remaining generator is $d=d_2$, a curve of the
standard pants decomposition. In the same marking,
\[
 w_2=a_1^2,\qquad
 w_{2i}=d_i d_{-i}\quad(i\geq2),\qquad
 w_{2i-1}^{2}=e_{i-1}\quad(i\geq2).
\]
All curves appearing on the right belong to that pants decomposition.
Since $4\rho$ is divisible by two, these identities express $K_F$ as
a product of powers of its mutually commuting pants twists. Thus
$K_F$ commutes with $d$. Both representatives take $d$ to $d^\lambda$.
Equality on the displayed generating set proves the assertion; in
genus one the chain generators already suffice. The same identities
survive the relation filling the marked point, so the argument also
applies to $\widehat\Gamma_{g,0}$.
\end{proof}

This is an equality on the overlapping families of full profinite
mapping class groups. A comparison with another completed-tower
construction additionally requires that construction's action formulas
or a stated uniqueness theorem for its objects and morphisms. The
existence of actions with the same genus-zero restriction alone does
not identify their chosen representatives on every object of a tower.

\section{The scope of inertia-preserving identifications}
For a generic surface $S$ of type $(g,n)$, namely
$g>0$, $3g-3+n>2$, and $(g,n)\ne(2,0)$, we use the established
identification
\[
 \operatorname{Out}^{*}(\check\Gamma(S))
       \cong\widehat{\operatorname{GT}}.
\]
Here the star requires preservation of the twist inertia subgroups.
This is \cite[Corollary 22]{H17}, using Lochak's compatible-pair
classification \cite[Definition 5.4, Proposition--Definition 5.6,
and Theorem 5.10(ii)]{Lochak2025}. It is a cited result about the
procongruence completion. Passing from nonseparating inertia to all
inertia is a further assertion.

Boggi states this further assertion as
$\operatorname{Aut}^{I_0}(\check\Gamma(S))=
\operatorname{Aut}^{I}(\check\Gamma(S))$ in
\cite[Theorem 3.12 and Corollary 3.15]{Boggi}.
The proof proceeds through the nonseparating curve complex and requires
recognition of the multicurves with connected complement, as used in
\cite[Remark 3.11]{Boggi}. Commutation of individual twists recognizes
disjointness. It does not by itself recognize connectedness of the
complement: a bounding pair and a pair with connected complement both
consist of disjoint nonseparating curves. Thus the latter recognition
step requires its own justification.

\begin{proposition}\label{prop:nonunique-stabilizer-extension}
Let $\sigma=\{\gamma_1,\ldots,\gamma_{g-1}\}$ be a cut system with
connected complement, and let $H$ be its oriented stabilizer in the
full profinite pure mapping class group. Let $H_i$ be the oriented
stabilizer of $\sigma\setminus\{\gamma_i\}$. The identity
of $H$ has more than one inertia-preserving extension to $H_i$.
\end{proposition}
\begin{proof}
The twist $\tau_{\gamma_i}$ is central in $H$, since every element
of the oriented stabilizer fixes its twist. Thus the identity and
$\operatorname{Ad}(\tau_{\gamma_i})$ restrict to the same automorphism
of $H$. Both are inertia-preserving automorphisms of $H_i$.
Choose a simple curve $\beta_i$ meeting $\gamma_i$ once and disjoint
from the other curves of the cut system. Its twist belongs to $H_i$.
On the symplectic plane generated by these two curves, their twist
matrices are
\[
 U=\begin{pmatrix}1&1\\0&1\end{pmatrix},\qquad
 V=\begin{pmatrix}1&0\\-1&1\end{pmatrix}.
\]
They do not commute, already after reduction modulo three. Thus
$\tau_{\gamma_i}$ is not central in $H_i$, and its inner automorphism
is not the identity. This finite symplectic quotient also shows that
the distinction survives profinite completion.
\end{proof}

Accordingly, uniqueness as a genuine automorphism in
\cite[Lemma 6.16]{Boggi} fails as written. This counterexample does
not refute a statement of uniqueness up to inner automorphisms.
There is a separate group identification issue in the proof of
\cite[Proposition 6.9]{Boggi}: the quotient
$B_{2g+1}/\langle z^2\rangle$ obtained by filling a distinguished
Weierstrass boundary point preserves that point. It therefore gives
the corresponding stabilizer in the closed hyperelliptic group,
whose index is $2g+2$, rather than the whole closed hyperelliptic
group. Neither observation is a counterexample to the earlier
nonseparating-inertia theorem; they concern distinct proof assertions.

\section{A global character from the two elliptic frames}
\label{sec:global-frame-character}

We define the frame signs, extend their sum to every multiplier, and
fix its arithmetic normalization. The construction uses the normalized
frame vanishing theorem of \cite[Theorem~3]{H17}. The extension and
the conjugation invariance are proved here.

\subsection{The frames and their normalization}

Let $\widehat B_3=\langle u,v\mid uvu=vuv\rangle$ be the profinite
braid group and put $c=(uv)^3$. Its central quotient is
\[
 \Gamma=\widehat B_3/\overline{\langle c\rangle}
      =\langle s,r\mid s^2=r^3=1\rangle,
 \qquad s=uvu,\quad r=uv,\quad \tau=sr=u^{-1}.
\]
The free pure subgroup $F=\ker(\Gamma\to S_3)$ has generators
$x=u^2=\tau^{-2}$ and $y=v^2=sxs$. We also use the subgroup
$\langle u^2,v^2\rangle\subset\widehat B_3$, which maps isomorphically
to $F$. A gentle element $g=(\lambda_g,f_g)$ acts on the braid
group by
\[
 \Phi_g(u)=u^{\lambda_g},\qquad
 \Phi_g(v)=f_g(u^2,v^2)^{-1}v^{\lambda_g}f_g(u^2,v^2),
 \qquad \Phi_g(c)=c^{\lambda_g}.
\]
These are automorphisms and induce the identity on the permutation
quotient \cite[Section 3, before Lemma~5]{H17}.

The finite-subgroup and centralizer theorem for the free profinite
product gives unique pure frames
\[
 \Phi_g(s)=\beta_gs\beta_g^{-1}\quad(g\in\GTz),
 \qquad
 \Phi_h(r)=A_hrA_h^{-1}\quad(\lambda_h=1),
 \qquad \beta_g,A_h\in F.
\]
These are the frames of \cite[Lemma~5]{H17}. For the cubic frame,
the multiplier-one condition fixes the conjugacy class of $r$ in
$\Gamma^{\mathrm{ab}}=C_6$, distinguishing it from $r^{-1}$.
The frames vary continuously: their graphs are closed in compact
products, and uniqueness makes the corresponding projections
homeomorphisms. Their composition identities are
\begin{equation}
 \beta_{gh}=\Phi_g(\beta_h)\beta_g,
 \qquad A_{hk}=\Phi_h(A_k)A_h
       \quad(\lambda_h=\lambda_k=1).
 \label{eq:frame-cocycles}
\end{equation}

Let $\theta$ interchange $x,y$. Comparison of the two expressions
for $\Phi_g(y)$ gives
$f_g=\theta(\beta_g)\beta_g^{-1}$. Indeed their possible discrepancy
centralizes $y^{\lambda_g}$; abelianization makes its exponent zero.
Thus the two coordinates of $[\beta_g^{-1}]\in F^{\mathrm{ab}}$
are equal. Define $\rho_2(g)$ by
\[
 [\beta_g^{-1}]=(\rho_2(g),\rho_2(g)).
\]
It is a continuous cocycle,
\[
 \rho_2(gh)=\rho_2(g)+\lambda_g\rho_2(h).
\]
Its restriction to $G_{\mathbb Q}$ is the Kummer cocycle of the
coherent positive roots of two \cite[Lemma~7]{H17}.
Write
\[
 \mathcal G=\GTz,\qquad
 N_0=\{g\in\mathcal G:\lambda_g=1,\ \rho_2(g)=0\},
 \qquad H=G_{\mathbb Q}.
\]
The cocycle makes $N_0$ closed and normal.

We will use the arithmetic decomposition
\begin{equation}
 \mathcal G=N_0H,\qquad
 N_0\cap H=G_{K_2},\qquad
 K_2=\mathbb Q(\mu_\infty,2^{1/\infty}).
 \label{eq:gentle-arithmetic-normalization}
\end{equation}
This is Proposition~\ref{prop:normalized-comparison}; its intersection formula follows from the two defining Galois actions.

\subsection{The signed matrices at level twelve}

Put
\[
 U=\begin{pmatrix}1&1\\0&1\end{pmatrix},\qquad
 L=\begin{pmatrix}1&0\\-1&1\end{pmatrix},\qquad
 S=ULU,\qquad R_+=UL,\qquad R=-R_+.
\]
The assignments $u\mapsto U$, $v\mapsto L$ define the signed
representation $\widetilde q_\ell:\widehat B_3\to
\operatorname{SL}_2(\mathbb Z_\ell)$ for every prime $\ell$.
Its restriction to $F$ is denoted $q_\ell$. Reducing the signed
representation modulo twelve gives $q_{12}$ on $F$ and a projective
representation $q$ of $\Gamma$.
The signed pure image is
\[
 E_{12}=V_4\times\operatorname{SL}_2(\mathbb F_3),
 \qquad V_4=\{I,U^2,L^2,U^2L^2\}\pmod4.
\]
The matrix $-I$ is not in $V_4$, so this image embeds in the
projective group modulo twelve. Its projective centralizers of
$R_+$ and $S$ are respectively
\begin{equation}
 \{I\}\times\bigl(\langle-I\rangle\times\langle R\rangle\bigr),
 \qquad
 \{I,U^2L^2\}\times\langle S\rangle.
 \label{eq:finite-frame-centralizers}
\end{equation}
These follow directly by multiplying the indicated matrices; they
are the finite groups in \cite[Lemma~12]{H17}.

\begin{lemma}\label{lem:frame-all-multipliers-congruence}
For every $g\in\mathcal G$ and every prime $\ell$ there is a unique
matrix
\[
 D_{g,\ell}=\begin{pmatrix}\lambda_{g,\ell}&b_{g,\ell}\\0&1\end{pmatrix}
\]
such that
$\widetilde q_\ell\Phi_g=\operatorname{Ad}(D_{g,\ell})\widetilde q_\ell$.
It depends continuously on $g$, and $D_{gh,\ell}=D_{g,\ell}D_{h,\ell}$.
For $h\in N_0$, its reduction modulo twelve is the identity.
\end{lemma}

\begin{proof}
Put $a=\lambda_{g,\ell}$ and
$M=\widetilde q_\ell(\Phi_g(v))=
\left(\begin{smallmatrix}\alpha&b\\d&2-\alpha\end{smallmatrix}\right)$.
It is conjugate to $L^a$, so its trace is two and its determinant
is one. The lower-left entry of the braid relation
$U^aMU^a=MU^aM$ gives $d(1+ad)=0$.
If $d=0$, the domain property of $\mathbb Z_\ell$, together with
the trace and determinant, makes $M$ upper unipotent. The images
of both braid generators would commute, contradicting surjectivity
onto $\operatorname{SL}_2(\mathbb Z_\ell)$. Hence $d=-a^{-1}$
and $b=a(\alpha-1)^2$. Direct multiplication gives
\[
 D_{g,\ell}=
  \begin{pmatrix}a&a(1-\alpha)\\0&1\end{pmatrix}.
\]
This proves existence and continuity. A matrix commuting with both
$U,L$ is scalar; the specified lower-right entry gives uniqueness.
Composition then gives the product identity.

For $\lambda_h=1$, write the reduction modulo twelve as $U^{t_h}$.
Modulo four the conjugator $f_h\in F'$ maps trivially because
$V_4$ is abelian, so $t_h=0\pmod4$. Modulo three,
$U^{-t_h}q_3(\beta_h)$ lies in the quaternion subgroup of
$\operatorname{SL}_2(\mathbb F_3)$: its projective class centralizes
the involution represented by $S$. In the quotient of order three,
the classes of $U$ and $L$ agree, whereas
$[\beta_h]=(-\rho_2(h),-\rho_2(h))$ gives class
$-4\rho_2(h)=-\rho_2(h)$. Thus
$t_h=-\rho_2(h)\pmod3$. This is the shear formula of
\cite[Lemma~13]{H17}; in particular $D_h=I\pmod{12}$ on $N_0$.
\end{proof}

Define the closed subgroup
\[
 H_0=\{h\in\mathcal G:\lambda_h=1,\ t_h=0\pmod{12}\}.
\]
It contains $N_0$. For $h\in H_0$, the two frames lie in
\eqref{eq:finite-frame-centralizers}, so define
\begin{equation}
 q_3(A_h)\in(-I)^{a_3(h)}\langle R\rangle,
 \qquad
 q_3(\beta_h)\in S^{a_2(h)}\langle-I\rangle,
 \qquad a_2(h),a_3(h)\in\mathbb F_2.
 \label{eq:frame-sign-definition}
\end{equation}
These are exactly the signs of \cite[Definition~14]{H17}.
The cocycles \eqref{eq:frame-cocycles}, the identity action on
$E_{12}$, and the abelian groups in
\eqref{eq:finite-frame-centralizers} make $a_2,a_3$ continuous
characters on $H_0$. Put $\delta=a_2+a_3$.

The sign convention can also be recorded by a single unit. Let
$\xi(h)\in(\mathbb Z/12\mathbb Z)^\times$ be determined by
\[
 \xi(h)\equiv(-1)^{a_2(h)}\pmod4,\qquad
 \xi(h)\equiv(-1)^{a_3(h)}\pmod3.
\]
For the quadratic Dirichlet character $\chi_{12}$ the complete table is
\[
 \begin{array}{c|rrrr}
  t&1&5&7&11\\ \hline
  \chi_{12}(t)&1&-1&-1&1
 \end{array}
 \qquad (-1)^{\delta(h)}=\chi_{12}(\xi(h)).
\]
All elements of $H_0$ have multiplier one. The unit $\xi$ records
the two frame signs separately from that multiplier.

\subsection{Conjugation by elements of every multiplier}

\begin{proposition}\label{prop:frame-sign-conjugation}
The restrictions of $a_2,a_3$, and hence $\delta$, to $N_0$ are
invariant under conjugation by $\mathcal G$.
\end{proposition}

\begin{proof}
Let $h\in N_0$, $g\in\mathcal G$, and $\eta=ghg^{-1}$.
Normality gives $\eta\in N_0$. All matrix calculations in this
proof are modulo three. Write $\overline D_g$ for the reduction
of $D_{g,3}$.

Applying $\Phi_\eta\Phi_g=\Phi_g\Phi_h$ to $s$ and comparing
unique pure frames gives
\[
 \beta_\eta=\Phi_\eta(\beta_g)^{-1}\Phi_g(\beta_h)\beta_g.
\]
Since $\overline D_\eta=I$,
\[
 q_3(\beta_\eta)=q_3(\beta_g)^{-1}\overline D_g
      q_3(\beta_h)\overline D_g^{-1}q_3(\beta_g).
\]
Thus the traces of their squares agree. The defining groups give
\[
 \operatorname{tr}(q_3(\beta_h)^2)=2(-1)^{a_2(h)},
\]
which proves invariance of $a_2$.

Put $\epsilon\in\{1,-1\}$ equal to $\lambda_g$ modulo three.
The torsion theorem gives
$\Phi_g(r)=\gamma_g r^\epsilon\gamma_g^{-1}$.
The exponent follows from the action on $\Gamma^{\mathrm{ab}}=C_6$.
Since $\Phi_g$ induces the identity on $S_3$, multiplying
$\gamma_g$ by a power of $r$ permits the form
$\gamma_g=B_gs^e$, with $B_g\in F$ and $e=0$ or one.
Here $e=0$ when $\epsilon=1$ and $e=1$ when $\epsilon=-1$.
Comparison of the two actions on $r$ gives
\begin{equation}
 A_\eta=\Phi_\eta(\gamma_g)^{-1}\Phi_g(A_h)\gamma_g.
 \label{eq:cubic-frame-conjugation}
\end{equation}
The right side is pure and conjugates $r$ to $\Phi_\eta(r)$,
so equality follows from uniqueness of the pure frame.

We check this equality in the signed representation. Lift
$\gamma_g$ to $\widetilde\gamma_g=B_g(uvu)^e$ in
$\widehat B_3$. Conjugation by $uvu$ interchanges $x,y$ and
normalizes the embedded free subgroup. Since $\lambda_\eta=1$,
\[
 \Phi_\eta(uvu)=\beta_\eta(uvu)\beta_\eta^{-1}.
\]
Both sides project to the same modular element and have square $c$.
A central discrepancy would have square one, and the central
procyclic group $\overline{\langle c\rangle}\simeq\widehat{\mathbb Z}$
is torsion free. This proves the equality.
It follows that the lifted right side of
\eqref{eq:cubic-frame-conjugation} belongs to the embedded free
subgroup: the two factors involving $uvu$ cancel through its
normalizing action. Its projection equals $A_\eta$, and the free
subgroup injects into the central quotient, so it is the signed
lift of $A_\eta$ itself.

Applying the signed matrix representation therefore gives
\[
 q_3(A_\eta)=
 \widetilde q_3(\widetilde\gamma_g)^{-1}\overline D_g
 q_3(A_h)\overline D_g^{-1}
 \widetilde q_3(\widetilde\gamma_g).
\]
The trace is unchanged. Every power of $R$ has trace two in
$\mathbb F_3$, so
$\operatorname{tr}(q_3(A_h))=2(-1)^{a_3(h)}$.
This proves invariance of $a_3$ as well.
\end{proof}

\begin{theorem}[Arithmetic normalization of the frame character]
\label{thm:global-frame-character}
There is a unique continuous character
\[
 \Delta:\mathcal G\longrightarrow\mathbb F_2
\]
which restricts to $\delta$ on $N_0$ and is zero on $G_{\mathbb Q}$.
For any $g\in\mathcal G$, choose $\sigma\in G_{\mathbb Q}$
with $(\lambda,\rho_2)(\sigma)=(\lambda,\rho_2)(g)$. Then
\[
 \Delta(g)=\delta(g\sigma^{-1}).
\]
This character vanishes on the full pentagon group.
\end{theorem}

\begin{proof}
By \cite[Theorem~3]{H17}, $\delta$ vanishes on the normalized
pentagon subgroup. In particular it vanishes on
$N_0\cap H=G_{K_2}$. Using
\eqref{eq:gentle-arithmetic-normalization}, define
$\Delta(n\sigma)=\delta(n)$ for $n\in N_0$, $\sigma\in H$.
Two decompositions differ by an element of $N_0\cap H$, so this
is well defined. Proposition~\ref{prop:frame-sign-conjugation}
proves its homomorphism law. The multiplication map
$N_0\times H\to\mathcal G$ is a quotient map, since it is a
continuous surjection from a compact space to a Hausdorff space.
Thus the character is continuous. The two specified restrictions
and the decomposition prove uniqueness, as well as the displayed
formula. Finally \cite[Theorem~2]{H17} decomposes the pentagon
group into its normalized subgroup and $G_{\mathbb Q}$, both of
which have zero image. Hence $\Delta$ vanishes on $\GT$.
\end{proof}

\section{The individual frame signs and the Kummer cocycle of three}
\label{sec:harmonic-kummer-comparison}

The two normalized frame signs can be identified separately on the
pentagon group. The comparison uses Anderson duplication, already a
consequence of the cited distribution and braid identities.

\begin{theorem}\label{thm:individual-frame-kummer}
For every $g\in N=\ker(\lambda,\rho_2)\subset\widehat{\mathrm{GT}}$,
\[
 a_2(g)=a_3(g)=\frac{\rho_3^I(g)}2\pmod2.
\]
In particular the character
$a_2+\rho_3^I/2\bmod2$ is zero on $N$. Both frame signs vanish
on the simultaneous radical-normalized subgroup $N_{\mathrm{rad}}$.
\end{theorem}
\begin{proof}
Let $g=(1,f)$ and put $F=\widehat F_2=\langle x,y\rangle$.
Write $c=[x,y]=xyx^{-1}y^{-1}$. The module $F'/F''$ is free of
rank one over
\[
 \mathcal A=\widehat{\mathbb Z}[[F^{\mathrm{ab}}]],
\]
with generator $[c]$; see \cite[\S2.3, p.~101]{Furusho2003}.
Let $X,Y$ denote the two conjugation operators and write
$[f]=A_f(X,Y)[c]$. The two-cycle identity gives
\begin{equation}
 A_f(X,Y)=A_f(Y,X),
 \label{eq:harmonic-A-symmetry}
\end{equation}
because interchange of $x,y$ sends $[c]$ to $-[c]$.
The action $x\mapsto x$, $y\mapsto f^{-1}yf$ gives
\[
 [\Phi_g(c)]=[c]+(X-1)(Y-1)[f].
\]
Thus its multiplier on this module is
\begin{equation}
 B_g(X,Y)=1+(X-1)(Y-1)A_f(X,Y).
 \label{eq:harmonic-beta-coefficient}
\end{equation}
For multiplier one this is also the adelic beta function, since
the normalizing cyclotomic factors in \cite[\S2.3]{Furusho2003}
are equal to one.

Work in $\mathcal O=\mathbb Z_2[\omega]$, where
$\omega^2+\omega+1=0$. Let $\nu_g$ be the triangle-group lattice
unit of \cite[Lemma 15]{H17}. Its relation to the beta multiplier is
\begin{equation}
 \nu_g=B_g(\omega,\omega),\qquad
 \nu_g\equiv(-1)^{a_2(g)}\pmod{4\mathcal O}.
 \label{eq:harmonic-triangle-beta}
\end{equation}
For the first equality, in the affine triangle quotient the two pure
generators have multiplier $\omega$ and $[x,y]$ is the translation
by $-3$. The action on translations is multiplication by $\nu_g$.
Applying the module action to $[x,y]$ therefore gives
$-3\nu_g=-3B_g(\omega,\omega)$, and $3$ is a unit in
$\mathcal O$. The second equality is the cited lattice calculation.

Since $\rho_2^I(g)=\rho_2(g)=0$, the Anderson duplication identity
at the torsion parameter $s=1/3$ gives
\begin{equation}
 B_g(\omega,\omega)=B_g(-1,\omega).
 \label{eq:harmonic-duplication}
\end{equation}
Indeed in the notation of \cite[Theorem 3.1]{Furusho2003} the
quotient of these two values is the third root of unity with exponent
$\Psi_2^{(0)}(g)=-\rho_2^I(g)$, which is zero here. The identity
is in the full adelic beta function; its evaluation in
$\mathcal O$ is therefore legitimate.

It remains to compute the right side modulo four. In the degree-three
cyclic-cover subgroup with basis $x^3,y_0,y_1,y_2$, write the
abelianized image of $f$ as
$c_0y_0+c_1y_1+c_2y_2$. The $x^3$ coefficient is zero.
The cyclic-cover symmetry of
Lemma~\ref{lem:cyclic-cover-symmetry} gives $c_1=c_2$;
abelianizing in $F$ gives $c_0+c_1+c_2=0$.
Consequently, for $b=\ell_3(f)/2=-\rho_3^I(g)/2$,
\[
 (c_0,c_1,c_2)=(2b,-b,-b).
\]
Map this cover homology to $\mathcal O$ by
$y_j\mapsto\omega^j$, $x^3\mapsto0$.
Conjugation by $x$ acts as multiplication by $\omega$ and
conjugation by $y$ acts trivially. The commutator $c$ maps to
$\omega-1$. Applying the map to $[f]=A_f[c]$ therefore gives
\begin{equation}
 (\omega-1)A_f(\omega,1)
   =2b-b\omega-b\omega^2=3b.
 \label{eq:harmonic-cover-coefficient}
\end{equation}

The evaluations at $(-1,\omega)$ and $(1,\omega)$ agree
modulo $2\mathcal O$. By \eqref{eq:harmonic-A-symmetry},
$A_f(1,\omega)=A_f(\omega,1)$. Using
\eqref{eq:harmonic-beta-coefficient} and
\eqref{eq:harmonic-cover-coefficient}, we obtain
\[
 \begin{split}
 \nu_g
 &=B_g(-1,\omega)\\
 &\equiv1-2(\omega-1)A_f(\omega,1)\\
 &=1-6b\equiv1+2b\pmod{4\mathcal O}.
 \end{split}
\]
Together with \eqref{eq:harmonic-triangle-beta}, this proves
$a_2(g)=b=\rho_3^I(g)/2\pmod2$.
The established vanishing $a_2+a_3=0$ on $N$
\cite[Theorem 3]{H17} gives the assertion for $a_3$.
Every rational-prime Kummer cocycle is zero on $N_{\mathrm{rad}}$,
so both signs vanish there.
\end{proof}

The comparison identifies these particular finite-frame coordinates.
It does not determine the residual quotient $B$ or its quotient $B/2B$.

\section{The full adelic image of the two elliptic frames}
\label{sec:adelic-cm-frames}

The two corrected elliptic frames detect no further relative abelian
characters on the profinite Grothendieck--Teichm\"uller group.  The
reason is an equality of full adelic images: the multiplier-one
subgroup and the cyclotomic Galois subgroup have the same image in the
product of the two CM tori.  We prove the integral comparison, identify
the one quadratic entanglement between the division fields, and obtain
a second normalization of the main character problem.

\subsection{Corrected frames and their integral targets}

Write $\mathcal G=\widehat{\mathrm{GT}}_0$ for the gentle group,
$G=\widehat{\mathrm{GT}}$, $H=G_{\mathbb Q}$, and
\[
 M_0=\ker(\lambda:\mathcal G\longrightarrow
                       \widehat{\mathbb Z}^{\times}),
 \qquad H^{\mathrm{cyc}}=G_{\mathbb Q^{\mathrm{ab}}}.
\]
We recall the normalization of the frames.  In the profinite braid
group put
\[
 \widehat B_3=\langle u,v\mid uvu=vuv\rangle,\qquad
 c=(uv)^3,\qquad s=uvu,\quad r=uv,
\]
\[
 \Gamma=\widehat B_3/\overline{\langle c\rangle}
             =\langle s,r\mid s^2=r^3=1\rangle.
\]
The free pure subgroup is
$F=\ker(\Gamma\to S_3)=\widehat F(x,y)$, where
$x=u^2$ and $y=v^2$.  The subgroup generated by these two elements
in $\widehat B_3$ identifies with $F$.  For $g\in\mathcal G$ its
braid action is
\[
 \Phi_g(u)=u^{\lambda_g},\qquad
 \Phi_g(v)=f_g^{-1}v^{\lambda_g}f_g,\qquad
 \Phi_g(c)=c^{\lambda_g},\qquad f_g\in F'.
\]
The unique pure frames satisfy
\[
 \Phi_g(s)=\beta_gs\beta_g^{-1},\qquad
 \Phi_h(r)=A_hrA_h^{-1}\quad(h\in M_0),
 \qquad \beta_g,A_h\in F.
\]
Their existence, uniqueness, and cocycles are as in
\cite[Lemmas~5 and~7]{H17}; in particular
\begin{equation}
 \beta_{gh}=\Phi_g(\beta_h)\beta_g,\qquad
 A_{gh}=\Phi_g(A_h)A_g\quad(g,h\in M_0),\qquad
 [\beta_g^{-1}]=(\rho_2(g),\rho_2(g)).
 \label{eq:adelic-frame-cocycles}
\end{equation}
The parity result gives $\rho_2(M_0)\subset2\widehat{\mathbb Z}$.

Set
\[
 U=\begin{pmatrix}1&1\\0&1\end{pmatrix},\quad
 L=\begin{pmatrix}1&0\\-1&1\end{pmatrix},\quad
 S=ULU=\begin{pmatrix}0&1\\-1&0\end{pmatrix},\quad
 R=-UL=\begin{pmatrix}0&-1\\1&-1\end{pmatrix}.
\]
Let $\widetilde q:\widehat B_3\to
\operatorname{SL}_2(\widehat{\mathbb Z})$ send $u,v$ to $U,L$,
and let $q$ be its restriction to $F$.  The intertwiner proved above
has the form
\begin{equation}
 D_g=\begin{pmatrix}\lambda_g&b_g\\0&1\end{pmatrix},\qquad
 \widetilde q\Phi_g=\operatorname{Ad}(D_g)\widetilde q,\qquad
 D_{gh}=D_gD_h.
 \label{eq:adelic-frame-intertwiner}
\end{equation}
Here all entries are adelic.  Its existence can also be seen directly,
prime by prime: if
$\widetilde q(\Phi_g(v))=\left(\begin{smallmatrix}
 a&b\\d&2-a\end{smallmatrix}\right)$, the braid equation gives
$d(1+\lambda_gd)=0$.  The zero possibility would make the two
braid images commute, contrary to their surjectivity onto
$\operatorname{SL}_2(\mathbb Z_\ell)$.  Thus
$d=-\lambda_g^{-1}$, $b=\lambda_g(a-1)^2$, and
$D_g=\left(\begin{smallmatrix}\lambda_g&\lambda_g(1-a)\\0&1
\end{smallmatrix}\right)$.  Its normalization makes it unique.

Put $O_3=\mathbb Z[\omega]$, where $\omega^2+\omega+1=0$,
and $O_4=\mathbb Z[i]$.  Embed these rings in matrices by
$\omega\mapsto R$ and $i\mapsto S$.  A bar denotes their quadratic
involution.  Define
\begin{equation}
 N_j=\{a\in(O_j\otimes\widehat{\mathbb Z})^{\times}:
                                      a\bar a=1\},\qquad
 N_j(4)=\{a\in N_j:a_2\equiv1\pmod4\}.
 \label{eq:adelic-cm-targets}
\end{equation}
The subscript $2$ in this formula means the component at the prime two.

\begin{lemma}\label{lem:adelic-corrected-frames}
The formulas
\begin{equation}
 T_g=q(A_g)^{-1}D_g,\qquad
 W_g=q(\beta_g)^{-1}D_g\qquad(g\in M_0)
 \label{eq:adelic-corrected-frames}
\end{equation}
define continuous homomorphisms
$T:M_0\to N_3(4)$ and $W:M_0\to N_4(4)$.
\end{lemma}

\begin{proof}
The centralizations are signed equalities.  Indeed the two lifts of
$\Phi_g(s)=\beta_gs\beta_g^{-1}$ in $\widehat B_3$ differ by
$c^e$.  Squaring, and using $s^2=c$ and $\Phi_g(c)=c$, gives
$2e=0$ in $\widehat{\mathbb Z}$, hence $e=0$.  Cubing the
corresponding identity for $r$, with $r^3=c$, similarly gives
$3e=0$.  Thus $T_g$ commutes with $R$ and $W_g$ commutes with $S$
as integral matrices, without a scalar ambiguity.

The centralizer rings of $R$ and $S$ in
$\operatorname{Mat}_2(\widehat{\mathbb Z})$ are respectively
$\widehat{\mathbb Z}[R]$ and $\widehat{\mathbb Z}[S]$.
For example $aI+bR$ has determinant $a^2-ab+b^2$, and
$aI+bS$ has determinant $a^2+b^2$.  Since $\det D_g=1$ on
$M_0$, this places $T_g,W_g$ in $N_3,N_4$.
Using \eqref{eq:adelic-frame-cocycles} and
\eqref{eq:adelic-frame-intertwiner} gives, with the order fixed,
\[
 q(A_{gh})^{-1}D_{gh}
 =q(A_g)^{-1}D_gq(A_h)^{-1}D_h=T_gT_h;
\]
the same calculation applies to $W$.

Modulo four the image of $F$ is the abelian four-group
\[
 V_4=\{I,U^2,L^2,U^2L^2\}.
\]
Consequently $q_4(f_g)=1$.  The braid covariance at $v$ makes
$D_g\bmod4$ commute with $L$, which forces the upper-right entry
of $D_g$ to be zero modulo four.  Since $q_4(A_g)\in V_4$
commutes with $R$ and $C_{V_4}(R)=\{I\}$, one has
$T_g\equiv I\pmod4$.  Finally the two abelian coordinates of
$\beta_g$ are $-\rho_2(g)$, both even.  The map $q_4$ factors
through their reductions modulo two, so $q_4(\beta_g)=I$ and
$W_g\equiv I\pmod4$ as well.  All maps used are continuous.
\end{proof}

\begin{lemma}\label{lem:adelic-kernel-normality}
The joint kernel of $T,W$ in $M_0$ is normal in $\mathcal G$.
\end{lemma}

\begin{proof}
The formula $q(\beta_g)^{-1}D_g$ defines a homomorphism to
$\operatorname{GL}_2(\widehat{\mathbb Z})$ for every $g\in\mathcal G$:
the calculation in the preceding proof uses the beta cocycle,
which is valid for every multiplier.  Thus its kernel is normal.

For the cubic frame, fix $g\in\mathcal G$, $h\in M_0$, and
$\eta=ghg^{-1}$.  Put $\epsilon\in\{1,-1\}$ equal to
$\lambda_g$ modulo three.  The torsion theorem gives a conjugator
$\gamma_g$ with
$\Phi_g(r)=\gamma_g r^\epsilon\gamma_g^{-1}$.
Because $\Phi_g$ induces the identity on $S_3$, this conjugator
may be chosen as $\gamma_g=B_gs^e$ with $B_g\in F$ and
$e\in\{0,1\}$.  Comparing
$\Phi_\eta\Phi_g(r)$ with $\Phi_g\Phi_h(r)$ and using the
uniqueness of the pure frame gives
\[
 A_\eta=\Phi_\eta(\gamma_g)^{-1}\Phi_g(A_h)\gamma_g.
\]
Lift $\gamma_g$ to $\widetilde\gamma_g=B_g(uvu)^e$ in
$\widehat B_3$.  Since $\lambda_\eta=1$, the equality
$\Phi_\eta(uvu)=\beta_\eta(uvu)\beta_\eta^{-1}$ is signed,
by the squaring argument above.  Also $uvu$ normalizes the embedded
free group.  Consequently the displayed expression, formed with
these braid lifts, belongs to that free group and is its unique
lift of $A_\eta$.  Writing
$Q_g=\widetilde q(\widetilde\gamma_g)$, covariance therefore yields
\[
 q(A_\eta)=D_\eta Q_g^{-1}D_\eta^{-1}
                 D_gq(A_h)D_g^{-1}Q_g.
\]
As $D_\eta=D_gD_hD_g^{-1}$, inversion and multiplication by
$D_\eta$ give
\begin{equation}
 T_\eta=(Q_g^{-1}D_g)\,T_h\,(Q_g^{-1}D_g)^{-1}.
 \label{eq:adelic-cubic-kernel-conjugation}
\end{equation}
Thus the kernel of $T$ in $M_0$ is also stable under every
conjugation by $\mathcal G$.  Intersecting the two kernels proves
the assertion.
\end{proof}

There are two parity maps on the target.  Let
$\pi_3=\sqrt{-3}$, an associate of $1-\omega$.  For $t\in N_3(4)$
and $w\in N_4(4)$ define
\begin{equation}
 (-1)^{\epsilon_3(t)}=t_3\bmod\pi_3\in\mathbb F_3^{\times},\qquad
 w_3\bmod3=S^a\quad\Longrightarrow\quad
 \epsilon_4(w)=a\bmod2.
 \label{eq:adelic-parities}
\end{equation}
The latter is well defined because the norm-one group in
$(O_4/3)^{\times}=\mathbb F_9^{\times}$ is cyclic of order four.
Both maps are surjective.  Write
\begin{equation}
 \mathcal E=\{(t,w)\in N_3(4)\times N_4(4):
                       \epsilon_3(t)=\epsilon_4(w)\}.
 \label{eq:adelic-equal-parity}
\end{equation}

\subsection{The CM division fields and the normalized images}

Let $K_j=O_j\otimes\mathbb Q$, and consider
\[
 E_3:Y^2=X^3+1,\qquad E_4:Y^2=X^3-X.
\]
Their adelic Tate representations are denoted $\rho_3,\rho_4$.
Over $K_j$ they take values in $(O_j\otimes\widehat{\mathbb Z})^
\times$.  We use the following precise CM input: the torsion
reciprocity theorem identifies the action modulo geometric
automorphisms with the ray class action
\cite[Theorem~12.27, printed p.~124]{MilneMF110}.
Since both rings have class number one, the ray class exact sequence
identifies the inverse limit with
$(O_j\otimes\widehat{\mathbb Z})^\times/O_j^\times$.
The action on the Tate module is abelian over $K_j$
\cite[Remark~12.12 and Proposition~12.13, printed p.~119]{MilneMF110}.

Put $\alpha=3^{1/4}$ and $r_0=(3/4)^{1/4}$, with positive real
roots, and define on $H^{\mathrm{cyc}}$
\[
 \kappa(\sigma)=\frac{\sigma(r_0)}{r_0}
               =\frac{\sigma(\alpha)}{\alpha}\in\{1,-1\}.
\]
The equality follows from $r_0=\alpha/\sqrt2$.

\begin{proposition}\label{prop:adelic-cm-images}
The full division fields are
\[
 \mathbb Q(E_j[\mathrm{tors}])=K_j^{\mathrm{ab}}.
\]
On $H^{\mathrm{cyc}}$, one has
\begin{equation}
 \rho_4(H^{\mathrm{cyc}})=N_4(4),\qquad
 \rho_3^{\circ}(H^{\mathrm{cyc}})=N_3(4),\qquad
 \rho_3^{\circ}=\kappa\rho_3.
 \label{eq:adelic-arithmetic-projections}
\end{equation}
The representations $\rho_3^{\circ}$ and $\rho_3$ have the same
kernel on $H^{\mathrm{cyc}}$.
\end{proposition}

\begin{proof}
First determine the images over the CM fields.  Put $\pi_4=1+i$.
The eight points of $E_4[\pi_4^3]$ are rational over $K_4$:
besides the four points of $E_4[2]$, they have coordinates
\[
 (i,\pm(1-i)),\qquad(-i,\pm(1+i)).
\]
For instance, with $[i](X,Y)=(-X,iY)$, the addition law gives
$(1+i)(i,1-i)=(1,0)$, then $(1+i)(1,0)=(0,0)$, and then zero.
The other three points give the same conclusion.  The eight points
exhaust the kernel because the endomorphism has degree eight.
The four roots of unity map bijectively to
$(O_4/\pi_4^3)^{\times}$.  Hence
\[
 C_4=\{u\in\widehat O_4^{\times}:u\equiv1\pmod{\pi_4^3}\}
       \longrightarrow\widehat O_4^{\times}/\mu_4
\]
is an isomorphism.  Rationality puts $\rho_4(G_{K_4})$ inside
$C_4$, and CM reciprocity makes its image modulo $\mu_4$ full.
Thus $\rho_4(G_{K_4})=C_4$.

For $E_3$, the two-torsion abscissae are
$-1,-\omega,-\omega^2$, and
$E_3[1-\omega]=\{O,(0,1),(0,-1)\}$.  These coprime kernels
make $E_3[2\pi_3]$ rational over $K_3$.
The reduction of $\mu_6$ is an isomorphism onto
\[
 (O_3/(2\pi_3))^{\times}
              \simeq\mathbb F_4^{\times}\times\mathbb F_3^{\times}.
\]
The same argument gives
\[
 \rho_3(G_{K_3})=C_3
       =\{u\in\widehat O_3^{\times}:u\equiv1\pmod{2\pi_3}\}.
\]

The CM theorem implies that the action on torsion points modulo
geometric automorphisms already cuts out $K_j^{\mathrm{ab}}$.
Since the full torsion action is
abelian over $K_j$, the full division field over $K_j$ is the same
field.  Moreover $K_3$ occurs among the two-torsion coordinates,
and $K_4$ occurs by the Weil pairing on four-torsion.  This proves
the displayed division-field identities over $\mathbb Q$.

Complex conjugation acts on $C_j$ by $u\mapsto\bar u$ and has
order two.  Thus the closed derived subgroup of the full Galois
image is $\{u/\bar u:u\in C_j\}$; this is already closed, as
the image of a compact group.  Kronecker--Weber identifies
$H^{\mathrm{cyc}}$ with the closed derived subgroup of $H$.
Consequently
\begin{equation}
 \rho_j(H^{\mathrm{cyc}})=\{u/\bar u:u\in C_j\}.
 \label{eq:adelic-ratio-image}
\end{equation}

At split primes the ratio map on units is
$(a,b)\mapsto(a/b,b/a)$ and is onto the local norm-one group.
At an unramified quadratic prime, Hilbert~90 supplies a preimage;
multiplication by a rational uniformizer power makes it a unit.
At the inert prime two for $K_3$, a norm-one unit $v\equiv1\pmod2$
has a unit preimage whose residue is fixed by the residue-field
involution.  Its residue therefore belongs to
$\mathbb F_2^\times=\{1\}$, so the preimage is also one modulo two.
At the ramified prime three, a Hilbert--90 preimage of a norm-one
unit $v\equiv1\pmod{\pi_3}$ has even $\pi_3$-valuation:
an odd valuation contributes the residue
$\pi_3/\bar\pi_3=-1$.  Multiplying by a rational power of three
and then by a sign makes it a unit congruent to one modulo $\pi_3$.
The converses follow by reduction.  Thus
\begin{equation}
 \rho_3(H^{\mathrm{cyc}})
 =\{v\in N_3:v_2\equiv1\pmod2,\quad
                     v_3\equiv1\pmod{\pi_3}\}=:M_3.
 \label{eq:adelic-unnormalized-cubic-image}
\end{equation}

For $K_4$ at two, a unit $u=a+bi\equiv1\pmod{\pi_4^3}$ has
$b$ even, so $u/\bar u\equiv1\pmod4$.  Conversely if
$v\bar v=1$ and $v\equiv1\pmod4$, then
$b_0=(1+v)/2$ is a unit congruent to one modulo two and
$b_0/\bar b_0=v$.  The two unit classes modulo $\pi_4^3$ that
are one modulo two are $1$ and $-1$.  Hence one of $b_0,-b_0$
belongs to $1+\pi_4^3O_{4,2}$.  Together with the other local
calculations, this proves $\rho_4(H^{\mathrm{cyc}})=N_4(4)$.

We now determine the cubic normalization without an unspecified
four-division comparison.  The norm-one units modulo four in
$O_3/4$ that are one modulo two are exactly $1,-1$:
write a unit as $1+2a$, and use
$\operatorname{Tr}_{\mathbb F_4/\mathbb F_2}(a)=0$.
Let $k=\mathbb Q^{\mathrm{ab}}$.  If
$Y^2=\prod_{i=1}^3(X-e_i)$ has all its roots in $k$, the halves
of $(e_i,0)$ are obtained from
\[
 a^2=e_i-e_j,\quad b^2=e_i-e_k,\qquad
 (X,Y)=(e_i+ab,\,ab(a+b)),
\]
with the possible signs.  Conversely the ratios $Y/(X-e_i)$
for the choices $(a,b)$ and $(-a,b)$ recover $a+b$ and $b-a$,
and hence $a,b$.  Thus its four-division field over $k$ is
generated by the square roots of the root differences.
For $E_3$ every nonzero root difference is a root of unity times
$\sqrt{-3}$.  Since $k$ contains every root of unity,
\begin{equation}
 k(E_3[4])=k(3^{1/4}).
 \label{eq:adelic-cubic-four-division}
\end{equation}
This extension is nontrivial: the real field
$\mathbb Q(3^{1/4})$ misses the nonreal conjugates of its
Eisenstein polynomial $X^4-3$, so it cannot lie in an abelian
extension of $\mathbb Q$.
It follows that the modulo-four sign of $\rho_3(\sigma)$ is
exactly $\kappa(\sigma)$.

Multiplication by this sign gives an isomorphism
$M_3\to N_3(4)$.  Its inverse sends $w\in N_3(4)$ to
$\varepsilon(w_3)w$, where
$\varepsilon(w_3)\in\{1,-1\}$ is its residue modulo $\pi_3$.
Indeed this makes the three-component one modulo $\pi_3$, while
the modulo-four sign at two becomes $\varepsilon(w_3)$.
The two maps are inverse homomorphisms.  This proves both the
normalized cubic image and the equality of kernels.
\end{proof}

\begin{lemma}\label{lem:adelic-cm-entanglement}
The intersection of the two full division fields is
\begin{equation}
 K_3^{\mathrm{ab}}\cap K_4^{\mathrm{ab}}
       =\mathbb Q^{\mathrm{ab}}(3^{1/4}).
 \label{eq:adelic-cm-intersection}
\end{equation}
In the representations of Proposition~\ref{prop:adelic-cm-images},
the joint image of $H^{\mathrm{cyc}}$ is exactly $\mathcal E$.
\end{lemma}

\begin{proof}
Put $L=K_3^{\mathrm{ab}}\cap K_4^{\mathrm{ab}}$.
It contains $\mathbb Q^{\mathrm{ab}}$ and is Galois over
$\mathbb Q$.  In $J=\operatorname{Gal}(L/\mathbb Q)$, the two
distinct index-two subgroups
$J_3=\operatorname{Gal}(L/K_3)$ and
$J_4=\operatorname{Gal}(L/K_4)$ are abelian.
Their intersection $Z$ is central, since $J_3,J_4$ generate $J$.
Choose $a\in J_3\setminus Z$ and $b\in J_4\setminus Z$.
Then $J=\langle Z,a,b\rangle$, $a^2,b^2\in Z$, and
$J'$ is generated by the central commutator $[a,b]$.
The identity $[a^2,b]=[a,b]^2=1$ makes $|J'|\leq2$.
Since $L$ contains the maximal abelian extension of $\mathbb Q$,
this proves $[L:\mathbb Q^{\mathrm{ab}}]\leq2$.

The splitting field $\mathbb Q(i,\alpha)$ of $X^4-3$ has dihedral
group of order eight.  Its subgroup over $\mathbb Q(i)$ is cyclic
of order four; its subgroup over $\mathbb Q(\sqrt{-3})$ is a
four-group.  It is therefore contained in both $K_j^{\mathrm{ab}}$.
The nontrivial extension in
\eqref{eq:adelic-cubic-four-division} gives the reverse bound and
proves \eqref{eq:adelic-cm-intersection}.

It remains to identify the common quadratic quotient in the stated
matrix coordinates.  By
\eqref{eq:adelic-unnormalized-cubic-image},
$\rho_3(\sigma)$ is one modulo $\pi_3$; hence
\[
 (-1)^{\epsilon_3(\rho_3^{\circ}(\sigma))}
                      =\kappa(\sigma).
\]
For $E_4$ its three-division polynomial is
$3X^4-6X^2-1$.  Choose a root $x$ with
$x^2=1+2\sqrt3/3$.  Then
\[
 \frac{x^2}{\sqrt3}=\frac{2+\sqrt3}{3}
     =\left(\frac{\sqrt6+\sqrt2}{2\sqrt3}\right)^2.
\]
Thus $3^{1/4}\in k(E_4[3])$.  The latter extension over $k$ is
cyclic of order four, since the reduction of $N_4(4)$ modulo
three is $\langle S\rangle$.  Its unique quadratic subextension
is consequently $k(3^{1/4})$, and
\[
 (-1)^{\epsilon_4(\rho_4(\sigma))}=\kappa(\sigma).
\]
The joint image of two Galois representations is the fibre product
of their separate images over the Galois group of the intersection
of their fields.  The separate images are full by
Proposition~\ref{prop:adelic-cm-images}; the intersection has order
two and the two quotient maps have just been computed.  The fibre
product is exactly \eqref{eq:adelic-equal-parity}.
\end{proof}

\subsection{A signed comparison on the Legendre family}

\begin{proposition}\label{prop:adelic-signed-path-comparison}
Choose the Tate-module bases by transport along paths to the
harmonic and equianharmonic fibres of the Legendre family.  For
$\sigma\in H^{\mathrm{cyc}}$ one then has the exact integral
identities
\begin{equation}
 W_\sigma=\rho_4(\sigma),\qquad
 T_\sigma=\rho_3^{\circ}(\sigma)
                       =\kappa(\sigma)\rho_3(\sigma).
 \label{eq:adelic-signed-path-comparison}
\end{equation}
In particular $(T,W)(H^{\mathrm{cyc}})=\mathcal E$.
\end{proposition}

\begin{proof}
Let $Y=\mathbb P^1\setminus\{0,1,\infty\}$ and consider the
Legendre family
\[
 \mathcal L_t:\quad Y_1^2=X_1(X_1-1)(X_1-t)
\]
with tangential base point at $t=0$ towards $1$.
Choose the integral vanishing-cycle basis in the peripheral
convention fixed by $x,y$.  The two local monodromies are
$U^2,L^2$: a loop about either collision of two ordered branch
points lifts to the square of the corresponding Dehn twist, and
the two vanishing cycles have intersection of absolute value one.
This identifies the geometric Tate-module representation with $q$
on $F$, at every prime.

At zero the special cubic is $Y_1^2=X_1^2(X_1-1)$, whose two
branches at the node have slopes $i,-i$.  Over
$\mathbb Q(i)((t))$ the reduction is split multiplicative.
Tate uniformization supplies a primitive cyclotomic line and a
free rank-one quotient in every integral Tate module; the only
additional character over $\mathbb Q((t))$ is the splitting
character of $\mathbb Q(i)$.
For an explicit uniformization over an arbitrary complete
nonarchimedean field, including the Galois compatibility and the
split multiplicative criterion, see
\cite[\S2, Theorem~2.2 and pp.~3--4]{PapikianTate}.
The integral filtration follows directly: in
$\overline K^{\times}/Q^{\mathbb Z}$, the $n$-torsion has kernel
$\mu_n$ and quotient $\mathbb Z/n$ under
$[z]\mapsto m\bmod n$ for $z^n=Q^m$.
The characteristic is zero, so this argument applies to every
prime, including two.

Both characters on the two graded pieces are trivial on
$H^{\mathrm{cyc}}$.  The tangential arithmetic matrix
$V_\sigma$ is therefore upper unipotent with both diagonal
entries one.  Functoriality gives
$q\Phi_\sigma=\operatorname{Ad}(V_\sigma)q$ on $F$.
Equation~\eqref{eq:adelic-frame-intertwiner} gives the same
covariance for $D_\sigma$.
An upper unipotent matrix commuting with $U^2$ and $L^2$ is the
identity: the second commutation gives twice its upper-right
entry equal to zero in each $\mathbb Z_\ell$.
Consequently
\begin{equation}
 V_\sigma=D_\sigma.
 \label{eq:adelic-tangential-matrix}
\end{equation}
This uniqueness uses the primitive integral filtration, so it
does not leave a sign undetermined at the prime two.

The geometric fundamental group of the quotient orbifold
$[Y/S_3]$ is $\Gamma$, with pure subgroup $F$.
The arithmetic action induced by the image of the chosen
tangential base point restricts to the canonical Ihara action
on $F$.  It is exactly $\Phi_\sigma$ on $\Gamma$.
Indeed two automorphisms of $\Gamma$ agreeing on $F$ differ,
on each group element, by an element of $C_\Gamma(F)$, and this
centralizer is trivial.  For an explicit argument, if $z$ centralizes
$F$, then $z^6\in F$ and $z^6$ lies in the centre of the nonabelian
free profinite group $F$, which is trivial.  Thus $z$ has finite
order.  A nontrivial torsion element of $\Gamma$ is conjugate into
one of its two finite factors, and its centralizer lies in that
factor \cite{HerfortRibes1985}.  Such a finite centralizer cannot
contain $F$, so $z=1$.

The transposition $t\mapsto1-t$ fixes $p_4=1/2$.
Choose an ordinary geometric path $\xi_4$ in $Y$ from the
tangential base point to $p_4$ so that its transported orbifold
isotropy is $s$.  Such a choice exists: every order-two subgroup
is conjugate into the order-two factor, and a conjugator preserving
the specified transposition in $S_3$ belongs to $F\langle s\rangle$.
It can therefore be replaced by a pure conjugator.
For the order-three isotropy use one of
$p_3=(1\pm\sqrt{-3})/2$ and choose an ordinary path $\xi_3$
whose transported positive generator is $r$.
The two points give the two possible pure conjugacy classes;
choose the one corresponding to $r$.  The same conjugator
argument, with centralizer $\langle r\rangle$ in $S_3$, then
makes the path adjustment pure.  Both chosen points and their
coordinate-permutation isotropies are fixed by $H^{\mathrm{cyc}}$.

Paths are composed as maps of fibre functors.  If $\xi$ is either
chosen path, set
\[
 c_\sigma=(\sigma\xi)^{-1}\xi\in F.
\]
For the isotropy $s_p$ at its endpoint, the identity
$s=\xi^{-1}s_p\xi$ gives
\[
 \Phi_\sigma(s)=(\sigma\xi)^{-1}s_p(\sigma\xi)
                        =c_\sigma s c_\sigma^{-1}.
\]
The analogous formula holds for $r$.  Uniqueness of the pure
frames yields $c_\sigma=\beta_\sigma$ or $A_\sigma$.

Now use the Tate sheaf on $Y$, where the Legendre family itself
is defined.  If $\rho_p(\sigma)$ is its matrix at the ordinary
point $p$, functoriality of parallel transport gives
\[
 q(\sigma\xi)=\rho_p(\sigma)q(\xi)D_\sigma^{-1}.
\]
In the transported basis put
$M_p(\sigma)=q(\xi)^{-1}\rho_p(\sigma)q(\xi)$.
Taking the inverse in the preceding formula and multiplying by
$q(\xi)$ proves
\begin{equation}
 q(c_\sigma)=D_\sigma M_p(\sigma)^{-1},\qquad
 q(c_\sigma)^{-1}D_\sigma=M_p(\sigma).
 \label{eq:adelic-path-order}
\end{equation}
The signed representation here is the one already fixed on $F$.
The calculation applies on every finite torsion sheaf, with
compatible transition maps.  It therefore proves an integral
adelic identity.  No descent of the signed Tate sheaf to
$[Y/S_3]$ is needed.

It remains to identify the two ordinary fibres.  At $p_4=1/2$,
the substitution $X_1=X+1/2$ gives
$Y_1^2=X^3-X/4$.  This is the quadratic twist of $E_4$ by
$1/2$, whose square root lies in $\mathbb Q^{\mathrm{ab}}$.
Its representation on $H^{\mathrm{cyc}}$ is therefore $\rho_4$.
At $p_3=(1+\sqrt{-3})/2$, substituting
$X_1=X+(1+p_3)/3$ gives
\[
 Y_1^2=X^3+\frac{\sqrt{-3}}9
             =X^3+\left(\frac1{\sqrt{-3}}\right)^3.
\]
This is the quadratic twist of $E_3$ by
$d=1/\sqrt{-3}$.  The quotient of a square root of $d$ by $r_0$
lies in $\mathbb Q^{\mathrm{ab}}$; its square is $-2i/3$, and
one may take $\sqrt{2/3}\,\zeta_8^{-1}$.
Its twisting character on $H^{\mathrm{cyc}}$ is consequently
$\kappa$.  The conjugate equianharmonic point changes $d$ by
$-1$ and gives the same character.

The transported CM algebras are the centralizers of $S$ and $R$
already computed in Lemma~\ref{lem:adelic-corrected-frames}.
Indeed Proposition~\ref{prop:adelic-cm-images} makes each local
arithmetic image open in its one-dimensional CM torus, so its
commutant over $\mathbb Q_\ell$ is the corresponding quadratic
algebra.  The matrices $S$ and $R$, which commute with that image,
generate its full quadratic integer ring.
At a prime, interchanging the two CM embeddings changes a unit
to its conjugate; it preserves the norm-one groups, their level-four
subgroups, and both parity maps.  Thus the CM image calculations
apply in these transported integral bases.
Equation~\eqref{eq:adelic-path-order} now gives
\eqref{eq:adelic-signed-path-comparison}, and
Lemma~\ref{lem:adelic-cm-entanglement} identifies the joint image.
\end{proof}

\subsection{Eliminating the entire adelic frame contribution}

Recall the global quadratic character $\Delta:\mathcal G\to
\mathbb F_2$ of Section~\ref{sec:global-frame-character}.
It is zero on $H$ and on $G$; on
$N_0=\ker(\lambda,\rho_2)$ it equals $a_3+a_2$.
The vanishing on $G$ uses the normalized frame theorem
\cite[Theorem~3]{H17}.  We require no stronger arithmetic input.

\begin{theorem}\label{thm:adelic-frame-elimination}
Put $G_1=G\cap M_0$.  Then
\begin{equation}
 (T,W)(G_1)=(T,W)(H^{\mathrm{cyc}})=\mathcal E.
 \label{eq:adelic-gt-arithmetic-image}
\end{equation}
In particular, the closed normal subgroup
\begin{equation}
 C_{\mathrm{CM}}=
 \ker\bigl((T,W):G_1\longrightarrow N_3(4)\times N_4(4)\bigr)
 \label{eq:adelic-cm-kernel}
\end{equation}
satisfies
\begin{equation}
 G=C_{\mathrm{CM}}H,\qquad
 C_{\mathrm{CM}}\cap H=
 G_{K_3^{\mathrm{ab}}K_4^{\mathrm{ab}}}.
 \label{eq:adelic-cm-normalization}
\end{equation}
The inclusion of $H$ induces the canonical isomorphism
\begin{equation}
 G/C_{\mathrm{CM}}\simeq
 \operatorname{Gal}(K_3^{\mathrm{ab}}K_4^{\mathrm{ab}}/\mathbb Q).
 \label{eq:adelic-cm-arithmetic-quotient}
\end{equation}
Every continuous abelian character of $G$ that is zero on $H$
and whose restriction to $G_1$ factors through $(T,W)$ is zero.
\end{theorem}

\begin{proof}
On $N_0$ the matrix $D_g$ is the identity modulo twelve.
The reductions of $q(A_g)$ and $q(\beta_g)$ therefore give
\[
 \epsilon_3(T_g)=a_3(g),\qquad
 \epsilon_4(W_g)=a_2(g).
\]
For the cubic group the first equality is reduction modulo
$\pi_3$: $R$ reduces to one and $-I$ to minus one.
For the harmonic group the second equality is exponent parity
in $\langle S\rangle$.  Inverting a frame does not change either
parity.  Thus
$(\epsilon_3\circ T)+ (\epsilon_4\circ W)=\Delta$ on $N_0$.

The arithmetic normalization by the Kummer cocycle of two gives
$M_0=N_0H^{\mathrm{cyc}}$: its value on $M_0$ is even, and its
image on $H^{\mathrm{cyc}}$ is exactly
$2\widehat{\mathbb Z}$.  On $H^{\mathrm{cyc}}$ both the sum of
the two parities and $\Delta$ are zero, by
Proposition~\ref{prop:adelic-signed-path-comparison} and the
arithmetic vanishing of $\Delta$.
All these maps are homomorphisms on $M_0$, so
\begin{equation}
 (\epsilon_3\circ T)+ (\epsilon_4\circ W)=\Delta
                          \quad\hbox{on }M_0.
 \label{eq:adelic-parity-delta}
\end{equation}
Since $\Delta|_G=0$, the image of $G_1$ lies in $\mathcal E$.
It contains the image of $H^{\mathrm{cyc}}$, which is already
all of $\mathcal E$.  This proves
\eqref{eq:adelic-gt-arithmetic-image}.

For $g\in G$, first choose $\sigma\in H$ with
$\lambda_\sigma=\lambda_g$.  Then $g\sigma^{-1}\in G_1$.
By \eqref{eq:adelic-gt-arithmetic-image}, choose
$h\in H^{\mathrm{cyc}}$ with the same corrected frame pair.
It follows that $g\sigma^{-1}h^{-1}\in C_{\mathrm{CM}}$, proving
$G=C_{\mathrm{CM}}H$.
The intersection formula follows from the exact signed comparison,
the equality of the normalized and unnormalized cubic kernels,
and the two full division-field identities.
Normality follows from Lemma~\ref{lem:adelic-kernel-normality}.
The compact surjection $H\to G/C_{\mathrm{CM}}$ has exactly the
kernel in \eqref{eq:adelic-cm-normalization}, proving
\eqref{eq:adelic-cm-arithmetic-quotient}.

Finally let $\chi$ be a character as in the statement.
Its restriction to $G_1$ is a character of the image
$\mathcal E$, and it is zero there because
$H^{\mathrm{cyc}}$ maps onto that image.  Thus $\chi|_{G_1}=0$.
Matching a multiplier by an element of $H$ then gives $\chi=0$
on $G$.
\end{proof}

In particular the same relative abelian quotient can be computed as
\[
 B\simeq C_{\mathrm{CM}}\big/
       \overline{[G,C_{\mathrm{CM}}]\,
                 G_{K_3^{\mathrm{ab}}K_4^{\mathrm{ab}}}}.
\]
For any prime $p$, its continuous $\mathbb F_p$-characters are the
$G$-invariant characters of $C_{\mathrm{CM}}$ that vanish on the
displayed Galois subgroup.  This follows from the general extension
criterion for a decomposition $G=NH$ with $N$ closed and normal.

This normalization concerns the full pair of CM frames.  Combining
it with simultaneous rational-radical normalization would require
the joint arithmetic image of the CM and radical representations;
the two separate surjectivity statements do not imply that joint
compatibility.

\section{The exact intersection of the radical and CM fields}
\label{sec:radical-cm-entanglement}

Put $k=\mathbb Q^{\mathrm{ab}}=\mathbb Q(\mu_\infty)$ and
\[
 K_3=\mathbb Q(\sqrt{-3}),\qquad K_4=\mathbb Q(i),\qquad
 F_{12}=K_3K_4=\mathbb Q(\zeta_{12}).
\]
Let $\mathcal P$ be the set of positive rational primes, and set
\[
 L_{\mathrm{rad}}=k(q^{1/n}:q\in\mathcal P,\ n\geq1),\qquad
 L_{\mathrm{CM}}=K_3^{\mathrm{ab}}K_4^{\mathrm{ab}}.
\]
All unqualified radicals are the positive real ones.  The presence
of all roots of unity makes these fields independent of that choice.

\begin{theorem}\label{thm:radical-cm-intersection}
One has
\begin{align}
 L_{\mathrm{rad}}\cap L_{\mathrm{CM}}
 &=L_{\mathrm{rad}}\cap F_{12}^{\mathrm{ab}}\notag\\
 &=k\bigl(q^{1/12}:q\in\mathcal P,\,
                         2^{1/8},3^{1/8}\bigr)=:I.
 \label{eq:radical-cm-intersection}
\end{align}
In particular
\begin{equation}
 \operatorname{Gal}(I/k)\simeq
 (\mathbb Z/12)^2\times
       \prod_{q\in\mathcal P\setminus\{2,3\}}\mathbb Z/6.
 \label{eq:radical-cm-common-group}
\end{equation}
The first two factors are indexed by $2,3$.  In these coordinates
conjugation by an element of $G_{\mathbb Q}$ is multiplication by
its cyclotomic character, reduced modulo the indicated orders.
\end{theorem}

\begin{proof}
We first compute the entire maximal subextension of $L_{\mathrm{rad}}$
that is abelian over $F_{12}$.  Recall the simultaneous radical
description
\[
 A_{\mathrm{rad}}=\operatorname{Gal}(L_{\mathrm{rad}}/\mathbb Q)
 =
 \left\{(u,(t_q)_q):
 \begin{array}{l}
 u\in\widehat{\mathbb Z}^{\times},\
 t_q\in\widehat{\mathbb Z},\\
 t_q\bmod2=\epsilon_q(u)\quad(q\in\mathcal P)
 \end{array}\right\},
\]
with multiplication $(u,t)(v,s)=(uv,t+us)$.
Here $\epsilon_q$ is the quadratic character of $\mathbb Q(\sqrt q)$,
written additively.  The translation subgroup over $k$ is
$T=\prod_q2\widehat{\mathbb Z}$.
This is the classical radical Galois description
\cite[\S2, Theorem~2.5]{MoreeStevenhagen}; the passage from finite
sets of primes to this product follows by compactness.
It also follows directly from the simultaneous Kummer-image lemma:
a positive rational radical lying in an abelian extension of
$\mathbb Q$ has its square in $\mathbb Q$, which is exactly the
displayed parity condition.

Write
\[
 U_{12}=\{u\in\widehat{\mathbb Z}^{\times}:u\equiv1\pmod{12}\},
 \qquad A_{12}=\operatorname{Gal}(L_{\mathrm{rad}}/F_{12}).
\]
Thus $A_{12}$ is the inverse image of $U_{12}$ in $A_{\mathrm{rad}}$.
Let $C_{12}=\overline{[A_{12},A_{12}]}$.  Since the multiplier quotient is
abelian, $C_{12}\subset T$.  For the convention
$[a,b]=aba^{-1}b^{-1}$, direct multiplication gives
\begin{equation}
 [(u,t),(v,s)]=
 \bigl(1,(1-v)t+(u-1)s\bigr).
 \label{eq:radical-cm-affine-commutator}
\end{equation}
We claim the exact simultaneous equality
\begin{equation}
 C_{12}=(24\widehat{\mathbb Z})_{\{2,3\}}\times
             \prod_{q\ne2,3}12\widehat{\mathbb Z}.
 \label{eq:radical-cm-derived-subgroup}
\end{equation}
The subscripts $\{2,3\}$ here index the rational-prime coordinates,
not their local components.

We make the distinction between these coordinates and the local
prime $\ell$ explicit in the proof.  Every closed subgroup of the
profinite abelian group $T$ is stable under
$\widehat{\mathbb Z}$-multiplication: approximate an element of
$\widehat{\mathbb Z}$ by integers and use closedness.
The idempotents of $\widehat{\mathbb Z}=\prod_\ell\mathbb Z_\ell$
therefore project $C_{12}$ onto its local components inside $C_{12}$ itself.
Consequently it suffices to determine each simultaneous local
subgroup $C_{12,\ell}$; there can be no additional entanglement between
distinct local primes.

Commuting a lift of $u\in U_{12}$ with a translation gives
$(u-1)T$.  At $\ell\geq5$ one may choose $u_\ell=2$ and all
other local components one.  Hence
$C_{12,\ell}=\prod_q\mathbb Z_\ell$.
At $\ell=3$, choosing $u_3=4$ gives
$\prod_q3\mathbb Z_3\subset C_{12,3}$.
Both multipliers in \eqref{eq:radical-cm-affine-commutator} are one
modulo three, proving the reverse inclusion.  Thus
$C_{12,3}=\prod_q3\mathbb Z_3$.
At $\ell=2$, choosing $u_2=5$ and using translations in
$\prod_q2\mathbb Z_2$ gives
\[
 \prod_q8\mathbb Z_2\subset C_{12,2}.
\]

It remains to determine the image modulo eight simultaneously
in every coordinate.  For $u\in U_{12}$ put
$a(u)=(u_2-1)/4\bmod2$.  Since $u_2\equiv1\pmod4$,
\[
 a(u)=\epsilon_2(u).
\]
Also $\epsilon_3(u)=0$, since $\sqrt3\in F_{12}$.
Reducing \eqref{eq:radical-cm-affine-commutator} modulo eight gives
in coordinate $q$
\begin{equation}
 4\bigl(\epsilon_2(u)\epsilon_q(v)
               -\epsilon_2(v)\epsilon_q(u)\bigr).
 \label{eq:radical-cm-mod-eight}
\end{equation}
It is zero for $q=2,3$, and is a multiple of four for every
other $q$.  This proves the necessary simultaneous upper bound.

For the reverse inclusion, the characters
$\epsilon_q|_{U_{12}}$ with $q\ne3$ are jointly independent.
Indeed a finite product would be trivial on $U_{12}$ precisely if
$\mathbb Q(\sqrt a)\subset F_{12}$, where $a$ is the corresponding
positive squarefree product of primes.  The three quadratic
subfields of $F_{12}$ are
$\mathbb Q(i)$, $\mathbb Q(\sqrt{-3})$, and $\mathbb Q(\sqrt3)$;
among positive squareclasses only $a=3$ occurs.  Excluding $3$
therefore gives independence.  Every finite evaluation map onto
the corresponding product of $\mathbb F_2$ is surjective.
Compactness then gives surjectivity onto the entire product.

Fix $u\in U_{12}$ with $\epsilon_2(u)=1$.  For an arbitrary family
$(b_q)_{q\ne2,3}\in\prod_{q\ne2,3}\mathbb F_2$, choose
$v\in U_{12}$ with
\[
 \epsilon_2(v)=0,\qquad \epsilon_q(v)=b_q\quad(q\ne2,3).
\]
Choose any lifts of $u,v$ to $A_{12}$.  Formula
\eqref{eq:radical-cm-mod-eight} now gives exactly the family
$4b_q$ in all the other coordinates, with zero at $2,3$.
This proves surjectivity onto the whole required product modulo
eight, not merely onto its individual projections.  Together
with $\prod_q8\mathbb Z_2\subset C_{12,2}$, it proves
\[
 C_{12,2}=(8\mathbb Z_2)_{\{2,3\}}\times
                       \prod_{q\ne2,3}4\mathbb Z_2.
\]
Combining the local equalities proves
\eqref{eq:radical-cm-derived-subgroup}.

The fixed field of $C_{12}$ is $L_{\mathrm{rad}}\cap F_{12}^{\mathrm{ab}}$.
In the radical coordinates, the subgroup
\eqref{eq:radical-cm-derived-subgroup} fixes precisely the
$24$th roots of $2,3$ and the $12$th roots of all other primes.
Since a $24$th root is generated by its eighth and twelfth powers
as root orders, this is the field $I$ in
\eqref{eq:radical-cm-intersection}.  More explicitly,
$q^{1/24}=(q^{1/8})/(q^{1/12})$.
The quotient of $2\widehat{\mathbb Z}$ by
$24\widehat{\mathbb Z}$ is $\mathbb Z/12$, and by
$12\widehat{\mathbb Z}$ is $\mathbb Z/6$.  This proves
\eqref{eq:radical-cm-common-group} and its conjugation action.

Each $K_j^{\mathrm{ab}}/F_{12}$ is abelian, so
$L_{\mathrm{CM}}\subset F_{12}^{\mathrm{ab}}$.  We have therefore
proved $L_{\mathrm{rad}}\cap L_{\mathrm{CM}}\subset I$.
For the reverse inclusion, all rational cube roots lie in
$K_3^{\mathrm{ab}}$, since $K_3$ contains $\mu_3$.
All rational fourth roots lie in $K_4^{\mathrm{ab}}$, since
$K_4$ contains $\mu_4$.

Let $\alpha=2^{1/8}$.  The field $\mathbb Q(i,\alpha)$ contains
\[
 \zeta_8=\frac{1+i}{\alpha^4},
\]
so it is the splitting field of $X^8-2$.  An automorphism fixing
$i$ sends $\alpha$ to $\zeta_8^a\alpha$ and sends $\zeta_8$ to
$\zeta_8^{1+4a}$.  Composition on the index $a\bmod8$ is
$a*b=a+b+4ab$, which is commutative.  Thus this extension is
abelian over $K_4$, and $2^{1/8}\in K_4^{\mathrm{ab}}$.

For the remaining radical put
$v=(2+\sqrt3)^{1/4}$ and $M=\mathbb Q(i,v)$.
The field $M$ contains $F_{12}$ and all conjugates
$i^a v^{\pm1}$, so it is Galois over $\mathbb Q$.
Every automorphism of $M$ fixing $\sqrt{-3}=i\sqrt3$ has the form
\[
 i\longmapsto i^\epsilon,\qquad
 v\longmapsto i^a v^\epsilon,
 \qquad a\in\mathbb Z/4,\quad\epsilon\in\{1,-1\}.
\]
Indeed the sign on $\sqrt3=v^4-2$ must equal the sign on $i$.
The product of two such automorphisms has exponent
$a\eta+\epsilon b$ and sign $\epsilon\eta$; this rule is
symmetric in $(a,\epsilon)$ and $(b,\eta)$.
Therefore $M/K_3$ is abelian and $v\in K_3^{\mathrm{ab}}$.

On the CM elliptic curve $E_4:Y^2=X^3-X$, the three-division
polynomial is $3X^4-6X^2-1$.  Choose a three-torsion point $(x,y)$
with $x^2=1+2\sqrt3/3$.  The curve equation gives
\[
 y^4=x^2(x^2-1)^2
       =\frac49\sqrt3(2+\sqrt3).
\]
All its torsion coordinates are abelian over $K_4$, since the
endomorphism by $i$ is defined there and the Tate-module action
commutes with that rank-one CM action
\cite[Proposition~12.13, printed p.~119]{MilneMF110}.
Thus $y\in K_4^{\mathrm{ab}}$.
As $\sqrt{2/3}\in k$, the element
\[
 z=\frac{y}{\sqrt{2/3}\,v}\in L_{\mathrm{CM}}
\]
satisfies $z^4=\sqrt3$ and hence $z^8=3$.
All eighth roots of unity already lie in $k$, so
$3^{1/8}\in L_{\mathrm{CM}}$.
This supplies every generator of $I$ and completes the proof.
\end{proof}

\subsection{The remaining compatibility condition on GT}

The arithmetic intersection theorem by itself does not combine
the two separate GT normalizations.  Here is the exact additional
condition it leaves.

Suppose the separate arithmetic quotient identifications are fixed:
\[
 G/N_{\mathrm{rad}}\simeq
        A_{\mathrm{rad}}=\operatorname{Gal}(L_{\mathrm{rad}}/\mathbb Q),
 \qquad
 G/C_{\mathrm{CM}}\simeq
        A_{\mathrm{CM}}=\operatorname{Gal}(L_{\mathrm{CM}}/\mathbb Q),
\]
with both identifications induced by the inclusion
$H=G_{\mathbb Q}\subset G$.
Restricting these quotient maps to $I$ gives two homomorphisms
$a,b:G\to A_I=\operatorname{Gal}(I/\mathbb Q)$, which agree on
$H$ and induce the same multiplier.  Their difference
\begin{equation}
 d(g)=a(g)b(g)^{-1}\in C=\operatorname{Gal}(I/k)
 \label{eq:radical-cm-compatibility-defect}
\end{equation}
is a continuous crossed homomorphism:
\[
 d(gh)=d(g)+\lambda_gd(h),\qquad d|_H=0.
\]
The addition uses the coordinates
\eqref{eq:radical-cm-common-group}.
The formula follows by writing $a(g)=d(g)b(g)$ and using the
scalar cyclotomic conjugation on $C$.

\begin{proposition}\label{prop:radical-cm-joint-gate}
The following are equivalent:
\begin{enumerate}
\item $G=(N_{\mathrm{rad}}\cap C_{\mathrm{CM}})H$;
\item the joint radical and CM image of $G$ equals that of $H$;
\item $a=b$, or equivalently $d=0$.
\end{enumerate}
Thus the remaining compatibility has exponent dividing twelve
over $k$: it consists of one cubic coordinate for every rational
prime, one quadratic coordinate for each prime other than $2,3$,
and one coordinate modulo four for each of $2,3$.
\end{proposition}

\begin{proof}
The first two assertions are equivalent by taking the common
kernel of the two quotient maps.
The joint Galois image is exactly
$A_{\mathrm{rad}}\times_{A_I}A_{\mathrm{CM}}$,
by the intersection theorem and the Galois correspondence.
The joint image of $G$ contains it, since $H\subset G$.
It is contained in it precisely when its two restrictions to
$A_I$ coincide.  This proves the equivalences.  The final list
is the primary decomposition of
\eqref{eq:radical-cm-common-group}.
\end{proof}

For example, on $G_1=\ker\lambda$, write the coordinate of the
CM restriction to $I$ as $\eta_q(g)$ in $\mathbb Z/n_q$, where
$n_2=n_3=12$ and $n_q=6$ for $q\ne2,3$.
Normalize it arithmetically by
$\eta_q(\sigma)=\kappa_q(\sigma)/2\bmod n_q$ for
$\sigma\in G_{\mathbb Q^{\mathrm{ab}}}$.
Then the required coordinate identities are precisely
\[
 \frac{\rho_q(g)}2=\eta_q(g)\pmod{n_q}
 \qquad(q\in\mathcal P,\ g\in G_1).
\]
They compare the rational Kummer cocycles with finite quotients
of the CM frames.  Separate surjectivity of the two quotient maps
does not prove these identities.

\section{Compatibility of the complete common coordinate of two}
\label{sec:eighth-two-compatibility}

Both the eighth-root and the cubic coordinates of $2$ in the
radical--CM intersection are compatible with the corrected CM
frames.  These assertions hold on the multiplier-one gentle group.
We first identify the harmonic torsion level by an exact conductor
computation and prove its congruence directly from the cycle
equations.  An abelianized calculation in the modular group then
gives the cubic coordinate.

\subsection{The relevant Gaussian ray fields}

Put $k=\mathbb Q^{\mathrm{ab}}$, $K_4=\mathbb Q(i)$,
$\mathfrak p=(1+i)$, and $E_4:Y^2=X^3-X$.

\begin{lemma}\label{lem:eighth-two-conductor}
For $n=2,4,8$, the field
$L_n=\mathbb Q(i,2^{1/n})$ is cyclic of degree $n$ over $K_4$.
Its conductor is respectively
$\mathfrak p^4,\mathfrak p^6,\mathfrak p^8$.
In particular
\[
 L_4\subset K_4(E_4[8]),\qquad
 L_8\subset K_4(E_4[16]).
\]
\end{lemma}

\begin{proof}
Write $n=2^r$, $\alpha=2^{1/n}$, and $F_n=\mathbb Q(\alpha)$.
For $n=2,4$, cyclicity follows because $K_4$ contains the required
roots of unity.  For $n=8$ the field $L_8$ contains
$\zeta_8=(1+i)/\alpha^4$, and hence is the splitting field of
$X^8-2$.  Automorphisms fixing $i$ are indexed by
$a\in\mathbb Z/8$, with $\alpha\mapsto\zeta_8^a\alpha$.
Their composition is
$a*b=a+b+4ab$ modulo eight.  The element indexed by $1$ has
order eight, so the group is cyclic.  The degrees are $n$:
$X^n-2$ is Eisenstein over $\mathbb Q$ and $F_n$ is real,
so $[F_n(i):K_4]=n$.

We compute the discriminants without assuming a ray conductor.
The polynomial $X^n-2$ is Eisenstein at two, and $\alpha$ is a
uniformizer of its totally ramified local extension.  Therefore
the local integer ring is $\mathbb Z_2[\alpha]$.
The polynomial discriminant is supported only at two, so
$\mathcal O_{F_n}=\mathbb Z[\alpha]$ globally, and
\begin{equation}
 |\operatorname{disc}(F_n/\mathbb Q)|
       =n^n2^{n-1}=2^{rn+n-1}.
 \label{eq:eighth-two-pure-discriminant}
\end{equation}

In the completion of $F_n$ at its unique prime over two, use
the valuation with $v(\alpha)=1$.  Define
\[
 a_2=1+\alpha,\qquad
 a_4=1+\alpha^2+\alpha^3,\qquad
 a_8=1+\alpha^4+\alpha^6+\alpha^7,
\]
where in each expression $\alpha^n=2$ for its indicated $n$.
Direct expansion gives
\[
 a_n^2+1\equiv\alpha^{2n-1}\pmod{\alpha^{2n}},
 \qquad v(a_n^2+1)=2n-1.
\]
The element
$b_n=(i+a_n)/\alpha^{n-1}$ satisfies
\begin{equation}
 b_n^2-a_n\alpha\,b_n+
              \frac{a_n^2+1}{\alpha^{2n-2}}=0.
 \label{eq:eighth-two-eisenstein}
\end{equation}
Both nonleading coefficients have valuation one, so this is an
Eisenstein quadratic polynomial.  It generates the relative local
integer ring.  Its discriminant has valuation two: the square
of its linear coefficient has valuation two, whereas four times
the constant coefficient has valuation $2n+1$.
Thus the relative discriminant of $L_n/F_n$ has exponent two
at the unique prime over two and has no odd prime factor.
The discriminant tower formula and
\eqref{eq:eighth-two-pure-discriminant} give
\[
 |\operatorname{disc}(L_n/\mathbb Q)|
       =2^{\,2(rn+n-1)+2}=2^{\,2n(r+1)}.
\]
Since $|\operatorname{disc}(K_4/\mathbb Q)|=4$ and
$N_{K_4/\mathbb Q}(\mathfrak p)=2$, a second application of
the tower formula yields
\begin{equation}
 \operatorname{disc}(L_n/K_4)=\mathfrak p^{\,2nr}.
 \label{eq:eighth-two-gaussian-discriminant}
\end{equation}

Let $c_2,c_4,c_8$ be the conductor exponents of faithful characters
of the three cyclic extensions.  Characters of the same order
have the same kernel and hence the same conductor.  The cyclic
tower $L_2\subset L_4\subset L_8$ and the conductor--discriminant
formula \cite[V, Theorem~3.27 and Remark~3.28, p.~165]{MilneCFT}
give
\[
 c_2=4,\qquad c_2+2c_4=16,\qquad
 c_2+2c_4+4c_8=48.
\]
Consequently $(c_2,c_4,c_8)=(4,6,8)$.
The conductors of the extensions themselves are the same as
those of their faithful characters.

The torsion reciprocity theorem for $E_4$ identifies its
torsion points modulo geometric automorphisms with the
corresponding Gaussian ray class fields
\cite[Theorem~12.27, printed p.~124]{MilneMF110}.
Thus the ray class field modulo $(8)=\mathfrak p^6$ is contained
in $K_4(E_4[8])$, and that modulo $(16)=\mathfrak p^8$ is
contained in $K_4(E_4[16])$.  The conductors just computed
give the two asserted inclusions.
\end{proof}

\begin{proposition}\label{prop:eighth-two-division-fields}
Over $k=\mathbb Q^{\mathrm{ab}}$ one has
\[
 k(E_4[8])=k(2^{1/4}),\qquad
 k(E_4[16])=k(2^{1/8}).
\]
\end{proposition}

\begin{proof}
The arithmetic CM image on $G_k$ is
$N_4(4)=\{w\in\mathbb Z_2[i]^\times:
w\bar w=1,\ w\equiv1\pmod4\}$ at the prime two.
This is the image calculated in the adelic frame comparison.
Write $w=a+bi$, where $a=1+4s$ and $b=4t$.
The norm equation gives
$s+2(s^2+t^2)=0$.  In particular $s$ is even and
$s\equiv2t\pmod4$, so
\begin{equation}
 w\equiv1+8t+4ti\pmod{16}.
 \label{eq:eighth-two-norm-reduction}
\end{equation}
Every class $t\bmod4$ occurs: the norm-one unit
$(1+2ti)/(1-2ti)$ has the required imaginary coefficient modulo
sixteen.  Thus the image modulo sixteen is cyclic of order four,
generated by $9+4i$, and its image modulo eight has order two.
It follows that
$[k(E_4[16]):k]=4$ and $[k(E_4[8]):k]=2$.
The two inclusions from Lemma~\ref{lem:eighth-two-conductor}
become inclusions of extensions of these same degrees, since
the radical Kummer image over $k$ is $2\widehat{\mathbb Z}$.
They are therefore equalities.
\end{proof}

\subsection{The cycle equations fix the harmonic frame modulo sixteen}

Use the signed matrices
\[
 U=\begin{pmatrix}1&1\\0&1\end{pmatrix},\quad
 L=\begin{pmatrix}1&0\\-1&1\end{pmatrix},\quad S=ULU,
 \qquad q(x)=U^2,\quad q(y)=L^2.
\]
Write $\gamma_3F=\overline{[F,[F,F]]}$ for the third lower
central subgroup of $F=\widehat F(x,y)$.

\begin{lemma}\label{lem:eighth-two-class-two-image}
Modulo sixteen the commutator subgroup of $q(F)$ is the central
cyclic group of order four generated by
\[
 C=[U^2,L^2]
   =\begin{pmatrix}13&8\\8&5\end{pmatrix}\pmod{16}.
\]
In particular $q_{16}(\gamma_3F)=1$.
\end{lemma}

\begin{proof}
Direct multiplication gives $C^2=9I$, $C^4=I$, and shows that
$C$ commutes with both $U^2$ and $L^2$ modulo sixteen.
In a group generated by two elements, the commutator subgroup
is the normal closure of their commutator.  Here that normal
closure is precisely the central group $\langle C\rangle$.
\end{proof}

\begin{lemma}\label{lem:eighth-two-normalized-gentle}
If $h=(1,f)$ is gentle and $\rho_2(h)=0$, then
\[
 \beta_h\in\gamma_3F,\qquad
 q_{16}(\beta_h)=D_h\bmod16=W_h\bmod16=I.
\]
\end{lemma}

\begin{proof}
In $F/\gamma_3F$, write $f=[x,y]^a$.
At multiplier one, the three-cycle equation makes $3a=0$,
because its three commutators have the same central image.
The commutator coordinate is in the torsion-free group
$\widehat{\mathbb Z}$, so $a=0$ and $f\in\gamma_3F$.
The equation $\rho_2(h)=0$ gives $\beta_h\in F'$.
Write its image modulo $\gamma_3F$ as $[x,y]^d$.
The interchange $\theta(x)=y$, $\theta(y)=x$ inverts this
central commutator, and the frame identity
$f=\theta(\beta_h)\beta_h^{-1}$ gives $2d=0$.
Therefore $d=0$, proving $\beta_h\in\gamma_3F$.

Lemma~\ref{lem:eighth-two-class-two-image} gives
$q_{16}(\beta_h)=I$.  The corrected frame
$W_h=q(\beta_h)^{-1}D_h$ centralizes $S$ as a signed matrix.
Consequently the upper-unipotent matrix $D_h\bmod16$ centralizes
$S$.  Comparing the two diagonal entries in $D_hS=SD_h$
forces its shear to be zero modulo sixteen.  Thus both $D_h$
and $W_h$ are the identity modulo sixteen.
\end{proof}

\begin{theorem}\label{thm:eighth-two-compatibility}
Let $M_0=\ker\lambda$ in the gentle group.  Define the continuous
homomorphism
$\nu_2:N_4(4)\to\mathbb Z/4$ by its arithmetic normalization:
\[
 \nu_2(\rho_4(\sigma))
       =\frac{\kappa_2(\sigma)}2\pmod4
       \qquad(\sigma\in G_{\mathbb Q^{\mathrm{ab}}}).
\]
This is well defined and has kernel the subgroup congruent to
one modulo sixteen.  For every $h\in M_0$,
\begin{equation}
 \nu_2(W_h)=\frac{\rho_2(h)}2\pmod4.
 \label{eq:eighth-two-frame-kummer}
\end{equation}
In particular the exceptional eighth-root coordinate of $2$
in the radical--CM compatibility defect is identically zero
on full GT.
\end{theorem}

\begin{proof}
Proposition~\ref{prop:eighth-two-division-fields} identifies
the quotient of the arithmetic harmonic image modulo sixteen
with $\operatorname{Gal}(k(2^{1/8})/k)$.
The Kummer action on the eighth root has even exponent on $G_k$,
so division by two identifies this group with $\mathbb Z/4$.
This proves the stated definition and kernel of $\nu_2$.

The normalization of the gentle group gives
$M_0=N_0G_k$, where $N_0=\ker(\lambda,\rho_2)$.
Explicitly, $\rho_2(h)$ is even, and the Kummer map
$G_k\to2\widehat{\mathbb Z}$ is onto; choose $\sigma\in G_k$
with $\kappa_2(\sigma)=\rho_2(h)$ and put $n=h\sigma^{-1}$.
Then $n\in N_0$.
By Lemma~\ref{lem:eighth-two-normalized-gentle},
$W_n\equiv I\pmod{16}$.
The corrected-frame homomorphism and its exact arithmetic
comparison give
\[
 \nu_2(W_h)=\nu_2(W_\sigma)
           =\nu_2(\rho_4(\sigma))
           =\kappa_2(\sigma)/2
           =\rho_2(h)/2\pmod4.
\]
This proves \eqref{eq:eighth-two-frame-kummer}.
It is the required compatibility on the multiplier-one subgroup.
Matching an arbitrary multiplier by a Galois element extends the
vanishing to the corresponding crossed-cocycle component on $G$.
\end{proof}

\subsection{The cubic coordinate of two}

Recall the modular presentation
$\Gamma=\langle s,r\mid s^2=r^3=1\rangle$ and its pure subgroup
$F=\ker(\Gamma\to S_3)$, with
$x=u^2$, $y=v^2$, $s=uvu$, $r=uv$, and $sr=u^{-1}$.
For $h\in M_0$ the positive peripheral automorphism satisfies
\[
 \Phi_h(s)=\beta_hs\beta_h^{-1},\qquad
 \Phi_h(r)=A_hrA_h^{-1},\qquad
 \Phi_h(sr)=sr,
 \qquad A_h,\beta_h\in F .
\]
Its corrected cubic frame is
\[
 T_h=q(A_h)^{-1}D_h,\qquad
 R=-UL=\begin{pmatrix}0&-1\\1&-1\end{pmatrix}.
\]
Thus $T_h$ centralizes $R$ as a signed matrix.

\begin{lemma}\label{lem:cubic-two-pure-frame}
For every $h\in M_0$, one has $A_h\in F'$.
\end{lemma}

\begin{proof}
Write the abelian group $F^{\mathrm{ab}}$ additively.
Expanding $\Phi_h(s)\Phi_h(r)=sr$ in the extension
$\Gamma/F'$ gives
\[
 (1-s)[\beta_h]+s(1-r)[A_h]=0.
\]
The harmonic normalization is
$[\beta_h]=-\rho_2(h)([x]+[y])$.
Conjugation by $s$ interchanges $[x]$ and $[y]$, so the first
term vanishes and $(1-r)[A_h]=0$.
Conjugation by $r$ sends
$x\mapsto y$ and $y\mapsto(xy)^{-1}$; its matrix on
$F^{\mathrm{ab}}=\widehat{\mathbb Z}^2$ is therefore
\[
 \begin{pmatrix}0&-1\\1&-1\end{pmatrix}.
\]
The determinant of $1-r$ is $3$.  Since multiplication by $3$
is injective on $\widehat{\mathbb Z}^2$, the adjugate identity
shows that $1-r$ is injective.  Hence $[A_h]=0$.
\end{proof}

The centralizer of $R$ in $\mathrm{SL}_2(\mathbb F_3)$ is
$\{\pm R^j:j\in\mathbb Z/3\}$.  Define
\[
 j_3:N_3(4)\longrightarrow\mathbb Z/3,\qquad
 j_3(t)=j\quad\hbox{if }t\bmod3=(-I)^eR^j .
\]
This is the quotient that kills the central sign.

\begin{proposition}\label{prop:cubic-two-compatibility}
For every $h\in M_0$,
\[
 j_3(T_h)=\rho_2(h)\pmod3.
\]
On the arithmetic subgroup, this quotient is the Kummer action
on $k(2^{1/3})/k$, with the indicated choice of exponent.
\end{proposition}

\begin{proof}
The abelianization of $\mathrm{SL}_2(\mathbb F_3)$ is cyclic of
order three.  Its quotient map $\ell$ can be normalized by
$\ell(U)=\ell(L)=1$; then
$\ell(S)=0$ and $\ell(R)=2$.
For completeness, its derived subgroup is the quaternion group
generated by $S$ and $USU^{-1}$.
These generators have square $-I$, anticommute, and generate
a normal subgroup of order eight.  Conjugation by $U$ permutes
its three cyclic subgroups of order four.  Commutators with $U$
generate this quaternion group, proving the abelianization
assertion.

Write $D_h=\left(\begin{smallmatrix}1&b_h\\0&1\end{smallmatrix}\right)$.
Since $q(x)=U^2$ and $q(y)=L^2$, the harmonic normalization gives
\[
 \ell(q(\beta_h))=-4\rho_2(h)=-\rho_2(h)\pmod3.
\]
The centralizer of $S$ in $\mathrm{SL}_2(\mathbb F_3)$ is
$\langle S\rangle$, so $\ell(W_h)=0$.  The identity
$W_h=q(\beta_h)^{-1}D_h$ now gives
$b_h=-\rho_2(h)$ modulo three.
By Lemma~\ref{lem:cubic-two-pure-frame},
$\ell(q(A_h))=0$, and consequently
\[
 2j_3(T_h)=\ell(T_h)=\ell(D_h)
       =-\rho_2(h)\pmod3 .
\]
Multiplying by $2$ yields the stated formula.

The arithmetic identification can also be seen directly from
$E_3:Y^2=X^3+1$.  The duplication formula
$X(2P)=9X(P)^4/(4Y(P)^2)-2X(P)$ shows that
its three-division polynomial is
$3X(X^3+4)$.  The nonzero three-torsion coordinates satisfy
either $X=0$, $Y=\pm1$, or
$X^3=-4$, $Y^2=-3$.
Since $k$ contains the roots of unity and $\sqrt{-3}$, this gives
\[
 k(E_3[3])=k(2^{1/3}).
\]
The signed arithmetic comparison is $T_\sigma=\kappa(\sigma)
\rho_3(\sigma)$, where $\kappa(\sigma)\in\{\pm1\}$.
Its reduction modulo three can therefore have an additional
central sign.  Passing to $j_3$ kills that sign, and gives the
cubic field just identified; no equality with the full corrected
modulo-three field is asserted.
On $\sigma\in G_k$, the formula already proved reads
$j_3(T_\sigma)=\kappa_2(\sigma)\pmod3$.
This fixes its normalization as the usual Kummer exponent.
\end{proof}

\begin{corollary}\label{cor:complete-two-compatibility}
There is a continuous homomorphism
\[
 \nu_{2,\mathrm{CM}}:N_3(4)\times N_4(4)
          \longrightarrow\mathbb Z/12
\]
specified by
\[
 \nu_{2,\mathrm{CM}}(t,w)\equiv\nu_2(w)\pmod4,\qquad
 \nu_{2,\mathrm{CM}}(t,w)\equiv2j_3(t)\pmod3 .
\]
For every $h\in M_0$,
\[
 \nu_{2,\mathrm{CM}}(T_h,W_h)=\frac{\rho_2(h)}2\pmod{12}.
\]
Arithmetically this is the common quotient
$\operatorname{Gal}(k(2^{1/24})/k)$, with Kummer exponent
$\kappa_2/2$.
In particular the entire order-twelve coordinate belonging to
$2$ in the radical--CM compatibility defect vanishes on full GT.
\end{corollary}

\begin{proof}
The Chinese remainder theorem defines the homomorphism.
The two preceding results give its asserted value because
$2$ is its own inverse modulo three.
The fields $k(2^{1/8})$ and $k(2^{1/3})$ have degrees four and
three over $k$, and together generate $k(2^{1/24})$.
Their arithmetic Kummer exponents agree with the two displayed
congruences.  Finally, the compatibility defect is a crossed
cocycle that vanishes on the arithmetic subgroup.  Matching
the multiplier of any full-GT element by an arithmetic element
reduces its value to the multiplier-one case proved here.
\end{proof}

\section{Finite tests for the radical--CM compatibility defect}
\label{sec:lr-relative-gate}

Write $G=\widehat{\mathrm{GT}}$, $H=G_{\mathbb Q}$,
$k=\mathbb Q^{\mathrm{ab}}$, and $D=N_{\mathrm{rad}}$.
We use the established individual arithmetic quotients, the exact
radical--CM intersection, the signed CM-frame comparison, the normalized
signed-congruence identity, and the full hexagonal comparison of
\cite[Proposition~18]{H17} as inputs. Thus $G=DH$, and on $D$ the
multiplier and every rational Kummer cocycle vanish. The crossed defect is
\[
 d:G\longrightarrow C=\operatorname{Gal}(I/k),\qquad
 d(gh)=d(g)+\lambda_gd(h),\quad d|_H=0,
\]
with $C\simeq(\mathbb Z/12)^2\times\prod_{q\ne2,3}\mathbb Z/6$.
Nothing below assumes $d=0$, joint factorization, or $B/2B=0$.
Complete prime-two compatibility is a separate input when that
coordinate is removed from the final list of gates.

\subsection{An obstruction to a group-theoretic shortcut}

\begin{proposition}\label{prop:lr-abstract-countermodel}
Separate arithmetic quotient identifications do not imply vanishing
of a cyclotomically crossed defect. Even a trivial ordinary relative
abelian quotient does not force its order-three component to vanish.
\end{proposition}
\begin{proof}
Let $U=C_2$ act by negation on $V=\mathbb F_3$, and put
\[
 G_*=(V\oplus V)\rtimes U,\qquad
 H_*=\{(x,x,u):x\in V,\ u\in U\},\qquad A_*=V\rtimes U.
\]
The two projections $a_*,b_*:G_*\to A_*$ restrict to isomorphisms
on $H_*$. Hence $G_*=(\ker a_*)H_*=(\ker b_*)H_*$. Nevertheless
$d_*(x,y,u)=x-y$ is a nonzero crossed homomorphism, trivial on $H_*$.
Since $[G_*,G_*]=V\oplus V$, the abelianization of $G_*$ is $U$,
already covered by $H_*$. Its ordinary relative abelian quotient
is therefore zero. This is a countermodel to the proposed abstract
inference, not a construction of an element of GT.
\end{proof}

There is an exact cohomological reduction. Put
$G^+=\ker(\lambda\bmod3)$ and $H^+=H\cap G^+$, and choose complex
conjugation $c\in H$. Then $G=G^+\rtimes\langle c\rangle$.
Crossed homomorphisms to $\mathbb F_3(1)$ trivial on $H$ correspond
to homomorphisms
\[
 \chi:G^+\longrightarrow\mathbb F_3,\qquad
 \chi|_{H^+}=0,\qquad \chi(cgc^{-1})=-\chi(g).
\]
Restriction gives one direction; the inverse is $d(gc^e)=\chi(g)$.
Thus the relevant group is the minus eigenspace in the relative
abelianization of an index-two subgroup. Ordinary transfer applies
$1+c$ and kills this eigenspace; it does not show that it is zero.
Complex conjugation is not central in $G$, so the argument for a
central element acting nontrivially on coefficients cannot be used.
Odd conjugation weights all agree with the cyclotomic weight modulo
three and supply no weight mismatch.

By contrast, the two-primary image is the subgroup $d_2(D)$, of
exponent at most four, because $G=DH$. The action on
$d_2(D)/2d_2(D)$ is trivial, so $d_2$ induces ordinary relative
quadratic characters. A nonzero group of exponent four has a nonzero
quotient by twice itself. This observation does not prove that the
resulting characters vanish.

\subsection{The integral cubic-frame comparison}

Put $N=\ker(\lambda,\rho_2)$ and $\nu_g=B_g(\omega,\omega)$,
where $\omega^2+\omega+1=0$. If
$[f_g]=A_{f_g}(X,Y)[x,y]$ in $F'/F''$, then, for multiplier one,
\[
 B_g(X,Y)=1+(X-1)(Y-1)A_{f_g}(X,Y).
\]

\begin{proposition}\label{prop:lr-integral-cubic-beta}
For $g\in N$, an integral identification of the quadratic orders
identifies $T_g$ with either $(-1)^{a_2(g)}\nu_g$ or
$(-1)^{a_2(g)}\overline{\nu_g}$. For $g\in D$ the sign is one.
No uniform orientation is needed for cube membership or for the
scalar modulo-eight bit.
\end{proposition}
\begin{proof}
We specify the groups in the full comparison of
\cite[Proposition~18]{H17}. Put
\[
 \Gamma=B_3/\langle(uv)^3\rangle\simeq C_2*C_3,\qquad
 \Gamma_{1,1}=B_3/\langle(uv)^6\rangle,\qquad
 \Gamma_{1,2}=B_4/\langle(abc)^4\rangle,
\]
where $a,b,c$ are the standard chain generators of $B_4$.
The forgetting map
$\operatorname{forg}:\widehat\Gamma_{1,2}\to
\widehat\Gamma_{1,1}$ sends $(a,b,c)$ to $(u,v,u)$.
Its kernel is the completed punctured-torus group $\widehat\Pi$,
free on
\[
 \xi_1=ca^{-1},\qquad
 \xi_2=abac^{-1}b^{-1}a^{-1}.
\]
On $\widehat\Pi^{\mathrm{ab}}=\widehat{\mathbb Z}^2$,
$a,c$ act by $U$ and $b$ by $L$.

For $g=(1,f_g)$ let $\mathcal A_g$ be the braid action on
$\widehat\Gamma_{1,2}$:
\[
 \mathcal A_g(a)=a,\quad
 \mathcal A_g(b)=f_g(a^2,b^2)^{-1}b f_g(a^2,b^2),\quad
 \mathcal A_g(c)=P_g^{-1}cP_g,\quad
 P_g=f_g(ba^2b,c^2).
\]
Let $\widetilde A_g=A_g(a^2,b^2)$ be the evaluation of the
pure cubic conjugator $A_g\in\widehat F(x,y)$, and put
$\Psi_g=\operatorname{Ad}(\widetilde A_g^{-1})\mathcal A_g$.
Writing $J_3=\operatorname{forg}^{-1}(\langle uv\rangle)$,
the cited proposition gives an isomorphism
$\varphi:\widehat\Gamma\to J_3$, induced by a discrete
isomorphism and carrying $\overline{\Gamma'}$ to
$\widehat\Pi$, and an $h\in J_3$ such that
\[
 \Psi_g\varphi=\operatorname{Ad}(h)\varphi\Phi_g .
\]
Here $\overline{\Gamma'}$ is the closure of the commutator
subgroup of $\Gamma$, whose abelianization is the hexagonal
lattice $\widehat{\mathbb Z}[\omega]$.
We take this identity on the whole integral homology.

Let $M_g$ be the action of $\mathcal A_g$ on
$\widehat\Pi^{\mathrm{ab}}$. Equivariance of the forgetting map
and the established normalized signed-congruence identity imply
that $M_g$ commutes with $\mathrm{SL}_2(\widehat{\mathbb Z})$.
The centralizer of its two elementary unipotents consists of
scalars. Moreover $P_g\in\widehat\Pi$: after forgetting,
its two arguments are $(uv)^3u^{-2}$ and $u^2$, which commute,
and $f_g$ is a commutator word. Since
$ba^2b=(ab)^3a^{-2}$ and $c$ commutes with $a$, the element
$P_g$ commutes with $a$. Thus
$\mathcal A_g(\xi_1)=P_g^{-1}\xi_1P_g$, so $M_g$ fixes this
primitive homology class and $M_g=I$.

The corrected action on homology is consequently
$q(A_g)^{-1}=T_g$, since $D_g=I$ on $N$. The other side is
\[
 R_+^j C\nu_gC^{-1},\qquad j\in\mathbb Z/6,
\]
where $R_+=UL$ and the discrete homology identification $C$ sends the order-six
lattice rotation to $R_+$ or $R_+^{-1}$.
Lemma~15 of \cite{H17} gives
$\nu_g\equiv(-1)^{a_2(g)}\pmod4$.
Since $T_g\equiv I\pmod4$ and $R_+$ has exact order six modulo
four, $j=3a_2(g)\pmod6$. The two embeddings of the quadratic
order differ by its involution, proving the assertion.
On $D$, the established sign identity is
$a_2(g)=\rho_3^I(g)/2\bmod2=0$.

The identification of the lattice scalar with $B_g(\omega,\omega)$
is integral: $[x,y]$ maps to the translation $-3$ in the affine
triangle quotient, so $-3\nu_g=-3B_g(\omega,\omega)$.
Multiplication by three is injective in
$\widehat{\mathbb Z}[\omega]$, including at three, and can be cancelled.
\end{proof}

\subsection{The order-three gate}

\begin{lemma}\label{lem:lr-cubic-cm-quotient}
The maximal exponent-three subextension of $K_3^{\mathrm{ab}}/k$
is $k(q^{1/3}:q\text{ rational prime})$. The order-three part
of $d$ is therefore zero if and only if
$T_g\in N_3(4)^3$ for every $g\in D$.
\end{lemma}
\begin{proof}
Let $A=\operatorname{Gal}(K_3^{\mathrm{ab}}/K_3)$ and let $c$
act by complex conjugation. The full group over $\mathbb Q$ is
$A\rtimes\langle c\rangle$, whose closed derived subgroup is
$(c-1)A=\operatorname{Gal}(K_3^{\mathrm{ab}}/k)$.
On the pro-three part the idempotents $(1\pm c)/2$ split $A$,
and $(c-1)A$ is precisely its minus part.
Every cubic character of this subgroup thus extends to $A$,
with $c$ acting by negation on the cyclic cubic quotient.

Kummer theory writes the corresponding field as
$K_3(\alpha^{1/3})$, where the class of $\alpha$ in
$K_3^\times/(K_3^\times)^3$ is fixed by conjugation: inversion
on the Galois quotient is already accounted for by the action
on $\mu_3$. A fixed class comes from $\mathbb Q$ because
$[N_{K_3/\mathbb Q}(\alpha)]=[\alpha]^2$, and squaring again
recovers $[\alpha]$ in this exponent-three group.
Conversely rational cube roots give such extensions.
Character separation proves the simultaneous field statement.
The normalized cubic torus is an isomorphic model of the Galois
group over $k$; its quadratic sign correction has no effect on
the quotient by cubes. On $D$ the radical restriction is trivial.
\end{proof}

\begin{lemma}\label{lem:lr-wild-three-cube}
For every $g\in D$, the three-adic component of $\nu_g$ is a cube
in the norm-one units of $\mathbb Z_3[\omega]$.
\end{lemma}
\begin{proof}
Put $\pi=\omega-1$. The class-two three-cycle relation gives
$3A_{f_g}(1,1)=0$, hence $A_{f_g}(1,1)=0$.
Symmetry gives a three-adic expansion
\[
 A_{f_g}(1+S,1+T)=a(S+T)+O((S,T)^2).
\]
In the degree-three cyclic-cover homology modulo three one has
$Y=1$, $X^3=1$, and $[x,y]\mapsto(X-1)y_0$.
Since $(X-1)^3=0$ in $\mathbb F_3[C_3]$, the image of $f_g$ is
$a(X-1)^2y_0=a(1+X+X^2)y_0$.
Its $y_0$ coefficient is $a$, proving
$\rho_3^I(g)\equiv-a\pmod3$. This is zero on $D$. Thus
\[
 \nu_g=1+\pi^2A_{f_g}(\omega,\omega)
      =1+2a\pi^3+O(\pi^4)\equiv1\pmod{\pi^4}.
\]
The unit has norm one by
Proposition~\ref{prop:lr-integral-cubic-beta}, and it is one
modulo three by \cite[Lemma~15]{H17}.
The logarithm identifies the norm-one units in $1+\pi^3O_3$
with $3\sqrt{-3}\,\mathbb Z_3$. Congruence modulo $\pi^4$
forces the logarithm into $9\sqrt{-3}\,\mathbb Z_3$, whose
valuation is five. This is three times the previous lattice;
exponentiation proves cube membership. The possible $\mu_3$
factor was already excluded by congruence to one modulo three.
\end{proof}

\begin{theorem}\label{thm:lr-tame-cubic-gate}
The order-three component of $d$ vanishes if and only if, for
all $g\in D$ and primes $\ell\ge5$,
\begin{equation}
 B_{g,\ell}(\omega,\omega)^{(\ell-\chi_{-3}(\ell))/3}=1
 \quad\text{in the residue norm-one torus}.
 \label{eq:lr-tame-cubic-test}
\end{equation}
Use a primitive cube root in $\mathbb F_\ell$ in the split
case and in $\mathbb F_{\ell^2}$ in the inert case.
\end{theorem}
\begin{proof}
Cubing is bijective on the pro-two factor $N_{3,2}(4)$.
The three-component is handled by the preceding lemma.
For $\ell\ge5$, local reduction has a pro-$\ell$ kernel and a
cyclic norm-one image of order $\ell-\chi_{-3}(\ell)$.
Cubing is bijective on the kernel, and its image on the residue
group is exactly the kernel of the displayed power.
Apply Proposition~\ref{prop:lr-integral-cubic-beta} and
Lemma~\ref{lem:lr-cubic-cm-quotient}. Inversion preserves the tests.
For the product, the cube image is the product of the local
cube images, so no simultaneous condition has been discarded.
\end{proof}

The Anderson relation $(A_\ell)$ is sufficient for
\eqref{eq:lr-tame-cubic-test}. In its beta form
\cite[Remark~2.4]{Furusho2003}, specialize at $s=1/3$ and
reduce in characteristic $\ell$. The $\ell$th roots in the
denominator all reduce to one, where $B_g(1,\omega)=1$.
The numerator contains $B_g(\omega,\omega)$ exactly
$(\ell-\chi_{-3}(\ell))/3$ times; all other factors are boundary
values equal to one. The arithmetic right side is one on $D$.
This implication does not assert $(A_\ell)$ for all of GT.

\subsection{A formal obstruction to removing the tame tests}

\begin{proposition}\label{prop:lr-formal-beta-obstruction}
Beta associativity, unitarity, all dyadic Anderson identities,
vanishing integral Kummer coordinates, and the logarithmic
derivatives of the Anderson identities do not, as formal integral
Alexander identities, imply all tests
\eqref{eq:lr-tame-cubic-test}.
\end{proposition}
\begin{proof}
Fix $\ell\ge5$. By the full affine radical image choose
$\sigma\in H^{\mathrm{cyc}}$ whose sole nonzero radical
translation is $\rho_\ell^I(\sigma)=2e_3$, with
$e_3\in\widehat{\mathbb Z}$ the coefficient-prime-three
idempotent. In particular $\rho_2^I(\sigma)=\rho_3^I(\sigma)=0$
and every $\ell$-adic Kummer coordinate is zero.
The arithmetic $(A_\ell)$ relation at $s=1/3$ gives
\[
 B_{\sigma,\ell}(\omega,\omega)^{
          (\ell-\chi_{-3}(\ell))/3}=\omega^{-2\ell}\ne1
 \quad\bmod\ell.
\]
Take the coefficient-$\ell$ component of $A_{f_\sigma}$ and
set every other coefficient-prime component to zero.
Its beta unit has component $B_{\sigma,\ell}$ at $\ell$ and
one elsewhere. This is integral in the completed group algebra.
Associativity, unitarity, and the dyadic identities hold
componentwise, because they hold for $\sigma$ and the identity.
Every cyclic-cover coefficient is a
$\widehat{\mathbb Z}$-linear functional on $F'/F''$.
Thus every Kummer coordinate of the patched data is zero.
The logarithmic derivatives here are taken in the coefficient-$p$
formal variable near each torsion specialization, for every
coefficient prime $p$. At $p=\ell$, write the variable as
$\zeta(1+S)$, with $\zeta$ of order prime to $\ell$.
The arithmetic Anderson discrepancy for $\sigma$ has principal
factor $(1+S)^{n\rho_{n,\ell}^I(\sigma)}=1$; its remaining
factor is constant torsion on that specialization.
Its logarithmic derivative is therefore zero. At every other
coefficient prime the patched unit is one. This proves all the
stated derivative identities while leaving the displayed
nontrivial constant $\omega^{-2\ell}$ intact.
The displayed tame test still fails.

This constructs formal Alexander data, not a lift satisfying the
full profinite pentagon. It also does not assert the full
nonabelian $(K_n)$ identities for a lift; these contain information
beyond the numerical Kummer and logarithmic-derivative consequences.
The example isolates a missing constraint in the proposed formal
inference without supplying a counterexample inside GT.
\end{proof}

\subsection{The quartic coordinates}

\begin{lemma}\label{lem:lr-quartic-cm-quotient}
The maximal elementary-two subextension of
$K_4^{\mathrm{ab}}/k$ is
\[k(q^{1/4}:q\text{ rational prime}).\]
\end{lemma}
\begin{proof}
Use CM reciprocity in the primary-unit realization
\[
 \operatorname{Gal}(K_4^{\mathrm{ab}}/K_4)
 \simeq\{u\in\widehat O_4^\times:
                    u_2\equiv1\pmod{(1+i)^3}\},
\]
whose commutator image is $N_4(4)$.
The local two-component is isomorphic by logarithm to
$4i\mathbb Z_2$ and has one quadratic character.
Each odd component has one quadratic character, since its
residue norm-one group is cyclic of order
$q-\chi_{-4}(q)$ and its principal kernel is pro-$q$.

The character of $2^{1/4}$ is supported only at two and is
nontrivial over $k$, so it spans the first character space.
For odd $q$, $K_4(q^{1/4})/K_4$ has conductor supported at
two and $q$, and its tame character at $q$ has order four.
On norm-one units it restricts to the unique nontrivial
quadratic character: writing such a unit as $u/\bar u$,
complex conjugation inverts the Kummer character, so its value
is the square of the order-four tame character of $u$.
The local ratio map is surjective.
Thus the restriction is the tame $q$-character plus a possible
two-component character. Conductor support and the primary-unit
form of reciprocity exclude every other local component.
These characters, together with the character of $2^{1/4}$,
span the continuous quadratic character space of the product.
Character separation proves the result. The spanning argument
uses local support, not a comparison of infinite dimensions.
\end{proof}

On $N$ one has $W_{g,2}\equiv I\pmod8$. Indeed
$\beta_g\in F'$ and $D_g=I$, and modulo eight the commutator
subgroup generated by $U^2,L^2$ is $\{I,5I\}$.
The reduction $5I$ cannot lift to a norm-one element $a+bi$:
if $a\equiv5$ and $b\equiv0\pmod8$, then
$a^2+b^2\equiv9\pmod{16}$. This observation is subsumed by
the separate stronger prime-two compatibility theorem.
Consequently the local two-component of $W_g$ is a square.
At three, the established sign identity on $D$ places
$W_{g,3}\bmod3$ in $\{I,-I\}$, the square subgroup of the
cyclic order-four residue torus.
Lemma~\ref{lem:lr-quartic-cm-quotient} now gives the exact
simultaneous reduction: all rational quartic compatibilities
hold if and only if
\begin{equation}
 W_{g,\ell}^{(\ell-\chi_{-4}(\ell))/2}=1
 \quad\text{in the residue norm-one torus}
 \qquad(g\in D,\ \ell\ge5).
 \label{eq:lr-tame-quartic-test}
\end{equation}
Proposition~\ref{prop:tame-harmonic-word-gate} expresses these
tests directly as finite word evaluations by the normalized
harmonic identity. Their vanishing remains open.

\subsection{The exceptional eighth root of three}

Put $E_3:y^2=x^3+1$, $\alpha=3^{1/4}$, and
$v=(2+\sqrt3)^{1/4}$, using positive roots.

\begin{lemma}\label{lem:lr-cubic-eight-field}
One has $k(E_3[8])=k(\alpha,v)$.
The scalar automorphism $[-1]$ of this division field fixes $v$.
On $H^{\mathrm{cyc}}$, if
\[
 T_{\sigma,2}\equiv(1+4\tau(\sigma))I\pmod8,
 \qquad \tau(\sigma)\in\mathbb F_2,
\]
then $\sigma(v)/v=(-1)^{\tau(\sigma)}$.
\end{lemma}
\begin{proof}
Write $e_j=-\omega^j$, for $j=0,1,2$.
The abscissae of the non-two-torsion points in $E_3[4]$ are
$\omega^j(-1+\sqrt3)$ and $\omega^j(-1-\sqrt3)$.
The established four-division calculation is
$k(E_3[4])=k(\alpha)$.
To halve $P=(x_P,y_P)$ one adjoins square roots of $x_P-e_j$
with product of the prescribed sign $-y_P$.
Conversely these square roots belong to the field of a half
$Q=(x_Q,y_Q)$, because duplication gives
\[
 x_P-e_j=
 \left(\frac{(x_Q-e_j)^2-(e_j-e_{j+1})(e_j-e_{j+2})}
 {2y_Q}\right)^2.
\]
Indices are modulo three. This supplies both directions of
the square-root criterion inside the full division field.

For $x_P=-1+\sqrt3$, set $A=\sqrt3-1+\omega$.
Since $v^{-2}=\sqrt{2-\sqrt3}$, direct multiplication gives
\[
 A=\zeta_{24}^{\,5}\sqrt3\,v^{-2},\qquad
 \sqrt A=\zeta_{48}^{\,5}\alpha/v
 \quad\text{up to sign}.
\]
Hence $v\in k(E_3[8])$.
Conversely every square root needed to halve the displayed
four-torsion abscissae is a root of unity times
$\alpha$, $\alpha/v$, or $\alpha v$.
For the minus choice the relevant identity is
\[
 -\sqrt3-1+\omega=\zeta_{24}^{\,11}\sqrt3\,v^2.
\]
Conjugating and multiplying by roots of unity gives all cases.
Thus the halving formulas place all of $E_3[8]$ in $k(\alpha,v)$.

The two quadratic extensions $k(\alpha)$ and $k(v)$ are
nontrivial and distinct. The positive numbers $\alpha,v,v/\alpha$
each have a nonreal conjugate, whereas any real element of
$k$ generates a totally real abelian field.
For $v/\alpha$, the automorphism of
$\mathbb Q(i,\sqrt2,\sqrt3)$ fixing $\sqrt3$ and changing
the sign of $\sqrt2$ negates $(v/\alpha)^2$.
The same criterion, or the defining polynomials, handles
$\alpha$ and $v$. Hence $[k(\alpha,v):k]=4$.

Under $[-1]$, the four-torsion ordinate changes sign,
so $\alpha$ changes sign. The recovered square root in
the duplication formula also changes sign when $Q$ is
replaced by $-Q$. Their quotient $v$ is therefore fixed.
The unnormalized cubic norm-one image modulo eight,
one modulo two, consists of the scalar units $1,3,5,7$.
The scalar $-1=7$ fixes $v$, and the other nontrivial
automorphism fixing $E_3[4]$, the scalar $5$, changes its sign.
The normalization from the Tate representation to $T$
multiplies by a scalar sign, which fixes $v$.
Its image modulo eight is $\{I,5I\}$, proving the assertion.
\end{proof}

For $g\in D$ define $w_3(g),\tau_3(g)\in\mathbb F_2$ by
\[
 W_{g,3}\bmod3=(-I)^{w_3(g)},\qquad
 T_{g,2}\bmod8=(1+4\tau_3(g))I.
\]
The first uses the known vanishing of the individual frame
parity; the second follows from norm one and congruence one
modulo four. Both are homomorphisms on $D$.

\begin{proposition}\label{prop:lr-exceptional-three-character}
The character $\chi_3=w_3+\tau_3$ on $D$ is $G$-invariant
and zero on $D\cap H$. It extends uniquely to an ordinary
continuous quadratic character on $G$ that is zero on $H$.
Its vanishing is exactly the remaining compatibility for
$3^{1/8}$, once the known quartic coordinate of three is imposed.
Moreover
\[
 \tau_3(g)=\frac{B_{g,2}(\omega,\omega)-1}{4}\pmod2,
\]
where the right side is scalar, in
$\mathbb F_2\subset\mathbb F_4$.
\end{proposition}
\begin{proof}
The harmonic frame extends to a matrix representation on the
whole gentle group, so conjugation fixes its scalar classes
$\{I,-I\}$ modulo three. The full cubic-frame conjugation
formula fixes the scalar classes $\{I,5I\}$ modulo eight.
Thus both bits are $G$-invariant.

Use the arithmetic witness in the intersection theorem:
an ordinate $y$ of a three-torsion point of
$E_4:y^2=x^3-x$ satisfies
\[
 y^4=\frac49\sqrt3(2+\sqrt3),\qquad
 z=\frac{y}{\sqrt{2/3}\,v},\qquad z^8=3.
\]
For $\sigma\in D\cap H$, the harmonic residue frame sends
$y$ to $(-1)^{w_3(\sigma)}y$, while
Lemma~\ref{lem:lr-cubic-eight-field} gives the sign
$(-1)^{\tau_3(\sigma)}$ on $v$.
Since $\sigma$ fixes every rational radical, it fixes $z$,
so the sum of the two bits is zero.

The extension is $\chi_3(dh)=\chi_3(d)$ for $d\in D,h\in H$.
Invariance and vanishing on the intersection prove
well-definedness and multiplicativity. Continuity follows
from the compact surjection $D\times H\to G$.
On $D$, the radical coordinate and the quartic coordinate of
three are zero, and the sign on the eighth root is exactly
this sum, proving its interpretation as the remaining gate.

Finally the integral cubic-frame comparison identifies $T_g$
with $B_g(\omega,\omega)$ or its conjugate.
Modulo eight a norm-one unit that is one modulo four is
the scalar $1$ or $5$. Conjugation changes neither, proving
the beta formula. Its vanishing on all of $G$ is proved in
Section~\ref{sec:normalized-cubic-identity}.
\end{proof}

\subsection{Status}

After the separate prime-two compatibility theorem, joint
compatibility is reduced to the tame cubic tests
\eqref{eq:lr-tame-cubic-test}, the tame harmonic quadratic
tests \eqref{eq:lr-tame-quartic-test}, and the single relative
character $\chi_3$. The wild three-adic cubic part is eliminated.
The two subsequent frame comparisons also give $\chi_3=0$,
so only the tame tests remain.
The exceptional eighth root of three is represented by a
concrete finite combination of cubic eight-division and
harmonic three-division frames. The individual quotients,
transfer, and the formal beta identities considered here do
not force the remaining tests to vanish. Full joint
factorization remains unproved.

\section{A finite interchange map for the residual character at three}

This section compares the exceptional character $\chi_3$ with two
finite specializations of the harmonic and equianharmonic formulas of
Nakamura and Tsunogai.  An explicit homomorphism between finite matrix
groups places both comparisons in the same quaternion central sign.
Furusho's dyadic identities will identify the harmonic comparison
on full GT in Section~\ref{sec:harmonic-dyadic-comparison}. The
cubic comparison follows in Section~\ref{sec:normalized-cubic-identity}.

Write $G=\widehat{\mathrm{GT}}$ and let
$D=\ker\lambda\cap\bigcap_q\ker\rho_q^I$.  The pure harmonic and
cubic conjugators are denoted $\beta_g,A_g\in\widehat F(x,y)$,
respectively, where $x=u^2,y=v^2$ in the modular group.
We use
\[
 U=\begin{pmatrix}1&1\\0&1\end{pmatrix},\quad
 L=\begin{pmatrix}1&0\\-1&1\end{pmatrix},\quad
 S=ULU,\quad R=-UL.
\]
On $D$ the established signed normalization is $D_g=I$, and the
established finite frame targets give
\[
 q_3(\beta_g)^{-1}=(-I)^{w(g)},\qquad
 q_8(A_g)^{-1}=5^{\tau(g)}I,\qquad
 \chi_3(g)=w(g)+\tau(g).
\]
Here $q_n$ is the signed braid representation reduced modulo $n$.
All signs in this section are in $\mathbb F_2$.  In particular there is
no extra minus sign in the definition of the harmonic frame.

\begin{proposition}[Finite interchange map]\label{prop:finite-interchange}
Put $X=U^2\bmod8$ and
$P=\langle X,R\bmod8\rangle\subset\mathrm{SL}_2(\mathbb Z/8)$.
There is a surjection
\[
 \kappa:P\longrightarrow\mathrm{SL}_2(\mathbb F_3)
\]
such that
\[
 \kappa(X)=-S\bmod3,\qquad
 \kappa(R)=U\bmod3,\qquad
 \kappa(5I)=-I.
\]
The group $P$ has order $192$ and the kernel has order eight.
For every profinite commutator word $f\in\widehat F(x,y)'$,
\[
 \kappa\bigl(f(X,R)\bigr)=f(S,U)=f(U,S)^{-1}
\]
if $f$ satisfies the two-cycle relation.
\end{proposition}
\begin{proof}
Put $Y=L^2\bmod8$, $C=5I$ and $Z=-I$.  Direct matrix
multiplication gives
\[
 [X,Y]=C,\qquad X^4=Y^4=C^2=Z^2=I,
\]
and $X^2,Y^2,C,Z$ are central in $N=\langle X,Y,Z\rangle$.
Consequently every element of $N$ has the form
$X^aY^bC^eZ^d$ with $0\le a,b<4$ and $e,d\in\{0,1\}$.
These 64 forms are distinct.  Indeed, if one equals $I$, its two
off-diagonal entries give $2a=2b=0\bmod8$, so $a=b=0$;
then $5^e7^d=1\bmod8$ gives $e=d=0$.
This also proves that the displayed class-two presentation is a
presentation of $N$.

Conjugation by $R$ has the following action:
\[
 X\longmapsto Y,\qquad
 Y\longmapsto Z(XY)^{-1},\qquad C\longmapsto C,\qquad
 Z\longmapsto Z.
\]
Since $R^3=I$ and its reduction modulo two has order three,
whereas every element of $N$ is one modulo two, one has
$P=N\rtimes\langle R\rangle$, of order $64\cdot3$.
Here $Y=RXR^{-1}$ and
$Z=(RYR^{-1})XY$, so this semidirect product is generated by
$X,R$ as asserted.

In $\mathrm{SL}_2(\mathbb F_3)$ put
\[
 i=-S=\begin{pmatrix}0&2\\1&0\end{pmatrix},\qquad
 t=U,\qquad
 j=tit^{-1}=\begin{pmatrix}1&1\\1&2\end{pmatrix}.
\]
Then $i^2=j^2=[i,j]=-I$, $t^3=I$, and
$tjt^{-1}=ij$.  Thus $i,j$ generate the quaternion group and
$t$ acts on its three noncentral cyclic subgroups by a three-cycle.
The assignment
\[
 X\mapsto i,\quad Y\mapsto j,\quad C\mapsto-I,\quad Z\mapsto-I
\]
respects the presentation of $N$.  It also respects conjugation by
$R$, because
$(-I)(ij)^{-1}=ij$.  It therefore extends with $R\mapsto t$
to a homomorphism $P\to Q_8\rtimes C_3$
$=\mathrm{SL}_2(\mathbb F_3)$.  It is surjective, and its kernel
has order $192/24=8$.

Evaluation of a profinite word commutes with a homomorphism to a
finite group.  Multiplying either argument of a commutator word
by a central element does not change its evaluation, since the
two exponent sums are zero.  Hence
$\kappa(f(X,R))=f(-S,U)=f(S,U)$, and the two-cycle relation gives
the final equality.
\end{proof}

\begin{lemma}[A common finite evaluation]\label{lem:common-sign}
For every $g\in D$ the matrix $f_g(U,S)\bmod3$ is central.
Consequently
\[
 f_g(U,S)=(-I)^{\eta(g)}\quad\text{in }\mathrm{SL}_2(\mathbb F_3)
\]
defines a continuous homomorphism $\eta:D\to\mathbb F_2$.
\end{lemma}
\begin{proof}
Write $[f_g]=A_f(X,Y)[x,y]$ in the completed Alexander module.
The two-cycle relation gives $A_f(X,Y)=A_f(Y,X)$.
We recall precisely the cyclic-cover input already established for
full GT.  In the homology of the subgroup
$\ker(\widehat F(x,y)\to C_3)$ with $x\mapsto1,y\mapsto0$,
use the basis $x^3,y_j=x^jyx^{-j}$, $0\le j<3$.
The three $y_j$ coefficients of $f_g$ are
$(2b,-b,-b)$, where $b=-\rho_3^I(g)/2$.
Sending $y_j$ to $\omega^j$ gives
\[
 (\omega-1)A_f(\omega,1)=3b.
\]
On $D$, $b=0$.  Since $\omega-1$ is a unit in
$\mathbb Z_2[\omega]$, it follows that
$A_f(\omega,1)=A_f(1,\omega)=0$, also after reduction modulo two.

Now project $\mathrm{SL}_2(\mathbb F_3)$ to
$A_4=V\rtimes C_3$, where $V=Q_8/\{\pm I\}$ is identified
with the additive group of $\mathbb F_4$.
The image of $S$ belongs to $V$ and the image of $U$ acts on
$V$ by a primitive cube root $\omega$ or $\omega^{-1}$.
Conjugation by $S$ is trivial on $V$.  Since this quotient is
metabelian, the image of $f_g(S,U)$ is
$A_f(1,\omega)[S,U]$, with either choice of the root; it is zero.
Thus $f_g(S,U)$, and hence $f_g(U,S)$ by the two-cycle relation,
lies in $\{\pm I\}$.

For $g\in D$, the free-group automorphism
$x\mapsto x$, $y\mapsto f_g^{-1}yf_g$ induces the identity under
the evaluation $x\mapsto U,y\mapsto S$, because the evaluated
conjugating word is central.  Applying this observation to the
composition law $f_{gh}=f_g\phi_g(f_h)$ proves that the evaluated
signs multiply.  Continuity follows from evaluation in the finite
group $\mathrm{SL}_2(\mathbb F_3)$.
\end{proof}

\begin{proposition}[An exact pair of finite comparison gates]
\label{prop:pair-of-gates}
For $g\in D$ define
\[
 \epsilon_h(g)=w(g)+\eta(g),\qquad
 \epsilon_c(g)=\tau(g)+\eta(g).
\]
These are continuous homomorphisms on $D$, and
\[
 \boxed{\quad \chi_3=\epsilon_h+\epsilon_c.\quad}
\]
More explicitly, the two signs are the following finite defects:
\[
 \begin{split}
 (-I)^{\epsilon_h(g)}
   &=q_3(\beta_g)^{-1}f_g(U,S)^{-1},\\
 (-I)^{\epsilon_c(g)}
   &=\kappa\!\left(q_8(A_g)^{-1}f_g(X,R)^{-1}\right).
 \end{split}
\]
In particular the two finite comparisons
\begin{equation}
 q_3(\beta_g)^{-1}=f_g(U,S),\qquad
 \kappa(q_8(A_g)^{-1})=\kappa(f_g(X,R))
 \label{eq:finite-harmonic-cubic-comparisons}
\end{equation}
would imply $\chi_3(g)=0$.  They are sufficient conditions;
vanishing of $\chi_3$ requires only equality of their two defects.
Both defects vanish on $D\cap G_{\mathbb Q}$.
\end{proposition}
\begin{proof}
The established frame homomorphisms make $w,\tau$ homomorphisms
on $D$, and Lemma~\ref{lem:common-sign} does the same for $\eta$.
The first finite formula follows directly from their definitions.
For the second, Proposition~\ref{prop:finite-interchange} gives
\[
 \kappa(q_8(A_g)^{-1})=(-I)^{\tau(g)},\qquad
 \kappa(f_g(X,R))=f_g(U,S)^{-1}=(-I)^{\eta(g)}.
\]
This also proves the displayed sum and the sufficient conditions.

For the arithmetic assertion use exactly the scope of
\cite[Theorem~2, pp.~199--200]{NT2002}.  Its $(\mathrm{GF}_0)$ and
$(\mathrm{HF}_1)$ formulas are stated for Galois elements.
When $\lambda=1$ and $\rho_2=\rho_3=0$, they reduce to
\[
 \beta_g^{-1}=f_g(u,s),\qquad A_g^{-1}=f_g(u^2,r).
\]
Here its harmonic parameter is $\beta^{-1}$.  Its cubic
parameter is $A^{-1}$, since $\operatorname{Ad}(r)$ cyclically
permutes $x,y,(xy)^{-1}$ and
$f_g=\operatorname{Ad}(r)(A_g)A_g^{-1}$.
These identifications also follow from the uniqueness of the
pure conjugators.  The equalities lift without a central error
from the modular quotient to the signed braid group: both
sides have degree zero in the braid abelianization, whereas
the central generator has degree six.
Under the signed representation $q(r)=UL=-R$; the central
minus sign disappears in a commutator word.
Reduction modulo three and eight therefore gives
\eqref{eq:finite-harmonic-cubic-comparisons} on the arithmetic
subgroup.  No use of those two formulas on all of GT is made.
\end{proof}

\subsection{Two relative quadratic characters}

Continue with the same notation:
$G=\widehat{\mathrm{GT}}$, $H=G_{\mathbb Q}$,
$D=\ker\lambda\cap\bigcap_q\ker\rho_q^I$, and
$f_g(U,S)=(-I)^{\eta(g)}$ for $g\in D$.
The two frame signs are $w,\tau$, and
$\epsilon_h=w+\eta$, $\epsilon_c=\tau+\eta$.
We use the previously established facts $G=DH$, the invariance
of $D$, and the signed frame conjugation laws.  The latter make
$w$ and $\tau$ invariant under $G$: on $D$ the finite frames are
scalar matrices, which every conjugation fixes.

\begin{lemma}\label{lem:eta-quotient-invariant}
The kernel of the evaluation
\[
 e:\widehat F(x,y)\longrightarrow Q=\mathrm{SL}_2(\mathbb F_3),
 \qquad e(x)=U,\quad e(y)=S,
\]
is invariant under every gentle automorphism.  Consequently $\eta$
is invariant under conjugation by all of $G$.
\end{lemma}
\begin{proof}
We first verify the precise finite presentation:
\begin{equation}\label{eq:finite-quotient-presentation}
Q=\langle x,y\mid x^3=y^4=(xy)^3=1,\ [x,y^2]=1\rangle.
\end{equation}
In the presented group, $z=y^2$ is central and has order at most
two.  Modulo $z$ one has
$\langle x,y\mid x^3=y^2=(xy)^3=1\rangle$, of order at most
twelve.  For an explicit bound, put
$b_j=x^jyx^{-j}$ for $j=0,1,2$.  They are involutions, their
product in cyclic order is one by $(xy)^3=1$, and hence two
of them commute and generate a group of order at most four.
This group is normalized by $x$, which cyclically permutes the
$b_j$.  Thus the whole quotient has order at most $4\cdot3$.
The group in \eqref{eq:finite-quotient-presentation} has order at most 24.  The matrices $U,S$ satisfy
all four relators.  They generate $\mathrm{SL}_2(\mathbb F_3)$:
$-S$ and its conjugate by $U$ generate the quaternion subgroup,
and $U$ has order three and permutes its three noncentral cyclic
subgroups.  This gives 24 distinct matrices, proving \eqref{eq:finite-quotient-presentation}.

Let $g=(\lambda,f)$ be a gentle element and use the convention
\[
 \phi_g(x)=x^\lambda,\qquad
 \phi_g(y)=f^{-1}y^\lambda f.
\]
The peripheral three-cycle condition says that $\phi_g(xy)$
is conjugate to $(xy)^\lambda$.  Therefore the images under
$e\phi_g$ satisfy the first three relators of \eqref{eq:finite-quotient-presentation}.
Moreover
\[
 e\phi_g(y^2)=e(f)^{-1}(S^2)^\lambda e(f)=-I,
\]
since $S^2=-I$ is central and $\lambda$ is odd.
This proves the fourth relator as well.  There is consequently
an induced endomorphism $\alpha_g$ of $Q$ with
$e\phi_g=\alpha_ge$.  Applying the same argument to $g^{-1}$
shows that $\alpha_g$ is an automorphism.
In particular every $\alpha_g$ fixes the central involution
$-I$, directly also visible in the displayed formula.

Put $s(g)=e(f_g)$.  The composition convention is
$f_{gh}=f_g\phi_g(f_h)$, so
\[
 s(gh)=s(g)\alpha_g(s(h)).
\]
For $d\in D$, Lemma~\ref{lem:common-sign} proves
$s(d)\in\{\pm I\}$.  Since $\lambda_d=1$, the formulas on
the two generators then give $\alpha_d=\mathrm{id}_Q$.
For $g\in G$ and $d\in D$, it follows that
\[
 \begin{split}
 s(gdg^{-1})
 &=s(g)\alpha_g(s(d))\alpha_{gd}(s(g^{-1}))\\
 &=s(g)s(d)\alpha_g(s(g^{-1}))=s(d).
 \end{split}
\]
The last equality uses centrality of $s(d)$ and
$s(g)\alpha_g(s(g^{-1}))=1$.
This proves the conjugation invariance of $\eta$.
\end{proof}

\begin{proposition}\label{prop:two-relative-extensions}
Each of $\epsilon_h,\epsilon_c$ has a unique extension to a
continuous homomorphism $G\to\mathbb F_2$ that is zero on $H$.
Their sum is the established residual relative character $\chi_3$.
The harmonic extension is proved to vanish in
Section~\ref{sec:harmonic-dyadic-comparison}. The cubic extension
and $\chi_3$ vanish by Section~\ref{sec:normalized-cubic-identity}.
\end{proposition}
\begin{proof}
By Lemma~\ref{lem:eta-quotient-invariant} and the scalar frame
conjugation laws, both homomorphisms on $D$ are invariant under
$G$.  Proposition~\ref{prop:pair-of-gates}, using only the Galois
scope of the published $(\mathrm{GF}_0)$ and $(\mathrm{HF}_1)$
formulas, proves that both are zero on $D\cap H$.

For either homomorphism $\epsilon$, define
\[
 \widetilde\epsilon(dh)=\epsilon(d),\qquad d\in D,\ h\in H.
\]
If $dh=d'h'$, then $(d')^{-1}d\in D\cap H$, so this is
well-defined.  Conjugation invariance and normality of $D$ give
\[
 \widetilde\epsilon(d_1h_1d_2h_2)
 =\epsilon(d_1)+\epsilon(h_1d_2h_1^{-1})
 =\epsilon(d_1)+\epsilon(d_2).
\]
Thus it is a homomorphism, zero on $H$.
The continuous surjection $D\times H\to G$ is a quotient map,
because its source is compact and its target is Hausdorff.
Its pullback is the continuous function $(d,h)\mapsto\epsilon(d)$,
which proves continuity.  The decomposition $G=DH$ proves
uniqueness.  The two extensions sum to $w+\tau$ on $D$ and
are zero on $H$, hence their sum is $\chi_3$ by the same uniqueness.
\end{proof}

\subsection{A limitation of a low-degree peripheral calculation}

The following completely explicit test records why checking the
modular order-three relation through degree four does not remove
the central ambiguity.  It is not a counterexample to that relation
in the full profinite group and is not a GT counterexample.

Let $\Pi=\langle a,b\rangle$ be the discrete free group, with
automorphisms
\[
 \theta(a)=a^{-1},\quad\theta(b)=b^{-1},\qquad
 \alpha(a)=a^{-1}b,\quad\alpha(b)=a^{-1}.
\]
These are the actions of the order-two and order-three modular
generators on the free subgroup $\Gamma'$ for the basis
$a=srsr^{-1}$, $b=sr^2sr^{-2}$.
For a candidate conjugator $\beta$, fixing $sr$ makes the
image of $r$ equal to $dr$, where
$d=\beta\theta(\beta)^{-1}$.  Its cube is one exactly when
$d\alpha(d)\alpha^2(d)=1$.

\begin{proposition}[A truncated peripheral counterexample]
\label{prop:truncated-counterexample}
In the mod-two Magnus ring
$\mathcal R=\mathbb F_2\langle A,B\rangle/(\text{degree}\ge5)$,
send $a$ to $1+A$ and $b$ to $1+B$.
The word $\beta=a^3bab^3$ has exponent sums $(4,4)$ and
satisfies
\[
 d\alpha(d)\alpha^2(d)=1\quad\text{in }\mathcal R.
\]
Nevertheless its image under $a\mapsto i,b\mapsto j$ in the
quaternion group is $-1$.  Thus this truncated norm equation
and the two exponent sums modulo four do not determine the
quaternion central sign.
\end{proposition}
\begin{proof}
The exponent sums are immediate, and
$i^3jij^3=(ij)^2=-1$.
Here are direct expansions verifying the norm equation.  Write
$d=1+D$, $\alpha(d)=1+E$, $\alpha^2(d)=1+F$ in $\mathcal R$.
Substitute $a^{-1}=1+A+A^2+A^3+A^4$ and the analogous
expression for $b^{-1}$.  One obtains
\[
\begin{split}
D={}&AAB+ABB+BAA+BBA\\
 &+AAAB+AABA+AABB+ABAA+ABBA+ABBB\\
 &+BAAA+BAAB+BABB+BBAA+BBAB+BBBA,\\
E={}&ABB+BBA+AABB+ABAB+ABBB+BAAB\\
 &+BABA+BABB+BBAB+BBBA,\\
F={}&AAB+BAA+AAAB+AABA+ABAA+ABAB\\
 &+ABBA+BAAA+BABA+BBAA.
\end{split}
\]
Every monomial in $D+E+F$ occurs twice, so the sum is zero.
Each of $D,E,F$ has degree at least three, hence all their
pairwise products vanish in $\mathcal R$.  Therefore
$(1+D)(1+E)(1+F)=1$.
The trivial word has the same exponent sums modulo four
and passes the same norm equation, but has quaternion image
$+1$.  This proves the claimed limitation.
\end{proof}

The displayed expansion proves only the limitation of this truncated
norm equation.  It verifies neither the full pentagon nor the full
nonabelian distribution relation, and does not produce a profinite
GT element.

\subsection{Two-dimensional specializations of relation $(\mathrm{III}^\prime)$}

Write
\[
 B_4=\langle a,b,c\mid aba=bab,\ bcb=cbc,\ ac=ca\rangle.
\]
In the notation of \cite[\S10, Theorem~20]{H17}, the two inputs
to the harmonic auxiliary word $g_g=\beta_g^{-1}$ are
\[
 x_{45}=(ab)^3,\qquad x_{51}=(bc)^3.
\]
On the subgroup $D$, $g_g\in\widehat F(x,y)'$.

\begin{lemma}\label{lem:iii-prime-two-dimensional}
Let $k$ be any field and let $\rho:B_4\to\mathrm{GL}_2(k)$ be
a homomorphism.  Then $\rho((ab)^3)$ and $\rho((bc)^3)$ commute.
For a finite field $k$ and $g\in D$, it follows that
\[
 g_g\bigl(\rho((ab)^3),\rho((bc)^3)\bigr)=1.
\]
Thus such a direct matrix specialization of $(\mathrm{III}')$
does not retain its harmonic auxiliary-word term.
\end{lemma}
\begin{proof}
The elementary braid identity says that $(ab)^3$ is central in
$\langle a,b\rangle$, and $(bc)^3$ is central in
$\langle b,c\rangle$.  Both displayed images therefore belong
to the centralizer of $B=\rho(b)$.

If $B$ is not scalar, its centralizer in $M_2(k)$ is the
commutative algebra $k[B]$.  To see this directly, choose a vector
$v$ such that $v,Bv$ are a basis.  Such a vector exists for every
nonscalar two-by-two matrix.  An operator commuting with $B$ is
determined by its value on $v$, and the operators $\alpha I+\beta B$
give every possible such value.  In particular the two centralizer
elements commute.

If $B=tI$, the braid relation with $A=\rho(a)$ gives
$tA^2=t^2A$, and invertibility implies $A=tI$.  The other braid
relation similarly gives $\rho(c)=tI$.  All images then commute.

Over a finite field the representation extends continuously to
the profinite completion.  Evaluation of a commutator proword on
two commuting elements is one.  Applying this to $g_g$ proves the
second assertion.
\end{proof}

This lemma only excludes this direct way of retaining the
$g_g$ term.  It does not assert that $(\mathrm{III}')$ has no
consequences for the finite defects, nor that higher-dimensional
representations or comparisons of covering groups cannot provide
the required equality.  No vanishing or nonvanishing statement for
the two defects follows from it.


\subsection{A genus-two quotient carrying the cubic comparison sign}

Use the finite-comparison notation
$U=\left(\begin{smallmatrix}1&1\\0&1\end{smallmatrix}\right)$,
$L=\left(\begin{smallmatrix}1&0\\-1&1\end{smallmatrix}\right)$,
$R=-UL$, $X=U^2$, $Y=L^2$, $C=5I$, and $Z=-I$ modulo eight.
Write $G_8=\mathrm{SL}_2(\mathbb Z/8)$ and
$P=\langle X,R\rangle$, of order $192$.
The already proved presentation has
\[
 N=\langle X,Y,C,Z\rangle,\quad |N|=64,\quad
 P=N\rtimes\langle R\rangle,
\]
with the unique forms $X^aY^bC^eZ^d$,
$0\leq a,b<4$, $0\leq e,d<2$, and
\[
 [X,Y]=C,\qquad RXR^{-1}=Y,\qquad
 RYR^{-1}=Z(XY)^{-1}.
\]
The elements $C,Z,X^2,Y^2$ are central in $N$.

In the target put
\[
 S=\begin{pmatrix}0&1\\-1&0\end{pmatrix},\quad
 T=\begin{pmatrix}1&1\\0&1\end{pmatrix},\quad
 A=I-S=\begin{pmatrix}1&2\\1&1\end{pmatrix}
 \quad\text{over }\mathbb F_3.
\]

\begin{proposition}\label{prop:genus-two-finite-quotient}
There is a surjection
\[
 \ell:G_8\longrightarrow\mathrm{GL}_2(\mathbb F_3),\qquad
 U\longmapsto A,\quad R\longmapsto T,
\]
whose kernel is the order-eight subgroup
\[
 K=\langle Z,X^2C,Y^2C\rangle.
\]
It kills $Z$, so the signed congruence representation induces
\[
 \bar\ell:\Gamma=C_2*C_3\longrightarrow
                 \mathrm{GL}_2(\mathbb F_3).
\]
For $s=uvu$, $r=uv$ and $u=r^{-1}s$, their images have orders
$2,3,8$, respectively.  On the subgroup $P$, this map agrees
with the previous interchange map on all commutator prowords in
$X,R$, although its value on $X$ is $S$ instead of $-S$.
\end{proposition}
\begin{proof}
The presentation of $P$ defines a character
$\chi:P\to\{\pm1\}$ by
\[
 \chi(X)=\chi(Y)=\chi(Z)=-1,\qquad
 \chi(C)=\chi(R)=1.
\]
For example the only nontrivial conjugation check is
$\chi(Y)=\chi(Z(XY)^{-1})=-1$; the remaining relations are
immediate.  Multiply the previously established interchange
homomorphism $\kappa:P\to\mathrm{SL}_2(\mathbb F_3)$ by this
central character.  The resulting surjection $\kappa'$ has
\[
 \kappa'(X)=S,\quad \kappa'(R)=T,\quad
 \kappa'(C)=-I,\quad\kappa'(Z)=I.
\]
Its kernel contains $K$.  The unique normal forms show that
$Z,X^2C,Y^2C$ are three independent commuting involutions, so
$|K|=8$.  Surjectivity and $|P|=192$ imply that this is the whole
kernel.  A character is trivial on commutators, giving the last
assertion about prowords.

The subgroup $P$ has index two in $G_8$, and $U^2=X\in P$.
The whole group is generated by $P$ and $U$: indeed
$L=U^{-1}ZR$.  Direct multiplication gives
\begin{equation}\label{eq:bolza-1}
UXU^{-1}=X,\qquad URU^{-1}=XYR^2.
\end{equation}
In the target, $A^2=S$, $A$ commutes with $S$, and
\begin{equation}\label{eq:bolza-2}
ATA^{-1}=S(TST^{-1})T^2.
\end{equation}
For clarity, both sides of \eqref{eq:bolza-2} are
$\left(\begin{smallmatrix}2&2\\1&0\end{smallmatrix}\right)$.
Since $X,R$ generate $P$, \eqref{eq:bolza-1}--\eqref{eq:bolza-2} show that conjugation by
$A$ on $\kappa'(P)$ agrees with conjugation by $U$ before
applying $\kappa'$.  Together with $A^2=\kappa'(U^2)$, these
identities give a well-defined homomorphism on the two cosets
$P\sqcup PU$, extending $\kappa'$ and sending $U$ to $A$.
Its image contains $\mathrm{SL}_2(\mathbb F_3)$ and the matrix
$A$ of determinant $-1$, hence is all of $\mathrm{GL}_2(\mathbb F_3)$.
Its kernel has order $384/48=8$ and is $K$.

The signed braid center maps to $Z$, so the map descends to
$\Gamma$.  There $r=uv$ maps to $T$, since $ZR$ maps to $T$;
and $s=ru$ maps to
\[
 TA=\begin{pmatrix}2&0\\1&1\end{pmatrix},
\]
an involution.  Finally $A^2=S$, $S^2=-I$ and $A^4=-I$,
so $A$ has order eight.  This proves the stated orders.
\end{proof}

\begin{corollary}\label{cor:genus-two-sign-cover}
The discrete kernel of $\bar\ell$ is the fundamental group of
a genus-two surface with six punctures.  Its compactification
carries a faithful deck action of $\mathrm{GL}_2(\mathbb F_3)$
of signature $(2,3,8)$.  The central involution is its
hyperelliptic involution.

On $D$ the induced action on this finite deck group is the
identity.  The pure harmonic conjugator maps to $I$, and the
pure cubic conjugator maps to $(-I)^{\tau(g)}$.
\end{corollary}
\begin{proof}
The subgroup has index $48$ in the modular group $\Gamma$.
It is torsion-free because the orders two and three of the
finite factors are preserved.  The cusp stabilizer has image
of order eight, so there are $48/8=6$ cusps.  Since
$\chi(\Gamma)=-1/6$, its Euler characteristic is $-8$.
Thus $2-2g-6=-8$ and $g=2$.  Filling the cusps gives the
signature $(2,3,8)$ action, by the same stabilizer counts.

The central matrix $-I=A^4$ lies in every conjugate of the
cusp stabilizer, so it fixes all six filled cusp points.
For an involution of a genus-two compact surface the
Riemann--Hurwitz formula is $2=4h-4+r$.  Since $r\geq6$,
one has $h=0$ and $r=6$.  This is the hyperelliptic involution.

On $D$ the established normalized signed covariance gives
$q_8\Phi_g=q_8$, hence identity on the deck group after $\ell$.
The established prime-two normalization gives
$q_8(\beta_g)=I$.  Also
$q_8(A_g)^{-1}=5^{\tau(g)}I=C^{\tau(g)}$, and
$\ell(C)=-I$.  Inversion leaves this sign unchanged, yielding
the asserted cubic conjugator value.
\end{proof}

This gives a concrete genus-two setting for the cubic finite
comparison. The equality $\tau=\eta$ requires the cubic branch-filling
comparison proved in Section~\ref{sec:normalized-cubic-identity}.

\paragraph{Identification with the classical Bolza cover.}
The genus-two cover in Proposition~\ref{prop:genus-two-finite-quotient}
is the classical Bolza cover.  Katz--Katz--Schein--Vishne
\cite[\S3 and Corollaries 13.3--13.4]{KKSV2016} describe its
index-two triangle-group inclusion and the corresponding
$\mathrm{SL}_2(\mathbb F_3)$ and $\mathrm{GL}_2(\mathbb F_3)$
quotients.  Their displayed matrices immediately before
Corollary 13.4 send the order-two and order-three generators to
\[
 s\longmapsto\begin{pmatrix}0&1\\1&0\end{pmatrix},\qquad
 r\longmapsto T=\begin{pmatrix}1&1\\0&1\end{pmatrix}.
\]
These are exactly conjugate to our pair:
\[
 T(TA)T^{-1}=\begin{pmatrix}0&1\\1&0\end{pmatrix},\qquad
 TTT^{-1}=T.
\]
Thus the kernels in the triangle group agree.  The compactification is the Bolza
curve $y^2=x^5-x$; the same model, its six hyperelliptic branch
points $0,\infty,\pm1,\pm i$, and its classical automorphism
group are recalled in \cite[\S4, equation (4.1)]{KRRBolza}.
That paper attributes the classification to Bolza's 1887 work.

Here we compute the factorization through
$\mathrm{SL}_2(\mathbb Z/8)$ with kernel
$\langle Z,X^2C,Y^2C\rangle$, together with the values of the
normalized frame conjugators. Section~\ref{sec:normalized-cubic-identity}
then identifies their cubic sign with $\eta$.


\section{The harmonic comparison from Furusho's dyadic identities}
\label{sec:harmonic-dyadic-comparison}

Let $g=(1,f)$ be an element of full profinite $\mathrm{GT}$ with
$\rho_2^I(g)=0$.  We use the established equality between this
Ihara coordinate and the harmonic coordinate $\rho_2$, so the
pure harmonic conjugator $\beta\in\widehat F(a,b)$ has
abelianization zero.  Write $\theta(a)=b$, $\theta(b)=a$.
The fixed convention for this conjugator is
\begin{equation}\label{eq:harmonic-word-H}
f(a,b)=\theta(\beta)\beta^{-1}.
\end{equation}
Relation $(\mathrm{IV})$ holds on full GT by
\cite[Theorem~1]{H17}, and $(K_2)$ holds there by
\cite[Theorem~8.1]{Enriquez2007}. We apply these geometric
identities to the harmonic frame.

\subsection{The precise dyadic word identities}

The kernel of $\widehat F(x,y)\to\mathbb Z/2$ sending
$x\mapsto1$ and $y\mapsto0$ is free on $xyx^{-1},x^2,y$.  Since $f$ has zero
exponent sums, there is a unique word $f^+(W,X,Y)$ such that
\begin{equation}\label{eq:harmonic-word-1}
f(x,y)=f^+(xyx^{-1},x^2,y).
\end{equation}
Furusho's general definition includes the factor
$y^{-2\Psi_2^{(0)}(g)}$ on the left.  Here that factor is one,
because $\Psi_2^{(0)}=-\rho_2^I=0$.
The exact $(K_2)$ relation, in the form
\cite[\S3, equation~(2), p.~102]{Furusho2003}, gives
\begin{equation}\label{eq:harmonic-word-2}
f^+(1,a,b)=f(a,b).
\end{equation}
We use the following two identities in their full nonabelian
domains:
\begin{equation}\label{eq:harmonic-word-3}
f^+(W,X,Y)\,
 f^+(X^{-1}W^{-1}Y^{-1},Y,X)=1,
\end{equation}
and, whenever $WXYZ=1$,
\begin{equation}\label{eq:harmonic-word-4}
f^+(W,X,Y)f^+(Z,W,X)
 f^+(Y,Z,W)f^+(X,Y,Z)=1.
\end{equation}
These are respectively Furusho~\cite[Lemmas 4.1 and 4.2,
pp.~105--106]{Furusho2003}, specialized to
$\lambda=1$ and $\Psi_2^{(0)}=0$.
Their stated hypotheses are $(\mathrm{IV})$, and
$(\mathrm{I}), (\mathrm{II}), (\mathrm{IV})$, respectively.
They are not restricted to Galois elements.  Equation \eqref{eq:harmonic-word-3} is
an identity in the free profinite group on $W,X,Y$; equation
\eqref{eq:harmonic-word-4} is an identity in the free group on three of the four
variables with their indicated product relation.  We are not
using just their Fox derivatives or abelianizations.

\begin{lemma}\label{lem:dyadic-branch-word}
Put $h(a,b)=f^+(a,1,b)$.  Then
\[
 f(a,b)=h(a,b)h(b,a)^{-1}
 \quad\text{in }\widehat F(a,b).
\]
\end{lemma}
\begin{proof}
Work with $abc=1$.  Substitute $(W,X,Y,Z)=(1,b,c,a)$ in
\eqref{eq:harmonic-word-4}, obtaining
\begin{equation}\label{eq:harmonic-word-5}
f^+(1,b,c)f^+(a,1,b)f^+(c,a,1)f^+(b,c,a)=1.
\end{equation}
By \eqref{eq:harmonic-word-2}, the first factor is $f(b,c)$, and the second is
$h(a,b)$.  Applying \eqref{eq:harmonic-word-3} to $(W,X,Y)=(b,1,a)$ gives
\[
 f^+(b,1,a)f^+(b^{-1}a^{-1},a,1)=1.
\]
Since $b^{-1}a^{-1}=c$, the third factor of \eqref{eq:harmonic-word-5} is therefore
$h(b,a)^{-1}$.
Applying \eqref{eq:harmonic-word-3} to $(W,X,Y)=(b,c,a)$ gives
\[
 f^+(b,c,a)f^+(c^{-1}b^{-1}a^{-1},a,c)=1.
\]
Here $c^{-1}b^{-1}a^{-1}=1$, so \eqref{eq:harmonic-word-2} and the two-cycle
identity make the fourth factor of \eqref{eq:harmonic-word-5} equal to
$f(a,c)^{-1}=f(c,a)$.
Consequently
\[
 f(b,c)h(a,b)h(b,a)^{-1}f(c,a)=1,
\]
whence
\[
 h(a,b)h(b,a)^{-1}=f(b,c)^{-1}f(c,a)^{-1}.
\]
The normalized three-cycle identity is
$f(a,b)f(c,a)f(b,c)=1$ for $abc=1$.
It identifies the right side with $f(a,b)$, in precisely
the displayed order.  This proves the lemma.
\end{proof}

\begin{lemma}\label{lem:unique-swap-frame}
The map $\gamma\mapsto\theta(\gamma)\gamma^{-1}$ on
$\widehat F(a,b)$ is injective.  In particular
\[
 \beta(a,b)=h(b,a).
\]
\end{lemma}
\begin{proof}
If $\theta(\gamma_1)\gamma_1^{-1}=
\theta(\gamma_2)\gamma_2^{-1}$, write
$\gamma_2=\gamma_1 k$.  Cancellation gives $\theta(k)=k$.
The semidirect product
$\widehat F(a,b)\rtimes\langle s\rangle$, where $s^2=1$
and $sas=b$, is the free profinite product
$\widehat{\mathbb Z}\amalg C_2$, generated by $a,s$.
The centralizer of $s$ is the factor $\langle s\rangle$,
by the centralizer theorem for a nonidentity factor element
\cite[Theorem~B, p.~155]{HerfortRibes1985}.
Its intersection with $\widehat F(a,b)$ is trivial, so $k=1$.

By Lemma~\ref{lem:dyadic-branch-word}, the word
$\gamma(a,b)=h(b,a)$ satisfies
$\theta(\gamma)\gamma^{-1}=f(a,b)$.
Equation \eqref{eq:harmonic-word-H} and injectivity now give $\gamma=\beta$.
\end{proof}

\begin{proposition}\label{prop:normalized-GF1-from-dyadic}
For every $g$ as above, in the full signed profinite braid group
on generators $u,v$, with $s=uvu$, one has
\[
 \beta_g(u^2,v^2)^{-1}=f_g(u^2,s).
 \tag{GF1$_0$}
\]
\end{proposition}
\begin{proof}
Equation \eqref{eq:harmonic-word-2} shows that the $X$- and $Y$-exponent sums
of $f^+$ vanish, because $f$ is a commutator word.
Equation \eqref{eq:harmonic-word-1} then shows that its $W$-exponent sum also
vanishes. In particular its $X$-exponent is zero.

In the signed braid group, put $c=(uv)^3$.
The braid relation gives $su^2s^{-1}=v^2$ and $s^2=c$,
and $c$ is central. Specializing \eqref{eq:harmonic-word-1} at $(x,y)=(s,u^2)$
therefore gives
\[
 f(s,u^2)=f^+(v^2,c,u^2)=f^+(v^2,1,u^2)
         =h(v^2,u^2)=\beta(u^2,v^2).
\]
The second equality holds because a central argument can
be collected into its total exponent, which here is zero.
This assertion for profinite words follows by continuity
from the corresponding identity for finite words.
The two-cycle identity now gives
$f(u^2,s)=f(s,u^2)^{-1}=\beta(u^2,v^2)^{-1}$,
as required.
\end{proof}

\begin{corollary}\label{cor:harmonic-defect-zero}
On $D=\ker\lambda\cap\bigcap_q\ker\rho_q^I$, the finite
harmonic frame sign satisfies $w=\eta$.  Thus the previously
constructed relative character $\epsilon_h=w+\eta$ is zero
on the whole of full GT, and the remaining exceptional
character is $\chi_3=\epsilon_c=\tau+\eta$.
The cubic comparison in Section~\ref{sec:normalized-cubic-identity}
will also make this character zero.
\end{corollary}
\begin{proof}
Use the signed matrices over $\mathbb F_3$
\[
 U=\begin{pmatrix}1&1\\0&1\end{pmatrix},\qquad
 S=\begin{pmatrix}0&1\\-1&0\end{pmatrix}.
\]
The earlier finite comparison proves
$f_g(U,S)=(-I)^{\eta(g)}$ on $D$.
Let $J=\operatorname{diag}(-1,1)$.  Then
$JUJ^{-1}=U^{-1}=U^2$ and $JSJ^{-1}=-S$.
A commutator proword is unaffected by multiplying either
input by a central element: its exponent sum in that input
is zero.  Hence
\[
 f_g(U^2,S)=f_g(U^2,-S)
          =Jf_g(U,S)J^{-1}=(-I)^{\eta(g)}.
\]
Proposition~\ref{prop:normalized-GF1-from-dyadic} identifies
the left side with $q_3(\beta_g)^{-1}=(-I)^{w(g)}$.
This proves $w=\eta$ and $\epsilon_h|_D=0$.
Its established extension is zero on $G_{\mathbb Q}$, and
$\widehat{\mathrm{GT}}=D G_{\mathbb Q}$, so it is zero
everywhere.  The final equality follows from the already
proved decomposition $\chi_3=\epsilon_h+\epsilon_c$.
\end{proof}

\section{The cubic branch filling and the normalized cubic identity}
\label{sec:normalized-cubic-identity}

We use the standard full genus-zero GT action and the exact
nonabelian $(K_3)$ identity. The rigidity input is
\begin{equation}\label{eq:cubic-normalizer-1}
\operatorname{Out}(\widehat\Gamma_{0,n})
 =\widehat{\mathrm{GT}}\times\mathfrak S_n\qquad(n\ge5),
\end{equation}
as proved by Hoshi--Minamide--Mochizuki
\cite[Definition~2.7 and Corollary~2.8(i), author pp.~31--32]{CubicHMM}.
Here the set of primes is the full set; their Definition~2.7
identifies the resulting GT group with the classical full
profinite GT group. The other topological inputs and their
precise scope are stated where they are used.

Schneps introduced special homomorphisms from cyclic automorphisms
of punctured spheres. Her subgroup $\widehat{\mathrm{GS}}^1_{0,0}$
is defined by two such compatibilities, and its elements satisfy
the normalized cubic relation $(\mathrm{NT}_2)$
\cite[Definition~1.2, p.~658; Theorems~3.1 and~3.4,
pp.~665 and~667]{Schneps2006Special}.
The argument here proves the cubic comparison on full GT with
multiplier one and vanishing $\rho_3^I$.


Throughout, $\Gamma_{0,n}$ denotes the pure mapping class
group of an oriented sphere with $n$ marked points. Brackets
around a set of marked points mean that its permutations are
allowed. A circumflex denotes full profinite completion.
Let $g=(1,f_g)$ belong to full profinite GT. Its standard genus-zero
action \cite[equation~(1.3), p.~657]{Schneps2006Special}
is represented by an automorphism $\alpha_g$
of each completed mixed mapping class group. It induces the
identity on the finite permutation quotient and the identity
on the abelianization of each pure subgroup. We distinguish
this action from its outer class on a pure subgroup.
The labelled forgetting maps are compatible with these actions
up to inner automorphisms by pure elements. One may choose
compatible based representatives for a single forgetting map
and obtain equality; changing such a choice adds a pure inner
automorphism. This is the usual labelled genus-zero action.
Below we keep these
pure inner corrections explicitly. On the four-point group it
has the usual outer class of
\begin{equation}\label{eq:cubic-normalizer-2}
\phi_g(x)=x,\qquad
 \phi_g(y)=f_g(x,y)^{-1}y f_g(x,y).
\end{equation}
The multiplier-one hypothesis is used here.

\subsection{The centralizer and its full completion}

Fix a primitive cube root $\omega$. On $\mathbb P^1$ mark
$0,\infty$ and three free orbits
\[
 D_{i,j}=\omega^j t_i,
 \qquad i\in\{1,2,3\},\quad j\in\mathbb Z/3,
 \qquad t_i\ne0,\quad t_i^3\ne t_k^3\ (i\ne k).
\]
Thus there are eleven marked points. Let $h(z)=\omega z$,
$\mathcal H=\langle h\rangle$, and let $\mathcal G$ be the subgroup of the
mixed mapping class group whose permutations fix $0,\infty$
individually and commute with the displayed permutation of
$h$. Put $\mathcal N=C_{\mathcal G}(\mathcal H)$. The cyclic quotient has coordinate
$w=z^3$, with its two branch points and three further marked
points. We first prove the discrete centralizer sequence
\begin{equation}\label{eq:cubic-normalizer-3}
1\longrightarrow \mathcal H\longrightarrow \mathcal N\longrightarrow
 Q_3\longrightarrow1,
 \qquad Q_3=\Gamma_{0,\{0\},\{\infty\},[3]}.
\end{equation}
The argument works with two or three free orbits. Delete all
marked points from the covering sphere and its quotient.
The resulting cover $S\to S/\mathcal H$ is unbranched, regular and
three-sheeted. Its monodromy character takes the values $1$
and $-1$ on positive loops around the two distinguished
branch labels and zero on the other punctures.

Every mapping class centralizing $h$ has a
representative commuting with $h$. Give $S$ its displayed
complex structure, in which $h$ is conformal, and choose the
extremal quasiconformal representative $n$ of that class.
Existence and uniqueness for a finitely punctured compact
surface follow from Gardiner--Hu, Section~5, Theorem~1,
and Section~8, Theorem~6
\cite[author pp.~14--15 and~22]{CubicGH}; the punctured
hypothesis is specified immediately before their Theorem~1.
Both $nh$ and $hn$ lie in the same homotopy class, and both
are extremal: any improvement, composed with $h^{-1}$ on
the appropriate side, would improve $n$. Uniqueness implies
that they differ by a conformal map homotopic to the identity.
Such a map fixes all punctures, hence at least three points
of the compactified sphere, and is the identity. Thus $nh=hn$.
This spells out in the punctured setting the extremal-map
argument of Birman--Hilden \cite[p.~433]{CubicBH}.

The representative $n$ descends to the punctured quotient.
If two such representatives are isotopic upstairs, their
quotient is fiber-preserving and isotopic to the identity.
Birman--Hilden Theorem~1
\cite[p.~424]{CubicBH}, applied to this unbranched cover,
makes it fiber-isotopic to the identity. Their descended
mapping classes consequently agree. Conversely every element
of $Q_3$ preserves the monodromy character and therefore
lifts; every lift commutes with $h$ and has exactly the
required label permutation. The lifts differ by $\mathcal H$.
Lifting an isotopy downstairs shows that the kernel is $\mathcal H$,
which proves \eqref{eq:cubic-normalizer-3}. This argument proves the required
liftability directly. It does not use Birman--Hilden's later
general lifting assertion in Lemma~5.1 or Theorem~5;
their erratum leaves Theorems~1--4 unaffected
\cite{CubicBHErr}.

We prove that $\mathcal N$ has its full profinite topology in the
ambient mapping class group. Let
$\Phi=\pi_1(S/\mathcal H)$ and $\Pi=\pi_1(S)$, with compatible
basepoints. Thus $\Pi\lhd\Phi$ is the kernel of the above
character $\chi:\Phi\to C_3$, and
$\widehat\Pi$ is open in the free profinite group
$\widehat\Phi$. Put
\[
 \mathcal L=
 \operatorname{Aut}_{\chi}(\widehat\Phi)
                 /\operatorname{Inn}(\widehat\Pi),
\]
where the subscript means preservation of $\chi$ pointwise.
Since $\widehat\Phi$ is finitely generated and
$\operatorname{Inn}(\widehat\Pi)$ is compact, these are
profinite groups. There is an exact sequence
\begin{equation}\label{eq:cubic-normalizer-4a}
1\longrightarrow C_3\longrightarrow\mathcal L
 \longrightarrow\operatorname{Out}_{\chi}(\widehat\Phi)
 \longrightarrow1.
\end{equation}
Indeed, its kernel is
$\operatorname{Inn}(\widehat\Phi)/
\operatorname{Inn}(\widehat\Pi)=\widehat\Phi/\widehat\Pi$;
the centre of a nonabelian free profinite group is trivial.

The lifted mapping classes give a homomorphism
$\mathcal N\to\mathcal L$. Choose a path upstairs from the basepoint
to its image under a commuting representative and project
that path downstairs. The resulting based automorphism of
$\Phi$ is well-defined modulo $\operatorname{Inn}(\Pi)$:
changing the path gives just such an inner automorphism, and
changing the representative is harmless by the preceding
fiber-isotopy argument. This construction is compatible with
composition. Its profinite extension fits with \eqref{eq:cubic-normalizer-3}, \eqref{eq:cubic-normalizer-4a}
and the ordinary outer action of $\widehat Q_3$ on
$\widehat\Phi$.

The completion of \eqref{eq:cubic-normalizer-3} has kernel exactly $\mathcal H$. This finite
subgroup is separated by its permutation of the upstairs
punctures; the universal property of completion identifies
$\widehat{\mathcal N}/\mathcal H$ with $\widehat Q_3$. Asada's faithfulness
theorem makes
$\widehat Q_3\to\operatorname{Out}(\widehat\Phi)$
injective in the pure case
\cite[introduction Theorem, p.~3, parts (i)--(ii)]{CubicAsada}.
The mixed case follows as well: a nontrivial puncture
permutation is visible on peripheral conjugacy classes,
and its permutation kernel is the pure subgroup. The map
$\mathcal H\to C_3$ in \eqref{eq:cubic-normalizer-4a} is an isomorphism, since a deck
transformation is represented by a projected path of the
corresponding monodromy. The two exact sequences therefore
make $\widehat{\mathcal N}\to\mathcal L$ injective.

Restriction $\mathcal L\to
\operatorname{Out}(\widehat\Pi)$ is injective. If a
representative restricts to an inner automorphism, adjust by
that inner automorphism so that it fixes $\widehat\Pi$
pointwise. For any $x\in\widehat\Phi$, comparison of its
conjugation on $\widehat\Pi$ then shows that
$x^{-1}\theta(x)$ centralizes $\widehat\Pi$.
This centralizer is trivial: it intersects $\widehat\Pi$
in the latter's trivial centre, so injects into the finite
group $\widehat\Phi/\widehat\Pi$; a free profinite group
has no nontrivial finite subgroup. Thus $\theta$ is the
identity. We have proved faithfulness of the full completed
$\mathcal N$-action on $\widehat\Pi$. This action factors through
the completion of the ambient mapping class group, so
$\widehat{\mathcal N}\to\widehat{\mathcal G}$ is injective. In particular
\begin{equation}\label{eq:cubic-normalizer-4}
\overline{\mathcal N}=\widehat{\mathcal N},\qquad
 \overline{\mathcal N}/\mathcal H=\widehat Q_3.
\end{equation}
This establishes the full profinite topology on $\mathcal N$.

Finally the Boggi--Zalesskii centralizer theorem gives
$C_{\widehat{\mathcal G}}(\mathcal H)=\overline{\mathcal N}$
\cite[Theorem~A(iii), p.~3; Section~6.1 and
Theorem~6.2(iii), pp.~16--17]{CubicBZ}.
The theorem applies to the present finite-index subgroup:
genus-zero mapping class groups are good and have finite
virtual cohomological type by the cited Section~6.1;
goodness and that finiteness property pass to finite-index
subgroups (by the finite-index cohomology comparison and
Shapiro's lemma). Residual finiteness is inherited as well.
Only finite $3$-subgroups are used.

\subsection{Correction and identification of the parameter}

Let $P=\Gamma_{0,11}$ and let
$\mathcal G_{\mathrm{cyc}}=P \mathcal H$ be the inverse image of the cyclic permutation
subgroup generated by $h$. Since $\alpha_g$ is the identity
on label permutations, $\alpha_g(\mathcal H)\subset\widehat{\mathcal G}_{\mathrm{cyc}}$.
The Boggi--Zalesskii finite-$3$-subgroup theorem
\cite[Theorem~A(i), p.~3; Theorem~6.2(i), p.~17]{CubicBZ}
applied to $\mathcal G_{\mathrm{cyc}}$ conjugates $\alpha_g(\mathcal H)$ into $\mathcal G_{\mathrm{cyc}}$.
The conjugator can be taken pure: write an initial conjugator
as $h^j p$ with $p\in\widehat P$, and absorb $h^j$ by
conjugating the resulting discrete subgroup. The resulting
generator has the same permutation of all eleven labels as
$h$, since the permutation quotient of $\mathcal G_{\mathrm{cyc}}$ is abelian.

Nielsen realization realizes this finite cyclic group by
conformal transformations of a punctured sphere. The required
finite-type, punctured scope is included in Kerckhoff's
extension to finitely generated Fuchsian groups with cusps
\cite[Section~4, pp.~255--256]{CubicKerckhoff}.
Putting its fixed points at $0,\infty$, the generator is
$z\mapsto\omega z$ or $z\mapsto\omega^{-1}z$.
In the first case its labelled configurations form the
connected space
\[
 \{(t_1,t_2,t_3)\in(\mathbb C^\times)^3:
                         t_i^3\ne t_k^3\text{ for }i\ne k\}.
\]
In the second case exchange the two fixed-point labels first.
It follows that there is an $e$ which is pure, or pure
followed by a transposition of $0,\infty$, and which fixes
every free-orbit label, such that
\begin{equation}\label{eq:cubic-normalizer-5}
\alpha'=\operatorname{Ad}(e^{-1})\alpha_g
       \quad\text{fixes }\mathcal H\text{ pointwise}.
\end{equation}
Here the assertion about the labels concerns the permutation
of $e$; its pure factor need not fix points pointwise as a
homeomorphism.

We verify invariance of the mixed groups before taking
outer classes. The permutation of $e$ is either the identity
or the single transposition $(0\ \infty)$; it fixes all nine
free labels. Either permutation centralizes the permutation
subgroup defining $\mathcal G$. Since $\alpha_g$ induces the identity
on permutations, $\alpha'$ preserves $\widehat{\mathcal G}$ and
induces the identity on its permutation quotient. Since it
fixes $\mathcal H$ pointwise, it preserves
$C_{\widehat{\mathcal G}}(\mathcal H)=\widehat{\mathcal N}$. Equation \eqref{eq:cubic-normalizer-4} therefore
gives an automorphism $\delta_3$ of $\widehat Q_3$.
It preserves its pure subgroup
$P_5=\widehat\Gamma_{0,5}$ and induces the identity on
$\widehat Q_3/P_5=\mathfrak S_3$.

Write $\vartheta=[\delta_3|_{P_5}]=(g',\sigma)$ by \eqref{eq:cubic-normalizer-1}.
For each $\pi\in\mathfrak S_3$, choose $z_\pi\in
\widehat Q_3$ over $\pi$. Conjugation by $z_\pi$ on
$P_5$ has outer class $(1,\pi)$ in \eqref{eq:cubic-normalizer-1}. As
$\delta_3(z_\pi)$ has the same permutation, the identity
\[
 \delta_3\operatorname{Ad}(z_\pi)\delta_3^{-1}
       =\operatorname{Ad}(\delta_3(z_\pi))
\]
restricted to $P_5$ gives
$\vartheta(1,\pi)\vartheta^{-1}=(1,\pi)$.
Hence $\sigma\pi\sigma^{-1}=\pi$ for every $\pi$.
The centralizer in $\mathfrak S_5$ of the full symmetric
group on the three orbit labels consists exactly of the
identity and the transposition of the other two labels.
Thus
\begin{equation}\label{eq:cubic-normalizer-6}
\sigma=1\quad\text{or}\quad\sigma=(0\ \infty).
\end{equation}
To forget the third orbit, first restrict $\widehat{\mathcal N}$
to the subgroup $\widehat{\mathcal N}^{\rm orb}$ fixing each free
orbit as a set; its quotient by $\mathcal H$ is $P_5$. This subgroup
is preserved by $\alpha'$, as just proved. Ambient forgetting
of all three labels in the third orbit maps it onto the
two-orbit centralizer $\widehat{\mathcal N}_2^{\rm orb}$.
On quotients by $\mathcal H$ this is the usual surjection
$P_5\to P_4=\widehat\Gamma_{0,4}$, and it is the
identity on $\mathcal H$. These facts prove surjectivity upstairs too.

Here are the inner corrections in this step. Choose standard
GT and permutation representatives on $P_5$. By definition
of an outer class, $\delta_3|_{P_5}$ is their product
followed by $\operatorname{Ad}(a_3)$ with $a_3\in P_5$.
Both possible $\sigma$ fix the deleted orbit label.
The standard representatives therefore preserve the
forgetting kernel, and so does $\delta_3$; the correction
becomes the image of $a_3$ in $P_4$. Thus the
induced automorphism $\delta_2$ has outer class
$(g',\sigma)$ on $P_4$.
On the ambient mixed group the same conclusion follows from
labelled forgetting compatibility. If a choice of based
representatives introduces a pure correction $c$, the
forgotten action is
$\operatorname{Ad}(\bar e^{-1}c)\alpha_g$.
Its effective conjugator $c^{-1}\bar e$ still fixes every
remaining free label and has exactly the same possible
branch transposition as $e$. Thus the label and homology
arguments below are unchanged. This notation also covers
all later forgetting maps: every extra based correction is
pure in its own target.

We make the resulting two-orbit calculation explicit. Scale
the first orbit to $\{1,\omega,\omega^2\}$, write the
second as $\{t,\omega t,\omega^2t\}$, and put $w=t^3$.
Then
\[
 \mathcal N_2/\mathcal H\cong\pi_1(\mathbb P^1_w\setminus\{0,1,\infty\})
       =F(x,y).
\]
The subgroup $\mathcal N_2^0$ fixing the chosen representatives $1,t$
fixes all six free-orbit labels and intersects $\mathcal H$ trivially.
Its image is
\begin{equation}\label{eq:cubic-normalizer-7}
H_3=\ker(F(x,y)\longrightarrow C_3),\qquad x\mapsto1,
 \quad y\mapsto0.
\end{equation}
It has basis $X=x^3$, $y_j=x^jyx^{-j}$, $j=0,1,2$.
The element $y_j$ is the positive simultaneous collision
multitwist pairing the three points of one free orbit with
the three points of the other, with relative offset $j$.
It is pure. Its abelian class $v_j$ in the ambient pure
eight-point group is fixed by $\alpha'$: $\alpha_g$ acts
as the identity on this abelianization and the only possible
permutation of $e$ exchanges $0,\infty$, fixing all the
paired labels.

The three $v_j$ are distinguished integrally by forgetting
maps. Retain $0,\infty$ and just one pair with offset $j$.
The class $v_j$ maps to the positive generator around their
collision in the four-point group, whereas the other two
classes map to zero. In particular it is nonzero and
primitive in this test quotient.

Put $\lambda'=\lambda(g')$. In \eqref{eq:cubic-normalizer-7}, an automorphism of
outer class $(g',\sigma)$ takes the conjugacy class of $y_j$
to that of
\begin{equation}\label{eq:cubic-normalizer-8}
y_{\varepsilon\lambda'j+c}^{\lambda'},
 \qquad \varepsilon=1\text{ if }\sigma=1,
 \quad\varepsilon=-1\text{ otherwise},
\end{equation}
with subscripts modulo three and a common offset $c$.
For the GT factor this follows directly from
$x\mapsto x^{\lambda'}$, $y\mapsto f_{g'}^{-1}y^{\lambda'}f_{g'}$;
the interchange of $0,\infty$ lifts to $t\mapsto1/t$ and
reflects the three collision points, and an inner correction
adds a common offset. Once the subscript in \eqref{eq:cubic-normalizer-8} has been
chosen, its conjugating element is in $\widehat H_3$, hence
is pure upstairs.
The preceding three abelian tests therefore give
\begin{equation}\label{eq:cubic-normalizer-9}
v_j=\lambda'v_{\varepsilon\lambda'j+c}
 \quad(j=0,1,2),\qquad
 \lambda'=1,\quad c=0,\quad\varepsilon=1.
\end{equation}
Thus $\sigma=1$. Moreover the two-orbit quotient
automorphism has the form
\begin{equation}\label{eq:cubic-normalizer-10}
\delta_2=\operatorname{Ad}(a)\phi_{g'},
                  \qquad a\in\widehat H_3,
\end{equation}
because its common offset is now the image of $a$ in $C_3$.

Forget all but $0,\infty,1,t$ upstairs. The resulting map
on \eqref{eq:cubic-normalizer-7} is exactly the standard cyclic-cover filling
\begin{equation}\label{eq:cubic-normalizer-11}
d_3:\widehat H_3\longrightarrow\widehat F(x,y),
 \qquad X\mapsto x,\quad y_0\mapsto y,
 \quad y_1,y_2\mapsto1.
\end{equation}
The exact $(K_3)$ identity says
$d_3(f_{g'})=y^{c_3(g')}f_{g'}$. In the automorphism formula
this extra power commutes with $y$, so
\begin{equation}\label{eq:cubic-normalizer-12}
d_3\phi_{g'}|_{\widehat H_3}=\phi_{g'}d_3.
\end{equation}
For clarity, \eqref{eq:cubic-normalizer-12} is checked on the four free generators:
on $X$ both sides give $x$; on $y_0$ both give
$f_{g'}^{-1}yf_{g'}$; and on $y_j$ for $j=1,2$ use
\[
 \phi_{g'}(y_j)=(x^j f_{g'}^{-1}x^{-j})\,y_j\,
                         (x^j f_{g'}x^{-j}),
\]
whose image under $d_3$ is one. The conjugating factors
belong to $H_3$ because $f_{g'}$ is a commutator.
Equation \eqref{eq:cubic-normalizer-10} now contributes the inner automorphism
$\operatorname{Ad}(d_3(a))$ in the target. Thus
\eqref{eq:cubic-normalizer-10}--\eqref{eq:cubic-normalizer-12} make the induced outer action downstairs
equal to that of $g'$. On the other hand, the ambient
forgetting compatibility and \eqref{eq:cubic-normalizer-5} make it the action of $g$,
possibly composed with the transposition of the chosen
labels $0,\infty$. That permutation is a single transposition on the
four chosen labels: the two free labels are individually
fixed. It is not a double transposition in the Klein four
kernel of the four-point permutation action. Equivalently,
in the coordinate $t$ it is $t\mapsto1/t$, which exchanges
the peripheral classes at zero and infinity. Both GT
factors and all pure inner corrections preserve the three
peripheral classes individually. The branch transposition
is therefore absent. Thus $e$ is pure and $g'=g$.

For completeness, the usual GT embedding into the outer
automorphism group of $\widehat F(x,y)$ is injective in this
comparison. An inner conjugation equating two standard
representatives with the same multiplier must centralize
$x$, so is a power of $x$. Comparing the conjugators of
$y$ and then taking their exponent sums makes that power
and the possible power of $y$ zero; the two commutator
parameters are equal.

\begin{lemma}[The other cubic filling]\label{lem:cubic-branch-filling}
Let $g=(1,f_g)\in\widehat{\mathrm{GT}}$.
Fill the punctures $0,\infty$ in the threefold cyclic cover
$\mathbb P^1_t\setminus\{0,\infty,1,\omega,\omega^2\}$.
The restriction of the standard action $\phi_g$ on the
base group, followed by this filling, induces on
$\pi_1(\mathbb P^1_t\setminus\{1,\omega,\omega^2\})$
the outer class of the same $g$.
\end{lemma}
\begin{proof}
Retain the three points $1,\omega,\omega^2$ of the first
orbit and the chosen point $t$ of the second, forgetting
$0,\infty$ and the other two points of the second orbit.
On $\mathcal N_2^0=H_3$ this is precisely the displayed branch
filling. The parameter comparison above gives
$\delta_2=\operatorname{Ad}(a)\phi_g$ with
$a\in\widehat H_3$, and the correction $e$ is pure.
Both corrections therefore become inner automorphisms in
the target four-point group. Ambient forgetting supplies
the standard outer action of $g$ there, proving the claim.
\end{proof}

\subsection{The exact cubic conjugator}

Let $\Gamma=\widehat{C_2*C_3}$ be the modular quotient of
the signed braid group on $u,v$, and put $r=uv$,
$a=u^2$, $b=v^2$. The subgroup $\mathcal E=\langle r,a\rangle$
has presentation
\[
 \mathcal E=\langle r,a\mid r^3=(ra)^3=1\rangle
       \cong C_3\amalg C_3.
\]
Its index-three subgroup $F=\langle a,b\rangle$ is free;
conjugation by $r$ cyclically permutes $a,b,(ab)^{-1}$.
It is the filled cover in Lemma~\ref{lem:cubic-branch-filling},
with $x\mapsto r$, $y\mapsto a$. These choices fix the
orientation and the order of the three peripheral loops.

\begin{proposition}\label{prop:normalized-HF1-from-normalizer}
If $g=(1,f_g)\in\widehat{\mathrm{GT}}$ and
$\rho_3^I(g)=0$, then its pure cubic conjugator satisfies
\begin{equation}\label{eq:cubic-normalizer-13}
\boxed{\quad A_g(u^2,v^2)^{-1}=f_g(u^2,uv)\quad}
\end{equation}
in the full signed profinite braid group.
\end{proposition}
\begin{proof}
The standard action $\phi_g$ descends to $\mathcal E$: it preserves
the two finite peripheral relations because it sends each
peripheral generator to a conjugate of itself. Denote this
action by $\psi$. Put $H_g=f_g(a,r)$. The two-cycle relation
gives
\begin{equation}\label{eq:cubic-normalizer-14}
\psi(r)=r,\qquad\psi(a)=H_g aH_g^{-1}.
\end{equation}
The braid action $\Phi_g$ preserves $\mathcal E$ and has
$\Phi_g(a)=a$, $\Phi_g(r)=A_g rA_g^{-1}$.
Lemma~\ref{lem:cubic-branch-filling} says that on the open
free subgroup $F$ one has
\begin{equation}\label{eq:cubic-normalizer-15}
\psi|_F=\operatorname{Ad}(k)\Phi_g|_F
                  \quad\text{for some }k\in F.
\end{equation}
The inner correction is in $F$, since the target four-point
group is pure and all corrections in the preceding lemma
were pure. Two automorphisms of $\mathcal E$ agreeing on $F$ agree
everywhere: their quotient fixes $F$ pointwise, and comparison
of conjugation on $F$ shows that its discrepancy at every
element centralizes $F$, whose centralizer is trivial.
Thus \eqref{eq:cubic-normalizer-15} holds on $\mathcal E$.

Applying it to $r$ gives $kA_g\in C_\mathcal E(r)=\langle r\rangle$.
Since $k,A_g\in F$ and $F\cap\langle r\rangle=1$, this
forces $k=A_g^{-1}$. Applying \eqref{eq:cubic-normalizer-15} to $a$ now gives
$A_gH_g\in C_F(a)=\overline{\langle a\rangle}$.
Write $A_gH_g=a^t$ with $t\in\widehat{\mathbb Z}$.

Both factors on the left have abelianization zero.
For $A_g$, this follows from
$f_g=\operatorname{Ad}(r)(A_g)A_g^{-1}$: on $F^{\rm ab}$
the order-three rotation minus the identity has determinant
three, and multiplication by three on $\widehat{\mathbb Z}$
is injective, so that rotation has no fixed vector. For $H_g$, use
the established integral cyclic-cover coefficient formula:
in the basis $X,y_0,y_1,y_2$ of $H_3$, the class of
$f_g(x,y)$ is
\[
 (0,2b_3,-b_3,-b_3),\qquad b_3=-\rho_3^I(g)/2.
\]
It is zero here, and the image under the branch filling is
$f_g(r,a)=H_g^{-1}$. Taking abelianizations of
$A_gH_g=a^t$ therefore gives $t=0$. This proves \eqref{eq:cubic-normalizer-13} in
the modular quotient.

Finally both sides of \eqref{eq:cubic-normalizer-13} have exponent sum zero in the
signed braid group: $A_g\in F'$ and $f_g$ is a commutator
word. The kernel of the modular quotient is generated
topologically by $(uv)^3$, whose exponent sum is six.
Multiplication by six on $\widehat{\mathbb Z}$ is injective,
so the possible central discrepancy is zero. This proves
the signed identity exactly.
\end{proof}

In the modular quotient the identity also gives the explicit
normalized cubic relation
\[
 f_g(a,b)=f_g(r,b)f_g(r,a)^{-1}.
\]
Indeed $f_g(a,b)=\operatorname{Ad}(r)(A_g)A_g^{-1}$ and
$\operatorname{Ad}(r)(a)=b$. Substituting
$A_g^{-1}=f_g(a,r)$ and using the two-cycle identity proves
the displayed formula. Its inverse is equation~(3.10) of
\cite[p.~667]{Schneps2006Special}; that source derives it on
$\widehat{\mathrm{GS}}^1_{0,0}$ by forgetting two strands.


\begin{corollary}\label{cor:cubic-defect-zero-normalizer}
On $D=\ker\lambda\cap\bigcap_q\ker\rho_q^I$ the finite
cubic comparison defect $\epsilon_c=\tau+\eta$ vanishes.
Together with the harmonic comparison, this gives
$\chi_3=0$ on full GT.
\end{corollary}
\begin{proof}
Reduce \eqref{eq:cubic-normalizer-13} modulo eight under the signed representation.
Writing $X=U^2$ and $R=-UL$, its right side is $f_g(X,R)$:
the central sign in $q(uv)=-R$ has no effect on a commutator
word. The previously proved finite interchange map sends
$q_8(A_g)^{-1}$ to $(-I)^\tau$ and $f_g(X,R)$ to
$f_g(U,S)^{-1}=(-I)^\eta$. Hence $\tau=\eta$.
The relative extension of $\epsilon_c$ is zero on the Galois
subgroup, and $\widehat{\mathrm{GT}}=D G_{\mathbb Q}$,
so it is zero everywhere. The already proved harmonic
equality $w=\eta$ now gives $\chi_3=w+\tau=0$.
\end{proof}

\section{The tame quartic gate as a finite word evaluation}

Put $G=\widehat{\mathrm{GT}}$ and
\[
 D=\ker\lambda\cap\bigcap_{q\ \mathrm{prime}}\ker\rho_q^I.
\]
Write $g=(1,f_g)$ for an element of $D$, and use the signed
matrices
\[
 U=\begin{pmatrix}1&1\\0&1\end{pmatrix},\qquad
 L=\begin{pmatrix}1&0\\-1&1\end{pmatrix},\qquad
 S=ULU=\begin{pmatrix}0&1\\-1&0\end{pmatrix}.
\]
For an odd prime $\ell$, a subscript $\ell$ on these matrices
will mean reduction in $\mathrm{SL}_2(\mathbb F_\ell)$.
Let
\[
 \chi_{-4}(\ell)=\left(\frac{-1}{\ell}\right),\qquad
 m_\ell=\frac{\ell-\chi_{-4}(\ell)}2
 =\begin{cases}(\ell-1)/2,&\ell\equiv1\pmod4,\\
               (\ell+1)/2,&\ell\equiv3\pmod4.
   \end{cases}
\]

\begin{proposition}\label{prop:tame-harmonic-word-gate}
For every $g\in D$ the corrected harmonic frame satisfies
\begin{equation}\label{eq:tame-harmonic-1}
W_g=f_g(U^2,S)\qquad\text{in }\mathrm{SL}_2(\widehat{\mathbb Z}).
\end{equation}
Consequently its residue belongs to the cyclic group
\[
 \mathcal T_\ell
 =\{aI+bS_\ell:a,b\in\mathbb F_\ell,\ a^2+b^2=1\},
 \qquad |\mathcal T_\ell|=\ell-\chi_{-4}(\ell).
\]
The simultaneous compatibility of the radical and CM actions
on every rational quartic root is equivalent to the following
family of finite word identities:
\begin{equation}\label{eq:tame-harmonic-2}
\boxed{\quad
 f_g(U_\ell^2,S_\ell)^{m_\ell}=I
 \qquad(g\in D,\ \ell\ge5\text{ prime}).\quad}
\end{equation}
For a fixed $g$, the same equivalence holds pointwise. This is
a reformulation of the remaining tame quartic gate; it does
not assert the identities in \eqref{eq:tame-harmonic-2}.
\end{proposition}
\begin{proof}
The harmonic frame is $W_g=q(\beta_g)^{-1}D_g$.
On $D$, the established signed normalization gives $D_g=I$.
Proposition~\ref{prop:normalized-GF1-from-dyadic} gives
\[
 \beta_g(u^2,v^2)^{-1}=f_g(u^2,uvu)
\]
in the signed profinite braid group. Applying its continuous
signed representation $q(u)=U$, $q(v)=L$ proves \eqref{eq:tame-harmonic-1}, and
evaluation of a profinite word commutes with reduction modulo
$\ell$.

The frame belongs to the integral norm-one torus of
$\mathbb Z[i]$, represented by $i\mapsto S$.
Reduction gives the displayed group $\mathcal T_\ell$.
When $\ell\equiv1\pmod4$, its two eigenspaces identify it
with $\mathbb F_\ell^\times$; when $\ell\equiv3\pmod4$,
it is the kernel of the norm
$\mathbb F_{\ell^2}^\times\to\mathbb F_\ell^\times$.
Thus it is cyclic of the stated even order in both cases.
Its square subgroup is exactly the kernel of the power
$z\mapsto z^{m_\ell}$. The kernel of reduction from the
integral local torus is pro-$\ell$, so squaring is an
automorphism on that kernel. Hence a local frame component
at odd $\ell$ is a square if and only if its residue passes
the displayed power test.

Lemma~\ref{lem:lr-quartic-cm-quotient} identifies the
maximal elementary-two subextension of
$\mathbb Q(i)^{\mathrm{ab}}/\mathbb Q^{\mathrm{ab}}$
with
\[
 \mathbb Q^{\mathrm{ab}}(q^{1/4}:q\text{ rational prime}).
\]
Accordingly the action on all these quartic roots is trivial
if and only if $W_g$ is a square in the normalized adelic norm-one
group $N_4(4)$. The radical action is already trivial for $g\in D$.
At two the proved prime-two compatibility places $W_{g,2}$
in the local square subgroup. At three the established
harmonic residue is $W_g\bmod3\in\{I,-I\}$, which is
the square subgroup of the residue torus of order four;
the pro-three reduction kernel contributes no obstruction.
Thus precisely the primes $\ell\ge5$ remain.

Finally, squaring in the product defining $N_4(4)$ has image
equal to the product of their square subgroups: local square
roots may be chosen independently, and their tuple is an
adelic element of the required product. This retains the
simultaneous condition, rather than only checking a separate
list of necessary conditions. Substitution of \eqref{eq:tame-harmonic-1} now proves
both the pointwise and the universal equivalences in \eqref{eq:tame-harmonic-2}.
\end{proof}

\section{The two first degeneracy correspondences}
The first degeneracy correspondence is already governed by relation (IV).
Put $B=\operatorname{PSL}_2(\mathbb Z)$ and write
\[
 u=\begin{pmatrix}1&1\\0&1\end{pmatrix},\qquad
 v=\begin{pmatrix}1&0\\-1&1\end{pmatrix}
\]
for their classes in $B$. The first map from $\Gamma_0(N)$ to $B$ is inclusion.
The second is conjugation by $D_N=\operatorname{diag}(N,1)$; denote it by $d_N$.
Thus $d_N(u)=u^N$ and $d_N(v^N)=v$. These maps extend to full profinite completions,
since each of their images has finite index in $B$.

\begin{proposition}\label{prop:degeneracy-two}
Let $F=(1,f)$ be a normalized element of $\widehat{\operatorname{GT}}$.
Let $\Phi_F$ be its standard action on $\widehat B$, so that
\[
 \Phi_F(u)=u,\qquad
 \Phi_F(v)=q^{-1}vq,\qquad q=f(u^2,v^2).
\]
Then $\Phi_F$ preserves the closure of $\Gamma_0(2)$, and its restriction
commutes with both degeneracy maps to $\widehat B$.
\end{proposition}
\begin{proof}
The standard index-three subgroup $\Gamma_0(2)$ is generated by $u,v^2$;
its presentation is $\langle u,w\mid(uw)^2=1\rangle$ with $w=v^2$.
This follows, for example, by applying the Schreier method to the action of
$B$ on $\mathbb P^1(\mathbb F_2)$: the stabilizer of infinity is $\Gamma_0(2)$,
and the two resulting generators can be chosen as $u,v^2$.
Because $q$ belongs to $\overline{\langle u^2,v^2\rangle}$, the formulas for
$\Phi_F$ take both generators into $\overline{\Gamma_0(2)}$.
The same argument for $F^{-1}$ gives equality of this subgroup with its image.
Commutation with inclusion is now its definition. For the other map,
\[
 d_2(q)=f(u^4,v)=f(u^2,v^2)=q.
\]
The middle equality follows by applying the interchange of $u,v$ to normalized
relation (IV), then applying $f(b,a)=f(a,b)^{-1}$ to both sides.
On the generator $u$ the two compositions both have value $u^2$.
On $v^2$ they both have value $q^{-1}vq$. Equality on these topological
generators proves the assertion.
\end{proof}

\section{Normalized congruence triviality and the level-three gate}
\label{sec:normalized-congruence}

The level-two identity has a consequence at every congruence level.
Use the signed representation
$\widetilde q:\widehat B_3\to\operatorname{SL}_2(\widehat{\mathbb Z})$
with $u\mapsto U=\left(\begin{smallmatrix}1&1\\0&1\end{smallmatrix}\right)$
and $v\mapsto V=\left(\begin{smallmatrix}1&0\\-1&1\end{smallmatrix}\right)$.

\begin{proposition}\label{prop:normalized-congruence-triviality}
Let $g=(1,f)\in\widehat{\mathrm{GT}}$ satisfy $\rho_2(g)=0$.
Then
\[
 f(U^2,V^2)=I
       \qquad\text{in }\operatorname{SL}_2(\widehat{\mathbb Z}).
\]
Consequently $\widetilde q\Phi_g=\widetilde q$ on $\widehat B_3$.
\end{proposition}
\begin{proof}
Normalized relation (IV), together with interchange of the two braid
generators and the two-cycle relation, gives
\[
 f(u^4,v)=f(u^2,v^2)
       \qquad\text{in }\widehat B_3.
\]
Fix a prime $\ell$ and put $Q=f(U^2,V^2)$ in
$\operatorname{SL}_2(\mathbb Z_\ell)$.
Conjugation by $D=\operatorname{diag}(2,1)$ in
$\operatorname{GL}_2(\mathbb Q_\ell)$ sends $U^2$ to $U^4$ and
$V^2$ to $V$. Therefore the displayed signed identity gives
$DQD^{-1}=Q$, with no scalar ambiguity. Since the two diagonal
entries of $D$ are distinct, $Q=\operatorname{diag}(a,a^{-1})$
for a unit $a\in\mathbb Z_\ell^\times$.

The images of the braid generators under
$\widetilde q_\ell\Phi_g$ are $U$ and
\[
 M=Q^{-1}VQ=\begin{pmatrix}1&0\\-a^2&1\end{pmatrix}.
\]
Their braid relation $UMU=MUM$, in its lower-left entry, is
$\gamma(1+\gamma)=0$ with $\gamma=-a^2$.
The coefficient ring is a domain and $a$ is a unit, so $a^2=1$.
Thus $Q=\pm I$ and $M=V$.

It remains to determine the sign using the three-cycle relation.
Put $A=U^2$, $B=V^2$, $Z=(AB)^{-1}$, and $R=UV$.
Direct multiplication gives
\[
 RAR^{-1}=B,\qquad RBR^{-1}=-Z,\qquad RZR^{-1}=-A.
\]
A commutator word is unchanged when one of its arguments is multiplied
by a central scalar. Since $Q$ is central, conjugating its defining
expression by $R$ and $R^2$ consequently gives
\[
 f(A,B)=f(B,Z)=f(Z,A)=Q.
\]
The multiplier-one three-cycle relation, with $ABZ=1$, now implies
$Q^3=I$. Together with $Q^2=I$, this forces $Q=I$.
This holds for every $\ell$, proving the signed adelic assertion.
Both braid generators now have unchanged images, which proves the
intertwining identity.
\end{proof}

In the following statements the ambient modular group is the
centerless group $B=\operatorname{PSL}_2(\mathbb Z)$ and its full
profinite completion, rather than the central braid extension. Put
\[
 \mathcal C=\ker\bigl(\widehat B\longrightarrow
      \operatorname{SL}_2(\widehat{\mathbb Z})/\{\pm I\}\bigr),
\]
where the scalar sign is the common sign in all components.

\begin{lemma}\label{lem:parabolic-centralizer}
The centralizer of $v$ in $\widehat B$ is
$\overline{\langle v\rangle}\simeq\widehat{\mathbb Z}$, and
it has trivial intersection with $\mathcal C$.
\end{lemma}
\begin{proof}
The free pure subgroup is
$\widehat F_2=\overline{\langle x\rangle}\amalg
\overline{\langle y\rangle}$, where $x=u^2,y=v^2$.
The centralizer theorem for a nonidentity element of a factor in a
free profinite product gives
$C_{\widehat F_2}(y)=\overline{\langle y\rangle}$
\cite[Theorem B, p.~155; Theorem B$'$, \S4, p.~159]{HerfortRibes1985}.
This theorem has no torsion assumption on the element of the factor.
If $c\in C_{\widehat B}(v)$, its image in $S_3$ centralizes the
transposition represented by $v$. Thus $cv^{-e}$ is pure for
some $e\in\{0,1\}$ and centralizes $y$. It lies in
$\overline{\langle y\rangle}$, proving the centralizer assertion.
Finally $v^t$ has congruence image represented by
$\left(\begin{smallmatrix}1&0\\-t&1\end{smallmatrix}\right)$.
It is a scalar sign only if $t=0$. This proves injectivity of the
procyclic subgroup in the congruence quotient and the last assertion.
\end{proof}

\begin{corollary}\label{cor:level-three-exact-gate}
For a normalized $g=(1,f)\in\widehat{\mathrm{GT}}$, put
$q=f(u^2,v^2)\in\widehat B$. The element $q$ lies in
$\mathcal C$, hence in $\overline{\Gamma_0(3)}$.
The automorphism $\Phi_g$ preserves that subgroup. For
$d_3:\overline{\Gamma_0(3)}\to\widehat B$ induced by
conjugation by $\operatorname{diag}(3,1)$, set
\[
 h_3=d_3(q)q^{-1}.
\]
Then $h_3\in\mathcal C$, and strict compatibility with both
level-three degeneracy maps is equivalent to the exact identity
\[
 d_3(q)=q.
\]
\end{corollary}
\begin{proof}
The preceding proposition shows that $q$ and the action of $g$
are trivial at every congruence level. Therefore $q$ belongs to
$\mathcal C$ and $\Phi_g$ preserves the level-three subgroup.
On its congruence image, $d_3$ is conjugation by the same diagonal
matrix: the equality holds on the discrete subgroup and extends
by continuity. Consequently $d_3(q)$ also has trivial congruence
image, and $h_3\in\mathcal C$.

The subgroup $\Gamma_0(3)$ is generated by $u,v^3$, and
$d_3(u)=u^3$, $d_3(v^3)=v$. The two proposed compositions agree
on $u$. On $v^3$ they agree exactly when
$h_3\in C_{\widehat B}(v)$. The preceding lemma and
$h_3\in\mathcal C$ make this equivalent to $h_3=1$.
\end{proof}

The corollary identifies the remaining full profinite identity. It does
not prove it: every congruence representation already kills $h_3$, so
congruence matrices alone cannot decide whether this element is trivial
in $\widehat B$.


\section{One local Brauer class for the two frame signs}
The two frame signs admit one local Brauer invariant which can be defined
from the traces of the two matrices. No choices of centralizer generators
or common matrix basis enter its definition.

\begin{proposition}\label{prop:local-brauer-frame}
Let
\[
 \mathcal D=\{(T,W)\in\operatorname{SL}_2(\mathbb Z_3)^2:
       \operatorname{tr}(T)\operatorname{tr}(W^2)\in\mathbb Z_3^\times\}.
\]
On this domain define
\[
 \mathfrak b(T,W)=
 \big[(\operatorname{tr}(T)\operatorname{tr}(W^2),3)_{\mathbb Q_3}\big]
       \in\operatorname{Br}(\mathbb Q_3)[2].
\]
Here brackets denote a quaternion algebra's Brauer class.
This invariant is unchanged by independent conjugation of $T$ and $W$.
Suppose their reductions have the form
\[
 \bar T\in(-I)^{\epsilon_3}\langle R_0\rangle,
 \qquad \bar W\in S_0^{\epsilon_2}\langle-I\rangle,
\]
where $R_0$ is the standard nonidentity unipotent matrix of order three
in $\operatorname{SL}_2(\mathbb F_3)$ and $S_0^2=-I$.
Then $(T,W)\in\mathcal D$ and
\[
 \mathfrak b(T,W)=(\epsilon_2+\epsilon_3)
                 [(-1,-3)_{\mathbb Q_3}].
\]
Its local Hilbert sign is $(-1)^{\epsilon_2+\epsilon_3}$.
\end{proposition}
\begin{proof}
Every power of $R_0$ has trace $2$ in $\mathbb F_3$, including its
zeroth power. Also $\bar W^2=(-I)^{\epsilon_2}$. Consequently
\[
 \operatorname{tr}(T)\equiv2(-1)^{\epsilon_3},\qquad
 \operatorname{tr}(W^2)\equiv2(-1)^{\epsilon_2}\pmod3.
\]
Both traces are units and their product has residue
$(-1)^{\epsilon_2+\epsilon_3}$. By Hensel's lemma, a unit of
$\mathbb Z_3$ is a square in $\mathbb Q_3$ precisely when its residue
is a square in $\mathbb F_3$. The odd-prime Hilbert formula gives
$(e,3)_3=(e/3)$ for a unit $e$. In particular $(-1,-3)_3=-1$,
so this quaternion algebra represents the unique nonzero class of
$\operatorname{Br}(\mathbb Q_3)[2]$. The formula follows.
Both traces are unchanged by independent conjugation, which proves
its invariance on the entire stated domain.
\end{proof}

For comparison with centralizer coordinates, suppose
\[
 T=aI+bR,\quad W=cI+dS,\qquad
 R^2+R+I=0,\quad S^2=-I,
 \qquad a^2-ab+b^2=c^2+d^2=1,
\]
where $R,S$ are the standard cubic and harmonic matrices over $\mathbb Z_3$.
The trace identities give
\[
 \operatorname{tr}(T)=2a-b\equiv2(a+b),\qquad
 \operatorname{tr}(W^2)=2(c^2-d^2)\pmod3.
\]
The norm equations imply that $a+b$ and $c^2-d^2$ are units. Thus
\[
 \mathfrak b(T,W)=\big[((a+b)(c^2-d^2),3)_{\mathbb Q_3}\big].
\]
This recovers the product of the two previously given local symbols,
while the trace expression makes the definition intrinsic to the ordered
pair of matrices.

The connection with the modular frame character is exact on its stated
domain. For $F\in H_0$, take the signed three-adic images of $A_F$ and
$\beta_F$ as $T$ and $W$. Their reductions satisfy the displayed conditions
with $\epsilon_3=a_3(F)$ and $\epsilon_2=a_2(F)$, by the definition of the
frame character. Therefore
\[
 \mathfrak b(T,W)=\delta(F)[(-1,-3)_{\mathbb Q_3}].
\]
Here $H_0=\{F:\lambda_F=1,\ t_F=0\pmod{12}\}$; this conclusion uses only
the finite frame conditions and does not require the uncorrected
three-adic matrices themselves to centralize the elliptic matrices.
For corrected frame pairs satisfying the centralizer norm forms above,
the same formula holds with their specified signed frame bits.

The class can be nonzero on the matrix domain: $T=-I,W=I$ gives its
nonzero value. This matrix calculation alone does not prove that such a
pair comes from any designated subgroup. On normalized elements of
$\widehat{\operatorname{GT}}$, the class vanishes by \cite[Theorem~3]{H17}. The algebra $(-1,-3)$ over $\mathbb Q$
has ramification set $\{3,\infty\}$, whereas $(2,3)$ has ramification set
$\{2,3\}$; these rational algebras cannot be identified merely because
their completions at three are isomorphic.

\section{The quartic Kummer layer over the Gaussian field}
The mod-four Kummer extension is a classical dihedral extension.
The following computation fixes its Hecke conductors and specifies the
additional analytic hypothesis that a square-root error term would use.
Let $K=\mathbb Q(i)$, let $\mathfrak p=(1+i)$, and put
$L=K(\alpha)$ with $\alpha=2^{1/4}>0$.
For a prime ideal $\mathfrak q\nmid2$ of $K$, use arithmetic Frobenius and define
\[
 \psi(\mathfrak q)=\frac{\operatorname{Frob}_{\mathfrak q}(\alpha)}{\alpha}
  =\left(\frac{2}{\mathfrak q}\right)_4.
\]
The last expression is the unique fourth root of unity whose reduction modulo
$\mathfrak q$ is $2^{(N\mathfrak q-1)/4}$. Extend multiplicatively to ideals
prime to two.

\begin{proposition}\label{prop:quartic-conductor}
The extension $L/K$ is cyclic of degree four. Its two faithful finite-order
Hecke characters are $\psi,\psi^{-1}$ and
\[
 \mathfrak f(\psi)=\mathfrak f(\psi^{-1})=\mathfrak p^6,
 \qquad \mathfrak f(\psi^2)=\mathfrak p^4.
\]
All three characters have trivial infinity type. The absolute discriminant of
$L$ is $2^{24}$. The representation
$\operatorname{Ind}_{G_K}^{G_{\mathbb Q}}\psi$ is odd and irreducible of
dimension two, has conductor $256$, and has determinant $\chi_{-4}$.
\end{proposition}
\begin{proof}
The polynomial $X^4-2$ is Eisenstein at two. Its real root field has degree four
and does not contain $i$, so $[L:\mathbb Q]=8$ and $[L:K]=4$.
The automorphism $r(\alpha)=i\alpha$, $r(i)=i$ generates
$\operatorname{Gal}(L/K)$; complex conjugation $s$ satisfies $srs=r^{-1}$.
This proves the dihedral description. Artin reciprocity identifies the
characters of this finite cyclic quotient with finite-order Hecke characters.
The displayed Frobenius formula follows by reducing $\alpha^{N\mathfrak q-1}$.
There is no ramification away from two, since the polynomial discriminant
of $X^4-2$ is a power of two and $K/\mathbb Q$ is unramified away from two.

We compute the different at two directly. Let $M=\mathbb Q(\zeta_8)=K(\sqrt2)$,
let $\varpi=1-\zeta_8$, and normalize the local valuation on $M$ by
$v_M(\varpi)=1$. Then $v_M(2)=4$ and $v_M(\sqrt2)=2$.
Put $D=\sqrt2/\varpi^2$. From
\[
 \sqrt2=\zeta_8+\zeta_8^{-1}
        =\frac{2-2\varpi+\varpi^2}{1-\varpi}
\]
we obtain $D\equiv1+\varpi\pmod{\varpi^2}$. For
$\beta=\alpha/\varpi$ and $\gamma=\beta-1$ we have
\[
 \gamma^2+2\gamma+(1-D)=0.
\]
This polynomial is Eisenstein over $M$ because $v_M(1-D)=1$.
Thus $L/M$ is totally ramified quadratic at two, $\gamma$ is a uniformizer,
and its different is generated by the derivative $2\gamma+2=2\beta$.
Since $\beta$ is a unit and $v_L(2)=8$, the different exponent is eight.
The discriminant of $M/\mathbb Q$ is $2^8$, from the cyclotomic polynomial
$X^4+1$. The discriminant formula in a tower therefore gives
$|d_L|=(2^8)^2\,2^8=2^{24}$.

Since $d_K=-4$, the relative discriminant ideals are
\[
 \mathfrak d_{M/K}=\mathfrak p^4,
 \qquad \mathfrak d_{L/K}=\mathfrak p^{16}.
\]
Indeed their norms are respectively $2^8/4^2=2^4$ and
$2^{24}/4^4=2^{16}$. The conductor-discriminant formula
\cite[Chapter V, Theorem~3.27, p.~165]{MilneCFT}
gives the conductor of the quadratic character $\psi^2$ as $\mathfrak p^4$.
Writing the common conductor exponent of the inverse faithful characters as
$a$, the same formula gives $2a+4=16$, hence $a=6$.

Conjugation by $s$ sends $\psi$ to $\psi^{-1}\ne\psi$, so induction is
irreducible. On $r$ its two eigenvalues are $i,-i$, and on $s$ they are
$1,-1$. Its determinant is consequently the character of $K/\mathbb Q$,
and it is odd. The induction formula for conductors gives
$|d_K|N_{K/\mathbb Q}\mathfrak f(\psi)=4\cdot2^6=256$.
\end{proof}

Hecke's theta construction associates to $\psi$ the weight-one form
\[
 \theta_\psi(z)=\sum_{\substack{\mathfrak a\subset\mathbb Z[i]\\
                     (\mathfrak a,2)=1}}
                 \psi(\mathfrak a)e^{2\pi i N\mathfrak a z}.
\]
It is the classical newform of level $256$ and character $\chi_{-4}$, recorded
as LMFDB $256.1.c.a$ \cite{LMFDB256}. Its first coefficients are
$q+q^9-2q^{17}-q^{25}+O(q^{33})$.

For an unramified prime ideal $\mathfrak q$ of $K$, specifying the value
$\psi(\mathfrak q)=i^a$ is a cyclic Chebotarev condition. A square-root
error term obtained from the explicit formula uses the named hypothesis
that every nontrivial zero of $L_K(s,\psi^j)$, $j=0,1,2,3$, has real part
$1/2$, where $j=0$ means $\zeta_K(s)$. The factorization
\[
 L_K(s,\psi^2)=L_{\mathbb Q}(s,\chi_8)
                    L_{\mathbb Q}(s,\chi_{-8}),\qquad
 \zeta_K(s)=\zeta_{\mathbb Q}(s)L_{\mathbb Q}(s,\chi_{-4})
\]
shows the new analytic input at this layer to be the quartic Hecke
$L$-function (the inverse character has the same $L$-function).
If extra cyclotomic congruences are imposed, the corresponding finite family
of Hecke twists $\psi^j(\chi\circ N_{K/\mathbb Q})$ must be included.
The existence of appropriate primes, and a terminating search for one, does
not require GRH. Nor does the algebraic decomposition itself require it.
A quantitative assertion must state its requested bound before a stronger
analytic hypothesis is described as necessary.

% Standalone research fragment: exact coefficient proof and arithmetic witness.
% The two global inputs are the normalized signed-congruence theorem and
% simultaneous rational Kummer surjectivity.  No coefficient identity below
% depends on the accompanying computer calculations.

\section{A cubic Magnus normalization beyond rational radicals}
\label{sec:cubic-magnus-normalization}

Put $G=\widehat{\mathrm{GT}}$, let $H\subseteq G$ be the Galois image, and
write $N=\ker(\lambda,\rho_2^I)$.  For a multiplier-one element
$g=(1,f)$, write its pro-$2$ Magnus expansion as
\[
 x=1+S,\qquad y=1+T,\qquad
 m_3(g)=[S^2T]f\in\mathbf Z_2.
\]
Here brackets around a word mean its coefficient, whereas group commutators
are $[a,b]=aba^{-1}b^{-1}$.  The notation $m_3$ separates this cubic
monodromy coordinate from the rational Kummer coordinate $\rho_3^I$.

We use two previously established global inputs.  First, for $g\in N$,
\begin{equation}\label{eq:magnus-normalized-congruence-input}
 f(U^2,V^2)=I,
 \qquad
 U=\begin{pmatrix}1&1\\0&1\end{pmatrix},\quad
 V=\begin{pmatrix}1&0\\-1&1\end{pmatrix}
\end{equation}
in $\mathrm{SL}_2(\widehat{\mathbf Z})$.  Second, over
$\mathbf Q^{\mathrm{ab}}$, the simultaneous rational Kummer map has image
$\prod_q2\widehat{\mathbf Z}$, with $q$ running over all rational primes.
Thus a fiber of a fixed cyclotomic multiplier is the full coset determined
by the quadratic Kummer characters.  These are the global inputs to the argument below, proved in
Proposition~\ref{prop:normalized-congruence-triviality} and
Theorem~\ref{thm:main-rational-normalization}.

\begin{lemma}[An integral alternating coefficient identity]
\label{lem:integral-alternating-cubic}
Let $f$ be a pro-$2$ commutator word satisfying
\[
 f(y,x)f(x,y)=1,\qquad
 f(z,x)f(y,z)f(x,y)=1,\qquad z=(xy)^{-1}.
\]
Set $a=[S^2T]f$ and $A=[S^2T^2S]f$.  Then
\begin{equation}\label{eq:alternating-cubic-identity}
 [STS]f=-2a,\qquad [ST^2S]f=a,\qquad
 [STSTS]f=-2a-4A.
\end{equation}
In particular, $[STSTS]f\equiv2a\pmod4$.
\end{lemma}

\begin{proof}
The degree-two part of $f$ is a multiple of $[S,T]$.  The degree-two
part of the three-cycle identity is three times that multiple, so it is
zero.  Thus $f\in\gamma_3$ in the pro-$2$ nilpotent quotient relevant
to the calculation.

For completeness, the low-weight collection calculation needed here is
given explicitly.  Let $L_w$ be the standard Lyndon group commutator
associated to a Lyndon word $w$ on $x<y$: $L_x=x$, $L_y=y$, and
$L_w=[L_u,L_v]$ for the standard factorization $w=uv$ with $v$ the
longest proper Lyndon suffix.  Up to weight five, collect $f$ as the
ordered product of these commutators, with coordinates
\[
\begin{gathered}
\begin{array}{c|rrrrr}
 w&xxy&xyy&xxxy&xxyy&xyyy\\
 \hline
 \text{coordinate}&a&d&u&v&w
\end{array}
\\[5pt]
\begin{array}{c|rrrrrr}
 w&xxxxy&xxxyy&xxyxy&xxyyy&xyxyy&xyyyy\\
 \hline
 \text{coordinate}&b&h&i&j&e&k.
\end{array}
\end{gathered}
\]
All coordinates belong to $\mathbf Z_2$.  Since every factor begins in
degree at least three, the product is linear in these coordinates
through degree five.  Write $I=f(y,x)f(x,y)-1$ and
$II=f(z,x)f(y,z)f(x,y)-1$.  Direct substitution of
$z=(1+T)^{-1}(1+S)^{-1}$ gives the following five coefficients:
\[
\begin{array}{c|l}
\text{coefficient}&\text{value}\\ \hline
{[S^2T]I}&a+d\\
{[S^3T]I}&u-w\\
{[S^3T^2]I}&h+j-u+w\\
{[S^3T]II}&a-2d+2u+v+2w\\
{[S^4T]II}&-a+3d-h-j-2u-v-3w.
\end{array}
\]
Each entry is zero.  The first three give $d=-a$, $w=u$, and
$h+j=0$; the last two then give
\[
 3a+4u+v=0,\qquad -4a-5u-v=0.
\]
Consequently $u=w=-a$ and $v=a$.  No division by two has occurred.
The same commutator expansion gives
\begin{equation}\label{eq:cubic-three-small-coefficients}
 [STS]f=-2a,\qquad
 [ST^2S]f=a,\qquad
 [STS^2]f=4a+3u=a.
\end{equation}
This supplies all the low-weight collection identities used below.

Let $C=[ST^2S^2]f$.  The group-ring antipode sends each generator to
its inverse and reverses products.  Since $f-1$ starts in degree three,
$f^{-1}=1-(f-1)$ through degree five.  The coefficient of $S^2T^2S$
in the antipode of $f$ is therefore
\[
 -A=-\bigl(C+[STS^2]f+[ST^2S]f+[STS]f\bigr)=-C.
\]
Thus $C=A$.

Finally, group-likeness of $f$ for
$\Delta(S)=S\otimes1+1\otimes S+S\otimes S$ and the analogous
formula for $T$ gives, on taking the coefficient of
$S\otimes ST^2S$,
\[
 0=2[S^2T^2S]f+[STSTS]f+2[ST^2S^2]f+2[ST^2S]f.
\]
The right-hand side counts the five insertions of the extra $S$ and
the two possible mergers with an existing $S$.  Substitution of
$C=A$ and \eqref{eq:cubic-three-small-coefficients} proves
\eqref{eq:alternating-cubic-identity}.
\end{proof}

\begin{proposition}[Divisibility and the cubic character]
\label{prop:cubic-magnus-eight}
One has
\[
 m_3(N)\subseteq8\mathbf Z_2.
\]
Consequently
\[
 \eta_3:N\longrightarrow\mathbf F_2,
 \qquad \eta_3(g)=m_3(g)/8\pmod2
\]
is an ordinary continuous $G$-invariant quadratic character.
\end{proposition}

\begin{proof}
Under the substitution in
\eqref{eq:magnus-normalized-congruence-input},
\[
 S=\begin{pmatrix}0&2\\0&0\end{pmatrix},\qquad
 T=\begin{pmatrix}0&0\\-2&0\end{pmatrix}.
\]
Consecutive equal letters vanish.  Every even-length surviving word
is diagonal, and every word of degree at least seven vanishes modulo
$128$.  The only contributions to the upper-right entry modulo $128$
are therefore
\[
 [STS]f\,(-8)+[STSTS]f\,32
 =16a+32(-2a-4A)\equiv80a\pmod{128}.
\]
The matrix identity makes this zero.  Since $80=16\cdot5$, it follows
that $a\in8\mathbf Z_2$.

On $\ker\lambda$, the cubic Magnus coordinate is additive under the
GT product.  Moreover, for $s\in G$ and $n\in\ker\lambda$,
\begin{equation}\label{eq:cubic-scalar-conjugation}
 m_3(sns^{-1})=\lambda_2(s)^3m_3(n).
\end{equation}
Indeed, a multiplier-one word has no part below degree three; the
GT product adds its degree-three terms, and conjugation acts on that
degree by the cube of the degree-one scalar.  Terms from the
conjugating word start in degree two and cannot change this leading
degree-three term.  These statements may equally be read from the
induced tangential automorphism in degree four.  Continuity follows
from the Magnus coefficient map.  Finally, every $2$-adic unit has
odd cube, so division by eight and reduction modulo two turn
\eqref{eq:cubic-scalar-conjugation} into invariance.
\end{proof}

We next give the arithmetic witness, including the small congruence
which detects its parity.

\begin{lemma}[A boundary parity congruence]
\label{lem:rho2-cubic-modfour}
For every multiplier-one gentle element,
\[
 \rho_2^I(g)\equiv-2m_3(g)\pmod4
\]
in the $2$-adic coordinate.
\end{lemma}

\begin{proof}
Use the affine matrix representation
\[
 x\longmapsto D=\begin{pmatrix}-1&0\\0&1\end{pmatrix},\qquad
 y\longmapsto E=\begin{pmatrix}1&1\\0&1\end{pmatrix}.
\]
In the index-two cyclic subgroup, $x^2$ maps to the identity, while
$Y_0=y$ and $Y_1=xyx^{-1}$ have translation coordinates $1$ and $-1$.
Since $f$ has zero total $y$-exponent, its translation is
$2\ell_2(f)=-2\rho_2^I(f)$.

On the Magnus side, $S=D-I=\operatorname{diag}(-2,0)$ and $T=E-I$.
Here $T^2=TS=0$.  Since the first nonzero term of $f$ has degree
three, its translation is
\[
 \sum_{n\ge2}(-2)^n[S^nT]f\equiv4m_3(f)\pmod8.
\]
The series converges $2$-adically.  Comparing translations and
dividing the congruence by $-2$ proves the claim.
\end{proof}

Let
\[
 D=N_{\mathrm{rad}}=\bigcap_q\ker(\rho_q^I)\cap\ker\lambda,
 \qquad
 K_{\mathrm{rad}}=D\cap H=G_L,
\]
where $L=\mathbf Q(\mu_\infty,q^{1/n}:q\text{ prime},n\ge1)$.
The all-prime normalization gives $G=DH$.

\begin{proposition}[An arithmetic value outside the rational Kummer sector]
\label{prop:cubic-arithmetic-witness}
There exists $k\in K_{\mathrm{rad}}$ such that
\[
 m_3(k)\in8\mathbf Z_2^\times.
\]
In particular, $\eta_3|_{K_{\mathrm{rad}}}\ne0$, so $\eta_3$ is
independent of every character generated by rational Kummer coordinates.
Moreover,
\[
 m_3(D)=m_3(K_{\mathrm{rad}})=8\mathbf Z_2.
\]
\end{proposition}

\begin{proof}
Simultaneous rational Kummer surjectivity supplies
$\sigma\in H^{\mathrm{cyc}}=G_{\mathbf Q^{\mathrm{ab}}}$ with
\[
 \rho_2^I(\sigma)=2\in\widehat{\mathbf Z},\qquad
 \rho_q^I(\sigma)=0\quad(q\text{ odd}).
\]
Lemma~\ref{lem:rho2-cubic-modfour} shows that $m_3(\sigma)$ is odd.
Choose $h\in H$ with $\lambda_2(h)=3$, and put
$t=\lambda(h)\in\widehat{\mathbf Z}^{\times}$.  For
$a\in\widehat{\mathbf Z}$, the power $\sigma^a$ means the value of
the continuous map $\widehat{\mathbf Z}\to
\overline{\langle\sigma\rangle}$ extending integer powers.
Define
\[
 k=(h\sigma h^{-1})\sigma^{-t}.
\]
Every rational Kummer coordinate of this element is zero:
on the multiplier-one subgroup the cocycles are homomorphisms,
and conjugation by $h$ multiplies each by $t$, exactly as the
profinite power $\sigma^t$ does.  Thus $k\in K_{\mathrm{rad}}$.
By additivity and \eqref{eq:cubic-scalar-conjugation},
\[
 m_3(k)=(3^3-3)m_3(\sigma)
        =24m_3(\sigma)\in8\mathbf Z_2^\times.
\]
Every rational Kummer character vanishes on $K_{\mathrm{rad}}$, while
$\eta_3(k)=1$, proving the independence assertion.

The image of the compact group $K_{\mathrm{rad}}$ under $m_3$ is a
closed additive subgroup of $\mathbf Z_2$.  It is contained in
$8\mathbf Z_2$ by Proposition~\ref{prop:cubic-magnus-eight}, and it
contains a generator of that closed subgroup.  It is therefore all of
$8\mathbf Z_2$.  The same conclusion follows for $D$ by inclusion.
\end{proof}

\begin{corollary}[A further proper normalization]
\label{cor:cubic-radical-normalization}
Define
\[
 D_3=N_{\mathrm{rad},3}=D\cap\ker m_3,
 \qquad K_{\mathrm{rad},3}=K_{\mathrm{rad}}\cap\ker m_3.
\]
Then $D_3$ is a proper closed normal subgroup of $G$, and
\[
 D=D_3K_{\mathrm{rad}},\qquad
 G=D_3H,\qquad D_3\cap H=K_{\mathrm{rad},3}.
\]
Writing
\[
 B=N/\overline{[G,N]K},\qquad K=N\cap H,
\]
the previously obtained radical normalization refines canonically to
\[
 B\simeq
 D/\overline{[G,D]K_{\mathrm{rad}}}
 \simeq
 D_3/\overline{[G,D_3]K_{\mathrm{rad},3}}.
\]
If $L_3=\overline{\mathbf Q}^{\,K_{\mathrm{rad},3}}$, then
$\operatorname{Gal}(L_3/L)\simeq\mathbf Z_2$.  This defines the
arithmetic extension by the explicit cubic monodromy coordinate; no
radical formula for $L_3$ is asserted.
\end{corollary}

\begin{proof}
Closedness and normality follow from continuity and
\eqref{eq:cubic-scalar-conjugation}.  Properness and the factorization
$D=D_3K_{\mathrm{rad}}$ follow from the equality of the two cubic
images.  The remaining factorization and intersection statements are
immediate from $G=DH$.

For the quotient assertion, one may use the following elementary
group argument.  If a profinite group factors as $G=EH$ with $E$
closed normal, then
\[
 E/\overline{[G,E](E\cap H)}
 \simeq G/\overline{[G,E]H}.
\]
The right-hand quotient is abelian because the image of $E$ is
central and it generates after the image of $H$ is killed.  If also
$E=E'(E\cap H)$ and $G=E'H$, then
$\overline{[G,E]H}=\overline{[G,E']H}$: after quotienting by the latter,
the image of $E'$ is central and $H$ is trivial.  Apply this with
$E=D$, $E'=D_3$, and combine with the earlier radical normalization.
Finally, the kernel of the surjective map
$m_3:G_L\to8\mathbf Z_2$ is exactly $G_{L_3}$.
\end{proof}

\begin{corollary}\label{cor:cubic-radical-joint-image}
The map
\[
 N\longrightarrow
 \left(\prod_{q\text{ odd prime}}2\widehat{\mathbf Z}\right)
       \times8\mathbf Z_2,
 \qquad n\longmapsto((\rho_q^I(n))_q,m_3(n))
\]
has kernel $D_3$ and is surjective; already its restriction to
$K=N\cap H$ is surjective.
Consequently the $G$-invariant continuous $\mathbf F_2$-characters
of $N$ which vanish on $D_3$ form exactly
\[
 \left(\bigoplus_{q\text{ odd prime}}\mathbf F_2\psi_q\right)
                \oplus\mathbf F_2\eta_3.
\]
For an odd prime $p$, no nonzero $G$-invariant
$\mathbf F_p$-character of $N$ vanishes on $D_3$.
\end{corollary}

\begin{proof}
Prescribe any rational Kummer vector and cubic value in the displayed
target.  First choose $h\in K$ with the prescribed Kummer vector.
Its cubic coordinate lies in $8\mathbf Z_2$.  The difference between
that coordinate and the prescribed one can be corrected by an element
of $K_{\mathrm{rad}}$, by
Proposition~\ref{prop:cubic-arithmetic-witness}; this changes none of
the Kummer coordinates.  The kernel is exactly the common kernel $D_3$.
Compactness identifies $N/D_3$ with the displayed product.

Every continuous character of this product into a finite group depends
on finitely many coordinates.  At two, each divided coordinate has a
one-dimensional character space, and conjugation multiplies the Kummer
coordinates by $\lambda$ and the cubic coordinate by $\lambda_2^3$.
Both scalars are odd, so every displayed quadratic character is invariant.
At an odd prime $p$, the pro-$2$ cubic factor has no nonzero character
into $\mathbf F_p$.  A Kummer character is multiplied by $\lambda_p$
under conjugation; choosing a multiplier with $\lambda_p\not\equiv1$
modulo $p$ forces every invariant coefficient to vanish.
\end{proof}

The character $\eta_3$ has nonzero arithmetic restriction.  The result
therefore enlarges the explicitly normalized arithmetic character
sector; it does not construct a nonzero character in the kernel of
restriction to the Galois subgroup.

% Complete research fragment. The finite coefficient proof is reproduced below.

\section{An exact weight-five normalization}
\label{sec:weight-five-normalization}

Let $G=\widehat{\mathrm{GT}}$, let $H$ be its Galois image, and write
the radical and cubic normalizations as
\[
 D=N_{\mathrm{rad}},\qquad
 D_3=D\cap\ker m_3,\qquad J_3=D_3\cap H.
\]
The established factorizations give $G=D_3H$.  All Magnus coefficients
in this section are $2$-adic.  For $g=(1,f)\in D_3$, the pro-$2$
Magnus expansion of $f$ starts in degree five: the degree-three and
degree-four Hall coordinates are respectively
$(m_3,-m_3)$ and $(-m_3,m_3,-m_3)$ by the preceding low-weight
calculation.  Define
\[
 m_5(g)=[S^4T]f,\qquad x=1+S,\quad y=1+T.
\]

\begin{theorem}\label{thm:weight-five-exact-image}
The exact integral images are
\[
 m_5(D_3)=m_5(J_3)=32\mathbf Z_2.
\]
In particular,
\[
 \eta_5:D_3\longrightarrow\mathbf F_2,
 \qquad \eta_5(g)=m_5(g)/32\pmod2,
\]
is a continuous $G$-invariant character with nonzero arithmetic
restriction.  The kernel of $m_5$,
\[
 D_5=D_3\cap\ker m_5,
 \qquad J_5=D_5\cap H=J_3\cap\ker m_5,
\]
is a proper closed normal subgroup satisfying
\[
 D_3=D_5J_3,\qquad G=D_5H.
\]
The relative abelian quotient refines canonically to
\[
 B\simeq D_3/\overline{[G,D_3]J_3}
       \simeq D_5/\overline{[G,D_5]J_5}.
\]
\end{theorem}

We first prove the upper bound.  The only global GT input needed
for this step is the normalized signed congruence identity
\begin{equation}\label{eq:weight-five-signed-input}
 f(U^2,V^2)=I,
 \qquad
 U=\begin{pmatrix}1&1\\0&1\end{pmatrix},\quad
 V=\begin{pmatrix}1&0\\-1&1\end{pmatrix}.
\end{equation}
The remaining input consists of the two-cycle and three-cycle
relations in the free pro-$2$ group.

\begin{lemma}[Finite integral coefficient relations]
\label{lem:weight-five-coefficients}
Suppose $f\in\gamma_5$ satisfies
\[
 f(y,x)f(x,y)=1,\qquad
 f(z,x)f(y,z)f(x,y)=1,\qquad z=(xy)^{-1}.
\]
Put
\[
 a=[S^4T]f,\quad
 c_5=[STSTS]f,\quad
 c_7=[STSTSTS]f,\quad
 c_9=[STSTSTSTS]f.
\]
Then
\begin{equation}\label{eq:weight-five-three-coefficients}
 c_5=-2a,\qquad c_7\equiv5a\pmod{16},\qquad
 c_9\equiv2a\pmod4.
\end{equation}
\end{lemma}

The explicit finite recursions and coefficient certificates proving
this lemma are given below.  In particular the third relation does
not omit a product of lower-weight terms: products of two
nonconstant factors begin in degree ten.

\begin{proposition}\label{prop:weight-five-upper-bound}
For $g\in D_3$, one has $m_5(g)\in32\mathbf Z_2$.
\end{proposition}

\begin{proof}
At the matrices in \eqref{eq:weight-five-signed-input}, consecutive
equal Magnus letters vanish.  Surviving even-length words are
diagonal.  The upper-right entry through degree nine is
\[
 32c_5-128c_7+512c_9.
\]
Write $c_7=5a+16r$.  Modulo $512$, the matrix identity gives
\[
 0\equiv32(-2a)-128(5a)=-64\cdot11a\pmod{512},
\]
so $a\in8\mathbf Z_2$.  The third relation in
\eqref{eq:weight-five-three-coefficients} now gives $c_9\in4\mathbf Z_2$.
Degree ten is diagonal, and every term of degree at least eleven
is divisible by $2048$.  Reducing the same entry modulo $2048$
therefore gives
\[
 0\equiv-64\cdot11a\pmod{2048}.
\]
Since $11$ is a $2$-adic unit, $a\in32\mathbf Z_2$.
\end{proof}

\subsection{Finite coefficient certificate}
\label{subsec:weight-five-certificate}

For a word $w$ in the alphabet $x<y$, write $[w]$ for its Magnus
coefficient after replacing $x,y$ by $S,T$.  A word is Lyndon if
it is lexicographically smaller than every nontrivial cyclic
rotation.  Define the standard Lyndon group words by
\[
 L_x=x,\quad L_y=y,\qquad L_w=[L_u,L_v]
\]
when $w=uv$ is the standard factorization with $v$ the longest
proper Lyndon suffix.  Let $h_w\in\mathbf Z_2$ be the corresponding
Hall coordinates of $f$.  Up to degree nine,
\begin{equation}\label{eq:weight-five-linear-collection}
 f=1+\sum_{5\le |w|\le9}h_w(L_w-1).
\end{equation}
This is an equality of truncated Magnus expansions; powers and
products have no quadratic contribution before degree ten.

Here are the complete recursions used to check the coefficient
certificates.  In the ring of noncommutative integer polynomials
truncated after degree eight, put
\[
 X=1+S,\quad Y=1+T,\quad Z=Y^{-1}X^{-1},\qquad
 (1+A)^{-1}=\sum_{j=0}^{8}(-A)^j.
\]
Form $L_w$ by the preceding commutator recursion and put
$F=1+\sum_{5\le|w|\le8}h_w(L_w-1)$.  Define the linear coefficient
forms
\[
 I_v=[v]\bigl(F(Y,X)+F(X,Y)-2\bigr),
\]
\[
 II_v=[v]\bigl(F(Z,X)+F(Y,Z)+F(X,Y)-3\bigr).
\]
All $I_v$ and $II_v$ vanish for an $f$ in the lemma.  The following
identities are identities of integer linear forms in the Hall
coordinates, with the indicated modulus; thus they may be checked
by the displayed finite multiplication and substitution recursions.

First,
\begin{align*}
 3(c_5+2a)={}&36I_{xxxxy}-18I_{xxxyy}+24I_{xxyxy}
             +66I_{xxxxyy}+24I_{xxxyxy}+6I_{xxyyxy}\\
             &+26II_{xxxyy}+66II_{xxxxxy}-12II_{xxxyxy}.
\end{align*}
This proves $c_5=-2a$ without division by a nonunit.

Second, modulo $16$ one has
\begin{equation}\label{eq:weight-five-c7-certificate}
 9(c_7-5a)\equiv\sum_v\alpha_v I_v+\sum_v\beta_v II_v,
\end{equation}
where the complete lists of nonzero coefficients are
\[
\begin{array}{c|r@{\qquad}c|r}
 v&\alpha_v&v&\beta_v\\ \hline
 xxxxy&-4&xxxyy&-5\\
 xxxyy&-1&xxxxxy&-7\\
 xxyxy&6&xxxyxy&-6\\
 xxxxyy&5&xxyxyy&8\\
 xxxyxy&-4&xxxxxyy&4\\
 xxyyxy&-1&xxxyxxy&4\\
 xxxxyyy&8&xxxyxyy&8\\
 xxxyxyy&8&xxxyyxy&8\\
 xxyxxyy&-4&xxxxxyyy&8\\
 xxyxyxy&8&xxxxyxyy&8\\
 xxxyyxyy&8&xxxyxyxy&8\\
 xxyxyyxy&8&xxyxxyxy&8.
\end{array}
\]
Since $9$ is a unit modulo $16$, this proves the second assertion
of \eqref{eq:weight-five-three-coefficients}.

For the degree-nine assertion, the only Hall coordinates that
can contribute modulo four are determined by a smaller matrix
recursion.  Work in $\mathbf Z[t]/(t^{10})$ and put
\[
 P_x=\begin{pmatrix}1&t\\0&1\end{pmatrix},\qquad
 P_y=\begin{pmatrix}1&0\\-t&1\end{pmatrix},\qquad
 P_w=P_uP_vP_u^{-1}P_v^{-1}.
\]
The inverse is the usual determinant-one matrix inverse, followed
by truncation.  Then $d_w=[t^9](P_w)_{12}$ is precisely
$[STSTSTSTS]L_w$.  The full table relevant modulo four is
\[
\begin{array}{c|r|r|l}
 |w|&\text{number of Lyndon words}&\gcd_w d_w&
       \{w:d_w\not\equiv0\pmod4\}\\ \hline
 5&6&2&xxxxy,\ xxxyy,\ xxyxy\quad(d_w\equiv2)\\
 6&9&2&xxxyxy\quad(d_w\equiv2)\\
 7&18&4&\varnothing\\
 8&30&8&\varnothing\\
 9&56&16&\varnothing.
\end{array}
\]
For additional explicit checks, the unabridged degree-nine
coefficient vectors in weights five and six, in lexicographic
Lyndon order, are
\[
 (-2,-22,18,4,0,0),\qquad
 (0,-8,14,4,8,-4,0,0,0).
\]
The recursion specifies all the remaining entries; their
divisibility by four is the only property used here.

There is also the following four-row identity of Hall coefficient
forms modulo two:
\begin{equation}\label{eq:weight-five-c9-parity-certificate}
 h_{xxxyy}+h_{xxyxy}+h_{xxxyxy}
 \equiv I_{xxxxyyy}+I_{xyxxxyyy}+II_{xxyxxy}+II_{xyxyyyy}
 \pmod2.
\end{equation}
Since $h_{xxxxy}=a$, the table, linear collection
\eqref{eq:weight-five-linear-collection}, and
\eqref{eq:weight-five-c9-parity-certificate} give
\[
 c_9\equiv
 2(h_{xxxxy}+h_{xxxyy}+h_{xxyxy}+h_{xxxyxy})
 \equiv2a\pmod4.
\]
This completes the finite certificate and the proof of
Lemma~\ref{lem:weight-five-coefficients}.

\subsection{An arithmetic generator of the exact image}

The arithmetic lower bound uses a weight-five homomorphism on
the entire multiplier-one subgroup.  For such an element, write
\[
 a_3=[S^2T]f,\qquad b_5=[S^4T]f,
 \qquad \tau_5(f)=12b_5-11a_3\in\mathbf Z_2.
\]

\begin{lemma}\label{lem:tau-five}
The map $\tau_5$ is a continuous homomorphism on $\ker\lambda$,
and
\[
 \tau_5(hzh^{-1})=\lambda_2(h)^5\tau_5(z)
 \qquad(h\in G,\ z\in\ker\lambda).
\]
On $D_3$, it is equal to $12m_5$.
\end{lemma}

\begin{proof}
Let $R(X,Y)$ be the Alexander coefficient of $f$ on $[x,y]$,
and let
\[
 B_f(X,Y)=1+(X-1)(Y-1)R(X,Y).
\]
The action on the rank-one Alexander module is semilinear.
Consequently, on the multiplier-one subgroup the cofactors
multiply, and conjugation by multiplier $t$ sends $B_f(X,Y)$
to $B_f(X^t,Y^t)$.  It follows that
\[
 [u^4v]\log B_f(e^u,e^v)
\]
is an additive weight-five $\mathbf Q_2$-valued homomorphism.
Only terms linear in $Y-1$ contribute to this coefficient, so
the logarithm can be omitted for this extraction.  The boundary
Magnus coefficients and the earlier identity $[S^3T]f=-a_3$ give
\[
 (X-1)R(X,1)
 =a_3(X-1)^2-a_3(X-1)^3+b_5(X-1)^4+O((X-1)^5).
\]
Now
\[
 [u^4](e^u-1)^2=\frac7{12},\quad
 [u^4](e^u-1)^3=\frac32,\quad
 [u^4](e^u-1)^4=1.
\]
Thus the extracted homomorphism is $b_5-11a_3/12$.
Multiplication by $12$ proves the asserted integral formula
and its exact transformation rule.  The last assertion uses
$a_3=0$ on $D_3$.
\end{proof}

\begin{proposition}\label{prop:weight-five-arithmetic-generator}
There exists $w\in J_3$ with $m_5(w)\in32\mathbf Z_2^\times$.
\end{proposition}

\begin{proof}
Simultaneous rational Kummer surjectivity over
$\mathbf Q^{\mathrm{ab}}$ supplies $\sigma\in H\cap\ker\lambda$
with
\[
 \rho_2^I(\sigma)=2\in\widehat{\mathbf Z},\qquad
 \rho_q^I(\sigma)=0\quad(q\text{ odd}).
\]
The established boundary congruence
$\rho_2^I\equiv-2m_3\pmod4$ makes $a_3(\sigma)$ odd.
Therefore $\tau_5(\sigma)=12b_5(\sigma)-11a_3(\sigma)$ is odd.

Choose $h\in H$ with $\lambda_2(h)=3$ and put
$t=\lambda(h)\in\widehat{\mathbf Z}^{\times}$.  For
$c\in\widehat{\mathbf Z}$ and $z\in H\cap\ker\lambda$, define
\[
 \mathcal D_c(z)=hzh^{-1}z^{-c}.
\]
The power $z^{-c}$ is the standard continuous profinite power
in the procyclic closure of $z$.

Every rational Kummer coordinate of $\mathcal D_t(\sigma)$
is zero, because its conjugation weight is $t$.  Set
\[
 w=\mathcal D_{t^3}\bigl(\mathcal D_t(\sigma)\bigr).
\]
All rational Kummer coordinates remain zero.  Its cubic
coordinate also vanishes, since conjugation on that coordinate
has weight $t_2^3$.  Hence $w\in J_3$.

Lemma~\ref{lem:tau-five} and the continuity of powers give the
exact equality
\[
 \tau_5(w)=(3^5-3^3)(3^5-3)\tau_5(\sigma)
          =51840\tau_5(\sigma).
\]
This equality holds for the exact homomorphism $\tau_5$ with the
stated conjugation weight.  Since $w\in D_3$,
\[
 m_5(w)=\frac{51840}{12}\tau_5(\sigma)
       =4320\tau_5(\sigma)
       =32\cdot135\,\tau_5(\sigma),
\]
which is a generator of $32\mathbf Z_2$.
\end{proof}

\begin{proof}[Proof of Theorem~\ref{thm:weight-five-exact-image}]
On $D_3$, $m_5$ is a continuous additive homomorphism, with
conjugation weight $\lambda_2^5$; this follows either from its
leading Magnus degree or from Lemma~\ref{lem:tau-five} and
the torsion-freeness of $\mathbf Z_2$.  The upper bound and
the arithmetic generator show that its compact images on
both $D_3$ and $J_3$ are exactly $32\mathbf Z_2$.  Dividing
by $32$ and reducing modulo two gives a $G$-invariant
character, since every $2$-adic unit has odd fifth power.

Equality of the images gives $D_3=D_5J_3$.  Normality follows
from the scalar conjugation rule, properness from the
arithmetic generator, and $G=D_5H$ follows from $G=D_3H$.
For the quotient assertion, if $G=EH$ with $E$ closed normal,
then
\[
 E/\overline{[G,E](E\cap H)}
 \simeq G/\overline{[G,E]H}.
\]
If $E=E'(E\cap H)$ and $G=E'H$, the closures
$\overline{[G,E]H}$ and $\overline{[G,E']H}$ coincide: in
the quotient by the latter, $H$ is trivial and the image
of $E'$ is central.  Apply this with $E=D_3$ and $E'=D_5$.
\end{proof}

If $J_3=G_{L_3}$ and $L_5=\overline{\mathbf Q}^{\,J_5}$,
then $\operatorname{Gal}(L_5/L_3)\simeq\mathbf Z_2$.
The character $\eta_5$ has nonzero arithmetic restriction on $D_3$.
The same witness shows precisely why it does not give another invariant
character on the original normalized subgroup.

\begin{corollary}\label{cor:weight-five-nonextension}
The character $\eta_5$ does not extend to a $G$-invariant continuous
$\mathbf F_2$-character of $D=N_{\mathrm{rad}}$, and hence does not
extend to one of $N$.  Nevertheless $D_5$ is a strictly smaller
normalization with the same relative abelian quotient.  The kernel of
restriction from $X_2$ to $H^1(D_5,\mathbf F_2)$ is exactly
\[
 \bigoplus_{q\text{ odd prime}}\mathbf F_2\psi_q
       \oplus\mathbf F_2\eta_3.
\]
The kernel of restriction from $H^1(D,\mathbf F_2)^G$ to
$H^1(D_5,\mathbf F_2)$ is exactly $\mathbf F_2\eta_3|_D$.
For every odd prime $p$, restriction from $H^1(N,\mathbf F_p)^G$
to $H^1(D_5,\mathbf F_p)$ is injective.
\end{corollary}

\begin{proof}
Use the elements $h,\sigma,t$ of
Proposition~\ref{prop:weight-five-arithmetic-generator}, and put
$z=\mathcal D_t(\sigma)\in D\cap H$.  Its witness is
\[
 w=hzh^{-1}z^{-t^3}=[h,z]z^{1-t^3}\in D_3\cap H,
 \qquad \eta_5(w)=1.
\]
Since $t$ is a profinite unit, $1-t^3\in2\widehat{\mathbf Z}$.
The continuous profinite power $z^{(1-t^3)/2}$ is defined, so
$w$ belongs to the closed subgroup generated by $[G,D]$ and the
squares in $D$.  Every invariant quadratic character of $D$ kills
that subgroup.  It therefore takes the value zero at $w$, which is
incompatible with extending $\eta_5$.  Restricting an invariant
character of $N$ to $D$ proves the second nonextension assertion.

If $\chi\in X_2$ vanishes on $D_5$, its restriction to $D_3$
factors through $D_3/D_5\simeq\mathbf Z_2$.  That restriction is
either zero or $\eta_5$.  The latter possibility has just been
excluded, so $\chi$ vanishes on $D_3$.  Conversely every character
vanishing on $D_3$ vanishes on $D_5$.  Corollary~\ref{cor:cubic-radical-joint-image}
identifies this common kernel as the displayed direct sum.
The same argument on $D$, using $D/D_3\simeq8\mathbf Z_2$,
gives the one-dimensional kernel there.

For odd $p$, a character of $D_3$ vanishing on $D_5$ must be zero,
because $D_3/D_5$ is pro-$2$.  A character of $N$ vanishing on
$D_3$ factors through the joint rational Kummer and cubic quotient.
The cubic factor is pro-$2$, and conjugation by multiplier $-1$
acts by inversion on every rational Kummer coordinate.  Invariance
therefore kills the remaining factor in $\mathbf F_p$.
\end{proof}

\section{The exact weight-seven normalization}
\label{sec:weight-seven-partial}

Let $G=\widehat{\mathrm{GT}}$, let $H$ be its Galois image, and use
the established normalization
\[
 D_5=N_{\mathrm{rad}}\cap\ker m_3\cap\ker m_5,
 \qquad J_5=D_5\cap H,
 \qquad G=D_5H.
\]
All Magnus coefficients below are $2$-adic, with $x=1+S$ and $y=1+T$.
We prove that the first possible nonzero degree on $D_5$ is seven and
write $m_7(g)=[S^6T]f$ for $g=(1,f)\in D_5$.

\begin{theorem}\label{thm:weight-seven-bounds}
Every word belonging to $D_5$ has zero Magnus components below degree
seven. The map $m_7:D_5\to\mathbf Z_2$ is a continuous homomorphism and
\[
 m_7(hgh^{-1})=\lambda_2(h)^7m_7(g)
 \qquad(h\in G,\ g\in D_5).
\]
Its arithmetic and full images are exactly
\begin{equation}\label{eq:weight-seven-image-interval}
 m_7(J_5)=m_7(D_5)=128\mathbf Z_2.
\end{equation}
The corresponding degree-seven Hall images both equal
$8\mathbf Z_2\mathbf L$, where $\mathbf L$ is the primitive
integral vector specified below.
\end{theorem}

\subsection{The first possible leading term}

Use standard Lyndon brackets on $x<y$, ordered first by degree and
then lexicographically, and let $h_w$ denote the corresponding integral
Hall coordinate. The standard factorization $w=uv$ uses the longest
proper Lyndon suffix, and the group word is $L_w=[L_u,L_v]$.

\begin{lemma}\label{lem:weight-seven-leading-spaces}
Suppose that $f\in\gamma_5$ satisfies the two-cycle and three-cycle
relations with multiplier one. Put $a=[S^4T]f$. Its Hall coordinates
in degrees five and six are
\begin{equation}\label{eq:weight-seven-lower-vector}
 a\left(1,-2,\tfrac32,2,\tfrac12,-1,
 -2,\tfrac72,-\tfrac32,-4,-\tfrac12,\tfrac32,
 \tfrac72,\tfrac32,-2\right).
\end{equation}
In particular, $a=0$ implies that $f$ starts in degree seven.

If $f$ starts in degree seven and satisfies the same two relations,
its degree-seven Hall vector has the form $a\mathbf v+b\mathbf w$,
where $a=[S^6T]f$, $b$ is the sixth degree-seven Hall coordinate, and
\begin{align}
\mathbf v={}&\left(1,-3,5,5,-2,0,-12,-5,-2,\tfrac23,
                  -5,5,8,3,-\tfrac{11}3,4,3,-1\right),
                  \label{eq:weight-seven-v}\\
\mathbf w={}&\left(0,0,0,0,0,1,1,0,0,\tfrac13,
                  -1,-1,-1,0,-\tfrac13,0,0,0\right).
                  \label{eq:weight-seven-w}
\end{align}
These are necessary leading-term conditions; they assert no lifting
of an arbitrary truncated solution to a profinite automorphism.
\end{lemma}

\begin{proof}
The finite identities in Section~\ref{subsec:weight-seven-leading-certificate}
give \eqref{eq:weight-seven-lower-vector} and
\eqref{eq:weight-seven-v}--\eqref{eq:weight-seven-w} directly from
coefficients of the two defining relations through degree nine.
They are exact rational identities between explicitly specified
integral coefficient forms. Consequently they apply to the
$2$-adic Hall coordinates of a word.

On $D_5$, the preceding cubic calculation first gives $f\in\gamma_5$,
and $m_5=0$. Equation~\eqref{eq:weight-seven-lower-vector} therefore
makes every degree-five and degree-six coordinate zero. This proves
the asserted absence of lower terms on $D_5$.
\end{proof}

\subsection{The additional leading pentagon relation}

The original pure-braid pentagon has the form
\[
 f(x_{12},x_{23}x_{24})f(x_{13}x_{23},x_{34})
 =f(x_{23},x_{34})f(x_{12}x_{13},x_{24}x_{34})f(x_{12},x_{23});
\]
see \cite[Remark~1.2, p.~98]{Furusho2003}.
For a word starting in degree seven, its degree-seven term $p$
therefore satisfies
\begin{equation}\label{eq:weight-seven-leading-pentagon}
\begin{split}
 p(t_{12},t_{23}+t_{24})+p(t_{13}+t_{23},t_{34})
 -p(t_{23},t_{34})\hspace{12mm}\\
 -p(t_{12}+t_{13},t_{24}+t_{34})-p(t_{12},t_{23})=0.
\end{split}
\end{equation}

\begin{lemma}\label{lem:weight-seven-integral-line}
For an element of $D_5$, the parameters in
\eqref{eq:weight-seven-v}--\eqref{eq:weight-seven-w} satisfy
\[
 16b=19a.
\]
Consequently $a\in16\mathbf Z_2$, and the integral leading line is
generated by the primitive vector
\begin{equation}\label{eq:weight-seven-primitive-vector}
\mathbf L=(16,-48,80,80,-32,19,-173,-80,-32,
           17,-99,61,109,48,-65,64,48,-16).
\end{equation}
The leading Hall vector of an element is $(a/16)\mathbf L$.
Moreover, its leading alternating Magnus coefficient is
$[STSTSTS]f=-2a$.
\end{lemma}

\begin{proof}
Evaluate \eqref{eq:weight-seven-leading-pentagon} in the infinitesimal
Artin action on the free Lie algebra on $X_1,\ldots,X_4$. Write
\[
 D_{ij}(X_i)=[X_i,X_j],\qquad
 D_{ij}(X_j)=[X_j,X_i],\qquad D_{ij}(X_k)=0\quad(k\ne i,j).
\]
The relations $[D_{ij},D_{kl}]=0$ for four distinct indices and
$[D_{ij},D_{ik}+D_{jk}]=0$ follow from Jacobi. A common change of sign
in this infinitesimal action does not change the required vanishing
in degree seven. No faithfulness assertion is needed.

Take the coefficient of $X_1^4X_2X_3X_2X_4$ in the image of $X_1$.
For $p=a\mathbf v+b\mathbf w$, the five terms, before attaching
their respective signs $+,+,-,-,-$, give
\[
\begin{array}{c|rr}
\text{term}&\text{coefficient of }a&\text{coefficient of }b\\ \hline
1&46/3&-16/3\\
2&-6&0\\
3&0&0\\
4&3&0\\
5&0&0.
\end{array}
\]
This table follows by the displayed formulas for $D_{ij}$ and the
recursion $[D,E](X_i)=D(E(X_i))-E(D(X_i))$, applied to the standard
Lyndon brackets. The pentagon thus gives $(19a-16b)/3=0$.

Both $a$ and $b$ are integral Hall coordinates, and $19$ is odd.
Hence $a\in16\mathbf Z_2$. Substituting $b=19a/16$ into
\eqref{eq:weight-seven-v}--\eqref{eq:weight-seven-w} gives
\eqref{eq:weight-seven-primitive-vector}; its coordinate $19$ proves
primitivity. Direct bracket expansion on the preceding plane gives
$[STSTSTS]p=(32/3)(a-b)$, which becomes $-2a$.
\end{proof}

\subsection{The exact beta coordinate and an arithmetic value}

For a multiplier-one element, put $a_j=[S^jT]f$ and define
\begin{equation}\label{eq:weight-seven-tau}
 \tau_7(g)=360a_6+900a_5+780a_4-239a_2\in\mathbf Z_2.
\end{equation}

\begin{lemma}\label{lem:weight-seven-beta}
The map $\tau_7$ is a continuous homomorphism on $\ker\lambda$, with
conjugation weight $\lambda_2^7$. On $D_5$ it equals $360m_7$.
\end{lemma}

\begin{proof}
Use the normalized Alexander beta cofactor
\[
 B_g(X,Y)=1+(X-1)(Y-1)R_g(X,Y),\qquad [f]=R_g(X,Y)[x,y].
\]
The action on the rank-one Alexander module has scalar
$S_\lambda(X)S_\lambda(Y)B_g(X,Y)$, where
$S_\lambda(X)=(X^\lambda-1)/(X-1)$. The $S_\lambda$ factors satisfy
the same semilinear multiplicative cocycle identity as that scalar;
cancelling them gives the corresponding identity for $B_g$.
Thus
\[
 \vartheta_7(g)=[u^6v]\log B_g(e^u,e^v)
\]
is an additive $\mathbf Q_2$-valued homomorphism on $\ker\lambda$,
with conjugation weight $\lambda_2^7$.

Since $B_g(X,1)=1$, differentiation at $v=0$ gives
\[
 \left.\partial_v\log B_g(e^u,e^v)\right|_{v=0}
   =(e^u-1)R_g(e^u,1)
   =\sum_{j\ge2}a_j(e^u-1)^j.
\]
The last equality follows from the affine specialization
$x=\operatorname{diag}(1+S,1)$, $y=I+TE_{12}$: the commutator has
translation $ST$, conjugation by $x$ multiplies it by $1+S$, and
conjugation by $y$ acts trivially on that translation module.
The coefficients of $u^6$ in $(e^u-1)^j$, for $j=2,\ldots,6$, are
\[
 \frac{31}{360},\quad\frac34,\quad\frac{13}{6},\quad\frac52,
 \quad1.
\]
Using the earlier identity $a_3=-a_2$ therefore gives
\[
 \vartheta_7=a_6+\frac52a_5+\frac{13}{6}a_4
                         -\frac{239}{360}a_2.
\]
Multiplication by $360$ proves the integral formula
\eqref{eq:weight-seven-tau}, continuity and the exact weight law.
On $D_5$, all the terms except $a_6$ vanish by
Lemma~\ref{lem:weight-seven-leading-spaces}.
\end{proof}

\begin{proposition}\label{prop:weight-seven-arithmetic-value}
There is $z_5\in J_5$ with
\[
 m_7(z_5)\in128\mathbf Z_2^\times.
\]
More explicitly, it can be chosen so that
$m_7(z_5)=128\cdot199017\,\tau_7(\sigma)$ for an arithmetic
$\sigma$ with $\tau_7(\sigma)$ odd.
\end{proposition}

\begin{proof}
Simultaneous rational Kummer surjectivity supplies
$\sigma\in H\cap\ker\lambda$ with $\rho_2^I(\sigma)=2$ in
$\widehat{\mathbf Z}$ and every odd-prime rational Kummer coordinate
zero. The established boundary congruence
$\rho_2^I\equiv-2m_3\pmod4$ makes $a_2(\sigma)$ odd.
Equation~\eqref{eq:weight-seven-tau} consequently makes
$\tau_7(\sigma)$ odd.

Choose $h\in H$ with $\lambda_2(h)=3$ and put
$t=\lambda(h)\in\widehat{\mathbf Z}^{\times}$. For fixed $h$, set
\[
 \mathcal D_c(z)=hzh^{-1}z^{-c},\qquad
 z_1=\mathcal D_t(\sigma),\quad
 z_3=\mathcal D_{t^3}(z_1),\quad
 z_5=\mathcal D_{t^5}(z_3).
\]
The powers are the continuous profinite powers in the procyclic
closure of the element. The first operation kills every rational
Kummer coordinate. The second kills $m_3$ and places $z_3$ in $J_3$.
The third kills $m_5$ by its established weight-five transformation
law on $D_3$. All preceding vanishings are preserved, so $z_5\in J_5$.

The exact homomorphism and weight law of $\tau_7$ give
\[
 \tau_7(z_5)=(3^7-3)(3^7-3^3)(3^7-3^5)\tau_7(\sigma).
\]
Since $\tau_7=360m_7$ on $D_5$, division by $360$ gives
\[
 m_7(z_5)=\frac{2184\cdot2160\cdot1944}{360}\tau_7(\sigma)
         =128\cdot199017\,\tau_7(\sigma).
\]
The last two factors are odd. The compact arithmetic image
therefore contains $128\mathbf Z_2$; the reverse inclusion is
proved below.
\end{proof}

\begin{proof}[Proof of Theorem~\ref{thm:weight-seven-bounds}]
The lower-degree vanishing and integral leading line were proved in
Lemmas~\ref{lem:weight-seven-leading-spaces} and
\ref{lem:weight-seven-integral-line}. The beta-coordinate identity
gives additivity and the scalar conjugation rule for $m_7$; division
by the nonzero integer $360$ is legitimate in the torsion-free group
$\mathbf Z_2$.

The arithmetic value and closedness show that both images contain
$128\mathbf Z_2$. Theorem~\ref{thm:weight-seven-exact-image}, proved
below from two full-pentagon coefficients and the signed matrix
identity, gives the reverse inclusion. Since the leading vector is
$(m_7/16)\mathbf L$, both Hall images are exactly
$8\mathbf Z_2\mathbf L$.
\end{proof}

\subsection{Finite identities for the leading spaces}
\label{subsec:weight-seven-leading-certificate}

Here are finite recursions and explicit matrix identities that verify
Lemma~\ref{lem:weight-seven-leading-spaces}. They specify the
coefficient calculation without a dependence on an implementation.
Work in noncommutative integer polynomials truncated after degree
nine. Put $X=1+S$, $Y=1+T$, $Z=Y^{-1}X^{-1}$ and
$(1+A)^{-1}=\sum_{j=0}^{9}(-A)^j$. Form the standard Lyndon group
words by the preceding commutator recursion.

For $d=5$ or $7$, let the columns be the Lyndon words of lengths
$d,\ldots,9$, ordered first by length and then lexicographically.
There are $119$ columns for $d=5$ and $104$ for $d=7$. Set
\[
 F_d=1+\sum_{d\le |w|\le9}h_w(L_w-1).
\]
All coefficients through degree nine are linear in the $h_w$,
because products of two nonconstant factors start in degree at
least ten. Define the rows of the integer coefficient matrix $A_d$
by
\[
 I_v=[v]\bigl(F_d(Y,X)+F_d(X,Y)-2\bigr),
\]
\[
 II_v=[v]\bigl(F_d(Z,X)+F_d(Y,Z)+F_d(X,Y)-3\bigr).
\]
The rows vanish on the coordinates of any word satisfying
the two-cycle and three-cycle relations.

Encode a word $v$ of length $n$ by
$c(v)=2^n+\operatorname{bin}(v)$, where $S$ is the binary digit $0$
and $T$ is $1$. The row codes for $I_v$ and $II_v$ are respectively
$c(v)$ and $1024+c(v)$. Let $R_5,R_7$ be the row lists displayed
below, let $C_5=(1,\ldots,96)$ and $C_7=(1,\ldots,82)$, and put
$B_d=A_d[R_d,C_d]$.

Let $\mathbf u$ be the fifteen-entry vector in
\eqref{eq:weight-seven-lower-vector}. In the $119$-dimensional
row space let $T_5$ have the fourteen rows
$2(e_j-u_je_1)$, $2\le j\le15$. In the $104$-dimensional row
space let $T_7$ have the sixteen rows
$3(e_j-v_je_1-w_je_6)$ for $1\le j\le18$, $j\ne1,6$.
Both matrices have integral entries. The finite identities are
\begin{equation}\label{eq:weight-seven-finite-matrix-identity}
 \det B_5=-6603882696,\qquad \det B_7=42661728,
 \qquad
 T_d[:,C_d]B_d^{-1}A_d[R_d,:]=T_d\quad(d=5,7).
\end{equation}
Every entry in these identities is specified by the integer
recursions and the row lists. The inverses are taken over
$\mathbf Q$, and the identities are exact, not congruences.
Applying them to $A_dh=0$ gives $T_dh=0$ and hence the two leading
vector formulas. This is a proof of necessary vanishings and
leading relations; it supplies no lifting assertion for the
remaining coordinates.

% The explicit row lists are appended immediately below.

\[
\begin{split}
R_{5}=\bigl( &33,35,37,65,67,69,77,129,131,133,135,137,\\
 &139,141,147,149,257,259,261,263,265,267,269,275,\\
 &277,281,283,285,293,301,513,515,517,519,521,523,\\
 &525,527,529,531,533,535,537,539,541,547,549,551,\\
 &553,555,557,563,565,569,587,589,595,597,1059,1089,\\
 &1093,1099,1153,1155,1157,1161,1163,1165,1281,1283,1285,1287,\\
 &1289,1291,1299,1301,1317,1325,1537,1539,1541,1545,1547,1549,\\
 &1553,1555,1557,1559,1561,1563,1571,1573,1577,1579,1581,1589\bigr).
\end{split}
\]

\[
\begin{split}
R_{7}=\bigl( &129,131,133,135,137,139,141,147,149,257,259,261,\\
 &263,265,267,269,275,277,281,283,285,293,301,513,\\
 &515,517,519,521,523,525,527,529,531,533,535,537,\\
 &539,541,547,549,551,553,555,557,563,565,569,587,\\
 &589,595,597,1155,1157,1161,1281,1283,1285,1287,1289,1291,\\
 &1299,1301,1317,1325,1537,1539,1541,1545,1547,1549,1553,1555,\\
 &1557,1559,1561,1563,1571,1573,1577,1579,1581,1589\bigr).
\end{split}
\]

\subsection{The exact weight-seven image}
\label{subsec:weight-seven-exact}

Two coefficients of the full pentagon, together with the signed
congruence identity, determine the exact image. The coefficient calculation uses the terms of
weights eight and nine arising from the ordinary pure-braid action.

\begin{lemma}[A matrix lower-central estimate]
\label{lem:weight-seven-matrix-filtration}
Let $Q$ be the closed subgroup of $\mathrm{SL}_2(\mathbf Z_2)$
generated by $U^2,V^2$, where
\[
 U=\begin{pmatrix}1&1\\0&1\end{pmatrix},\qquad
 V=\begin{pmatrix}1&0\\-1&1\end{pmatrix}.
\]
Let $\mathcal B_r$ be the principal congruence group modulo $2^r$,
and let $\mathcal A_r$ consist of determinant-one matrices whose
diagonal entries are $1$ modulo $2^r$ and whose off-diagonal
entries are zero modulo $2^{r+1}$. Then
\[
 \gamma_{2k}(Q)\subseteq\mathcal A_{3k-1},\qquad
 \gamma_{2k+1}(Q)\subseteq\mathcal B_{3k+1}\quad(k\ge1).
\]
In particular, $\gamma_7(Q)\subseteq\mathcal B_{10}$,
$\gamma_9(Q)\subseteq\mathcal B_{13}$ and
$\gamma_{10}(Q)\subseteq\mathcal A_{14}$.
\end{lemma}

\begin{proof}
For $M=(\alpha,\beta;\gamma,\delta)$ of determinant one,
\[
 [M,U^2]=
 \begin{pmatrix}
 1-2\alpha\gamma&2(\alpha^2-1)+4\alpha\gamma\\
 -2\gamma^2&1+2\alpha\gamma+4\gamma^2
 \end{pmatrix}.
\]
It follows for $r\ge1$ that
$[\mathcal B_r,U^2]\subseteq\mathcal A_{r+1}$ and
$[\mathcal A_r,U^2]\subseteq\mathcal B_{r+2}$.
Exchanging the coordinates gives the same assertions for $V^2$.
The determinant identity also gives
$\alpha-\delta\equiv0\pmod{2^{r+1}}$ in $\mathcal A_r$.
The conjugation formulas therefore show that $U^2,V^2$ normalize
$\mathcal A_r$; they also normalize $\mathcal B_r$.
Hence the two generator calculations imply
\[
 [\mathcal B_r,Q]\subseteq\mathcal A_{r+1},\qquad
 [\mathcal A_r,Q]\subseteq\mathcal B_{r+2}.
\]
Starting with $Q\subseteq\mathcal B_1$ and iterating proves the
stated lower-central estimates.
\end{proof}

\begin{theorem}\label{thm:weight-seven-exact-image}
One has
\[
 m_7(D_5)=m_7(J_5)=128\mathbf Z_2.
\]
Consequently, for $D_7=D_5\cap\ker m_7$ and $J_7=D_7\cap H$,
\[
 D_5=D_7J_5,\qquad G=D_7H,\qquad
 B\simeq D_7/\overline{[G,D_7]J_7}.
\]
The subgroup $D_7$ is proper, closed and normal in $G$, and
\[
 \eta_7:D_5\longrightarrow\mathbf F_2,
 \qquad \eta_7(g)=m_7(g)/128\pmod2
\]
is a continuous $G$-invariant character with nonzero arithmetic
restriction.
\end{theorem}

\begin{proof}
It remains to prove $m_7(D_5)\subseteq128\mathbf Z_2$, since an
arithmetic value of exact valuation seven was already constructed.
Use the integral degree-seven Hall coordinates $a,b$, for which
the leading pentagon gives $a=16A$, $b=19A$ with $A\in\mathbf Z_2$.
Define
\[
 M_f=2^{-10}\bigl(f(U^2,V^2)\bigr)_{12}\pmod{16}.
\]
The matrix lower-central estimates give
$\gamma_7(Q)\subseteq\mathcal B_{10}$ and
$\gamma_{10}(Q)\subseteq\mathcal A_{14}$.
Thus $M_f$ is an integral linear form modulo sixteen in the Hall
coordinates of weights seven through nine. Higher factors are
invisible modulo $2^{14}$, and products and powers of the remaining
factors are linear at that precision.

Let $P_\alpha$ and $P_\beta$ be the full-pentagon coefficient forms
defined below. They are integral and vanish for every GT
word. The finite coefficient identity is
\begin{equation}\label{eq:weight-seven-upper128-certificate}
 6(a+b)-8P_\alpha-4P_\beta-5M_f
       \equiv\sum_r e_rE_r\pmod{16},
\end{equation}
where the $E_r$ are coefficients of the two-cycle and three-cycle
relations through degree eleven. All these forms, and every
coefficient $e_r$, are specified below.

For an element of $D_5$, the defining equations make
$P_\alpha=P_\beta=E_r=0$. The normalized signed congruence identity
makes $M_f=0$. Hence $6(a+b)\equiv0\pmod{16}$, so
$a+b\equiv0\pmod8$. Since $a+b=35A$ and $35$ is odd, $A$ is
divisible by eight. Therefore $a\in128\mathbf Z_2$.

The arithmetic construction and compactness now give equality of
both images. The scalar conjugation law proves normality and
invariance, while the arithmetic value proves properness and
nonvanishing. Equality of the two images gives $D_5=D_7J_5$ and
then $G=D_7H$. The relative quotient identification follows from
the preceding centralization argument.
\end{proof}

\paragraph{The full-pentagon coefficient forms.}
Let $\Sigma_i$ be the ordinary Artin automorphism of the free group
on $x_1,x_2,x_3,x_4$:
\[
 \Sigma_i(x_i)=x_ix_{i+1}x_i^{-1},\qquad
 \Sigma_i(x_{i+1})=x_i,
\]
with the other generators fixed. Its inverse sends $x_i$ to
$x_{i+1}$ and $x_{i+1}$ to $x_{i+1}^{-1}x_ix_{i+1}$.
Put
\[
 \mathcal T_{ij}=(\Sigma_{j-1}\cdots\Sigma_{i+1})
                 \Sigma_i^2
                 (\Sigma_{j-1}\cdots\Sigma_{i+1})^{-1}.
\]
These are the images of the pure generators in the pentagon
convention of \cite[Section~5, equation~(17) and the following displayed
formula]{Enriquez2007}; see also \cite[Section~1]{NS2000}.
Use the following five pairs and signs:
\[
\begin{array}{c|c|c|r}
i&\mathcal U_i&\mathcal V_i&\varepsilon_i\\ \hline
1&\mathcal T_{12}&\mathcal T_{23}\mathcal T_{24}&1\\
2&\mathcal T_{13}\mathcal T_{23}&\mathcal T_{34}&1\\
3&\mathcal T_{23}&\mathcal T_{34}&-1\\
4&\mathcal T_{12}\mathcal T_{13}&\mathcal T_{24}\mathcal T_{34}&-1\\
5&\mathcal T_{12}&\mathcal T_{23}&-1.
\end{array}
\]
Products mean composition of automorphisms in the written order.
Write $x_j=1+X_j$ and let $\Delta_{\mathcal U}=\mathcal U-\mathrm{Id}$
act linearly on the completed noncommutative Magnus algebra.
Every pure braid acts trivially in degree one, so each $\Delta$
raises degree by at least one.

For a standard Lyndon group word, let
$Q_w(S,T)=L_w(1+S,1+T)-1$, truncated after degree nine. Define
\begin{equation}\label{eq:weight-seven-pentagon-functional}
 P_\xi(h)=[\xi]\sum_{i=1}^5\varepsilon_i
       \sum_{7\le |w|\le9}h_w
       Q_w(\Delta_{\mathcal U_i},\Delta_{\mathcal V_i})(X_1),
\end{equation}
for the two words
\[
 \alpha=X_1^6X_2X_3X_2X_4,
 \qquad
 \beta=X_1^4X_2X_3X_2X_4X_2X_4.
\]
Multiplication in the polynomial $Q_w$ means composition of the
operators. Formula~\eqref{eq:weight-seven-pentagon-functional}
defines integral linear forms in the Hall coordinates.

We verify their precise scope. A Hall factor of weight at least
seven has Artin operator equal to the identity through degree
increase six. Products of two nonconstant such factors have
degree increase at least fourteen. Thus, through output degree
ten on $X_1$, the group pentagon is exactly the signed sum in
\eqref{eq:weight-seven-pentagon-functional}, linear in the Hall
coordinates. Hall weights at least ten have no effect at this
output degree. Moreover, every operator monomial that can contribute
has at least seven factors $\Delta$. A piece of one factor that
raises degree by four would give total degree increase at least
ten, hence output degree at least eleven. It therefore suffices
to expand the four base Artin images through degree four and use
their full substitution operators at the required precision.
This retains all terms needed from the full pure-braid action.
Since the full pentagon is an identity, $P_\alpha=P_\beta=0$
for the word of every element under consideration.

\paragraph{Complete finite data for the congruence.}
Form the integer Magnus polynomials through degree eleven, using
$(1+A)^{-1}=\sum_{j=0}^{11}(-A)^j$ and the standard group
commutator recursion. The $389$ Hall columns have weights seven
through eleven. Set
\[
 F=1+\sum_{7\le |w|\le11}h_w(L_w-1),\qquad Z=Y^{-1}X^{-1},
\]
and define
\[
 I_v=[v]\bigl(F(Y,X)+F(X,Y)-2\bigr),\qquad
 II_v=[v]\bigl(F(Z,X)+F(Y,Z)+F(X,Y)-3\bigr).
\]
They are linear, because products of two nonconstant factors begin
in degree fourteen. To compute $M_f$, form the two-by-two matrix
words $P_x=U^2$, $P_y=V^2$,
$P_w=P_uP_vP_u^{-1}P_v^{-1}$ modulo $2^{14}$; its coefficient on
$h_w$ is $(P_w)_{12}/2^{10}$ modulo sixteen. The higher Hall
columns have zero contribution as above. The two pentagon forms
are given by \eqref{eq:weight-seven-pentagon-functional}, with
the same Hall column order.

For the following table, encode $I_v$ by
$r=2^{|v|}+\operatorname{bin}(v)$ and $II_v$ by
$r=4096+2^{|v|}+\operatorname{bin}(v)$, with $S=0$, $T=1$.
The complete $244$ nonzero coefficients $e_r$ in
\eqref{eq:weight-seven-upper128-certificate} are listed below.
The asserted congruence is an identity of integral linear forms
on all $389$ Hall columns, so every entry is checkable from the
displayed finite recursions.


{\small
\[
\begin{array}{rr@{\qquad}rr@{\qquad}rr@{\qquad}rr}
r&e_r&r&e_r&r&e_r&r&e_r\\ \hline
129&-2&258&-1&259&3&261&4\\
262&-2&268&2&518&-6&519&-1\\
521&-6&523&6&524&-5&526&-3\\
534&7&1036&-1&1038&-5&1039&-4\\
1042&-5&1044&1&1047&-4&1051&-4\\
1052&-7&1054&-4&1057&-5&1058&-3\\
1059&-5&1063&8&1067&8&1068&5\\
1069&8&1070&-5&1074&-6&1076&8\\
1084&-4&1086&4&1094&-3&1107&4\\
1108&-6&1109&-4&1118&-4&1123&4\\
1144&-7&1180&-6&1208&-2&1212&4\\
1280&-4&1584&-4&2072&-3&2076&-7\\
2078&-4&2079&-4&2084&-6&2086&4\\
2087&8&2089&-5&2090&-5&2091&4\\
2094&4&2096&5&2097&-3&2098&-1\\
2100&-6&2101&8&2103&-4&2108&-4\\
2110&-4&2114&-1&2115&-5&2117&8\\
2119&8&2121&-3&2122&-3&2124&8\\
2126&4&2129&5&2130&-7&2131&8\\
2132&8&2133&8&2137&8&2139&4\\
2140&-7&2145&5&2148&-6&2150&-4\\
2151&8&2155&8&2157&8&2158&-4\\
2162&-3&2163&-4&2164&3&2165&-4\\
2168&8&2172&4&2174&4&2177&-4\\
2181&-3&2188&2&2190&-6&2195&8\\
2196&-1&2198&-4&2200&8&2209&-2\\
2212&5&2214&-6&2217&4&2219&4\\
2221&-4&2226&-2&2227&4&2228&-1\\
2229&-4&2232&-1&2238&-4&2242&-7\\
2244&-2&2260&7&2261&8&2268&4\\
2276&-6&2278&-4&2310&1&2314&-5\\
2332&3&2348&8&2360&-2&2379&8\\
\end{array}
\]
}

{\small
\[
\begin{array}{rr@{\qquad}rr@{\qquad}rr@{\qquad}rr}
r&e_r&r&e_r&r&e_r&r&e_r\\ \hline
2380&6&2382&4&2388&2&2390&-4\\
2404&-2&2406&8&2412&8&2416&1\\
2424&4&2458&8&2470&4&2512&5\\
2544&8&2616&6&2672&7&2800&4\\
2816&6&3184&4&3248&8&3264&2\\
3376&-4&3472&4&3520&-4&3552&4\\
3840&-2&4354&3&4355&-4&4357&8\\
4358&8&4364&-4&4614&-3&4615&-1\\
4617&-6&4619&1&4620&4&4622&4\\
4630&3&5132&4&5134&-6&5135&4\\
5138&8&5139&-7&5140&2&5148&4\\
5150&4&5153&-2&5154&2&5155&7\\
5170&4&5172&-6&5190&-3&5204&-6\\
5212&8&5240&-6&5304&4&6168&-4\\
6174&8&6175&-4&6180&-4&6182&3\\
6183&-4&6185&-6&6186&3&6188&4\\
6192&4&6193&6&6194&-4&6195&4\\
6196&4&6197&4&6201&-4&6204&8\\
6206&-4&6211&-4&6213&1&6214&-3\\
6215&4&6217&-2&6218&-5&6219&4\\
6220&2&6222&-4&6225&3&6226&-3\\
6227&4&6228&2&6233&8&6244&-4\\
6258&8&6260&4&6268&4&6275&8\\
6277&1&6284&2&6286&4&6292&4\\
6294&8&6305&-1&6308&1&6310&-4\\
6322&4&6324&-4&6328&2&6338&-2\\
6348&4&6356&8&6392&-2&6410&7\\
6428&-6&6476&-6&6484&2&6500&-6\\
6912&4&7280&4&7344&4&7472&-4\\
\end{array}
\]
}

\begin{corollary}\label{cor:weight-seven-nonextension}
The character $\eta_7$ has no $G$-invariant continuous extension to
$D_3$, $N_{\mathrm{rad}}$ or $N$.
Moreover, writing $\psi_q=\rho_q^I/2\pmod2$ for the divided
odd-prime Kummer characters on $N$, one still has
\[
 \ker\bigl[H^1(N,\mathbf F_2)^G\longrightarrow
                         H^1(D_7,\mathbf F_2)\bigr]
 =\left(\bigoplus_{q\text{ odd prime}}\mathbf F_2\psi_q\right)
       \oplus\mathbf F_2\eta_3.
\]
On $D_3$, the corresponding kernel is exactly $\mathbf F_2\eta_5$.
\end{corollary}

\begin{proof}
Use the arithmetic elements constructed above, for which
\[
 z_3\in J_3,\qquad
 z_5=h z_3h^{-1}z_3^{-t^5}\in J_5,
 \qquad \eta_7(z_5)=1.
\]
For every continuous character $\chi$ to $\mathbf F_2$ and every
profinite exponent $c$, continuity of powers gives
$\chi(z^c)=(c\bmod2)\chi(z)$. If $\chi$ is $G$-invariant on
$D_3$, then
\[
 \chi(z_5)=\chi(z_3)-(t_2^5\bmod2)\chi(z_3)=0,
\]
because $t_2$ is odd. Thus its restriction to $D_5$ cannot be
$\eta_7$. Restricting an invariant character from either larger
group to $D_3$ proves the other nonextension assertions.

Now suppose that a $G$-invariant character of $N$ vanishes on
$D_7$. Its restriction to $D_5$ factors through
$D_5/D_7=128\mathbf Z_2$, so is zero or $\eta_7$.
Nonextension excludes the latter, and hence it vanishes on $D_5$.
Its restriction to $D_3$ is similarly zero or $\eta_5$.
The weight-five arithmetic word
$z_3=h z_1h^{-1}z_1^{-t^3}$, with $z_1\in N_{\mathrm{rad}}$,
satisfies $\eta_5(z_3)=1$, whereas the same odd-power calculation
makes every invariant character of $N$ zero on $z_3$.
Consequently the character vanishes on $D_3$.
The established cubic joint quotient
\[
 N/D_3\simeq
 \left(\prod_{q\text{ odd prime}}2\widehat{\mathbf Z}\right)
       \times8\mathbf Z_2
\]
now gives exactly the displayed direct sum. Conversely every one
of those characters already vanishes on $D_3$ and hence on $D_7$.
For a character defined on $D_3$, the first nonextension argument
again forces vanishing on $D_5$; the quotient
$D_3/D_5=32\mathbf Z_2$ then gives precisely $\mathbf F_2\eta_5$.
\end{proof}

If $J_5=G_{L_5}$ and $L_7=\overline{\mathbf Q}^{\,J_7}$, the exact
image also gives $\operatorname{Gal}(L_7/L_5)\simeq\mathbf Z_2$.
This specifies the extension by the weight-seven monodromy
coordinate; no expression by rational radicals is asserted.

\section{An exact weight-eight commutator normalization}
\label{sec:weight-eight-normalization}

Put $M=\ker\lambda$, and write
\[
 \begin{gathered}
 D=N_{\mathrm{rad}},\qquad D_3=D\cap\ker m_3,
 \qquad D_5=D_3\cap\ker m_5,\\
 D_7=D_5\cap\ker m_7,\qquad J_7=D_7\cap H.
\end{gathered}
\]
The preceding normalizations give $G=D_7H$. All Magnus expansions
below are in the free pro-$2$ group, with $x=1+S$ and $y=1+T$.
For $g=(1,f)\in D_7$, define the mixed coefficient
\[
 m_8(g)=[S^4TST^2]f.
\]
The depth-one coefficient $[S^7T]f$ is zero; the first available
degree-eight direction is an Ihara bracket.

\begin{theorem}\label{thm:weight-eight-exact-image}
Every word belonging to $D_7$ starts in degree eight. Its homogeneous
part of degree eight lies in a primitive rank-one integral Lie module,
and
\[
 m_8(D_7)=m_8(J_7)=\mathbf Z_2.
\]
Set $D_8=D_7\cap\ker m_8$ and $J_8=D_8\cap H$. Then $D_8$ is a
proper closed normal subgroup of $G$,
\[
 D_7=D_8J_7,\qquad G=D_8H,\qquad
 B\simeq D_8/\overline{[G,D_8]J_8}.
\]
Every word belonging to $D_8$ starts in degree at least nine.
The continuous character $\eta_8=m_8\bmod2$ is $G$-invariant on
$D_7$ and has nonzero arithmetic restriction. It has no
$G$-invariant extension to $D$ or to $N=\ker(\lambda,\rho_2)$.
\end{theorem}

We give the integral leading calculation and an arithmetic
commutator. No existence claim is drawn from a truncated system of
necessary equations.

\subsection{The integral leading line}

For a Lyndon word $w$ in $S<T$, let $L_w$ be its standard Lie
bracket, using the longest proper Lyndon suffix. Write
\[
 \psi_3=L_{SST}-L_{STT},
\]
and, in the degree-five Lyndon order, put
\[
 \psi_5=L_{SSSST}-2L_{SSSTT}
       +\tfrac32 L_{SSTST}+2L_{SSTTT}
       +\tfrac12 L_{STSTT}-L_{STTTT}.
\]
For a Lie element $P$, define the derivation
$\partial_P(S)=0$, $\partial_P(T)=[T,P]$, and use the Ihara bracket
\[
 \{P,Q\}=[P,Q]+\partial_P(Q)-\partial_Q(P).
\]
Let
\[
 \mathcal L_8=2\{\psi_3,\psi_5\}.
\]
Its integral Lyndon coordinates are as follows:
\[
\begin{array}{cr@{\qquad}cr@{\qquad}cr}
 w&[L_w]\mathcal L_8&w&[L_w]\mathcal L_8&w&[L_w]\mathcal L_8\\ \hline
 SSSSSSST&0&SSSTSTTT&-6&SSTTSTST&21\\
 SSSSSSTT&0&SSSTTSST&4&SSTTSTTT&3\\
 SSSSSTST&-4&SSSTTSTT&-3&SSTTTSTT&-15\\
 SSSSSTTT&0&SSSTTTST&-6&SSTTTTST&-3\\
 SSSSTSST&-2&SSSTTTTT&0&SSTTTTTT&0\\
 SSSSTSTT&7&SSTSSTST&-5&STSTSTTT&8\\
 SSSSTTST&17&SSTSSTTT&3&STSTTSTT&-11\\
 SSSSTTTT&0&SSTSTSTT&3&STSTTTTT&-4\\
 SSSTSSTT&-1&SSTSTTST&6&STTSTTTT&-10\\
 SSSTSTST&-12&SSTSTTTT&7&STTTTTTT&0
\end{array}
\]
In particular $\mathcal L_8$ is primitive. Direct associative
expansion gives
\begin{equation}\label{eq:weight-eight-mixed-coefficient}
 [S^4TST^2]\mathcal L_8=7,
 \qquad [S^7T]\mathcal L_8=0.
\end{equation}
The bracket formula and the displayed finite table specify this
identity without a choice of a group lift of either rational Lie
normalization.

\begin{lemma}\label{lem:weight-eight-leading-line}
Suppose $f\in\gamma_8$ satisfies the two-cycle and three-cycle
relations
\[
 f(y,x)f(x,y)=1,\qquad
 f(z,x)f(y,z)f(x,y)=1,\qquad z=(xy)^{-1}.
\]
Then its homogeneous part of degree eight is
\[
 \operatorname{in}_8(f)=b\mathcal L_8\quad\text{for a unique }b\in\mathbf Z_2.
\]
If $h_w$ denotes a degree-eight Lyndon coordinate, then
$b=-h_{SSSTSSTT}$ and $m_8(f)=7b$.
\end{lemma}

\begin{proof}
The following degree-nine certificate proves the necessary rational
line. Its coordinate $-1$ then proves the stated integral assertion.
Collect $f$ modulo $\gamma_{10}$ in the standard Lyndon group
commutators $C_w$, with integral coordinates $h_w$. Explicitly, for the
standard factorization $w=uv$, put
\[
 C_S=x,\qquad C_T=y,\qquad C_w=C_uC_vC_u^{-1}C_v^{-1}.
\]
Set
\[
 F=1+\sum_{8\le |w|\le9}h_w(C_w-1).
\]
There are $30+56=86$ columns, in order of weight and then lexicographic
order. The expression for $F$ is exact through degree nine, since
products of two nonconstant factors have degree at least sixteen.
All the following polynomial operations are truncated after degree
nine. Put $X=1+S$, $Y=1+T$, $Z=Y^{-1}X^{-1}$, and form
\[
 I_v=[v]\bigl(F(Y,X)+F(X,Y)-2\bigr),\qquad
 II_v=[v]\bigl(F(Z,X)+F(Y,Z)+F(X,Y)-3\bigr).
\]
These are integral linear forms on the $86$ Hall coordinates and
vanish for the given word.

Encode $I_v$ by $2^{|v|}+\operatorname{bin}(v)$ and $II_v$ by
$1024+2^{|v|}+\operatorname{bin}(v)$, with $S=0$, $T=1$. Let $A_R$
be the $66$ by $86$ matrix consisting of the following rows, in the
written order:
\[
\begin{aligned}
R=( &257,259,261,263,265,267,269,275,277,281,283,285,293,301,\\
    &513,515,517,519,521,523,525,527,529,531,533,535,537,539,\\
    &541,547,549,551,553,555,557,563,565,569,587,589,595,597,\\
    &1281,1287,1291,1299,1301,1317,1537,1539,1541,1545,\\
    &1547,1549,1553,1555,1557,1559,1561,1563,1571,1573,\\
    &1577,1579,1581,1589).
\end{aligned}
\]
Let $C$ be its first $66$ columns and let $A_{R,C}$ denote the
corresponding square matrix. If $\ell_j$ is the $j$th entry of the
thirty-coordinate table above, ordered lexicographically by its word,
define $T$ to have the $29$ rows
\[
 T_j=e_j+\ell_j e_9\quad(1\le j\le30,\ j\ne9),
\]
where the $e_j$ are rows of length $86$. The finite identities are
\begin{equation}\label{eq:weight-eight-leading-certificate}
 \det A_{R,C}=-12312=-2^3\,3^4\,19,
 \qquad
 T_C A_{R,C}^{-1}A_R=T.
\end{equation}
Here every matrix entry is determined by the displayed group
commutator recursion and the finite geometric expansion
$(1+A)^{-1}=\sum_{j=0}^{9}(-A)^j$. Thus
\eqref{eq:weight-eight-leading-certificate} is an explicit finite
rational matrix identity, not a rank assumption. Multiplying
$A_Rh=0$ by $T_C A_{R,C}^{-1}$ gives
$h_j=-\ell_jh_9$ for all $j\le30$. This proves the claimed line
over $\mathbf Q_2$. Since $h_9\in\mathbf Z_2$ and $\ell_9=-1$,
its scalar $b=-h_9$ is integral. The mixed coefficient is then
$7b$ by \eqref{eq:weight-eight-mixed-coefficient}.
\end{proof}

The earlier leading degree-seven calculation shows that a word in
$D_5$ has its degree-seven part determined by $m_7$. Therefore every
word in $D_7$ starts in degree eight and satisfies the preceding
lemma. In particular its degree-eight part vanishes exactly when
$m_8$ vanishes. This upper-lattice argument uses only the two-cycle
and three-cycle identities once the preceding normalization has
been made; it uses no new pentagon coefficient.

\subsection{A primitive arithmetic commutator}

Recall the three homomorphisms on all of $M$, together with their
scalar conjugation laws:
\[
 \begin{array}{c|c|c}
  \text{homomorphism}&\text{formula}&\text{conjugation weight}\\ \hline
  m_3 &[S^2T]f&\lambda_2^3\\
  \tau_5&12[S^4T]f-11m_3&\lambda_2^5\\
  \tau_7&360[S^6T]f+900[S^5T]f+780[S^4T]f-239m_3
            &\lambda_2^7.
 \end{array}
\]
They are the coefficients of the logarithmic Alexander beta
cocycle: $m_3=c_3$, $\tau_5=12c_5$, and $\tau_7=360c_7$, where
$c_j=[u^{j-1}v]\log B(e^u,e^v)$. Thus these formulas apply before
any radical normalization. On $D_3$ one has $\tau_5=12m_5$, and
on $D_5$ one has $\tau_7=360m_7$.

\begin{lemma}\label{lem:weight-eight-arithmetic-commutator}
One has $\overline{[M,M]}\subseteq D_7$. There are
$\sigma\in H\cap M$ and $h\in H$ such that, with
$s=h\sigma h^{-1}$,
\[
 w_8=[\sigma,s]\in J_7,
 \qquad m_8(w_8)=63\,m_3(\sigma)\,\tau_5(\sigma)
                 \in\mathbf Z_2^\times.
\]
Moreover $w_8=[\sigma,k]$ for some $k\in D\cap H$.
\end{lemma}

\begin{proof}
Every rational-prime Kummer cocycle is a homomorphism on $M$.
Consequently $[M,M]\subseteq D$. The homomorphism $m_3$ makes
these commutators belong to $D_3$; the homomorphism $\tau_5$, which
there equals $12m_5$, makes them belong to $D_5$; and $\tau_7=360m_7$
on $D_5$ makes them belong to $D_7$. All the kernels are closed,
which proves the first assertion.

The exact arithmetic Kummer image supplies
$\sigma\in G_{\mathbf Q^{\mathrm{ab}}}\subset H\cap M$ with
\[
 \rho_2(\sigma)=2\in\widehat{\mathbf Z},\qquad
 \rho_q(\sigma)=0\quad(q\text{ odd prime}).
\]
Write $a=m_3(\sigma)$. The established congruence
$\rho_2\equiv-2m_3\pmod4$ makes $a$ odd, and the formula for
$\tau_5$ makes $\tau_5(\sigma)$ odd as well.
Choose $h\in H$ with $\lambda_2(h)=3$, and put $t=\lambda(h)$.
Profinite powers are defined by the continuous map
$\widehat{\mathbf Z}\to\overline{\langle\sigma\rangle}$ sending
$1$ to $\sigma$. Hence both
\[
 k=s\sigma^{-t}\quad\text{and}\quad r=s\sigma^{-27}
\]
are elements of $H\cap M$. Conjugation has weight $\lambda$
on every Kummer homomorphism, so $k\in D\cap H$. Since powers of
$\sigma$ commute with $\sigma$,
\[
 [\sigma,k]=[\sigma,s]=[\sigma,r]=w_8.
\]
The weight-three law gives $m_3(r)=(3^3-27)a=0$. The low-degree
cycle calculation, whose degree-three and degree-four Hall
coordinates are $(m_3,-m_3)$ and $(-m_3,m_3,-m_3)$, therefore shows
that $r$ starts in degree five. Its leading part is a multiple of
$\psi_5$, by the preceding integral degree-five calculation. The
weight-five law gives its multiplier explicitly:
\[
 [S^4T]f_r=\frac{3^5-27}{12}\tau_5(\sigma)
       =18\tau_5(\sigma).
\]
The leading part of $\sigma$ is $a\psi_3$.

For completeness, with the normalized action
$\Phi_f(x)=x$, $\Phi_f(y)=f^{-1}yf$, multiplication of parameters
is $f\star g=f\Phi_f(g)$. Its leading Lie bracket is
$[P,Q]+\partial_P(Q)-\partial_Q(P)$, since the first-order change
in $T$ is $[T,P]$. Thus the leading part of this group
commutator is
\[
 \operatorname{in}_8(w_8)
   =18a\tau_5(\sigma)\{\psi_3,\psi_5\}
   =9a\tau_5(\sigma)\mathcal L_8.
\]
Taking the coefficient in
\eqref{eq:weight-eight-mixed-coefficient} gives the claimed odd
value $63a\tau_5(\sigma)$.
\end{proof}

\begin{proof}[Proof of Theorem~\ref{thm:weight-eight-exact-image}]
The leading coefficient $m_8$ is an additive continuous map on
$D_7$. This follows either by multiplying the filtered normalized
automorphisms or by noting that substitutions affect a leading
word only above degree eight. Conjugation acts on it by
\[
 m_8(gng^{-1})=\lambda_2(g)^8m_8(n)
 \quad(g\in G,\ n\in D_7).
\]
Indeed the associated graded action on the two free generators is
scalar $\lambda_2$, and conjugation of a tangential derivation of
degree eight has weight $\lambda_2^8$.

The image of $J_7$ is a compact additive subgroup of $\mathbf Z_2$
containing the unit given by the arithmetic commutator. It is
therefore all of $\mathbf Z_2$; the same holds for $D_7$.
The scalar law proves that $D_8$ is normal and that $\eta_8$ is
invariant. The unit proves properness and arithmetic nonvanishing.
Equality of the arithmetic and full images gives $D_7=D_8J_7$, hence
$G=D_8H$. The leading-line lemma shows that words in $D_8$ start
in degree at least nine.

The relative abelian quotient is unchanged by this replacement.
Explicitly, a continuous $G$-invariant homomorphism from $D_8$ to
a profinite abelian group, trivial on $J_8$, extends uniquely to
$G=D_8H$ by sending $dh$ to the value at $d$. The intersection
condition proves well-definedness, and the compact surjection
$D_8\times H\to G$ proves continuity. This is exactly the universal
property of $D_8/\overline{[G,D_8]J_8}$ and of the preceding quotient
$B$.

Finally let $\chi:D\to\mathbf F_2$ be $G$-invariant. The arithmetic
word satisfies $w_8=[\sigma,k]$ with $k\in D$, so
\[
 \chi(w_8)=\chi(\sigma k\sigma^{-1})-\chi(k)=0,
 \qquad \eta_8(w_8)=1.
\]
Thus $\chi$ cannot extend $\eta_8$. An extension to $N$ would
restrict to one on $D$, so it is impossible as well.
\end{proof}

\begin{corollary}\label{cor:weight-eight-original-character-kernel}
Let $E$ be any closed $G$-stable subgroup containing $D$ and contained
in $M$, and let $A$ be a finite abelian group with trivial $G$-action.
Then
\[
 \{\chi\in\operatorname{Hom}_{\mathrm{cts}}(E,A)^G:
                         \chi(D_8)=0\}
 =
 \{\chi\in\operatorname{Hom}_{\mathrm{cts}}(E,A)^G:
                         \chi(D_7)=0\}.
\]
In particular,
\[
 \ker\!\left(H^1(N,\mathbf F_2)^G
                 \longrightarrow H^1(D_8,\mathbf F_2)\right)
 =\left(\bigoplus_{q\text{ odd prime}}\mathbf F_2\psi_q\right)
       \oplus\mathbf F_2\eta_3.
\]
\end{corollary}

\begin{proof}
Every invariant homomorphism on $E$ kills $w_8=[\sigma,k]$, since
$k\in D\subseteq E$. If it also kills $D_8$, its restriction to
$D_7$ factors through $D_7/D_8\simeq\mathbf Z_2$. The image of
$w_8$ topologically generates this quotient, so the restriction is
zero. The reverse inclusion is immediate. Apply this assertion
to $E=N$, $A=\mathbf F_2$, and use the previously proved kernel
formula for restriction to $D_7$.
\end{proof}

\subsection{The extension threshold of the eighth character}
\label{subsec:eta-eight-extension-threshold}

Complex conjugation gives an arithmetic obstruction to extending the
new character through the preceding normalizations.

\begin{proposition}\label{prop:eta-eight-no-extension-d3}
The character $\eta_8$ on $D_7$ extends to the invariant character
\[
 \alpha_8:D_5\longrightarrow\mathbf F_2,\qquad
 \alpha_8(1,f)=[S^4TST^2]f\pmod2,
\]
but has no $G$-invariant extension to $D_3$. More precisely, no
nonzero linear combination of $\eta_7$ and $\alpha_8$ on $D_5$
extends to a $G$-invariant character of $D_3$.
\end{proposition}


\begin{proof}
The exact weight-seven theorem puts the degree-seven part of every
$D_5$ word in $8\mathbf Z_2\mathcal L_7$, where $\mathcal L_7$
is the primitive vector of \eqref{eq:weight-seven-primitive-vector}.
Thus its Magnus expansion over $\mathbf F_2$ starts in degree eight.
The first nonzero term is additive under the normalized GT product,
and conjugation acts by $\lambda_2^8$, which is one modulo two.
This follows by substituting the linear terms in the GT automorphism:
higher terms raise degree, and a homogeneous ambiguity commuting with
$T$ would be a pure power of $T$, absent from a commutator word.
Hence $\alpha_8$ is an invariant character and restricts to $\eta_8$.

On $M=\ker\lambda$ use the additive coordinates
\[
 a=m_3,\qquad T_5=12[S^4T]f-11a,\qquad
 T_7=360[S^6T]f-1020[S^4T]f+661a.
\]
Their conjugation weights are $\lambda_2^3,\lambda_2^5,\lambda_2^7$.
Together with the rational Kummer homomorphisms they imply
\[
 [M,M]\subset D_7,\qquad [M,D_7]\subset D_8.
\]
For the first inclusion all four families vanish on commutators;
then $T_5=12m_5$ and $T_7=360m_7$, and $\mathbf Z_2$ has no torsion.
For the second inclusion use the weight-eight law on $D_7$ and
$\lambda(M)=1$. Therefore $M/D_8$ has nilpotency class at most two,
and its commutator subgroup is contained in $D_7/D_8\simeq\mathbf Z_2$.
The previously proved commutator formula is
\begin{equation}\label{eq:eta-eight-commutator-pairing}
 m_8([z,z'])=\frac7{24}\bigl(a(z)T_5(z')-a(z')T_5(z)\bigr)
 \qquad(z,z'\in M).
\end{equation}

Let $c\in H$ be complex conjugation, and choose $h\in H$ with
$\lambda_2(h)=3$. Put $t=\lambda(h)$ and $\sigma=[h,c]$.
Since $c^2=1$, $\lambda(c)=-1$ and $\rho_2(c)=0$, one has
\[
 c\sigma c^{-1}=\sigma^{-1},\qquad
 \rho_2(\sigma)=2\rho_2(h).
\]
The established parity law makes $\rho_2(h)$ odd at two, and
$\rho_2(\sigma)\equiv-2m_3(\sigma)\pmod4$ makes
\[
 A=a(\sigma),\qquad B=T_5(\sigma),\qquad C=T_7(\sigma)
\]
all odd. The last two assertions also follow directly from their
formulas above. The odd-prime Kummer coordinates of $\sigma$ need
not vanish.

For $e\in\widehat{\mathbf Z}$ put
$D_e(z)=hzh^{-1}z^{-e}$, and define
\[
 z_1=D_t(\sigma),\qquad z_3=D_{t^3}(z_1),\qquad
 z_5=D_{t^5}(z_3).
\]
The weight laws give $z_1\in D\cap H$, $z_3\in D_3\cap H$ and
$z_5\in D_5\cap H$. Moreover
\[
 m_7(z_5)=\frac{(3^7-3)(3^7-3^3)(3^7-3^5)}{360}C
        =128\cdot199017\,C,
\]
so $\eta_7(z_5)=1$. Every invariant quadratic character on $D_3$
kills $z_5$, since conjugation fixes its value on $z_3$ and
$t^5$ is odd at two.

A second arithmetic element is $v=cz_3c^{-1}z_3$.
Every rational Kummer coordinate and each of $a,T_5,T_7$ changes
sign under $c$, so $v\in D_7$. Every invariant character on $D_3$
also kills $v$. We compute $m_8(v)$ in the class-two quotient.
Put $X=\sigma$, $Y=h\sigma h^{-1}$, $Z=h^2\sigma h^{-2}$.
The identity $chc^{-1}=X^{-1}h$ gives
\[
 cXc^{-1}=X^{-1},\quad
 cYc^{-1}=X^{-1}Y^{-1}X,\quad
 cZc^{-1}=X^{-1}Y^{-1}Z^{-1}YX.
\]
For profinite exponents $p,r$, set $W=ZY^{-p}(YX^{-p})^{-r}$.
Class-two collection gives
\begin{equation}\label{eq:eta-eight-conjugation-norm}
 (cWc^{-1})W=[X,Y]^{pr(p-1)-p-r}
                  [X,Z]^{1-pr}[Y,Z]^{1+p+r}.
\end{equation}
Here is a direct verification. Encode the normal form
$X^aY^bZ^c[X,Y]^d[X,Z]^e[Y,Z]^f$ by $(a,b,c,d,e,f)$.
Multiplication is
\[
 (a,b,c,d,e,f)(A,B,C,D,E,F)
 =(a+A,b+B,c+C,d+D-Ab,e+E-Ac,f+F-Bc),
\]
and its $n$th power has coordinates
\[
 (na,nb,nc,nd-\tbinom n2ab,
             ne-\tbinom n2ac,nf-\tbinom n2bc).
\]
The normal form of $W$ is
\[
 (pr,-p-r,1,\ p^2r+pr(r-1)/2,-pr,p+r).
\]
Conjugation by $c$ sends a coordinate tuple to
$(-a,-b,-c,d+b,e+c,f+c)$. Multiplying these two tuples gives
\eqref{eq:eta-eight-conjugation-norm}. Integer exponent identities
extend to profinite exponents by continuity, including the integral
binomial polynomial $n(n-1)/2$.

For $W=z_3$, the exponents are $p=t,r=t^3$. Only after applying
$m_8$ to the central commutator group do their two-adic values become
$3,27$. Thus the exponents in \eqref{eq:eta-eight-conjugation-norm}
are $132,-80,31$ in that coordinate. Formula
\eqref{eq:eta-eight-commutator-pairing} gives
\[
 m_8([X,Y])=63AB,\quad m_8([X,Z])=17010AB,\quad
 m_8([Y,Z])=413343AB.
\]
Consequently
\[
 m_8(v)=(132\cdot63-80\cdot17010+31\cdot413343)AB
       =11461149\,AB
\]
is odd. Hence $(\eta_7(v),\alpha_8(v))=(0,1)$.

Suppose an invariant character on $D_3$ restricts to
$\mu\eta_7+\nu\alpha_8$ on $D_5$. Evaluation at $v$ forces
$\nu=0$; evaluation at $z_5$ then forces $\mu=0$.
The value $\alpha_8(z_5)$ is unnecessary. Finally, an extension
of $\eta_8$ to $D_3$ would restrict on $D_5$ to
$\alpha_8+\epsilon\eta_7$, because $D_5/D_7\simeq128\mathbf Z_2$.
Both possibilities have been excluded.
\end{proof}

\begin{corollary}\label{cor:d3-character-kernel-d8}
One has
\[
 \ker\!\left(H^1(D_3,\mathbf F_2)^G
                    \longrightarrow H^1(D_8,\mathbf F_2)\right)
       =\mathbf F_2\eta_5.
\]
\end{corollary}
\begin{proof}
Let $\chi$ be an invariant character on $D_3$ that vanishes on $D_8$.
Since $D_7/D_8\simeq\mathbf Z_2$ through $m_8$, its restriction
to $D_7$ is zero or $\eta_8$. Subtracting the corresponding multiple
of $\alpha_8$ from $\chi|_{D_5}$ makes it vanish on $D_7$.
It then factors through $D_5/D_7\simeq128\mathbf Z_2$, so
$\chi|_{D_5}=\mu\eta_7+\nu\alpha_8$.
The proposition forces both coefficients to be zero.
Thus $\chi$ factors through $D_3/D_5\simeq32\mathbf Z_2$ and is
a multiple of $\eta_5$. Conversely $\eta_5$ is invariant on $D_3$
and vanishes on $D_5$, hence on $D_8$.
\end{proof}


\section{The exact weight-nine arithmetic normalization}
\label{sec:weight-nine-bounds}

The next odd coefficient gives another exact arithmetic quotient.
Use $G=D_8H$, $J_7=D_7\cap H$ and $J_8=D_8\cap H$.
For $g=(1,f)\in D_7$, put
\[
 m_9(g)=[S^8T]f.
\]
We prove that its image on each of $D_7,D_8,J_7,J_8$ is
$512\mathbf Z_2$. Killing this coefficient on $D_8$ therefore
preserves the relative abelian obstruction and leaves only words
starting in degree ten. The proof gives a uniform integral bound for odd coefficients,
constructs an arithmetic value, and then determines the leading line.

\subsection{An integral measure for all odd coefficients}
\label{subsec:integral-beta-measure}

Put $M=\ker\lambda$ and $N=\ker(\lambda,\rho_2^I)$.
For $g\in M$ write
\[
 B_g(X,Y)=1+(X-1)(Y-1)R_g(X,Y),\qquad
 q_g(X)=(X-1)R_g(X,1)=\sum_{n\ge1}a_n(g)(X-1)^n,
\]
where $a_n(g)=[S^nT]f_g$.  All coefficients lie in $\mathbf Z_2$.
The equality with the single-$T$ Magnus coefficients follows by taking
the metabelian image: the ending-$T$ part of $[x,y]$ is $S$, conjugation
by $x$ multiplies it by $X$, and a product of two commutator terms has
at least two occurrences of $T$.

The beta cocycle and its symmetry follow from Ihara's factorization
\cite[Proposition~1.6.1]{Ihara1999}, recalled in
\cite[Section~2.3, p.~101]{Furusho2003}. The boundary reflection is
an elementary consequence of the two cycle relations:
\[
 B(X,Y)B(XY,Z)=B(Y,Z)B(X,YZ),\quad B(X,Y)=B(Y,X),
 \quad q(X^{-1})=q(X).
\]
Here is a direct check of that reflection. The metabelian forms of
relations (I) and (II), at multiplier one, are
\[
 R(X,Y)=R(Y,X),\qquad
 R(X,Y)+(XY)^{-1}R(Y,(XY)^{-1})
       +Y^{-1}R((XY)^{-1},X)=0.
\]
The factors in the second formula are the commutator coefficients
of $[y,(xy)^{-1}]$ and $[(xy)^{-1},x]$, respectively.
Specializing at $Y=1$ and at $X=1$, and using symmetry, gives
$(X-1)R(X,X^{-1})=0$. Since the one-variable power-series ring is
a domain, $R(X,X^{-1})=0$ and $R(X^{-1},1)=-XR(X,1)$.
Multiplication by $X^{-1}-1$ proves $q(X^{-1})=q(X)$.

For $g\in N$, the normalized duplication equation is
\begin{equation}
 B_g(X,X)=B_g(-1,X).\label{eq:measure-duplication}
\end{equation}
The geometric input is \cite[Theorem~3.1 and the displayed
$(A_2^0)$, p.~102]{Furusho2003}: relations (I), (II), (IV) and
$(K_2)$ imply $B(X,X)=X^{-2\Psi_2}B(-1,X)$.
The convention there is $\Psi_2=-\rho_2^I$
\cite[p.~100, footnote]{Furusho2003}.
The full profinite $(K_2)$ relation follows from
\cite[Section~8, Theorem~8.1, equation~(50), author p.~30]{Enriquez2007}.
Relation (IV) on full $G$ is the input cited in Section~1.

\begin{lemma}[Integral boundary measure]
\label{lem:integral-boundary-measure}
For every $g\in N$ there is an integral $\mathbf Z_2$-valued measure
$\nu_{g,+}$ on $1+4\mathbf Z_2$, of total mass zero, such that
\begin{equation}
 \int t^{2k}\,d\nu_{g,+}(t)
  =(1-2^{2k})\tau_{2k+1}(g)\qquad(k\ge1),
 \label{eq:integral-boundary-moments}
\end{equation}
where the normalization is exactly
\[
 c_{2k+1}(g)=[u^{2k}]q_g(e^u),\qquad
 \tau_{2k+1}(g)=\frac{(2k)!}{2}c_{2k+1}(g).
\]
The same formula defines an integral continuous homomorphism
$\tau_{2k+1}:M\to\mathbf Z_2$ with conjugation weight
$\lambda_2^{2k+1}$. Its value modulo two is independent of $k$:
\[
 \tau_{2k+1}(g)\equiv\tau_3(g)=a_2(g)\pmod2
 \qquad(g\in M).
\]
\end{lemma}
\begin{proof}
The isomorphism
$\mathbf Z_2[[\mathbf Z_2]]\simeq\mathbf Z_2[[X-1]]$ gives an integral
measure $\mu_g$ with transform $q_g(X)=\int X^t\,d\mu_g(t)$.
One can verify the isomorphism directly: the ideals
$(2^r,X^{2^s}-1)$ and $(2^r,(X-1)^n)$ give cofinal topologies.
For fixed $r,n$, the coefficients below degree $n$ of $X^{2^s}-1$
are divisible by $2^r$ for large $s$; conversely,
$(X-1)^{2^s}$ is divisible by two modulo $X^{2^s}-1$.
The finite group-ring coefficients give the measure on residue classes.
Continuous functions are integrated by uniform locally constant
approximation, so pointwise divisibility passes to the integral.

Write $D_Y\log B=B^{-1}Y\partial B/\partial Y$; no analytic
logarithm is intended.  Differentiating the beta cocycle in $Z$ at
$Z=1$ gives
\[
 D_Y\log B(X,Y)=q(XY)-q(Y).
\]
Symmetry gives the analogous $X$ derivative.  Differentiating
\eqref{eq:measure-duplication}, with the chain rule on the diagonal,
therefore gives
\begin{equation}
 q(X)+q(-X)=2q(X^2).\label{eq:boundary-doubling-measure}
\end{equation}
Substitution $X\mapsto-X$ is legitimate: $X-1$ becomes
$-2-(X-1)$, and every resulting coefficient converges $2$-adically.

Boundary reflection makes $\mu_g$ invariant under $t\mapsto-t$.
The transforms $q(-X)$ and $q(X^2)$ correspond respectively to
$(-1)^t\mu_g$ and to the pushforward $[2]_*\mu_g$.
Equation~\eqref{eq:boundary-doubling-measure} and the torsion-freeness
of the measure module imply
\[
 \mu_g|_{2\mathbf Z_2}=[2]_*\mu_g.
\]
If $\nu_g=\mu_g|_{\mathbf Z_2^\times}$ and
$\nu_{g,+}=\mu_g|_{1+4\mathbf Z_2}$, it follows that
\[
 2\int t^{2k}\,d\nu_{g,+}
 =\int t^{2k}\,d\nu_g
 =(1-2^{2k})\int t^{2k}\,d\mu_g.
\]
The first equality uses negation symmetry, which exchanges the two
odd residue classes modulo four.  Expansion of $e^{tu}$ shows that
$\tau_{2k+1}=\tfrac12\int t^{2k}\,d\mu_g$, proving
\eqref{eq:integral-boundary-moments}; its factor $1-2^{2k}$ is a unit.
Also $\mu_g(\mathbf Z_2)=q_g(1)=0$, hence
$\nu_g(\mathbf Z_2^\times)=0$ and $\nu_{g,+}(1+4\mathbf Z_2)=0$.

The integral measure $\mu_g$ and its negation symmetry exist on all
of $M$, using only the integral boundary series and reflection.
Modulo two, the even residue classes contribute zero to
$\int t^{2k}\,d\mu_g$, while the two odd residue classes modulo four
have equal mass. The moment is therefore even, so its half is
integral. Moreover $t^{2k}-t^2$ is divisible by four for every
$t\in\mathbf Z_2$. Taking the integral and dividing by two proves
the asserted congruence; the identity $\tau_3=a_2$ follows from
$a_1=0$. Each moment is a finite rational linear combination of
boundary coefficients, hence is continuous.

On $M$ the beta multipliers multiply.  Their conjugation law is
$B_{hgh^{-1}}(X,Y)=B_g(X^{\lambda_2(h)},Y^{\lambda_2(h)})$,
as follows from the semilinear action on the metabelian commutator
module, with the multipliers of $h$ and $h^{-1}$ cancelling.
Thus $q$ is additive on $M$ and
$q_{hgh^{-1}}(X)=\lambda_2(h)q_g(X^{\lambda_2(h)})$.
The coefficient of $u^{2k}$ has the asserted weight.
\end{proof}

\begin{lemma}[Uniform odd-coordinate upper bound]
\label{lem:uniform-odd-coordinate-bound}
If $g\in N$ and $a_2(g)=a_4(g)=\cdots=a_{2r-2}(g)=0$, then
\[
 a_{2r}(g)\in2^{2r+1}\mathbf Z_2.
\]
\end{lemma}
\begin{proof}
Reflection implies that $q(1+S)$ is a power series in
$S^2/(1+S)$.  This follows recursively by cancelling the first
nonzero even-degree term of a reflection-invariant series; an odd
first degree would have its coefficient equal to its negative.
Consequently the assumed lower even coefficients make $q$ start in
$S^{2r}$.  The measure in Lemma~\ref{lem:integral-boundary-measure}
has zero moments in degrees $0,2,\ldots,2r-2$.
For $t=2z+1$ put
\[
 P_r(t)=\prod_{j=0}^{r-1}\bigl(t^2-(2j+1)^2\bigr)
 =2^{2r}(2r)!\binom{z+r}{2r}.
\]
The binomial polynomial is $\mathbf Z_2$-valued, by continuity from
integers.  Hence
\[
 (1-2^{2r})\tau_{2r+1}
 =\int t^{2r}\,d\nu_{g,+}
 =\int P_r(t)\,d\nu_{g,+}
 \in2^{2r}(2r)!\mathbf Z_2.
\]
As $\tau_{2r+1}=\frac{(2r)!}{2}a_{2r}$, cancellation in
$\mathbf Q_2$ gives the asserted divisibility.
\end{proof}



\subsection{An additive coefficient and an arithmetic value}

For $g\in M=\ker\lambda$, write $a_j=[S^jT]f$. The logarithmic
Alexander beta coefficient
\[
 c_9(g)=[u^8v]\log B_g(e^u,e^v)
\]
is a continuous homomorphism on $M$ and has conjugation weight
$\lambda_2^9$. The boundary formula for the beta cocycle gives
\[
 c_9=[u^8]\sum_{j\ge1}a_j(e^u-1)^j.
\]
Using $a_1=0$, $a_3=-a_2$, and the finite exponential expansion gives
\begin{equation}\label{eq:weight-nine-beta-formula}
 \begin{split}
 \tau_9:=20160c_9={}&20160a_8+70560a_7+95760a_6\\
                  &+63000a_5+20412a_4-2771a_2.
 \end{split}
\end{equation}
This is an integral homomorphism on all of $M$. Explicitly, the
coefficients of $a_2,a_4,a_5,a_6,a_7,a_8$ in $c_9$ are
\[
 -\frac{2771}{20160},\quad\frac{81}{80},\quad
 \frac{25}{8},\quad\frac{19}{4},\quad\frac72,\quad1.
\]
These numbers can also be checked from
$[u^8](e^u-1)^j=j!\genfrac{\{}{\}}{0pt}{}{8}{j}/8!$.
On $D_7$ the lower depth-one terms vanish: the possible degree-eight
bracket has no depth-one component. Consequently $c_9=m_9$ already
on $D_7$, and in particular on $D_8$.

\begin{proposition}\label{prop:weight-nine-arithmetic-value}
There is $w_9\in J_8$ such that
\[
 m_9(w_9)=512\cdot12746612385\,\tau_9(\sigma),
 \qquad \tau_9(\sigma)\in\mathbf Z_2^\times,
\]
where $\sigma$ is the arithmetic element used in the preceding
normalizations. Thus $512\mathbf Z_2\subseteq m_9(J_8)$.
\end{proposition}

\begin{proof}
Choose $\sigma\in H\cap M$ with $\rho_2(\sigma)=2$ and all the
odd-prime Kummer coordinates zero, and choose $h\in H$ with
$\lambda_2(h)=3$. Put $t=\lambda(h)$ and
\[
 \Delta_j(z)=hzh^{-1}z^{-t^j},\qquad
 z_1=\Delta_1(\sigma),\quad z_3=\Delta_3(z_1),\quad
 z_5=\Delta_5(z_3),\quad z_7=\Delta_7(z_5).
\]
All powers are the continuous profinite powers. The previously
proved weight-one, three, five and seven conjugation laws show
successively that $z_1\in D\cap H$, $z_3\in J_3$, $z_5\in J_5$,
and $z_7\in J_7$. By the weight-nine law,
\[
 c_9(z_7)=\frac{\prod_{j\in\{1,3,5,7\}}(3^9-3^j)}{20160}
                          \tau_9(\sigma)
        =512\cdot12746612385\,\tau_9(\sigma).
\]
The integer multiplying $512$ is odd; the full numerator divided
by $20160$ is $2^9 3^{14}\cdot5\cdot13\cdot41$.
Also $m_3(\sigma)$ is odd, and all other coefficients in
\eqref{eq:weight-nine-beta-formula} are even. Thus
$\tau_9(\sigma)$ is odd.

Let $w_8\in J_7$ be the primitive arithmetic commutator of the
preceding section. Its $m_8$ value is a unit. Choose
$u\in\widehat{\mathbf Z}$ whose $2$-adic component is
$m_8(z_7)/m_8(w_8)$, and set
\[
 w_9=z_7w_8^{-u}.
\]
Then $w_9\in J_8$. Since $w_8$ is a commutator of two elements of
$M$ and $c_9$ is a homomorphism on $M$, one has $c_9(w_8)=0$.
Therefore $m_9(w_9)=c_9(w_9)=c_9(z_7)$, as claimed. Compactness
of the arithmetic image supplies the entire ideal
$512\mathbf Z_2$.
\end{proof}

\subsection{The full leading line}

Use the standard Lyndon Lie and group commutators, with $S<T$ and
the longest proper Lyndon suffix. The degree-nine leading Lie
module of multiplier-one GT words that start in degree nine is contained
in $\mathbf Z_2\mathcal L_9$, where the primitive vector
$\mathcal L_9$ has the following coordinates:

{\small
\[
\begin{array}{cr@{\quad}cr@{\quad}cr}
w&[L_w]\mathcal L_9&w&[L_w]\mathcal L_9&w&[L_w]\mathcal L_9\\ \hline
SSSSSSSST&144&SSSTSTSST&-1118&SSTSTTTTT&-3686\\
SSSSSSSTT&-576&SSSTSTSTT&-5303&SSTTSSTTT&-1013\\
SSSSSSTST&1512&SSSTSTTST&-2951&SSTTSTSTT&2059\\
SSSSSSTTT&1344&SSSTSTTTT&4353&SSTTSTTST&5340\\
SSSSSTSST&-1440&SSSTTSSTT&-682&SSTTSTTTT&-3340\\
SSSSSTSTT&-346&SSSTTSTST&-11276&SSTTTSTST&4388\\
SSSSSTTST&-6322&SSSTTSTTT&479&SSTTTSTTT&1868\\
SSSSSTTTT&-2016&SSSTTTSST&-1762&SSTTTTSTT&3558\\
SSSSTSSST&360&SSSTTTSTT&-8109&SSTTTTTST&2434\\
SSSSTSSTT&-1598&SSSTTTTST&-6663&SSTTTTTTT&576\\
SSSSTSTST&3214&SSSTTTTTT&-1344&STSTSTSTT&3299\\
SSSSTSTTT&-2337&SSTSSTSTT&-89&STSTSTTTT&-5446\\
SSSSTTSST&2722&SSTSSTTST&-5496&STSTTSTTT&-1768\\
SSSSTTSTT&4724&SSTSSTTTT&-1949&STSTTTSTT&6452\\
SSSSTTTST&8823&SSTSTSSTT&291&STSTTTTTT&1368\\
SSSSTTTTT&2016&SSTSTSTST&1041&STTSTTTTT&2088\\
SSSTSSSTT&600&SSTSTSTTT&7472&STTTSTTTT&936\\
SSSTSSTST&2148&SSTSTTSTT&-4352&STTTTTTTT&-144\\
SSSTSSTTT&3259&SSTSTTTST&-2416&&\\
\end{array}
\]
}

Its first coordinate is $144$, and its nineteenth coordinate is
$3259$, an odd integer. Hence, if the leading vector is
$A\mathcal L_9$, then $A\in\mathbf Z_2$ and $m_9=144A$.

Here is complete finite data for the necessary leading line.
Collect group commutators of weights nine and ten, giving $155$
columns. Form $I_v,II_v$ as in the preceding section, now through
degree ten. Encode them by $2^{|v|}+\operatorname{bin}(v)$ and
$2048+2^{|v|}+\operatorname{bin}(v)$, respectively.
For the pentagon rows, let $d_{ij}$ be the derivation on the four
free letters $X_1,\ldots,X_4$ defined by
\[
 d_{ij}(X_i)=[X_i,X_j],\qquad
 d_{ij}(X_j)=[X_j,X_i],\qquad d_{ij}(X_k)=0\ (k\ne i,j).
\]
Use the five pairs and signs
\[
 (d_{12},d_{23}+d_{24},+),\quad
 (d_{13}+d_{23},d_{34},+),\quad
 (d_{23},d_{34},-),\quad
 (d_{12}+d_{13},d_{24}+d_{34},-),\quad
 (d_{12},d_{23},-).
\]
For a homogeneous Lie polynomial $P$, let $\mathcal P(P)$ be the
signed sum of its five evaluations on these pairs. This is the
leading Artin image of the pentagon, up to the common sign $(-1)^9$;
therefore it vanishes for every GT word under consideration.
Let $P_j$ be the coefficient on $X_1$ at the following five words,
respectively:
\[
 1111112324,\quad1111123324,\quad1111124324,\quad
 1111232324,\quad1111232424,
\]
where a digit $i$ denotes $X_i$. Give $P_j$ the row code $4096+j$.
Their entries on the degree-ten Hall columns are zero. The
coefficients on the degree-nine columns are determined by
$L_w(\mathcal U,\mathcal V)(X_1)$, with Lie bracket interpreted as
the commutator of derivations.

Let $A_R$ consist of the following $121$ rows, in the written order:
\[
\begin{aligned}
R=(&513,515,517,519,521,523,525,527,529,531,533,535,537,539,\\
&541,547,549,551,553,555,557,563,565,569,587,589,595,597,\\
&1025,1027,1029,1031,1033,1035,1037,1039,1041,1043,1045,1047,1049,1051,\\
&1053,1059,1061,1063,1065,1067,1069,1073,1075,1077,1079,1081,1083,1085,\\
&1093,1095,1097,1099,1101,1107,1109,1113,1115,1117,1125,1129,1131,1133,\\
&1139,1171,1173,1197,1203,1237,2563,2565,2569,2577,2583,2587,2603,2605,\\
&2613,3073,3075,3077,3079,3081,3083,3085,3089,3091,3093,3095,3097,3099,\\
&3101,3107,3109,3113,3115,3117,3119,3121,3123,3141,3145,3147,3149,3157,\\
&3159,3161,3163,3221,3243,3253,4097,4098,4101).
\end{aligned}
\]
Let $C$ be its first $121$ columns, and write
$q_j=\ell_j/144$, where $\ell_j$ is the $j$th coordinate of
$\mathcal L_9$ in lexicographic order. Define $T$ by its $55$ rows
$e_j-q_je_1$ for $2\le j\le56$. The exact finite identities are
\[
 \det A_{R,C}=22883210311680,\qquad
 T_C A_{R,C}^{-1}A_R=T.
\]
All entries are specified by the finite group-commutator and
geometric-series recursions through degree ten, and the displayed
derivations. Thus the identities give $h_j=q_jh_1$ for every
leading vector of a GT word. In particular it is a rational multiple
of $\mathcal L_9$, and its nineteenth coordinate makes the multiple
integral. This proves the asserted upper lattice without claiming
that arbitrary solutions of the necessary equations lift.

\subsection{The exact weight-nine normalization}
\label{subsec:weight-nine-exact}

The integral boundary measure forces the ninth coefficient into
its arithmetic lattice.

\begin{theorem}\label{thm:weight-nine-exact-image}
One has
\[
 m_9(D_8)=m_9(J_8)=512\mathbf Z_2.
\]
The same equality holds with $D_8,J_8$ replaced by $D_7,J_7$.
For $D_9=D_8\cap\ker m_9$ and $J_9=D_9\cap H$,
\[
 D_8=D_9J_8,\qquad G=D_9H,\qquad
 B\simeq D_9/\overline{[G,D_9]J_9}.
\]
The subgroup $D_9$ is proper, closed and normal in $G$, and every
word belonging to it starts in degree at least ten. Moreover,
\[
 (m_8,m_9):D_7\longrightarrow\mathbf Z_2\times512\mathbf Z_2
\]
is surjective already on $J_7$ and has kernel $D_9$.
\end{theorem}

\begin{proof}
The arithmetic word in Proposition~\ref{prop:weight-nine-arithmetic-value}
gives a value of exact valuation nine on $J_8$.
For the upper bound, use Lemma~\ref{lem:uniform-odd-coordinate-bound}
with $r=4$. On $D_7$ the coefficients $a_2,a_4,a_6$ vanish, so
$m_9=a_8\in512\mathbf Z_2$. Concretely, the polynomial
\[
 (t^2-1)(t^2-9)(t^2-25)(t^2-49)
 =t^8-84t^6+1974t^4-12916t^2+11025
\]
is divisible by $2^{15}$ on odd $t$. Its lower moments vanish,
and its eighth moment is $(1-2^8)\tau_9$, with
$\tau_9=20160m_9$ and $v_2(20160)=6$. This is exactly the required
valuation-nine bound. Compactness and the arithmetic witness
therefore prove equality of the two images.

The coefficient $m_9=c_9$ is already an additive, weight-nine
homomorphism on $D_7$. Given any element of $D_7$, multiply it by
a suitable profinite power of the primitive arithmetic
weight-eight commutator $w_8$ to make $m_8$ zero. This does not
change $c_9$, because every logarithmic beta homomorphism vanishes
on $w_8$. Consequently the same exact image equality holds on
$D_7$ and $J_7$. The two arithmetic words $w_8,w_9$ have respective
coordinate values
\[
 (m_8(w_8),0),\qquad(0,m_9(w_9)),
\]
where the first is a unit and the second has exact valuation
nine. They prove the asserted joint surjectivity.

The scalar conjugation laws imply normality of $D_9$. The
arithmetic value proves properness, and equality of the two images
gives $D_8=D_9J_8$, hence $G=D_9H$. The same compact gluing and
centralization argument as in the preceding normalizations gives
the displayed description of $B$. The leading-line calculation
shows that $m_9=0$ kills the whole degree-nine term, so the words
of $D_9$ start in degree at least ten.
\end{proof}

\subsection{The ninth character and its extension obstruction}

\begin{proposition}\label{prop:eta-nine-nonextension}
The character
\[
 \eta_9:D_7\longrightarrow\mathbf F_2,\qquad
 \eta_9(g)=m_9(g)/512\pmod2
\]
is continuous and $G$-invariant. Its restriction to $D_8$ has
nonzero arithmetic restriction and has no $G$-invariant extension
to $D_5$.
\end{proposition}

\begin{proof}
Integrality on $D_7$ and invariance follow from the exact image
and the weight-nine law. Recall the arithmetic construction
\[
 z_7=h z_5h^{-1}z_5^{-t^7}\in J_7,
 \qquad z_5\in J_5,
 \qquad w_9=z_7w_8^{-u}\in J_8,
\]
where $u_2=m_8(z_7)/m_8(w_8)$ and $m_8(w_8)$ is a unit.
The invariant character $\alpha_8$ on $D_5$ vanishes on $z_7$,
since $t$ is odd. Its restriction to $D_7$ is $m_8\bmod2$, so
$u_2$ is even.

Let $\chi:D_5\to\mathbf F_2$ be $G$-invariant. The same odd-power
argument gives $\chi(z_7)=0$. Also
$\chi(w_8^{-u})=-u_2\chi(w_8)=0$. Hence $\chi(w_9)=0$, whereas
$\eta_9(w_9)=1$ by its exact arithmetic valuation. This rules out
an invariant extension from $D_8$ to $D_5$.
\end{proof}

\begin{corollary}\label{cor:weight-nine-character-kernels}
The invariant quadratic-character kernels on restriction to $D_9$
are as follows:
\[
\begin{array}{c|l}
 \text{domain}&\text{kernel}\\ \hline
 D_8&\mathbf F_2\eta_9\\
 D_7&\mathbf F_2\eta_8\oplus\mathbf F_2\eta_9\\
 D_5&\mathbf F_2\eta_7\oplus\mathbf F_2\alpha_8\\
 D_3&\mathbf F_2\eta_5\\
 N&\displaystyle\left(\bigoplus_{q\text{ odd prime}}\mathbf F_2\psi_q\right)
          \oplus\mathbf F_2\eta_3.
\end{array}
\]
\end{corollary}

\begin{proof}
The first two rows follow from the exact scalar and joint quotient
maps. A character on $D_5$ killed on $D_9$ restricts on $D_8$ to
zero or $\eta_9$; the latter is impossible by the preceding
proposition. Its restriction therefore already kills $D_8$,
whose kernel on $D_5$ is the displayed span of $\eta_7,\alpha_8$.
No nonzero combination of those two characters extends to $D_3$,
so the $D_3$ kernel is unchanged. Similarly any invariant character
on $N$ killed on $D_9$ must kill $D_8$, since otherwise its
restriction to $D_5$ would extend $\eta_9$. The preceding kernel
formula for $N$ then gives the last row.
\end{proof}


\section{An exact weight-ten commutator normalization}
\label{sec:weight-ten-normalization}

Use $G=D_9H$ and $J_9=D_9\cap H$. For $g=(1,f)\in D_9$, define
\[
 m_{10}(g)=[S^5TST^3]f.
\]
This mixed coefficient detects the next even-degree bracket.

\begin{theorem}\label{thm:weight-ten-exact-image}
The full and arithmetic images are
\[
 m_{10}(D_9)=m_{10}(J_9)=\mathbf Z_2.
\]
For $D_{10}=D_9\cap\ker m_{10}$ and $J_{10}=D_{10}\cap H$,
\[
 D_9=D_{10}J_9,\qquad G=D_{10}H,\qquad
 B\simeq D_{10}/\overline{[G,D_{10}]J_{10}}.
\]
The subgroup $D_{10}$ is proper, closed and normal, and its words
start in degree at least eleven.
\end{theorem}

\subsection{The primitive leading bracket}

Let $L_w$ denote the standard Lyndon Lie brackets in $S<T$, and set
\[
 \psi_3=L_{SST}-L_{STT},\qquad
 \mathcal L_{10}=\{\psi_3,\mathcal L_7\},
\]
where $\mathcal L_7$ is the primitive degree-seven vector displayed
above. As before,
\[
 \{P,Q\}=[P,Q]+\partial_P(Q)-\partial_Q(P),\qquad
 \partial_P(S)=0,\quad\partial_P(T)=[T,P].
\]
All coefficients of $\mathcal L_{10}$ are integral. Two exact
coefficient calculations are
\begin{equation}\label{eq:weight-ten-unit-coefficient}
 [L_{SSSSSTSTTT}]\mathcal L_{10}=-141,
 \qquad [S^5TST^3]\mathcal L_{10}=-141.
\end{equation}
In particular it is primitive over $\mathbf Z_2$.

\begin{lemma}\label{lem:weight-ten-leading-line}
For every multiplier-one GT word starting in degree ten, the
homogeneous term of degree ten is $A\mathcal L_{10}$ with
$A\in\mathbf Z_2$. It vanishes exactly when $m_{10}$ vanishes.
\end{lemma}

\begin{proof}
The following explicit finite certificate proves the necessary
rational line. Collect the $99+186=285$ standard Lyndon group
commutators of weights ten and eleven. Through degree eleven
write
\[
 F=1+\sum_{10\le|w|\le11}h_w(C_w-1).
\]
Products of two nonconstant factors start in degree twenty, so
this expression is exact at the indicated precision. Define the
cycle forms
\[
 I_v=[v]\bigl(F(Y,X)+F(X,Y)-2\bigr),\qquad
 II_v=[v]\bigl(F(Z,X)+F(Y,Z)+F(X,Y)-3\bigr),
\]
where $X=1+S$, $Y=1+T$ and $Z=Y^{-1}X^{-1}$. Encode the rows by
$2^{|v|}+\operatorname{bin}(v)$ and
$4096+2^{|v|}+\operatorname{bin}(v)$, respectively.

Adjoin one leading Artin pentagon row $P$, with code $8192$.
It is the coefficient of the output word
\[
 11123234324
\]
on $X_1$. Explicitly, use the infinitesimal derivations
$d_{ij}(X_i)=[X_i,X_j]$, $d_{ij}(X_j)=[X_j,X_i]$, with the other
letters fixed. Evaluate the degree-ten Lie polynomial on the
five pairs
\[
 (d_{12},d_{23}+d_{24},+),\quad
 (d_{13}+d_{23},d_{34},+),\quad
 (d_{23},d_{34},-),\quad
 (d_{12}+d_{13},d_{24}+d_{34},-),\quad
 (d_{12},d_{23},-),
\]
and take the indicated coefficient of the signed sum. This is a
necessary coefficient of the full pentagon: a word starting in
degree ten contributes first in Artin output degree eleven, so
only the leading derivations enter. The common sign relating the
ordinary pure Artin action to the displayed $d_{ij}$ is raised to
the even power ten. The row is zero on all degree-eleven columns.

Let $A_R$ be the following $222$ rows, in the written order:
\[
\begin{aligned}
R=(&1025,1027,1029,1031,1033,1035,1037,1039,1041,1043,1045,1047,1049,1051,\\
&1053,1059,1061,1063,1065,1067,1069,1073,1075,1077,1079,1081,1083,1085,\\
&1093,1095,1097,1099,1101,1107,1109,1113,1115,1117,1125,1129,1131,1133,\\
&1139,1171,1173,1197,1203,1237,2049,2051,2053,2055,2057,2059,2061,2063,\\
&2065,2067,2069,2071,2073,2075,2077,2079,2081,2083,2085,2087,2089,2091,\\
&2093,2095,2097,2099,2101,2103,2105,2107,2109,2115,2117,2119,2121,2123,\\
&2125,2127,2129,2131,2133,2135,2137,2139,2141,2147,2149,2151,2153,2155,\\
&2157,2161,2163,2165,2167,2169,2185,2187,2189,2195,2197,2199,2201,2203,\\
&2205,2211,2213,2215,2217,2219,2221,2227,2229,2233,2247,2249,2251,2253,\\
&2259,2261,2265,2267,2277,2281,2341,2347,2349,2355,2357,2379,2381,2387,\\
&2389,5121,5127,5131,5133,5139,5141,5145,5155,5157,5161,5167,5169,5189,\\
&5193,5207,5211,5291,5301,6145,6147,6149,6151,6153,6155,6157,6159,6161,\\
&6163,6165,6167,6169,6171,6177,6179,6181,6183,6185,6187,6189,6191,6193,\\
&6195,6197,6199,6211,6213,6215,6217,6219,6221,6225,6227,6229,6231,6233,\\
&6235,6237,6245,6249,6251,6253,6281,6283,6285,6291,6293,6295,6297,6299,\\
&6307,6309,6313,6315,6317,6325,6345,6347,6357,6437,6443,8192).
\end{aligned}
\]
Let $C$ be the first $222$ columns. Write $\ell_j$ for the
$99$ degree-ten Lyndon coordinates of the explicitly defined
bracket $\mathcal L_{10}$. Its twelfth coordinate is $-141$.
Define the $98$ target rows, of length $285$, by
\[
 T_j=e_j-\frac{\ell_j}{\ell_{12}}e_{12}
       \qquad(1\le j\le99,\ j\ne12).
\]
The exact finite identities are
\[
 \det A_{R,C}=-63256059346747392\equiv37\pmod{101},
 \qquad T_C A_{R,C}^{-1}A_R=T.
\]
Every entry is specified by standard group commutators, geometric
inversion through degree eleven, and the displayed derivation
recursion. Thus the identity forces
$h_j=(\ell_j/\ell_{12})h_{12}$ for every GT word. Its rational
leading scalar is $A=h_{12}/(-141)$, which belongs to
$\mathbf Z_2$ because $141$ is odd. Equation
\eqref{eq:weight-ten-unit-coefficient} gives $m_{10}=-141A$.
This proves the assertion without inferring a lift from the
necessary equations.
\end{proof}

\subsection{An arithmetic unit}

\begin{lemma}\label{lem:weight-ten-arithmetic-unit}
There is an arithmetic commutator $w_{10}\in J_9$ such that
\[
 m_{10}(w_{10})=-141\cdot729\,m_3(\sigma)\tau_7(\sigma)
                                  \in\mathbf Z_2^\times.
\]
Moreover $w_{10}=[\sigma,k]$ for some $k\in D\cap H$.
\end{lemma}

\begin{proof}
Use the established arithmetic element $\sigma\in H\cap M$ with
$\rho_2(\sigma)=2$ and every odd-prime Kummer coordinate zero.
Both $a=m_3(\sigma)$ and $\tau_7(\sigma)$ are odd. Choose
$h\in H$ with $\lambda_2(h)=3$, set $t=\lambda(h)$, and put
\[
 r_3=h\sigma h^{-1}\sigma^{-27},\qquad
 r_5=h r_3h^{-1}r_3^{-243}.
\]
These are multiplier-one arithmetic elements. The first has
$c_3=m_3=0$, so its parameter starts in degree five. The second
has $c_3=c_5=0$, so its parameter starts in degree seven. These
vanishings use the cycle calculations before imposing any radical
normalization.

The weight-seven beta law gives
\[
 c_7(r_5)=\frac{(3^7-3^5)(3^7-3^3)}{360}\tau_7(\sigma)
        =16\cdot729\tau_7(\sigma).
\]
For a word starting in degree seven, $c_7=m_7$. The degree-seven
leading relation, which holds for every multiplier-one GT word
starting in that degree, therefore gives
\[
 \operatorname{in}_7(f_{r_5})=729\tau_7(\sigma)\mathcal L_7.
\]
No normalized weight-seven divisibility claim is being applied to
$r_5$. Its parameter is allowed to have a primitive leading term.
The leading parameter of $\sigma$ is $a\psi_3$. Hence the
commutator
\[
 w_{10}=[\sigma,r_5]
\]
starts in degree ten and has leading term
\[
 729a\tau_7(\sigma)\{\psi_3,\mathcal L_7\}
       =729a\tau_7(\sigma)\mathcal L_{10}.
\]
Every Kummer homomorphism on $M$ vanishes on this commutator. Its
lower Magnus coefficients also vanish, so $w_{10}\in J_9$.
The unit assertion follows from
\eqref{eq:weight-ten-unit-coefficient}.

For the last assertion, the all-prime conjugation law gives
\[
 \rho_q(r_5)=(t-243)(t-27)\rho_q(\sigma)
\]
for every rational prime $q$. Put $c=(t-243)(t-27)$ and
$k=r_5\sigma^{-c}$. The continuous profinite power is defined for
this $c\in\widehat{\mathbf Z}$, and all Kummer coordinates of $k$
are zero. Thus $k\in D\cap H$. Since powers of $\sigma$ commute
with $\sigma$, one has $[\sigma,k]=[\sigma,r_5]=w_{10}$.
\end{proof}

\begin{proof}[Proof of Theorem~\ref{thm:weight-ten-exact-image}]
The leading coefficient $m_{10}$ is an additive continuous
homomorphism on $D_9$, and conjugation acts by
$\lambda_2^{10}$. The arithmetic image contains the unit furnished
by the preceding lemma. Being compact, it is all of
$\mathbf Z_2$, as is the full image. Equality gives
$D_9=D_{10}J_9$ and $G=D_{10}H$; the usual compact gluing argument
identifies the relative quotient $B$ as displayed. The scalar law
proves normality and the unit proves properness. The leading-line
lemma shows that the words of $D_{10}$ start in degree at least
eleven.
\end{proof}

\subsection{The tenth character}

\begin{proposition}\label{prop:eta-ten-extension-kernels}
The invariant character $\eta_{10}=m_{10}\pmod2$ on $D_9$ extends
invariantly to $D_8$ by
\[
 \alpha_{10}(1,f)=[S^5TST^3]f\pmod2,
\]
but has no invariant extension to $D$ or $N$. The original kernel is
\[
 \ker\!\left(H^1(N,\mathbf F_2)^G
                  \longrightarrow H^1(D_{10},\mathbf F_2)\right)
 =\left(\bigoplus_{q\text{ odd prime}}\mathbf F_2\psi_q\right)
                  \oplus\mathbf F_2\eta_3.
\]
The joint raw-coefficient map $(m_9,m_{10})$ on $D_8$ has image
$512\mathbf Z_2\times\mathbf Z_2$ already on $J_8$, with kernel
$D_{10}$. Its invariant quadratic-character kernel on restriction
to $D_{10}$ is $\mathbf F_2\eta_9\oplus\mathbf F_2\alpha_{10}$.
\end{proposition}
\begin{proof}
Indeed the exact weight-nine theorem puts every degree-nine leading
vector on $D_8$ in $32\mathbf Z_2\mathcal L_9$. Modulo two, such
words therefore start in degree at least ten, making the displayed
coefficient additive and invariant under the odd multiplier.

In the other direction, $\eta_{10}$ has no invariant extension to
$D$ or $N$. The arithmetic witness is $w_{10}=[\sigma,k]$ with
$k\in D$, so every such invariant character kills it, whereas
$\eta_{10}(w_{10})=1$. Consequently the original character kernel
is still
\[
 \ker\!\left(H^1(N,\mathbf F_2)^G
                  \longrightarrow H^1(D_{10},\mathbf F_2)\right)
 =\left(\bigoplus_{q\text{ odd prime}}\mathbf F_2\psi_q\right)
                  \oplus\mathbf F_2\eta_3.
\]
To see the last assertion, restrict a character in this kernel to
$D_9$. It is zero or $\eta_{10}$, since $D_9/D_{10}\simeq\mathbf Z_2$;
the second case cannot extend to $N$. The previously proved kernel
on $D_9$ then applies.

Finally, extend $m_{10}$ to $D_8$ by the same raw coefficient.
The map $(m_9,m_{10})$ is additive there, since all
nonconstant product and substitution corrections start above
degree ten. The arithmetic witnesses $w_9,w_{10}$ show that its
image is $512\mathbf Z_2\times\mathbf Z_2$, with kernel $D_{10}$.
Thus the invariant quadratic-character kernel on $D_8$ is
$\mathbf F_2\eta_9\oplus\mathbf F_2\alpha_{10}$.

\end{proof}

\subsection{A group-theoretic extension test in weight ten}
\label{subsec:eta-ten-extension-threshold}

\begin{proposition}\label{prop:eta-ten-no-extension-d7}
The character $\eta_{10}$ on $D_9$ extends invariantly to $D_8$
by $\alpha_{10}$, but has no $G$-invariant extension to $D_7$.
Consequently it has no such extension to $D_5$, $D_3$, $D$ or $N$.
\end{proposition}

\begin{proof}
The extension to $D_8$ was proved above. We give an obstruction
using only class-two commutator identities and the previously
established leading terms.

For $r\ge3$, let $\mathcal F^r$ be the subgroup of
$M=\ker\lambda$ whose pro-$2$ Magnus parameter starts in degree
at least $r$. The filtered normalized group law gives
\[
 [\mathcal F^r,\mathcal F^s]\subseteq\mathcal F^{r+s}.
\]
One has $M=\mathcal F^3$ and
$[M,M]\subseteq D_7\subseteq\mathcal F^8$. Thus
$[M,[M,M]]\subseteq\mathcal F^{11}$, and the quotient
$M/\mathcal F^{11}$ has nilpotence class at most two.
Moreover,
\[
 [\mathcal F^5,\mathcal F^5]\subseteq\mathcal F^{11}.
\]
Indeed its possible degree-ten leading term is the Ihara bracket
of two multiples of $\psi_5$, and this is zero. Vanishing over
$\mathbf Q_2$ implies vanishing in the integral leading module,
which is torsion-free.

Keep the arithmetic $\sigma,h$ of the weight-ten construction, and
put
\[
 a_0=\sigma,\qquad a_1=h\sigma h^{-1},\qquad
 a_2=h^2\sigma h^{-2},\qquad c_{ij}=[a_i,a_j].
\]
The weight-three law shows that
\[
 r=a_1a_0^{-27},\qquad q=a_2a_0^{-729}
\]
belong to $\mathcal F^5$. Therefore $[r,q]=1$ in
$M/\mathcal F^{11}$. All commutators are central in this quotient,
and expanding this identity gives
\[
 1=[a_1a_0^{-27},a_2a_0^{-729}]
   =c_{12}c_{01}^{729}c_{02}^{-27}.
\]
Hence $c_{12}=c_{02}^{27}c_{01}^{-729}$ there.

The arithmetic weight-eight and weight-ten witnesses are
\[
 w_8=c_{01},\qquad
 w_{10}=[a_0,(a_2a_1^{-27})(a_1a_0^{-27})^{-243}].
\]
The same class-two calculation gives
\[
 w_{10}=c_{02}c_{01}^{-270}\pmod{\mathcal F^{11}}.
\]
Consequently the word
\[
 z=h w_8h^{-1}w_8^{-6561}=c_{12}c_{01}^{-6561}
\]
satisfies
\[
 z=c_{02}^{27}c_{01}^{-7290}
   =w_{10}^{27}\pmod{\mathcal F^{11}}.
\]
Since $z\in J_7$ and its parameter agrees through degree ten with
that of $w_{10}^{27}$, one has $z\in J_9$ and
\[
 m_{10}(z)=27m_{10}(w_{10}),\qquad\eta_{10}(z)=1.
\]
On the other hand, every $G$-invariant character
$\chi:D_7\to\mathbf F_2$ satisfies
\[
 \chi(z)=\chi(h w_8h^{-1})-6561\chi(w_8)=0,
\]
because $6561$ is odd. Thus such a character cannot restrict to
$\eta_{10}$ on $D_9$. Any extension to one of the larger listed
subgroups would restrict to an extension on $D_7$, so those are
also impossible.
\end{proof}

\begin{corollary}\label{cor:weight-ten-character-kernels}
The invariant quadratic-character kernels on restriction to
$D_{10}$ are
\[
\begin{array}{c|l}
 \text{domain}&\text{kernel}\\ \hline
 D_9&\mathbf F_2\eta_{10}\\
 D_8&\mathbf F_2\eta_9\oplus\mathbf F_2\alpha_{10}\\
 D_7&\mathbf F_2\eta_8\oplus\mathbf F_2\eta_9\\
 D_5&\mathbf F_2\eta_7\oplus\mathbf F_2\alpha_8\\
 D_3&\mathbf F_2\eta_5\\
 N&\displaystyle\left(\bigoplus_{q\text{ odd prime}}\mathbf F_2\psi_q\right)
                    \oplus\mathbf F_2\eta_3.
\end{array}
\]
\end{corollary}

\begin{proof}
The first two rows follow from the exact scalar and joint maps.
For any larger listed domain, a character killed on $D_{10}$
restricts on $D_9$ to zero or $\eta_{10}$. The latter is impossible
by the preceding proposition. Such a character therefore already
kills $D_9$, and the previously proved weight-nine kernel table
applies.
\end{proof}


\section{All odd coefficients and the dyadic metabelian quotient}
\label{sec:uniform-dyadic-normalization}

The integral measure makes it possible to remove every dyadic
metabelian coefficient at once, while preserving the arithmetic
factorization. This is a reduction of the relative obstruction to a
smaller group; it does not determine that obstruction.

Use the dyadic beta multiplier $B_g(X,Y)$ and the integral
homomorphisms $\tau_{2r+1}$ of
Lemma~\ref{lem:integral-boundary-measure}. Put
\[
 C_0=D,\qquad
 C_r=D\cap\bigcap_{j=1}^{r}\ker\tau_{2j+1},\qquad
 J(C_r)=C_r\cap H,\qquad C_\infty=\bigcap_{r\ge0}C_r.
\]
The subgroups $C_1,C_2,C_3$ are $D_3,D_5,D_7$, respectively.
After this point the $C_r$ impose only odd metabelian coordinates;
the even mixed coefficients are treated separately below.

\begin{theorem}\label{thm:uniform-odd-arithmetic-image}
For every integer $r\ge1$, the coefficient
\[
 m_{2r+1}(g)=[S^{2r}T]f_g\qquad(g\in C_{r-1})
\]
is a continuous homomorphism of conjugation weight
$\lambda_2^{2r+1}$, with
\[
 m_{2r+1}(C_{r-1})=m_{2r+1}(J(C_{r-1}))
                 =2^{2r+1}\mathbf Z_2.
\]
Its kernel is $C_r$. In particular each inclusion
$C_r\subset C_{r-1}$ is proper, all these subgroups are closed and
normal in $G$, and
\[
 G=C_rH\quad(r\ge0),\qquad G=C_\infty H.
\]
Moreover
\[
 C_\infty=\{g\in D:B_g(X,Y)=1\},\qquad
 B\simeq C_\infty/
       \overline{[G,C_\infty](C_\infty\cap H)}.
\]
Here $B_g$ is its dyadic power-series component; no assertion about
other components of the adelic beta function is made.
\end{theorem}

\begin{proof}
Boundary reflection makes $q_g(e^u)$ even. Since it has constant
term zero, vanishing of $\tau_3,\ldots,\tau_{2r-1}$ is equivalent
to vanishing of the coefficients of $q_g(1+S)$ below $S^{2r}$.
On $C_{r-1}$ one therefore has
\[
 \tau_{2r+1}=\frac{(2r)!}{2}\,m_{2r+1}.
\]
The asserted additivity and weight law follow, since multiplication
by a nonzero integer on $\mathbf Z_2$ is injective.
Lemma~\ref{lem:uniform-odd-coordinate-bound} gives the upper bound
$m_{2r+1}\in2^{2r+1}\mathbf Z_2$.

Choose the arithmetic elements already used for the finite
normalizations: $\sigma\in H\cap\ker\lambda$ with
$\rho_2^I(\sigma)=2$ and all odd-prime rational Kummer coordinates
zero, and $h\in H$ with $t=\lambda(h)$ and $t_2=3$.
Then $m_3(\sigma)$ is odd. Lemma~\ref{lem:integral-boundary-measure}
gives
\[
 \tau_{2r+1}(\sigma)\equiv m_3(\sigma)\equiv1\pmod2
 \quad\text{for every }r\ge1.
\]
For a positive odd integer $j$ define
\[
 \Delta_j(z)=hzh^{-1}z^{-t^j},\qquad
 v_r=\Delta_{2r-1}\Delta_{2r-3}\cdots\Delta_3\Delta_1(\sigma).
\]
The exponent $t^j$ lies in $\widehat{\mathbf Z}^{\times}$, and the
power is taken in the closed procyclic subgroup generated by $z$.
The operator $\Delta_1$ kills every rational Kummer coordinate,
so $v_1\in D\cap H$. Each later operator kills the indicated
odd-weight homomorphism and preserves its predecessors' kernels.
Thus $v_r\in J(C_{r-1})$.

The exact value of its next coordinate is
\begin{equation}\label{eq:uniform-odd-arithmetic-witness}
 m_{2r+1}(v_r)=
 \frac{2\prod_{j=0}^{r-1}(3^{2r+1}-3^{2j+1})}{(2r)!}
                         \tau_{2r+1}(\sigma).
\end{equation}
For $k\ge1$, one has
$v_2(3^{2k}-1)=3+v_2(k)$.
For odd $k$, factor $9^k-1$ by $9-1$ and note that the remaining
sum is odd. Repeatedly squaring then adds one to the valuation,
because every $9^k+1$ has valuation one. Consequently the numerator
product in \eqref{eq:uniform-odd-arithmetic-witness}, before its
factor two, has valuation
\[
 \sum_{k=1}^{r}(3+v_2(k))=3r+v_2(r!).
\]
The denominator satisfies
$v_2((2r)!)=r+v_2(r!)$, by separating its even and odd factors.
The value in \eqref{eq:uniform-odd-arithmetic-witness} therefore has
valuation exactly $2r+1$. The closed arithmetic image contains all
profinite powers of this element, hence the whole ideal
$2^{2r+1}\mathbf Z_2$. Together with the upper bound, this proves
the two exact image equalities and identifies the kernel.

The weight laws prove normality, and the nonzero arithmetic value
proves properness. Equality of the full and arithmetic images gives
\[
 C_{r-1}=C_rJ(C_{r-1}).
\]
Starting with $G=DH$, induction gives $G=C_rH$ for all $r$.
For any fixed $g\in G$, the sets
\[
 \{a\in H:ga^{-1}\in C_r\}
\]
are nonempty, closed and nested in the compact group $H$.
Their intersection is nonempty. This proves $G=C_\infty H$.

Finally, vanishing of all $\tau_{2r+1}$ is equivalent to $q_g=0$.
The beta differential identity then gives
$B_g^{-1}Y\partial_YB_g=0$. In characteristic zero this makes
$B_g$ independent of $Y$, and $B_g(X,1)=1$ makes it identically
one. The converse follows by differentiation.
The relative-normalization argument used for
Corollary~\ref{cor:cubic-radical-normalization} applies to the
closed normal subgroup $C_\infty$ and the factorization just
proved, giving the displayed canonical quotient for $B$.
\end{proof}

\begin{proposition}\label{prop:uniform-odd-character-kernels}
For $s\ge0$, the kernel of restriction
\[
 H^1(C_s,\mathbf F_2)^G\longrightarrow
 H^1(C_\infty,\mathbf F_2)^G
\]
is the one-dimensional space generated by
\[
 \zeta_s(g)=m_{2s+3}(g)/2^{2s+3}\pmod2.
\]
For the original normalized group $N$, the kernel of restriction
to $C_\infty$ is
\[
 \left(\bigoplus_{q\text{ odd prime}}\mathbf F_2\psi_q\right)
                     \oplus\mathbf F_2\eta_3.
\]
Each of these kernels has injective restriction to the corresponding
arithmetic subgroup.
\end{proposition}

\begin{proof}
The character $\zeta_s$ is defined and invariant by the exact image
and weight law, and it vanishes on $C_{s+1}$. Its arithmetic value
on $v_{s+1}$ is one.
Let $\chi$ be an invariant character on $C_s$ vanishing on
$C_\infty$. Continuity makes $\ker\chi$ open. Compactness and the
nested intersection imply $C_r\subseteq\ker\chi$ for some
$r>s$: otherwise the compact sets $C_r\setminus\ker\chi$ would
have nonempty intersection.

Suppose $r\ge s+2$ and $\chi$ kills $C_r$. Its restriction to
$C_{r-1}$ is zero or the unique nonzero character of the quotient
$C_{r-1}/C_r\simeq\mathbf Z_2$. But
\[
 v_r=h v_{r-1}h^{-1}v_{r-1}^{-t^{2r-1}},
 \qquad v_{r-1}\in C_{r-2}\subseteq C_s.
\]
Invariance and the odd exponent give $\chi(v_r)=0$, whereas the
nonzero quotient character has value one by
\eqref{eq:uniform-odd-arithmetic-witness}. The restriction must
therefore be zero. Descending induction shows that $\chi$ kills
$C_{s+1}$, proving the first assertion.

For a character on $N$, the same compactness and descent argument
works until $C_1=D_3$: every $v_{r-1}$ used with $r\ge2$ lies
in $D\subseteq N$. Thus the kernel on $N$ is exactly its already
computed kernel on restriction to $D_3$, namely the displayed
rational Kummer span and $\eta_3$. The converse is immediate from
$C_\infty\subseteq D_3$. The arithmetic value of $\zeta_s$ and
the previously proved independence of the rational Kummer
characters and $\eta_3$ prove the last assertion.
\end{proof}

\begin{corollary}\label{cor:dyadic-metabelian-deep-gate}
Put
\[
 E=C_\infty\cap D_{10},\qquad J_E=E\cap H.
\]
Then $E$ is a proper closed normal subgroup of $D_{10}$ and of
infinite index there, and
\[
 G=EH,\qquad B\simeq E/\overline{[G,E]J_E}.
\]
Every pro-$2$ Magnus word on $E$ starts in degree at least eleven,
and its dyadic beta multiplier is identically one.
The kernel of restriction from $H^1(N,\mathbf F_2)^G$ to
$H^1(E,\mathbf F_2)$ is exactly
\[
 \left(\bigoplus_{q\text{ odd prime}}\mathbf F_2\psi_q\right)
                   \oplus\mathbf F_2\eta_3.
\]
\end{corollary}

\begin{proof}
We first verify the factorization without assuming any joint-image
statement. Since $C_\infty\subseteq D_7$, the coordinate $m_8$
is defined there. The arithmetic word $w_8$ of the eighth
normalization is a commutator of multiplier-one elements. The
beta multiplier is a homomorphism on $M=\ker\lambda$, so
$B_{w_8}=1$. Hence $w_8\in C_\infty\cap H$, and its $m_8$
value is a unit. Correcting by its powers gives
\[
 C_\infty=(C_\infty\cap D_8)(C_\infty\cap H).
\]
Since $m_9=0$ on $C_\infty$, the first factor lies in $D_9$.
The arithmetic word $w_{10}$ of the tenth normalization is also
a commutator of multiplier-one elements, so belongs to
$C_\infty\cap J_9$ and has an odd $m_{10}$ value. Correcting
by its powers gives
\[
 C_\infty\cap D_8=E(C_\infty\cap D_8\cap H).
\]
Together with $G=C_\infty H$, these equations prove $G=EH$.
Normality and the relative quotient follow as before.

To prove strictness, take the word $v_5\in C_4\cap H$ in
\eqref{eq:uniform-odd-arithmetic-witness}. It has
$m_{11}(v_5)\in2^{11}\mathbf Z_2^{\times}$.
First correct its eighth coordinate by a power of $w_8$.
The result lies in $C_4\cap D_8\subseteq D_9$.
Then correct its tenth coordinate by a power of $w_{10}$.
The result $v$ belongs to $D_{10}\cap H$.
Both correcting words have beta multiplier one, so they leave
$\tau_{11}$, and therefore the nonzero eleventh depth-one
coefficient, unchanged. Thus $v\notin E$.
In fact $D_{10}\subseteq C_4$, so the uniform bound and the
corrected arithmetic word give the exact images
\[
 m_{11}(D_{10})=m_{11}(D_{10}\cap H)=2^{11}\mathbf Z_2.
\]
This additive coordinate vanishes on $E$, proving infinite index. The final assertions follow
from $E\subseteq D_{10}$ and its definition.

For the last kernel assertion, every invariant character $\chi$ on
$N$ kills both correcting words. Indeed,
\[
 w_8=[\sigma,h\sigma h^{-1}]=[\sigma,k_8],\qquad
 k_8=(h\sigma h^{-1})\sigma^{-t}\in D\subseteq N;
\]
the rational Kummer coordinates of $k_8$ cancel by their exact
weight-one law. Invariance under $\sigma\in G$ gives
$\chi(w_8)=0$. The tenth construction gives
$w_{10}=[\sigma,k_{10}]$ for $k_{10}\in D$ and hence
$\chi(w_{10})=0$ as well. The two correction factorizations in
this proof now imply that $\chi$ vanishes on $E$ if and only if
it vanishes on $C_\infty$. Proposition~\ref{prop:uniform-odd-character-kernels}
gives the asserted kernel.
\end{proof}

\section{A primitive mixed coordinate in degree eleven}
\label{sec:eleven-mixed-mixed}

We remove a primitive mixed coefficient while preserving the arithmetic
factorization. Use the groups $G$, $H$, $N$, $D$, and $D_{10}$ defined above.
Write $M=\ker\lambda$ and
\[
 C_\infty=\{g\in D:B_g(X,Y)=1\},\qquad
 E=C_\infty\cap D_{10},\qquad J(E)=E\cap H.
\]
Here $B_g$ is the dyadic beta multiplier. The preceding normalizations
give $G=EH$, and every Magnus series $f_g$ with $g\in E$ has no
nonconstant term below degree eleven. We use these established facts
and the exact degree-eight leading lattice below. The argument uses one coordinate of the degree-eleven leading space.

\begin{theorem}\label{thm:eleven-mixed-mixed-reduction}
The coordinate
\[
 b_{11}(g)=[S^6TST^3]f_g\qquad(g\in E)
\]
is a continuous additive homomorphism satisfying
\[
 b_{11}(aga^{-1})=\lambda_2(a)^{11}b_{11}(g)
 \quad(a\in G,\ g\in E).
\]
Its full and arithmetic images are both saturated:
\[
 b_{11}(E)=b_{11}(J(E))=\mathbf Z_2.
\]
Consequently
\[
 E'=\ker(b_{11}|_E)
\]
is a proper closed normal subgroup of $G$, and
\[
 E/E'\simeq\mathbf Z_2,\qquad G=E'H,
 \qquad
 B\simeq E'/\overline{[G,E'](E'\cap H)}.
\]
For $X_2=H^1(N,\mathbf F_2)^G$, its restriction kernel remains
\[
 \ker\bigl(X_2\longrightarrow H^1(E',\mathbf F_2)\bigr)
 =\left(\bigoplus_{q\text{ odd prime}}\mathbf F_2\psi_q\right)
       \oplus\mathbf F_2\eta_3.
\]
The invariant character $b_{11}\bmod2$ on $E$ is nonzero on $J(E)$
and has no $G$-invariant extension to $N$.
\end{theorem}

\begin{proof}
We first give the integral coefficient calculation that supplies an
arithmetic unit. For Lie polynomials use
\[
 P_3=[S,[S,T]]+[T,[S,T]],\qquad
 \{A,B\}=[A,B]+\partial_A B-\partial_B A,
\]
where $\partial_A(S)=0$ and $\partial_A(T)=[T,A]$.
This is the Ihara bracket in the convention of the GT product.
Let $L_8$ be the integer Lie polynomial for the already established
degree-eight leading line, normalized so that
\[
 [S^4TST^2]L_8=7.
\]
For specificity it is completely specified by its coordinates in the
lexicographically ordered Lyndon basis of weight eight, with $S<T$:
\begin{align*}
(&0,0,-4,0,-2,7,17,0,-1,-12,-6,4,-3,-6,0,\\
 &-5,3,3,6,7,21,3,-15,-3,0,8,-11,-4,-10,0).
\end{align*}
The Lyndon bracket of a nonletter word is $[L_u,L_v]$, where $v$
is its longest proper Lyndon suffix and the word is $uv$.
This finite recursion gives, in particular,
\begin{equation}\label{eq:eleven-mixed-four-coefficients}
 [S^4TST^2]L_8=7,\qquad [S^5TST]L_8=-4,\qquad
 [S^5T^3]L_8=[S^6T^2]L_8=0.
\end{equation}
An element $z\in D_7$ therefore has degree-eight term
$(m_8(z)/7)L_8$.

Put $\omega=S^6TST^3$. Expanding
\[
 P_3=S^2T-2STS+TS^2+2TST-T^2S-ST^2
\]
shows that $[\omega]P_3L_8=[\omega]L_8P_3=0$: the required
prefix or suffix of length three is absent from $P_3$. Moreover
\[
 [T,P_3]=-2TSTS+T^2S^2+3T^2ST-T^3S-3TST^2
             -S^2T^2+2STST+ST^3.
\]
Replacing one $T$ in a word of $L_8$ by this expression gives
\[
 [\omega]\partial_{P_3}L_8
     =2[S^5T^3]L_8-2[S^6T^2]L_8=0.
\]
For $\partial_{L_8}P_3$, the only surviving contributions are
$-S^2L_8T$ from $S^2T$ and $+SL_8T^2$ from $-ST^2$.
Thus the final, negative derivation term contributes
\[
 -[\omega]\partial_{L_8}P_3
    =[S^4TST^2]L_8-[S^5TST]L_8=7+4=11.
\]
The four Ihara summands have coefficients $(0,0,0,11)$, proving
\begin{equation}\label{eq:eleven-mixed-bracket-unit}
 [S^6TST^3]\{P_3,L_8\}=11.
\end{equation}

Choose the arithmetic elements from the earlier exact
normalizations: $\sigma\in H\cap M$ with
$\rho_2^I(\sigma)=2$ and all odd-prime rational Kummer coordinates
zero, and $h\in H$ with $t=\lambda(h)$ and $t_2=3$.
Let
\[
 a=m_3(\sigma),\qquad u=\tau_5(\sigma).
\]
The established integral moment parity gives $a,u\in\mathbf Z_2^\times$.
The arithmetic word
\[
 w_8=[\sigma,h\sigma h^{-1}]
\]
belongs to $D_7\cap H$ and satisfies the exact earlier formula
\begin{equation}\label{eq:eleven-mixed-old-eight}
 m_8(w_8)=63au.
\end{equation}
Recall that all rational Kummer coordinates and the dyadic beta
multiplier are homomorphisms on $M$, the latter into a commutative
group. In particular $B_{w_8}=1$.

Define the new arithmetic commutator
\[
 w_{11}=[\sigma,w_8].
\]
The first nonconstant terms of $f_\sigma$ and $f_{w_8}$ have
degrees three and eight. The filtered GT commutator has no term
below degree eleven, and its term in degree eleven is their
Ihara bracket. Its rational Kummer coordinates and beta multiplier
are trivial, since both inputs lie in $M$. Hence $w_{11}\in E\cap H$.
Equations \eqref{eq:eleven-mixed-bracket-unit} and
\eqref{eq:eleven-mixed-old-eight} now give the exact value
\begin{equation}\label{eq:eleven-mixed-arithmetic-unit}
 b_{11}(w_{11})
  =\frac{11}{7}a\,m_8(w_8)
  =99a^2u\in\mathbf Z_2^\times.
\end{equation}
There are no contributions from higher terms of either input to
degree eleven.

On a subgroup whose words start in degree eleven, the degree-eleven
part of the GT product is the sum of the two leading parts.
Conjugation by an arbitrary element of $G$ multiplies this
homogeneous part by $\lambda_2^{11}$: on the associated graded
algebra both generators are multiplied by $\lambda_2$, and all
conjugating corrections have higher degree. This proves the
claimed homomorphism and weight laws. Coefficient extraction is
continuous. The closed arithmetic image contains the image of the
procyclic group generated by $w_{11}$, which is all of $\mathbf Z_2$
by \eqref{eq:eleven-mixed-arithmetic-unit}. This proves both image
equalities and the normality and properness of $E'$.

For each $e\in E$, choose $r\in\widehat{\mathbf Z}$ whose
2-adic component is $b_{11}(e)/b_{11}(w_{11})$. The legitimate
profinite power $w_{11}^r$ lies in $J(E)$ and has the same
coordinate as $e$. Thus $e w_{11}^{-r}\in E'$, proving
$E=E'J(E)$ and hence $G=E'H$.

For completeness, let $K=N\cap H$ and
$L=\overline{[G,E']}$. Since $N=E'K$, the image of $E'$ is
central in $G/L$, while $[H,K]\subseteq K$. Consequently
\[
 \overline{[G,N]K}=LK.
\]
The product is closed because both factors are compact and $L$
is normal. Also $E'\cap LK=L(E'\cap H)$. The inclusion of $E'$
therefore induces the stated canonical presentation of $B$.

Finally let $\chi\in X_2$. Since $w_8\in N$ and $\sigma\in G$,
invariance gives
\[
 \chi(w_{11})=\chi(\sigma w_8\sigma^{-1})-\chi(w_8)=0.
\]
If $\chi$ vanishes on $E'$, the same correction by $w_{11}^r$
shows that it vanishes on $E$. The converse is immediate, so the
restriction kernels to $E$ and $E'$ coincide. The previously proved
kernel for $E$ is exactly the displayed rational Kummer span plus
$\mathbf F_2\eta_3$. On the other hand $b_{11}\bmod2$ is invariant
and takes value one on $w_{11}$, proving both its arithmetic
nonvanishing and its nonextension to $N$.
\end{proof}

\begin{remark}\label{rem:eleven-mixed-scope}
This theorem removes a saturated arithmetic coordinate from the
beta-trivial group. It does not produce a nonzero character of the
relative quotient $B$. It also does not assert that the words of
$E'$ start in degree twelve: that stronger conclusion requires a
separate proof that $b_{11}$ detects the whole degree-eleven
beta-trivial leading space.
\end{remark}


\section{Questions remaining}
The group $B$ of Theorem~\ref{thm:main-structural}, its maximal pro-$2$ quotient and its maximal pro-$p$ quotients for odd $p$ are not determined here. The equality $\GT=\GQ$ is therefore not decided. The rational Kummer characters form a subspace on which restriction is injective, but neither its complement in $X_2$ nor the full image in $Y_2$ is determined. Enriquez's distribution theorem makes all multiplicative defects of these cocycles zero on $\GT$. The remaining character obstruction can be computed on $E'=\ker(b_{11}|_{C_\infty\cap D_{10}})$, a proper subgroup of $E=C_\infty\cap D_{10}$. On $E'$ all rational Kummer cocycles, the Magnus terms through degree ten, the entire dyadic beta multiplier, and the mixed coordinate $b_{11}$ vanish. Vanishing of the other degree-eleven terms has not been proved. It can alternatively be computed on $C_{\mathrm{CM}}$, with Galois intersection $G_{K_3^{\ab}K_4^{\ab}}$. For the radical and CM normalizations, the exact joint-image gate consists of the tame cubic tests and the tame harmonic quadratic tests. The wild cubic, complete prime-two and exceptional prime-three conditions have been proved. Combining the Magnus and CM normalizations also requires their joint arithmetic image. The explicit determinant character $\psi_{17,11}$ of Corollary~\ref{cor:concrete-radical-anderson-character} belongs to this obstruction space, but its vanishing or nonvanishing remains undecided. Likewise, strict level-three compatibility is reduced to $d_3(q)=q$ in the full profinite modular group; all its congruence images are already equal.

\section*{Method and records}
The mathematics and text were produced by GPT-6 Astra Pro (ChatGPT Pro mode). The quotient, character extension, parity, dimension, simultaneous rational-radical normalization, normalized congruence, individual-frame comparison and conductor arguments received independent same-model re-proofs. The cubic and weight-five divisibilities, exact arithmetic images and further normalizations received separate statement-only re-proofs, including integral coefficient calculations, as well as targeted checks. The prime-two compatibility, exact radical--CM intersection, and exceptional-three division field and relative character also received complete independent same-model re-proofs. The weight-eight exact image and further normalization received a complete statement-only re-proof by the same model; its sharper extension threshold received a second complete statement-only proof using complex conjugation. For weight seven, the initial independent reconstruction proved a weaker upper bound; a subsequent complete targeted check verified the exact bound and the printed full-pentagon certificate by a separate implementation. The finite relative comparison characters received a complete separate statement-only re-proof, followed by a targeted comparison. The normalized dyadic harmonic comparison received a complete separate statement-only re-proof. The exact weight-nine normalization received a complete separate statement-only re-proof, using an integral measure for the upper bound and a separate leading-space calculation. The weight-ten normalization and its extension threshold received a complete separate statement-only re-proof, with a different arithmetic witness and leading-space certificate. The uniform dyadic normalization received a complete proof-informed reconstruction and targeted check, followed by a complete separate statement-only re-proof; the existence of a measure approach had been announced, but its argument had not been disclosed. The primitive mixed degree-eleven reduction received a complete targeted proof check and a separate coefficient recomputation; its statement-only check is pending. The cubic comparison received a complete targeted proof and primary-source check; a separate statement-only attempt recovered its algebraic reduction but did not complete the geometric comparison. The finite CM tests, adelic CM comparison and finite-cover character arguments received targeted proof checks and, where indicated, exact coefficient verification. Each check is recorded with its scope; the effort setting of the same-model checks is not exposed. The cited background note records checks by Claude Fable 5.2 and Claude Opus 5.5 at effort max. These are AI checks; no human has checked this paper. The text and original programs are dedicated to the public domain under CC0 1.0. The workbench is \url{\workbenchurl}; the permanent records are at \url{\recordurl}.

\begin{thebibliography}{99}
\bibitem{Drinfeld1990} V. G. Drinfeld, On quasitriangular quasi-Hopf algebras and on a group that is closely connected with $\operatorname{Gal}(\overline{\Q}/\Q)$, Leningrad Math. J. \textbf{2} (1991), 829--860.
\bibitem{Ihara1999} Y. Ihara, On beta and gamma functions associated with the Grothendieck--Teichm\"uller group, in \emph{Aspects of Galois Theory}, London Math. Soc. Lecture Note Ser. 256, Cambridge University Press, 1999, 144--179, Proposition~1.6.1 and \S1.8.
\bibitem{IchimuraKaneko1989} H. Ichimura and M. Kaneko, On the universal power series for Jacobi sums and the Vandiver conjecture, J. Number Theory \textbf{31} (1989), 312--334, Section~1.2 and Theorem~1, pp.~314--315; \href{https://www2.math.kyushu-u.ac.jp/~mkaneko/papers/04universal_power.pdf}{author's scan}.
\bibitem{Anderson1989} G. W. Anderson, The hyperadelic gamma function, Invent. Math. \textbf{95} (1989), 63--131, Corollary 8.6.3; the beta formulation used here is stated in \cite[Remark 2.4, p.~101]{Furusho2003}.
\bibitem{HerfortRibes1985} W. Herfort and L. Ribes, Torsion elements and centralizers in free products of profinite groups, J. reine angew. Math. \textbf{358} (1985), 155--161, Theorem B, p.~155, and Theorem B$'$ in \S4, p.~159.
\bibitem{Furusho2003} H. Furusho, Geometric and arithmetic subgroups of the Grothendieck--Teichm\"uller group, Math. Res. Lett. \textbf{10} (2003), 97--108, \S2.2, p.~100.
\bibitem{Enriquez2007} B. Enriquez, Quasi-reflection algebras and cyclotomic associators, Selecta Math. (N.S.) \textbf{13} (2007), 391--463, \href{https://arxiv.org/abs/math/0408035v2}{arXiv:math/0408035v2}, \S8, Theorem 8.1, equation (50) and Remark 8.6, author pp.~30--32.
\bibitem{MoreeStevenhagen} P. Moree and P. Stevenhagen, Computing higher rank primitive root densities, Acta Arith. \textbf{163} (2014), 15--32, \href{https://arxiv.org/abs/1203.4313v2}{arXiv:1203.4313v2}, \S2, Theorem 2.5, author pp.~4--6.
\bibitem{NS2000} H. Nakamura and L. Schneps, On a subgroup of the Grothendieck--Teichm\"uller group acting on the tower of profinite Teichm\"uller modular groups, Invent. Math. \textbf{141} (2000), 503--560, Theorems 1.3--1.4 and Definition 5.1.
\bibitem{HLS} A. Hatcher, P. Lochak and L. Schneps, On the Teichm\"uller tower of mapping class groups, J. reine angew. Math. \textbf{521} (2000), 1--24.
\bibitem{LS2006} P. Lochak and L. Schneps, Open problems in Grothendieck--Teichm\"uller theory, in \emph{Problems on Mapping Class Groups and Related Topics}, Proc. Sympos. Pure Math. 74, American Mathematical Society, 2006, 165--186, \S1.
\bibitem{Lochak2025} P. Lochak, Automorphisms of profinite and procongruence curve complexes and the Grothendieck--Teichm\"uller group, Algebra i Analiz \textbf{37} (2025), no.~2, author version, Definition 5.4 and Theorem 5.10, pp.~50, 60--62.
\bibitem{LNS} P. Lochak, H. Nakamura and L. Schneps, The Grothendieck--Teichm\"uller group acts on the profinite mapping class group of genus $g$, \href{https://arxiv.org/abs/2602.12462}{arXiv:2602.12462}, Theorem 1.1.
\bibitem{Boggi} M. Boggi, Automorphisms of profinite mapping class groups, \href{https://arxiv.org/abs/2011.15075v6}{arXiv:2011.15075v6}, 22 April 2026, Theorem 3.12, Corollary 3.15, Remark 3.11, Proposition 6.9 and Lemma 6.16.
\bibitem{H17} \emph{Relations $(\mathrm{III}')$ and $(\mathrm{IV})$ of Nakamura and Schneps hold on the profinite Grothendieck--Teichm\"uller group, with an application to the two elliptic points of the modular curve}, version 2, 24 November 2026, \href{https://hexagonmath.org/2611.00940v2}{hexagon:2611.00940v2}, Theorems 1--3 and 20, Corollaries 21--22, Lemmas 5, 7 and 15, Propositions 11 and 18.
\bibitem{Monoid} \emph{The profinite Grothendieck--Teichm\"uller monoid is a group, on a question of Lochak and Schneps}, 25 November 2026, \href{https://hexagonmath.org/2611.01024}{hexagon:2611.01024}, Theorems A and E.
\bibitem{MilneCFT} J. S. Milne, \emph{Class Field Theory}, version 4.03, 2020, Chapter V, Theorem 3.27, p.~165, \url{https://www.jmilne.org/math/CourseNotes/CFT.pdf}.
\bibitem{LMFDB256} The LMFDB Collaboration, Classical modular form $256.1.c.a$, \url{https://www.lmfdb.org/ModularForm/GL2/Q/holomorphic/256/1/c/a/}, accessed 25 November 2026.
\bibitem{SchmidtWingberg} A. Schmidt and K. Wingberg, Shafarevich's theorem on solvable groups as Galois groups, \href{https://arxiv.org/abs/math/9809211}{arXiv:math/9809211}, Theorem 1, p.~1.
\bibitem{MilneMF110}
J. S. Milne, \emph{Modular Functions and Modular Forms},
version 1.10, 1997, Remark~12.12, Proposition~12.13,
and Theorem~12.27, printed pp.~119 and~124,
\url{https://www.jmilne.org/math/CourseNotes/MF110.pdf}.

\bibitem{PapikianTate}
M. Papikian, \emph{Rigid analytic geometry and abelian varieties},
notes dated 13 June 2004, \S2, Theorem~2.2 and pp.~3--4,
\url{https://math.stanford.edu/~vakil/snowbird/mihranjun21.pdf}.

\bibitem{Marin2008}
I. Marin, Characters of the Grothendieck--Teichm\"uller group through
rigidity of the Burau representation, Ann. Fac. Sci. Toulouse Math.
(6) \textbf{17} (2008), 765--780, \S6.3, p.~780,
\url{https://www.numdam.org/item/10.5802/afst.1202.pdf}.

\bibitem{NT2002}
H. Tsunogai and H. Nakamura, Harmonic equations in the
Grothendieck--Teichm\"uller group, RIMS K\^oky\^uroku
\textbf{1267} (2002), 197--208, Theorem~2, pp.~199--200,
\url{https://www.kurims.kyoto-u.ac.jp/~kyodo/kokyuroku/contents/pdf/1267-18.pdf}.


\bibitem{KKSV2016}
K. Katz, M. Katz, M. M. Schein and U. Vishne,
\emph{Bolza quaternion order and asymptotics of systoles along
congruence subgroups}, Experimental Mathematics \textbf{25}
(2016), no.~4, 399--415,
doi:10.1080/10586458.2015.1073642.
Author version: \url{https://u.math.biu.ac.il/~scheinm/text61.pdf};
arXiv:1405.5454.
\bibitem{KRRBolza}
V. Koziarz, C. Rito and X. Roulleau,
\emph{The Bolza curve and some orbifold ball quotient surfaces},
author PDF, \S4, equation (4.1),
\url{https://www.math.u-bordeaux.fr/~vkoziarz/Bolza.pdf}.

\bibitem{CubicHMM} Y.~Hoshi, A.~Minamide and S.~Mochizuki,
\emph{Group-theoreticity of numerical invariants and distinguished
subgroups of configuration space groups}, Kodai Math. J.
\textbf{45} (2022), 295--348. Definition~2.7 and
Corollary~2.8(i), author pp.~31--32.
\href{https://www.kurims.kyoto-u.ac.jp/~motizuki/Numerical\%20invariants\%20and\%20distinguished\%20subgroups\%20of\%20configuration\%20spaces.pdf}{Author PDF}.
\bibitem{CubicBZ} M.~Boggi and P.~Zalesskii,
\emph{Finite subgroups of the profinite completion of good groups},
arXiv:2406.08639v2, 28 October 2024. Theorem~A(i),(iii),
p.~3; Section~6.1 and Theorem~6.2(i),(iii), pp.~16--17.
\href{https://arxiv.org/pdf/2406.08639v2}{Version 2 PDF}.
\bibitem{CubicAsada} M.~Asada,
\emph{The faithfulness of the monodromy representations associated
with certain families of algebraic curves}, RIMS preprint 1179
(1997). Introduction Theorem, p.~3, parts (i)--(ii), summarizing
Theorems~2 and~3A; the identification with full profinite
mapping class groups is given there.
\href{https://www.kurims.kyoto-u.ac.jp/preprint/file/RIMS1179.ps.gz}{Original preprint}.
\bibitem{CubicBH} J.~S.~Birman and H.~M.~Hilden,
\emph{On isotopies of homeomorphisms of Riemann surfaces},
Ann. of Math. \textbf{97} (1973), 424--439.
Theorem~1, p.~424, and the extremal-map proof on p.~433.
\href{https://www.math.columbia.edu/~jb/bir-hilden-annals.pdf}{Author-hosted scan}.
\bibitem{CubicBHErr} J.~S.~Birman and H.~M.~Hilden,
\emph{Erratum}, 2017, one page, concerning Lemma~5.1 and
Theorem~5 of the preceding article; Theorems~1--4 are unaffected.
\href{https://www.math.columbia.edu/~jb/bh-erratum.pdf}{Author-hosted erratum}.
\bibitem{CubicKerckhoff} S.~P.~Kerckhoff,
\emph{The Nielsen realization problem}, Ann. of Math.
\textbf{117} (1983), 235--265. Section~4, pp.~255--256,
extends the realization argument to finite-type groups with cusps.
\href{https://people.math.harvard.edu/~ctm/home/text/others/kerckhoff/nielsen/nielsen.pdf}{Primary article scan}.
\bibitem{CubicGH} F.~P.~Gardiner and J.~Hu,
\emph{A short course on Teichm\"uller's theorem},
in \emph{Proceedings of the Year on Teichm\"uller Theory},
HRI, Allahabad, vol.~10 (2009), 195--228.
Section~5, Theorem~1 and its punctured-surface hypothesis,
author pp.~14--15; Section~8, Theorem~6, author p.~22.
\href{https://userhome.brooklyn.cuny.edu/gardiner/A\%20short\%20course\%20on\%20Teichmuller\%27s\%20theorem.pdf}{Author PDF}.
\bibitem{Schneps2006Special} L.~Schneps,
Automorphisms of curves and their role in Grothendieck--Teichm\"uller
theory, Math. Nachr. \textbf{279} (2006), 656--671,
Definition~1.2, p.~658, Theorem~3.1, p.~665, and
Theorem~3.4 with equations~(3.10)--(3.11), p.~667.
\href{https://webusers.imj-prg.fr/~leila.schneps/automorphisms.pdf}{Author-hosted journal PDF}.
\end{thebibliography}
\end{document}
